\documentclass[11pt]{article}
\usepackage[a4paper,margin=18mm]{geometry}
\usepackage[no-math]{fontspec}
\setmainfont{Latin Modern Roman}
\usepackage{amsmath,amssymb,amsthm,xfrac,xspace,hyperref,graphicx,comment}
\excludecomment{DefTralics}

\providecommand{\bibliofont}{\small}
\newcommand{\ie}{i.e.,\xspace}
\newcommand{\eg}{e.g.,\xspace}
\newcommand{\etc}{etc.\xspace}
\newcommand{\cf}{cf.\xspace}
\newcommand{\resp}{resp.\xspace}
\newcommand{\smaller}{\small}
\newcommand{\Subsection}{\subsection}
\newcommand{\Subsubsection}{\subsubsection}
\newcommand{\skpt}{\smallskip}
\emergencystretch=3em
\numberwithin{equation}{section}\providecommand{\og}{«\,}\providecommand{\fg}{\,»}

% Begin literal retained input author-preamble.tex
\setcounter{tocdepth}{2}

\usepackage[all,cmtip]{xy}\let\labelstyle\textstyle
\RequirePackage[mathscr]{eucal}
\let\mathcal \EuScript
\let\cal\mathcal

\multlinegap0pt

\newcommand{\citePR}{\cite{P-R08}}
\begin{DefTralics}
\newcommand{\citePR}{[28]}
\end{DefTralics}

\let\oldenumerate\enumerate
\renewcommand{\enumerate}{\oldenumerate\itemindent0pt}

\newenvironment{enumeratei}
{\bgroup\def\theenumi{\roman{enumi}}\def\theenumii{\arabic{enumii}}\oldenumerate}
{\endenumerate\egroup}

\newcommand{\Changel}{\mathcode`l="7160}
\newcommand{\Changelback}{\mathcode`l="716C}
\newcommand{\NoChangel}[1]{%
\expandafter\let\csname old\string#1\endcsname=#1
\let#1=\relax
\newcommand{#1}{\mathcode`l="716C\csname old\string#1\endcsname\mathcode`l="7160 }%
}

\NoChangel{\log}
\NoChangel{\ln}
\NoChangel{\lim}
\NoChangel{\limsup}
\NoChangel{\liminf}
\NoChangel{\varlimsup}
\NoChangel{\varliminf}
\NoChangel{\varinjlim}
\NoChangel{\varprojlim}

\let\epsilon\varepsilon
\let\<\langle
\let\>\rangle
\let\sep\mid
\let\oldS\S
\renewcommand{\S}{\oldS\,}
\newcommand{\quadet}{\quad\text{et}\quad}

\newcommand\mto{\mathchoice{\longmapsto}{\mapsto}{\mapsto}{\mapsto}}
\def\To#1{\mathchoice{\xrightarrow{\textstyle\kern4pt#1\kern3pt}}{\stackrel{#1}{\longrightarrow}}{}{}}
\let\oldbigsqcup\bigsqcup
\renewcommand{\bigsqcup}{\mathop{\textstyle\oldbigsqcup}\displaylimits}

\newcommand\Cf{Cf.\kern.3em}

\newcommand{\gf}{{}^g\!f}
\newcommand{\nf}{{}^n\!f}
\newcommand{\rmod}{\textup{mod}\,}

\newcommand{\frb}{\mathfrak{b}}
\newcommand{\frg}{\mathfrak{g}}
\newcommand{\frk}{\mathfrak{k}}
\newcommand{\frl}{\mathfrak{l}}\NoChangel{\frl}
\newcommand{\mfrm}{\mathfrak{m}}
\newcommand{\fro}{\mathfrak{o}}
\newcommand{\frp}{\mathfrak{p}}
\newcommand{\mfrq}{\mathfrak{q}}
\newcommand{\frt}{\mathfrak{t}}
\newcommand{\fru}{\mathfrak{u}}
\newcommand{\frz}{\mathfrak{z}}
\newcommand{\frV}{\mathfrak{V}}
\newcommand{\frY}{\mathfrak{Y}}

\newcommand{\AD}{\mathrm{AD}}
\newcommand{\Ann}{\mathrm{Ann}}
\newcommand{\App}{\mathrm{App}}
\newcommand{\rc}{\mathrm{c}}
\newcommand{\Cl}{\mathrm{Cl}}\NoChangel{\Cl}
\newcommand{\comp}{\mathrm{comp}}
\newcommand{\tcomp}{\textup{-comp}}
\newcommand{\cusp}{\mathrm{cusp}}
\newcommand{\diag}{\mathrm{diag}}
\newcommand{\rell}{\mathrm{ell}} \NoChangel{\rell}
\newcommand{\Ell}{\mathrm{Ell}} \NoChangel{\Ell}
\newcommand{\rendo}{\mathrm{endo}}
\newcommand{\ent}{\mathrm{ent}}
\newcommand{\Fac}{\mathrm{Fac}}
\newcommand{\Fil}{\mathrm{Fil}} \NoChangel{\Fil}
\newcommand{\Gr}{\mathrm{Gr}}
\newcommand{\Hyp}{\textup{Hyp}}
\newcommand{\Imm}{\mathrm{Imm}}
\newcommand{\Ind}{\mathrm{Ind}}
\newcommand{\Int}{\mathrm{Int}}
\newcommand{\inst}{\mathrm{inst}}
\newcommand{\Irr}{\mathrm{Irr}}
\newcommand{\nil}{\mathrm{nil}} \NoChangel{\nil}
\newcommand{\Nil}{\mathrm{Nil}} \NoChangel{\Nil}
\newcommand{\nr}{\mathrm{nr}}
\newcommand{\trel}{\textup{-rel}} \NoChangel{\trel}
\newcommand{\reg}{\mathrm{reg}}
\newcommand{\res}{\mathrm{res}}
\newcommand{\SI}{\mathrm{SI}}
\newcommand{\rss}{\mathrm{ss}}
\newcommand{\st}{\mathrm{st}}
\newcommand{\Temp}{\mathrm{Temp}}
\newcommand{\tn}{\mathrm{tn}}
\newcommand{\tu}{\mathrm{tu}}
\newcommand{\tors}{\mathrm{tors}}

\DeclareMathOperator{\ad}{ad}
\DeclareMathOperator{\Ad}{Ad}
\DeclareMathOperator{\Aut}{Aut}
\DeclareMathOperator{\Gal}{Gal}\NoChangel{\Gal}
\DeclareMathOperator{\im}{Im}\let\Im\im
\DeclareMathOperator{\rI}{I}
\DeclareMathOperator{\rIE}{I\cE}
\DeclareMathOperator{\rIH}{I\cH}
\DeclareMathOperator{\Ker}{Ker}
\DeclareMathOperator{\mes}{mes}
\DeclareMathOperator{\Norm}{Norm}
\DeclareMathOperator{\Norme}{Norme}
\DeclareMathOperator{\transfert}{transfert}
\DeclareMathOperator{\val}{val} \NoChangel{\val}

\newcommand{\bmd}{\boldsymbol{d}}
\newcommand{\bmf}{\boldsymbol{f}}
\newcommand{\gbmf}{{}^g\!\boldsymbol{f}}
\newcommand{\bmu}{\boldsymbol{u}}
\newcommand{\bmv}{\boldsymbol{v}}

\newcommand{\bA}{\boldsymbol{A}}
\newcommand{\bD}{\boldsymbol{D}}
\newcommand{\bG}{\boldsymbol{G}}
\newcommand{\bl}{\boldsymbol{l}}
\newcommand{\bL}{\boldsymbol{L}}
\newcommand{\bM}{\boldsymbol{M}}
\newcommand{\bP}{\boldsymbol{P}}
\newcommand{\bQ}{\boldsymbol{Q}}
\newcommand{\bU}{\boldsymbol{U}}

\newcommand{\cA}{\mathcal{A}}
\newcommand{\cD}{\mathcal{D}}
\newcommand{\cE}{\mathcal{E}}
\newcommand{\cF}{\mathcal{F}}
\newcommand{\cG}{\mathcal{G}}
\newcommand{\cH}{\mathcal{H}}
\newcommand{\cL}{\mathcal{L}}
\newcommand{\cM}{\mathcal{M}}
\newcommand{\cN}{\mathcal{N}}
\newcommand{\cO}{\mathcal{O}}
\newcommand{\cP}{\mathcal{P}}
\newcommand{\cT}{\mathcal{T}}
\newcommand{\cX}{\mathcal{X}}

\newcommand{\CC}{\mathbb{C}}
\newcommand{\NN}{\mathbb{N}}
\newcommand{\ZZ}{\mathbb{Z}}
\newcommand{\QQ}{\mathbb{Q}}
\newcommand{\RR}{\mathbb{R}}

\theoremstyle{plain}
\newtheorem{thm}{Théorème}
\newtheorem{lem}{Lemme}
\newtheorem*{thm*}{Théorème}
\newtheorem*{lem*}{Lemme}
\newtheorem*{prop*}{Proposition}
\newtheorem*{cor*}{Corollaire}

\theoremstyle{definition}
\newtheorem*{rem*}{Remarque}
\newtheorem*{rems*}{Remarques}


% End literal retained input author-preamble.tex

\renewcommand{\proofname}{Démonstration}
\renewcommand{\refname}{Références}
\begin{document}
\begin{center}{\Large Stable item 2042}\end{center}
\paragraph{Source and attribution.}
Jean-Loup Waldspurger,\par\emph{Représentations et quasi-caractères de niveau $0$ ; endoscopie}, Journal de l’École polytechnique — Mathématiques 8 (2021), publisher version, DOI 10.5802/jep.145. \href{https://jep.centre-mersenne.org/articles/10.5802/jep.145/}{Publisher article}. Authors' text: \href{https://creativecommons.org/licenses/by/4.0/}{CC BY 4.0}.
\paragraph{Editorial problem boundary.}
Prove only the independent Proposition in Section 3.6 below, under the exact original large-residue-characteristic, facet fixed-point, compact Kottwitz-grade, cuspidal and Levi assumptions. The explicit orbital-distribution pairing, its vanishing alternatives and every Haar-volume normalisation are retained, including the conjugate on the normaliser-translated function. The entire new constant-term, root-group/segment and coset-measure core is retained. Endoscopic-transfer and later formal corollary claims are context only. Author prose remains in French.
\paragraph{Difficulty (editorial): H3.}
The definitions of parahoric cuspidal functions and Levi induction do not alone identify this orbital-distribution pairing or its exact volume factors. The author proves a constant-term calculation by facet segments and root-group cuspidal vanishing, then performs two nested coset decompositions with their Haar measures. Those geometric and measure comparisons produce the normaliser-indexed pairing, including its conjugation and the stated vanishing alternatives.
\paragraph{Editorial transcription note.}
Original author language, mathematics and ordering are preserved. Publisher front matter is replaced by this independent article wrapper. Bibliography is recovered from matching cached publisher HTML, including its TeX formulas. No new proof or mathematical repair is supplied. Surrounding original results are retained for conservative context and are not extra selected corpus claims.
\clearpage

% Begin literal retained input author-body.tex
\Changel

\vspace*{-1.2\baselineskip}\vskip0pt
\enlargethispage{.3\baselineskip}%
\section*{Introduction}

Soit $F$ un corps local non archimédien de caractéristique nulle et soit $G$ un groupe réductif connexe défini sur $F$. On s'intéresse ici aux représentations admissibles et irréductibles de $G(F)$ dans des espaces vectoriels complexes. On note $\Irr(G)$ l'ensemble des classes d'isomorphismes de ces représentations. Dans le cas où $G$ est un groupe classique, on sait associer à toute $\pi\in \Irr(G)$ un paramètre de Langlands, cela grâce aux résultats de Harris et Taylor, Henniart, Arthur et Mok. Dans le cas général, l'existence de ce paramètre reste conjecturale.
Moy et Prasad ont défini le niveau de $\pi$, \cf \cite[Th.\,5.2]{MP1}. Notons $\Irr(G)^0$ le sous-ensemble des représentations de niveau~$0$. Dans le cas où $G$ est adjoint, Lusztig a défini dans \cite{Lu} un paramétrage du sous-ensemble de $\Irr(G)^0$ formé des représentations de réduction unipotente, qui est un bon candidat pour être celui de Langlands. L'hypothèse que $G$ est adjoint a été récemment levée par Solleveld, \cf \cite{S}. On peut espérer que leurs méthodes permettent de paramétrer $\Irr(G)^0$ tout entier.
Hormis le cas des groupes classiques, savoir quelles conditions caractérisent le paramétrage de Langlands n'est pas clair, du moins pour l'auteur. Il~y~a en tout cas une condition minimale: le paramétrage doit satisfaire à des conditions de compatibilité à l'endoscopie. D'où une question préalable que l'on formule ici en termes vagues: le niveau~$0$ se conserve-t-il par endoscopie? Dans cet article, nous répondons positivement à cette question, pourvu que la caractéristique résiduelle $p$ de $F$ soit grande relativement à $G$ (ou que $G$ soit petit relativement à $p$, c'est une question de point de vue).

Pour tout ensemble $X$, notons $\CC[X]$ l'espace vectoriel complexe de base $X$. En~associant à toute représentation son caractère-distribution, on définit une injection $\Theta:\CC[\Irr(G)] \to \rI(G)^*$, où $\rI(G)^*$ est l'espace des distributions sur $G(F)$ invariantes par conjugaison. Définissons le projecteur de Bernstein $p^0:\CC[\Irr(G)]\to \CC[\Irr(G)^0]$ qui annule toute représentation irréductible qui n'est pas de niveau~$0$.

Supposons d'abord que $G$ est quasi-déployé. On sait définir le sous-espace $\SI(G)^*\subset \rI(G)^*$ des distributions stables. On note $\CC[\Irr(G)]^{\st}$ le sous-espace des $\pi\in \CC[\Irr(G)]$ telles que $\Theta_{\pi}\in \SI(G)^*$ (ici, $\pi$ n'est plus irréductible, c'est une combinaison linéaire à coefficients complexes de représentations irréductibles). Une conséquence de l'arti\-cle~\cite{A3} est qu'il existe une projection naturelle $p^{\st}:\CC[\Irr(G)]\to \CC[\Irr(G)]^{\st}$. Notre premier résultat est le

\begin{thm}\label{th:1}
Supposons que $G$ soit quasi-déployé et que l'hypothèse \ref{HypG} soit satisfaite. Alors on~a l'égalité $p^{\st}\circ p^0=p^0\circ p^{\st}$.
\end{thm}

\Cf le corollaire \ref{lecasquasideploye}. L'hypothèse \ref{HypG}, précisément énoncée en \ref{legroupeG}, est l'hypothèse sur $p$ évoquée plus haut. Grosso-modo, elle suppose $p\geq c(G)\val_{F}(p)$, où $c(G)$ est un entier dépendant de $G$ et $\val_{F}$ est la valuation usuelle de $F$.

Revenons au cas d'un groupe $G$ quelconque. Soit $\bG'$ une donnée endoscopique elliptique de $G$, \cf \ref{donneesendoscopiques}. Une telle donnée est un triplet dont l'un des termes est un groupe endoscopique $G'$, qui est quasi-déployé sur $F$. En général, on doit fixer des données auxiliaires pour définir un facteur de transfert. Pour simplifier l'introduction, supposons que cela ne soit pas nécessaire et fixons un facteur de transfert défini sur un sous-ensemble de $G'(F)\times G(F)$. Toujours en conséquence de \cite{A3}, il y a alors un homomorphisme de transfert spectral
\[
\transfert:\CC[\Irr(G')]^{\st}\to \CC[\Irr(G)].
\]

\begin{thm}\label{th:2}
Supposons que l'hypothèse \ref{HypendoG} soit satisfaite. Alors on~a l'égalité $p^0\circ \transfert=\transfert\circ p^0$.
\end{thm}

\Cf le théorème \ref{letheoreme}. L'hypothèse \ref{HypendoG} renforce \ref{HypG}, \cf \ref{Hypendo}, et implique que $(\Hyp)(G')$ est satisfaite. Évidemment, le premier $p^0$ vit sur $G$ et le second sur $G'$. Remarquons que ce dernier conserve l'espace $\CC[\Irr(G')]^{\st}$ d'après le premier théorème.

L'application $p^0$ est un projecteur de Bernstein et il s'en déduit un tel projecteur sur divers objets, par exemple l'espace $C_{\rc}^{\infty}(G(F))$ des fonctions à valeurs complexes sur $G(F)$, localement constantes et à support compact. Les théorèmes ci-dessus ont des contreparties pour les espaces de fonctions, \cf \ref{derechef} et \ref{transfertdefonctions}.

La méthode utilisée pour prouver ces théorèmes est de caractériser les représentations de niveau~$0$ par le développement local de leurs caractères. Notons $\frg$ l'algèbre de Lie de $G$. On sait définir l'exponentielle $\exp$ qui envoie un voisinage de $0$ dans $\frg(F)$ sur un voisinage de $1$ dans $G(F)$. Sous l'hypothèse \ref{HypG}, ces voisinages peuvent être choisis les plus gros possibles. À savoir l'ensemble $\frg_{\tn}(F)$ des éléments topologiquement nilpotents dans $\frg(F)$ et l'ensemble $G_{\tu}(F)$ des éléments topologiquement unipotents dans $G(F)$. Il en est de même pour tout sous-groupe réductif connexe de~$G$.

Soit $D\in \rI(G)^*$ une distribution invariante et localement intégrable sur
$G(F)$, donc associée à une fonction $\theta_{D}$ définie presque partout sur $G(F)$, invariante par conjugaison et localement intégrable. Nous dirons que $D$ est un quasi-caractère si et seulement si $\theta_{D}$ satisfait à la condition suivante. Soit $x\in G(F)$ un élément semi-simple, notons~$G_{x}$ la composante neutre de son commutant dans $G$. Alors il existe un voisinage $\frV_{x}$ de~$0$ dans $\frg_{x}(F)$ de sorte que la fonction $X\mto \theta_{D}(x\, \exp(X))$ définie presque partout sur~$\frV_{x}$ soit combinaison linéaire de transformées de Fourier d'intégrales orbitales nilpotentes sur $\frg_{x}(F)$. Harish-Chandra a démontré que le caractère de toute représentation irréductible était un quasi-caractère (d'où la terminologie \og quasi-caractère\fg), \cf \cite[Th.\,16.2]{HC}. Appelons $p'$-élément un élément $\epsilon\in G(F)$ qui est semi-simple et satisfait à la condition suivante. Fixons une extension finie $F'$ de~$F$ contenant toutes les valeurs propres de l'opérateur $\ad(\epsilon)$ agissant dans $\frg$ et notons $\Sigma$ l'ensemble de ces valeurs propres. Notons aussi $\vert {\cdot}\vert _{F'}$ la valeur absolue usuelle de~$F'$. Alors, pour tout $\sigma\in \Sigma$ telle que $\vert \sigma\vert _{F'}=1$, $\sigma$ est une racine de l'unité d'ordre premier à $p$. Soit $D\in \rI(G)^*$ une distribution invariante localement intégrable. Nous dirons que $D$ est un quasi-caractère de niveau~$0$ si et seulement si, pour tout $p'$-élément $\epsilon\in G(F)$, la fonction $X\mto \theta_{D}(\epsilon \exp(X))$ définie presque partout sur $\frg_{\epsilon,\tn}(F)$ est combinaison linéaire de transformées de Fourier d'intégrales orbitales nilpotentes sur $\frg_{\epsilon}(F)$. Autrement dit, le développement est valable sur le plus gros voisinage possible dans $\frg_{\epsilon}(F)$.\enlargethispage{\baselineskip}%

\begin{thm}\label{th:3}
Supposons l'hypothèse \ref{HypG} satisfaite. Soit \hbox{$\pi\!\in\!\CC[\Irr(G)]$}. Alors $\pi\in \CC[\Irr(G)^0]$ si et seulement si $\Theta_{\pi}$ est un quasi-caractère de niveau~$0$.
\end{thm}

\Cf \ref{caracterisation}. Le fait que le caractère d'une représentation de niveau~$0$ se développe sur de gros voisinages se trouve déjà dans la littérature, par exemple dans des articles de S.\,DeBacker, J.-L.\,Kim, F.\,Murnaghan, \cf \cite[Th.\,3.5.2]{D}, \cite[Th.\,5.3.1]{KM}. Mais je ne crois pas que l'on y trouve la réciproque.

Il s'avère que cette caractérisation des représentations de niveau~$0$ \og passe bien\fg à l'endoscopie et cela nous permet d'en déduire les théorèmes \ref{th:1} et \ref{th:2}.

Pour démontrer le théorème \ref{th:3}, on doit décrire l'espace des quasi-caractères de niveau~$0$. Pour que cette introduction reste d'une longueur raisonnable, donnons un simple exemple. Soit $\cF$ un sommet de l'immeuble de Bruhat-Tits du groupe adjoint de~$G$. On lui associe un groupe parahorique $K_{\cF}^0\subset G(F)$. Notons $K_{\cF}^+$ son plus grand sous-groupe distingué pro-$p$-unipotent. On sait qu'il existe un groupe réductif connexe~$\bG_{\cF}$ défini sur le corps résiduel $k_{F}$ de $F$, de sorte que $K_{\cF}^0/K_{\cF}^+\simeq \bG_{\cF}(k_{F})$, \cf \cite[\S3.4]{T}. Soit $f:\bG_{\cF}(k_{F})\to \CC$ une fonction invariante par conjugaison et cuspidale. Par l'isomorphisme précédent, on la considère comme une fonction sur $K_{\cF}^0$, invariante par translations par $K_{\cF}^+$, et on l'étend en une fonction sur $G(F)$, nulle hors de $K_{\cF}^0$. Notons $A_{G}$ le plus grand sous-tore central de $G$ qui soit déployé sur $F$ et fixons des mesures de Haar sur $G(F)$ et $A_{G}(F)$. On définit la distribution $D_{f}$ qui, à~une fonction $\varphi\in C_{\rc}^{\infty}(G(F))$, associe\vspace*{-3pt}
\[
D_{f}(\varphi)=\int_{A_{G}(F)\backslash G(F)}\int_{G(F)}\varphi(g^{-1}xg)f(x)\,dx\,dg.
\]
L'intégrale n'est pas absolument convergente mais converge dans l'ordre indiqué. On~montre que $D_{f}$ est un quasi-caractère de niveau~$0$, \cf \ref{unpremiertheoreme}. On généralise cette construction dans deux directions. D'une part en considérant non pas des groupes parahoriques mais des fixateurs de points de l'immeuble, qui donnent naissance à des groupes non connexes sur $k_{F}$. D'autre part, en induisant de telles distributions définies sur des groupes de Levi de $G$. On montre alors que tout quasi-caractère de niveau~$0$ est obtenu par cette construction convenablement généralisée, \cf \ref{undeuxiemetheoreme}. Ce qui relie les distributions ci-dessus aux représentations de niveau~$0$ est la formule calculant le carac\-tère d'une telle représentation que l'on~a obtenue dans un article antérieur, \cf \cite{W2}.

Signalons un résultat curieux. Les constructions ci-dessus conduisent naturellement à la définition d'un sous-espace de l'espace des quasi-caractères de niveau~$0$, noté $D^G[\cD_{\cusp}(G)]$ dans l'article. Ce sous-espace est relié à celui des caractères des représentations de niveau~$0$ qui sont elliptiques au sens d'Arthur, \cf \cite{A2}. Mais il ne lui est pas égal. On montre en \ref{stabilite} et \ref{finale} qu'il possède néanmoins les mêmes propriétés magiques de ces espaces de caractères, démontrées par Arthur dans \cite{A3}. À savoir que, dans le cas où $G$ est quasi-déployé, leur stabilité se lit sur les éléments elliptiques du groupe et que, dans le cas général, le transfert entre de tels quasi-caractères se lit lui aussi sur les éléments elliptiques des groupes en question.

L'hypothèse \ref{HypG} est utilisée de deux façons. D'une part, elle entraîne diverses propriétés concernant le groupe $G$. Par exemple, il existe une extension $F'$ de $F$ de degré fini et premier à $p$ telle que $G$ soit déployé sur $F'$. Ou bien, l'ordre du groupe des composantes connexes du centre de $G$ est premier à $p$. Pour ces propriétés, une hypothèse plus faible serait suffisante. D'autre part, comme on l'a déjà dit, l'hypothèse \ref{HypG} entraîne que l'exponentielle est définie sur les plus gros voisinages possibles. On peut envisager de remplacer l'exponentielle par un substitut, comme dans \cite[\S1.8]{KV}, ce qui permettrait d'améliorer grandement l'hypothèse sur $p$. Cela semble possible mais pas très facile et on~a préféré y renoncer. La certitude est qu'une restriction sur $p$ est nécessaire: les résultats seraient faux pour $G=\mathrm{SL}_{p}$.

Les liens entre niveau, paramétrage et endoscopie ont fait l'objet de plusieurs travaux récents. Citons \cite{Oi} et aussi \cite{Lanard1} et \cite{Lanard2}, où l'auteur étudie des sous-catégories de celle des représentations de niveau~$0$.

\subsubsection*{Remerciements}
Je remercie vivement les rapporteurs dont les remarques ont permis la correction d'un nombre à peine dénombrable d'erreurs qui se trouvaient dans la première version de l'article.

\section{Les données}

\subsection{Le corps local}

Soit $F$ un corps local non archimédien de caractéristique nulle. On note $\fro_{F}$ son anneau d'entiers, $\fro_{F}^{\times}$ le groupe des unités, $ \frp_{F}$ l'idéal maximal, $k_{F}=\fro_{F}/\frp_{F}$ le corps résiduel, $p$ la caractéristique de $k_{F}$, $\vert .\vert _{F}$ la valeur absolue usuelle, $\val_{F}$ la valuation.

Fixons une clôture algébrique $\bar{F}$ de $F$, \resp $\bar{k}_{F}$ de $k_{F}$. Notons $F^{\nr}$ la plus grande extension non ramifiée de $F$ contenue dans $\bar{F}$. On note $\Gamma_{F}$ le groupe de Galois de $\bar{F}/ F$ et $I_{F}$ son sous-groupe d'inertie, c'est-à-dire le groupe de Galois de $\bar{F}/F^{\nr}$.

\subsection{Notations diverses}

Quand un groupe abstrait $H$ agit sur un ensemble $X$, on~note $X^{H}$ le sous-ensemble des points fixes. Si $Y$ est un sous-ensemble de $X$, on~note $\Norm_{H}(Y)$ le sous-ensemble des éléments de $H$ dont l'action conserve $Y$.

Si $X$ est un ensemble, on~note $\CC[X]$ le $\CC$-espace vectoriel de base $X$.

\subsection{Groupes algébriques}\label{groupesalgebriques}

Posons $K=F$ ou $K=k_{F}$. Soit $H$ un groupe algébrique défini sur $K$. On note $H^0$ sa composante neutre et $Z(H)$ le centre de $H$. Les sous-groupes algébriques de $H$ que l'on considérera seront implicitement supposés définis sur $K$, sauf mention explicite du contraire.

Supposons $H$ réductif et connexe. On appelle Levi de $H$ une composante de Levi d'un sous-groupe parabolique de $G$ (tous deux définis sur $K$ comme on vient de le~dire). On note $H_{\AD}=H/Z(H)$ le groupe adjoint. Pour $x\in H$, on~note $x_{\ad}$ son image dans $H_{\AD}$, $x_{\rss}$ la partie semi-simple de $x$, $Z_{H}(x)$ le commutant de $x$ dans $H$ et $H_{x}=Z_{H}(x)^0$.

On utilise les notations d'Arthur concernant les Levi et paraboliques: si $M$ est un Levi, $\cL(M)$ est l'ensemble des Levi contenant $M$ et $\cP(M)$ est l'ensemble des sous-groupes paraboliques de $H$ de composante de Levi $M$. Si $P$ est un sous-groupe parabolique de $H$, on~note $U_{P}$ son radical unipotent.

Soit $T$ un tore défini sur $K$. On note $X^{*}(T)$, \resp $X_{*}(T)$, les groupes de caractères algébriques de $T$, \resp de sous-groupes à un paramètre. Pour ces définitions, $T$ est vu comme un tore sur la clôture algébrique $\bar{K}$, c'est-à-dire que les caractères ou cocaractères ne sont pas forcément définis sur $K$. Si $K=F$, on~note $T(F)_{\rc}$ le plus grand sous-groupe compact de $T(F)$. La définition de $X^*(T)$ s'étend au cas où $T$ est un groupe diagonalisable, c'est-à-dire un sous-groupe algébrique d'un tore.

Soit $H$ un groupe réductif connexe défini sur $K$. On note $A_{H}$ le plus grand sous-tore de $Z(H)$ qui soit défini et déployé sur $K$. On pose $a_{H}=\dim(A_{H})$ et $\cA_{H}=X_{*}(A_{H})\otimes_{\ZZ}\RR$. On appelle sous-tore déployé maximal de $H$ un sous-tore de $H$ qui est déployé sur $K$ et est maximal parmi les sous-tores déployés sur $K$. Ces sous-tores déployés maximaux sont en bijection avec les groupes de Levi minimaux de $H$: à un sous-tore $A$ est associé son commutant $Z_{H}(A)$.

Soit $H$ un groupe réductif connexe défini sur $k_{F}$. On appelle espace tordu sous $H$ une variété algébrique $\tilde{H}$ définie sur $k_{F}$, munie de deux actions algébriques à droite et à gauche de $H$ commutant entre elles et telles que, pour chacune des actions, $\tilde{H}$ soit un espace principal homogène sous $H$. On impose de plus ici que $\tilde{H}(k_{F})\neq\emptyset$. Cette terminologie a été introduite par Labesse. Une partie de la théorie habituelle des groupes réductifs s'étend aux espaces tordus. Ainsi, on définit les notions de sous-espace parabolique ou d'espace de Levi. On envoie pour tout cela au premier chapitre de \cite{MW}.

\subsection{Le groupe \texorpdfstring{$G$}{G}}\label{legroupeG}

On fixe pour tout l'article un groupe réductif connexe $G$ défini sur $F$. Soit $N\geq1$ un entier tel qu'il existe un plongement de $G$ dans $\mathrm{GL}(N)$.
On impose l'hypothèse
\begin{equation}\tag*{$\textup{(\Hyp)}(G)$}\label{HypG}
p\geq (2+\val_{F}(p))N.
\end{equation}

Cela
entraîne qu'il existe une extension $F'$ de $F$ de degré premier à $p$ telle que $G$ soit déployé sur $F'$. Plus généralement, pour tout tore défini sur $F$ de dimension inférieure ou égale au rang de $G$, il existe une telle extension telle que le tore soit déployé sur~$F'$. L'hypothèse \ref{HypG} implique $(\Hyp)(H)$ pour tout sous-groupe réductif connexe $H\subset G$ défini sur $F$.

On fixe un Levi minimal $M_{\min}$ de $G$. On pose simplement $A=A_{M_{\min}}$, $\cA=\cA_{M_{\min}}$, $\cL_{\min}=\cL(M_{\min})$. On définit $W^G=\Norm_{G(F)}(A)/M_{\min}(F)$.

Quand un objet a été défini relativement à notre groupe $G$, nous noterons souvent l'objet analogue défini relativement à un autre groupe $H$ en ajoutant un exposant~$H$ dans la notation. Par exemple, si $M\in \cL_{\min}$, on~note $\cL_{\min}^M$ le sous-ensemble des $L\in \cL_{\min}$ qui sont contenus dans $M$. Le groupe $W^G$ agit sur $\cL_{\min}$.
Pour $M\in \cL_{\min}$, on~note $N_{W^G}(M)$ le fixateur de $M$ dans $W^G$ et $W^G(M)=N_{W^G}(M)/W^M$.

On munit $G(F)$ d'une mesure de Haar. Plus généralement, pour tout sous-groupe fermé unimodulaire $J$ de $G(F)$, on suppose $J$ muni d'une mesure de Haar.

\section{L'immeuble de \texorpdfstring{$G$}{G}}

\subsection{Facettes, groupes parahoriques}\label{facettes}

On note $\Imm(G_{\AD})$ l'immeuble de Bruhat-Tits sur $F$ du groupe adjoint $G_{\AD}$. Cet ensemble est réunion d'appartements associés aux sous-tores déployés maximaux de $G$ ou encore aux Levi minimaux de $G$. On note $\App(A_{M})$ l'appartement associé à un tel Levi $M$. L'appartement $\App(A_{M})$ est un espace affine euclidien sous l'espace vectoriel réel $\cA_{M}/\cA_{G}$. Pour $x,y\in \App(A_{M})$, on~note $x-y$ l'élément de $\cA_{M}/\cA_{G}$ tel que $x=y+(x-y)$.

L'immeuble se décompose aussi en réunion disjointe de facettes, chaque facette étant contenue dans (au moins) un appartement. On note $\Fac(G)$ l'ensemble des facettes et $\Fac(G,A)$ le sous-ensemble des facettes contenues dans $\App(A)$. Le~grou\-pe~$G(F)$ agit sur l'immeuble. Pour $\cF\in \Fac(G)$, notons $K_{\cF}^{\dag}$ le stabilisateur de $\cF$ dans~$G(F)$.

Introduisons le groupe $\hat{G}$ dual de Langlands de $G$. Modulo le choix d'une paire de Borel épinglée de $\hat{G}$, ce groupe est muni d'une action de $\Gamma_{F}$. Le groupe $\Gamma_{F}/I_{F}$ agit sur $X^*(Z(\hat{G})^{I_{F}})$. Notons $\cN=X^{*}(Z(\hat{G})^{I_{F}})^{\Gamma_{F}/I_{F}}$ le sous-groupe des invariants. Kottwitz a défini un homomorphisme surjectif $w_{G}:G(F)\to \cN$, \cf \cite[\S7]{K}.

Soit $\cF\in \Fac(G)$. Notons $\cN(\cF)$ l'image de $K_{\cF}^{\dag}$ par $w_{G}$ et, pour tout $\nu\in \cN$, notons~$K_{\cF}^{\nu}$ l'ensemble des $g\in K_{\cF}^{\dag}$ tels que $w_{G}(g)=\nu$ (on~a donc $ K_{\cF}^{\nu}\neq\emptyset$ si et seulement si $\nu\in \cN(\cF)$). L'ensemble $K_{\cF}^{0}$ est le sous-groupe parahorique de $G(F)$ associé à $\cF$, \cf \cite[Prop.\,3]{HR}. Notons $K_{\cF}^{+}$ le plus grand sous-groupe distingué et pro-$p$-unipotent dans $K_{\cF}^0$. Bruhat et Tits ont associé à $\cF$ un schéma en groupes $\cG_{\cF}$ défini sur $\fro_{F}$. On a $\cG_{\cF}(\fro_{F})=K_{\cF}^{0}$. Notons $\bG_{\cF}$ la partie réductive de la fibre spéciale de $\cG_{\cF}$. C'est un groupe réductif connexe défini sur $k_{F}$ et $K_{\cF}^0/K_{\cF}^+$ est isomorphe à $\bG_{\cF}(k_{F})$, \cf \cite[Prop.\,3.7]{HV}. Pour $\nu\in \cN(\cF)$, il existe un espace tordu $\bG_{\cF}^{\nu}$ sous $\bG_{\cF}$ de sorte que $K_{\cF}^{\nu}/K_{\cF}^{+}$ s'identifie à $\bG_{\cF}^{\nu}(k_{F})$. Faute de référence, indiquons brièvement la construction de cet espace tordu. Rappelons la propriété suivante. Pour $g\in G(F)$, notons $\Ad(g)$ la conjugaison par $g$. Soit $\cF'$ l'image de $\cF$ par l'action de $g$ sur $\Imm(G_{\AD})$. Alors $\Ad(g)(K_{\cF}^0)=K_{\cF'}^0$ et l'application $\Ad(g)$ se descend en un isomorphisme algébrique $\overline{\Ad(g)}:\bG_{\cF}\to \bG_{\cF'}$. Fixons
un élément $x\in K_{\cF}^{\nu}$. En appliquant ce qui précède, on dispose d'un automorphisme algébrique $\overline{\Ad(x)}$ de $\bG_{\cF}$. On définit $\bG_{\cF}^{\nu}$ comme la variété algébrique $\bG_{\cF}$ munie des deux actions de $\bG_{\cF}$:
\begin{align*}
\bG_{\cF}\times \bG_{\cF}^{\nu}\times \bG_{\cF}&\to\bG_{\cF}^{\nu}\\(a,y,b)&\mto ay\,\overline{\Ad(x)}(b).
\end{align*}
Pour $g\in K_{\cF}^{\nu}$, on écrit $g=hx$ avec $h\in K_{\cF}^0$. Alors l'image de $g$ dans $\bG_{\cF}^{\nu}(k_{F})$ est égale à celle de $h$ dans $\bG_{\cF}(k_{F})=\bG_{\cF}^{\nu}(k_{F})$. À isomorphisme près, cette construction ne dépend pas du choix de l'élément $x$.

Le groupe $K_{\cF}^0$ fixe tout point de $\cF$. Pour $\nu\in \cN(\cF)$, il existe une permutation isométrique $\sigma_{\cF,\nu}$ de $\cF$ telle que tout élément de $K_{\cF}^{\nu}$ agisse dans $\cF$ par cette permutation. On note $\cF^{\nu}$ le sous-ensemble des points fixes de $\sigma_{\cF,\nu}$ dans $\cF$. Pour tout appartement contenant $\cF$, $\cF^{\nu}$ est l'intersection de $\cF$ avec un sous-espace affine de cet appartement. Soit $x\in \cF$. Alors le fixateur de $x$ dans $G(F)$ est la réunion des $K_{\cF}^{\nu}$ sur les $\nu\in \cN(\cF)$ tels que $\sigma_{\cF,\nu}$ fixe $x$.

On note $\Fac^*(G)$ l'ensemble des couples $(\cF,\nu)$, où $\cF\in \Fac(G)$ et $\nu\in \cN(\cF)$. On note $\Fac^*_{\max}(G)$ le sous-ensemble des $(\cF,\nu)$ tels que $\cF^{\nu}$ est réduit à un point. On note $\Fac^*(G,A)$ et $\Fac^*_{\max}(G,A)$ les sous-ensembles des $(\cF,\nu)$ tels que $\cF\subset \App(A)$.

\begin{rem*}
L'indice $\max$ est contestable puisqu'il s'agit de couples $(\cF,\nu)$ tels que~$\cF^{\nu}$ est de dimension minimale. Ce sont plutôt les groupes qui leur sont attachés qui sont maximaux. En tout cas, on conserve cette notation de \cite{W2} pour simplifier les références.
\end{rem*}

Du schéma en groupes $\cG_{\cF}$ se déduit une sous-$\fro_{F}$-algèbre de Lie $\frk_{\cF}$ de $\frg_{F}$. On note~$\frk_{\cF}^+$ son radical pro-$p$-nilpotent. Le quotient $\frk_{\cF}/\frk_{\cF}^+$ s'identifie à l'espace des points sur $k_{F}$ de la partie réductive de l'algèbre de Lie de $\cG_{\cF}$.

On peut aussi considérer l'immeuble étendu $\Imm(G)=\Imm(G_{\AD})\times \cA_{G}$, qui est muni d'une action de $G(F)$. Pour $x=(y,a)\in \Imm(G)$, avec $y\in \Imm(G_{\AD})$ et $a\in \cA_{G}$, le fixateur de $x$ dans $G(F)$ est le sous-groupe des éléments $g\in G(F)$ qui fixent $y$ et tels que $w_{G}(g)$ appartienne au sous-groupe de torsion $\cN_{\tors}$ de $\cN$.

\subsection{Description des facettes}\label{description}

Pour $a\in A(F)$, on~note $a_{\ZZ}$ l'élément de $ X_{*}(A)$ tel que $\<x^*,a_{\ZZ}\>=-\val_{F}(x^*(a))$ pour tout $x^*\in X^*(A)$. Le noyau de cet homomorphisme $a\mto a_{\ZZ}$ n'est autre que le plus grand sous-groupe compact $A(F)_{\rc}$ de $A(F)$. L'action sur l'immeuble du groupe $A(F)$ conserve $\App(A)$. Pour $a\in A(F)$ et $x\in \App(A)$, on~a $ax= x+a_{\ZZ}$, où ici $X_{*}(A)$ est vu comme un sous-groupe de $\cA$. L'action du groupe $\Norm_{G}(A)(F)$ conserve aussi l'appartement $\App(A)$. Notons $M_{\min}(F)_{\rc}$ l'unique sous-groupe compact maximal de $M_{\min}(F)$. L'action de $\Norm_{G}(A)(F)$ se quotiente en une action du groupe $\Norm_{G}(A)(F)/M_{\min}(F)_{\rc}$.

Notons $\Sigma$ l'ensemble des racines réduites de $A$ dans $G$. Fixons un sommet spécial de $\App(A)$ que l'on~note $0$ qui nous permet d'identifier $\App(A)$ et $\cA/\cA_{G}$. À tout $\alpha\in \Sigma$ est associé un sous-ensemble $\Gamma_{\alpha}$ de $\QQ$, qui est l'image réciproque dans $\QQ$ d'un sous-ensemble fini de $\QQ/\ZZ$. Pour $c\in \Gamma_{\alpha}$, on~note $c^+$ le plus petit élément de $\Gamma_{\alpha}$ strictement supérieur à $c$ et $c^-$ le plus grand élément de $\Gamma_{\alpha}$ strictement inférieur à $c$. On note $H_{\alpha,c}$ l'hyperplan affine de $\App(A)$ défini par l'équation $\alpha(x)=c$. Alors la décomposition en facettes de $\App(A)$ est définie par la famille d'hyperplans $(H_{\alpha,c})_{\alpha\in \Sigma, c\in \Gamma_{\alpha}}$. À toute facette $\cF\in \Fac(G,A)$ sont associés un sous-ensemble $\Sigma_{\cF}\subset \Sigma$ et, pour tout $\alpha\in \Sigma$, un élément $c_{\alpha,\cF}\in \Gamma_{\alpha}$ de sorte que $\cF$ soit le sous-ensemble des éléments $x\in \App(A)$ qui satisfont aux relations
\begin{enumerate}
\item\label{description1}
$\alpha(x)=c_{\alpha,\cF}$ pour tout $\alpha\in \Sigma_{\cF}$;

\item\label{description2}
$c_{\alpha,\cF}<\alpha(x)<c_{\alpha,\cF}^+$ pour tout $\alpha\in \Sigma-\Sigma_{\cF}$.
\end{enumerate}

Pour $\alpha\in \Sigma$, notons $U_{\alpha}$ le groupe radiciel associé à $\alpha$. À tout $c\in \Gamma_{\alpha}$ est associé un sous-groupe ouvert compact $U_{\alpha,c}$ de $U_{\alpha}(F)$ de sorte que les propriétés suivantes soient satisfaites:

\begin{enumerate}\setcounter{enumi}{2}
\item\label{description3}
si $c,c'\in \Gamma_{\alpha}$ et $c<c'$, alors $U_{\alpha,c}\subsetneq U_{\alpha,c'}$;

\item\label{description4}
quelle que soit $\cF\in \Fac(G,A)$, on~a $U_{\alpha}(F)\cap K_{\cF}^0=U_{\alpha,c_{\alpha,\cF}}$; on~a $U_{\alpha}(F)\cap K_{\cF}^+=U_{\alpha,c_{\alpha,\cF}^-}$ si $\alpha\in \Sigma_{\cF}$ et $U_{\alpha}(F)\cap K_{\cF}^+=U_{\alpha,c_{\alpha,\cF}}$ si $\alpha\not \in \Sigma_{\cF}$.
\end{enumerate}

Notons $M_{\min}(F)^{\star}=\{m\in M_{\min}(F)_{\rc}\sep w_{G}(m)=0\}$. C'est un sous-groupe d'indice fini de $M_{\min}(F)_{\rc}$. Pour toute facette $\cF\in \Fac(G,A)$, $K_{\cF}^0$ est le sous-groupe de $G(F)$ engendré par les $U_{\alpha}(F)\cap K_{\cF}^0$ et par le groupe $M_{\min}(F)^{\star}$. Les $\fro_{F}$-algèbres~$\frk_{\cF}$ et~$\frk_{\cF}^+$ admettent une description similaire. Il en résulte que
les groupes $K_{\cF}^0$ et les $\fro_{F}$-algèbres~$\frk_{\cF}$ et~$\frk_{\cF}^+$ caractérisent les facettes. C'est-à-dire
\begin{enumerate}\setcounter{enumi}{4}
\item\label{description5}
soient $\cF,\cF'\in \Fac(G)$; supposons $K_{\cF}^0=K_{\cF'}^0$, ou $\frk_{\cF}=\frk_{\cF'}$ ou $\frk^+_{\cF}=\frk^+_{\cF'}$; alors $\cF=\cF'$.
\end{enumerate}

Pour $\cF\in \Fac(G,A)$, il existe un unique Levi $M_{\cF}\in \cL_{\min}$ tel que $\cA_{M_{\cF}}/\cA_{G}$ soit l'espace vectoriel réel engendré par les $x-y$ pour $x,y\in \cF$. Pour $(\cF,\nu)\in \Fac^*(G,A)$, il existe un unique Levi $M_{\cF,\nu}\in \cL_{\min}$ tel que $\cA_{M_{\cF},\nu}/\cA_{G}$ soit l'espace vectoriel réel engendré par les $x-y$ pour $x,y\in \cF^{\nu}$. On a évidemment $M_{\cF}\subset M_{\cF,\nu}$.

\subsection{Groupes de Levi}\label{groupesdelevi}
Soit $M\!\!\in\!\cL_{\min}$. Considérons le sous-ensemble $\Imm(G_{\AD},M)$ de $\Imm(G_{\AD})$ réunion des appartements $\App(A_{M'})$ associés aux Levi minimaux $M'$ contenus dans $M$. Le groupe $\cA_{M}$ agit sur chacun de ces appartements et ces actions se recollent en une action sur la réunion. De plus, $\Imm(G_{\AD},M)$ est conservé par l'action de $M(F)$. Le quotient de $\Imm(G_{\AD},M)$ par l'action de $\cA_{M}$, muni de son action de $M(F)$, s'identifie canoniquement à l'immeuble $\Imm(M_{\AD})$. En particulier, $\App(A)/\cA_{M}$ s'identifie à l'appartement $\App^M(A)$ associé à $M_{\min}$ dans l'immeuble $\Imm(M_{\AD})$. On note $p_{M}:\App(A)\to \App^M(A)$ cette projection.

L'action sur l'immeuble du groupe $\Norm_{G}(M)(F)$ conserve $\Imm(G_{\AD},M)$ et se descend en une action sur $\Imm(M_{\AD})$. Cette action permute les éléments de $\Fac(M)$. Pour $\cF_{M}\in \Fac(M)$, on~note $K_{\cF_{M}}^{\dag,G}$ le sous-groupe des éléments de $\Norm_{G}(M)(F)$ qui conservent $\cF_{M}$. Il contient $K_{\cF_{M}}^{\dag}$ comme sous-groupe distingué.

Notons $\Norm_{G}(M,A)$ le normalisateur commun de $M$ et $A$ dans $G$. L'action du groupe $\Norm_{G}(M,A)(F)$ dans $\Imm(M_{\AD})$ conserve $\App^M(A)$.

Pour $\cF\in \Fac(G,A)$, il existe une unique facette notée $\cF^{M}\in \Fac(M,A)$ telle que $p_{M}(\cF)\subset \cF^{M}$. Si $\cF$ est décrit par les relations \eqref{description1} et \eqref{description2} de \ref{description}, $\cF^{M}$ est l'ensemble des $x\in \App^M(A)$ qui satisfont aux relations
\begin{enumerate}
\item\label{groupesdelevi1}
$\alpha(x)=c_{\alpha,\cF}$ pour tout $\alpha\in \Sigma_{\cF}\cap \Sigma^M$;

\item\label{groupesdelevi2}
$c_{\alpha,\cF}<\alpha(x)<c_{\alpha,\cF}^+$ pour tout $\alpha\in \Sigma^M-\Sigma_{\cF}\cap \Sigma^M$.
\end{enumerate}

De l'inclusion $Z(\hat{G})\to Z(\hat{M})$ se déduit un homomorphisme naturel $\cN^M\to \cN$. Le diagramme
\[
\xymatrix{
M(F)\ar[r]\ar[d]_{w_{M}}&G(F)\ar[d]^{w_{G}}\\
\cN^M\ar[r]&\cN
}
\]
est commutatif. Posons $M_{\ad}=M/Z(G)$. On a aussi un homomorphisme naturel $\cN^M\to \cN^{M_{\ad}}$. Notons $\cN_{G\tcomp}^M$ l'image réciproque dans $\cN^M$ du sous-groupe de torsion de $\cN^{M_{\ad}}$. Alors l'homomorphisme $\cN^M\to \cN$ se restreint en un homomorphisme injectif $\cN^M_{G\tcomp}\to \cN$, \cf \cite[6(1)]{W2}. On identifie $\cN^M_{G\tcomp}$ à son image dans $\cN$. On note $\Fac^*_{G\tcomp}(M)$ le sous-ensemble des $(\cF_{M},\nu)\in \Fac^*(M)$ tels que $\nu\in \cN^M_{G\tcomp}$ (avec les variantes $\Fac^*_{\max,G\tcomp}(M)$ etc.). Remarquons que, pour $(\cF_{M},\nu)\in \Fac^*_{G\tcomp}(M)$ et pour $n\in K_{\cF_{M}}^{\dag,G}$, la conjugaison par $n$ conserve l'ensemble $K_{\cF_{M}}^{\nu}$.

\section{Espaces de fonctions et de distributions}

\subsection{Les espaces \texorpdfstring{$\rI(G)$ et $\rI(G)^*$}{GetG}}\label{lesespaces}

Le groupe $G(F)$ agit sur l'espace $C_{\rc}^{\infty}(G(F))$ par conjugaison: pour $g\in G(F)$ et $f\in C_{\rc}^{\infty}(G(F))$, $\gf$ est la fonction $x\mto f(g^{-1}xg)$. On note $\rI(G)$ le quotient de $C_{\rc}^{\infty}( G(F))$ par le sous-espace complexe engendré par les $\gf-f$ pour $f\in C_{\rc}^{\infty}(G(F))$ et $g\in G(F)$.

On note $G_{\reg}$ le sous-ensemble des éléments fortement réguliers de $G$. Soit $f\in C_{\rc}^{\infty}(G(F))$. Pour $x\in G_{\reg}(F)$, on définit l'intégrale orbitale
\[
I^G(x,f)=D^G(x)^{1/2}\int_{A_{G_{x}}(F)\backslash G(F)}f(g^{-1}xg)\,dg,
\]
où $D^G$ est le discriminant de Weyl usuel. L'espace $\rI(G)$ est aussi le quotient de $C_{\rc}^{\infty}(G(F))$ par le sous-espace des $f\in C_{\rc}^{\infty}(G(F))$ telles que $I^G(x,f)=0$ pour tout $x\in G_{\reg}(F)$.

On appelle distribution invariante sur $G(F)$ une forme linéaire sur $C_{\rc}^{\infty}(G(F))$ qui se quotiente en une forme linéaire sur $\rI(G)$. On identifie une telle distribution à la forme linéaire quotient. On note $\rI(G)^*$ l'espace des distributions invariantes sur $G(F)$.

Soient $M$ un Levi de $G$ et $D^M\in \rI(M)^*$. On définit une distribution $\Ind_{M}^G(D^M)\in \rI(G)^*$ de la façon suivante. On fixe un sous-groupe parabolique $P\in \cP(M)$. Comme on l'a dit, des mesures de Haar sont fixées sur $M(F)$, $U_{P}(F)$ et $G(F)$. Il s'en déduit ce que l'on peut appeler une pseudo-mesure invariante à droite sur $P(F)\backslash G(F)$. Précisément, c'est une forme linéaire non pas sur $C_{\rc}^{\infty}(P(F)\backslash G(F))$ mais sur l'espace des fonctions localement constantes $\varphi:G(F)\to \CC$ qui satisfont à la relation $\varphi(mug)=\delta_{P}(m)\varphi(g)$ pour tous $m\in M(F)$, $u\in U_{P}(F)$ et $g\in G(F)$, où $\delta_{P}$ est le module usuel. Cette pseudo-mesure est caractérisée par l'égalité
\[
\int_{G(F)}f(g)\,dg\,=\int_{P(F)\backslash G(F)}\int_{M(F)}\int_{U_{P}(F)}f(mug)\,du\,dm\,dg
\]
pour tout $f\in C_{\rc}^{\infty}(G(F))$.
On définit la fonction $f_{U_{P}}$ sur $M(F)$ par
\[
f_{U_{P}}(m)=\delta_{P}(m)^{1/2}\int_{U_{P}(F)}f(mu)\,du,
\]
puis la distribution $\Ind_{M}^G(D^M)$ par l'égalité
\[
\Ind_{M}^G(D^M)(f)=\int_{P(F)\backslash G(F)}D^M(({^g}f)_{U_{P}})\,dg.
\]
Cela ne dépend pas du choix de $P$.

On dit que $x\in G_{\reg}(F)$ est elliptique si le commutant $T$ de $x$ dans $G$ est un tore elliptique modulo $Z(G)$, c'est-à-dire si $A_{T}=A_{G}$. On note $G_{\rell}(F)$ le sous-ensemble des éléments elliptiques dans $G_{\reg}(F)$ (la notation étant un peu abusive: il n'y a pas de sous-ensemble algébrique $G_{\rell}$ de $G$). On dit qu'une fonction $f\in C_{\rc}^{\infty}(G(F))$ est cuspidale si et seulement si les intégrales orbitales de $f$ sont nulles en tout point $x\in G_{\reg}(F)-G_{\rell}(F)$. Autrement dit si $f$ est annulée par $\Ind_{M}^G(D^M)$ pour tout Levi propre $M$ et tout $D^M\in \rI(M)^*$. On note $C_{\cusp}(G(F))$ l'espace des fonctions cuspidales et $\rI_{\cusp}(G)$ son image dans $\rI(G)$.

On dit que $f$ est très cuspidale si et seulement si, pour tout sous-groupe parabolique propre $P$ de $G$, de composante de Levi $M$, la fonction $f_{U_{P}}$ est nulle. En fait, $\rI_{\cusp}(G)$ est aussi l'image dans $\rI(G)$ de l'espace des fonctions très cuspidales, \cf \cite[Lem.\,2.7]{W3}. Pour une fonction $f$ très cuspidale, on définit une distribution $D_{f}\in \rI(G)^*$ par l'égalité
\[
D_{f}(\varphi)=\int_{A_{G}(F)\backslash G(F)}\int_{G(F)}\varphi(x^{-1}gx)f(g)\,dg\,dx.
\]
La double intégrale n'est pas absolument convergente mais converge dans cet ordre, \cf \cite[Lem.\,9]{W2}.

\subsection{Filtrations}\label{filtrations}

On fixe un ensemble de représentants $\underline{\cL}_{\min}\subset \cL_{\min}$ des classes de conjugaison de Levi de $G$. Pour tout entier $n\in \ZZ$, notons $\underline{\cL}_{\min}^n$ le sous-ensemble des $M\in \underline{\cL}_{\min}$ tels que $ a_{M}=n$ (on peut évidemment se limiter aux $n$ appartenant à l'intervalle $\{a_{G},\dots,a_{M_{\min}}\}$).

On définit une filtration sur $\rI(G)$ de la façon suivante. Pour tout $n \in \ZZ$, notons $\Fil^n\rI(G)$ l'image dans $\rI(G)$ du sous-espace des $f\in C_{\rc}^{\infty}(G(F))$ qui satisfont à la condition: pour tout Levi $M$ tel que $a_{M}>n$ et tout $m\in G_{\reg}(F)\cap M(F)$, on~a $I^G(m,f)=0$. On a
\[
\Fil^{a_{G}-1}\rI(G)=\{0\}\subset \Fil^{a_{G}}\rI(G)=\rI_{\cusp}(G)\subset\cdots\subset \Fil^{a_{M_{\min}}}\rI(G)=\rI(G),
\]
et, en posant $\Gr^n\rI(G)=\Fil^n\rI(G)/\Fil^{n-1}\rI(G)$, on a
\[
\Gr^n\rI(G)\simeq \bigoplus_{M\in \underline{\cL}_{\min}^n}\rI_{\cusp}(M)^{ W^G(M)},
\]
\cf \ref{legroupeG} pour la définition du groupe $W^G(M)$, qui agit naturellement sur $\rI_{\cusp}(M)$.

Notons $\Ann^n\rI(G)^*$ l'annulateur de $\Fil^{n-1}\rI(G)$ dans $\rI(G)^*$. On a
\[
\Ann^{a_{M_{\min}}+1}\rI(G)^*=\{0\}\subset \Ann^{a_{M_{\min}}}\rI(G)^*\subset\cdots\subset \Ann^{a_{G}}\rI(G)^*=\rI(G)^*,
\]
et, en posant $\Gr^n\rI(G)^*=\Ann^n\rI(G)^*/\Ann^{n+1}\rI(G)^*$, on a
\begin{equation}\tag{1}\label{filtrations1}
\Gr^n\rI(G)^*\simeq \bigoplus_{M\in \underline{\cL}_{\min}^n}\rI_{\cusp}(M)^{* W^G(M)}.
\end{equation}
On vérifie que $\Ann^n\rI(G)^*$ est le sous-espace de $\rI(G)^*$ engendré par les distributions induites $\Ind_{M}^G(\rI(M)^*)$ pour les Levi $M$ de $G$ tels que $a_{M}\geq n$. Pour $M\in \underline{\cL}_{\min}^n$ et $d^M\in \rI(M)^*$, l'image de $\Ind_{M}^G(d^M)$ dans $\Gr^n\rI(G)^*$ est nulle dans les composantes $\rI_{\cusp}(M')^{* W^G(M')}$ de \eqref{filtrations1} pour $M'\neq M$ et est l'image naturelle de $d^M$ dans $\rI_{\cusp}(M)^{* W^G(M)}$ (c'est-à-dire la restriction de $d^{M}$ à $I_{\cusp}(M)^{ W^G(M)}$).

\subsection{Éléments compacts et topologiquement unipotents}\label{elementscompacts}
Soit $x\in G(F)$. On dit que $x$ est compact si et seulement s'il est contenu dans un sous-groupe compact de $G(F)$, c'est-à-dire si l'adhérence $\overline{x^{\ZZ}}$ du groupe $x^{\ZZ}$ engendré par $x$ est compacte. On dit qu'il est compact $\rmod Z(G)$ si et seulement si l'image $x_{\ad}$ de $x$ dans $G_{\AD}(F)$ est compacte. On dit que $x$ est topologiquement unipotent si et seulement si $\lim_{n\to \infty}x^{p^n}=1$

Fixons un sous-tore maximal $T$ de $G$ contenant la partie semi-simple $x_{\rss}$ de $x$ et fixons une extension finie $F'$ de $F$ telle que $T$ soit déployé sur $F'$. Alors:
\begin{itemize}
\item
$x$ est compact si et seulement si $ \chi(x_{\rss})\in \fro^{\times}_{F'}$ pour tout $\chi\in X^*(T)$; $x$ est compact $\rmod Z(G)$ si et seulement si $ \chi(x_{\rss})\in \fro^{\times}_{F'}$ pour tout $\chi\in X^*(T)$ tel que la restriction de $\chi$ à $Z(G)$ soit triviale; $x$ est topologiquement unipotent si et seulement si $\chi(x_{\rss})\in 1+\frp_{F'}$ pour tout $\chi\in X^*(T)$.
\end{itemize}

D'autre part, $x$ est compact $\rmod Z(G)$ si et seulement s'il existe $\cF\in \Fac(G)$ tel que $x\in K_{\cF}^{\dag}$. L'ensemble des éléments compacts, \resp compacts $\rmod Z(G)$, est un sous-ensemble ouvert et fermé de $G(F)$ invariant par conjugaison. Soit $M$ un Levi de~$G$, soit $(\cF_{M},\nu)\in \Fac^*(M)$ et soit $x\in K_{\cF_{M}}^{\nu}$. Alors $x$ est compact $\rmod Z(G)$ si et seulement si $\nu\in \cN^M_{G\tcomp}$.

On note $G_{\comp}(F)$ l'ensemble des éléments de $G(F)$ qui sont compacts $\rmod Z(G)$ (il serait plus correct de le noter $G_{\comp\rmod Z(G)}$...). On note $G_{\tu}(F)$ l'ensemble des éléments topologiquement unipotents de $G(F)$. On note $\rI(G)^*_{\comp}$ l'ensemble des distributions dont le support est contenu dans $G_{\comp}(F)$.

\subsection{Facettes, fonctions et distributions}\label{facettesfonctions}
Soit $(\cF,\nu)\in \Fac_{\max}^*(G)$. On note $C_{\cusp}(\bG^{\nu}_{\cF})$ l'espace des fonctions sur $\bG^{\nu}_{\cF}(k_{F})$ qui sont invariantes par conjugaison par $\bG_{\cF}(k_{F})$ et qui sont cuspidales. On munit cet espace du produit hermitien défini par
\[
\<f,f'\>=\vert \bG_{\cF}(k_{F})\vert ^{-1}\sum_{x\in \bG^{\nu}_{\cF}(k_{F})}\bar{f}(x)f'(x).
\]
Il est défini positif. Soit $f\in C_{\cusp}(\bG^{\nu}_{\cF})$. L'espace $\bG_{\cF}^{\nu}(k_{F})$ s'identifie à $K_{\cF}^{\nu}/K_{\cF}^+$. On~identifie $f$ à une fonction sur ce quotient $K_{\cF}^{\nu}/K_{\cF}^+$, on la relève en une fonction sur~$K_{\cF}^{\nu}$ puis on l'étend en une fonction sur $G(F)$ par $0$ hors de $K_{\cF}^{\nu}$. On note~$f_{\cF}$ la fonction ainsi obtenue. Elle est très cuspidale, \cf \cite[Lem.\,10]{W2}. On déduit de $f_{\cF}$ une distribution $D_{f_{\cF}}$, \cf \ref{lesespaces}, que l'on~note simplement $D_{f}$. Le support de cette distribution est contenu dans $G_{\comp}(F)$.

Notons $\boldsymbol{\cD}_{\cusp}(G)$ le sous-ensemble des familles
\[
(f_{\cF,\nu})_{(\cF,\nu)\in \Fac^*_{\max}(G)}\in \prod_{(\cF,\nu)\in \Fac^*_{\max}(G)}C_{\cusp}(\bG^{\nu}_{\cF})
\]
qui satisfont à la condition: pour tout $\nu\in \cN$, l'ensemble des $(\cF,\nu')\in \Fac^*_{\max}(G)$ tels que $\nu'=\nu$ et $f_{\cF,\nu'}\neq0$ est fini.

Pour $\bmf=(f_{\cF,\nu})_{(\cF,\nu)\in \Fac^*_{\max}(G)}\in \boldsymbol{\cD}_{\cusp}(G)$, la somme $\sum_{\cF,\nu}D_{f_{\cF,\nu}}$ est définie. \hbox{En~effet}, chaque $D_{f_{\cF,\nu}}$ est à support dans l'ouvert fermé $w_{G}^{-1}(\nu)$ et, pour tout $\nu$, il n'y a qu'un nombre fini de $(\cF,\nu')$ tels que $\nu'=\nu$ et $f_{\cF,\nu'}\neq0$. Pour tout compact~$C$ de~$G(F)$, il n'y a donc qu'un nombre fini de $(\cF,\nu)$ tels que le support de $D_{f_{\cF,\nu}}$ coupe~$C$. On pose $D_{\bmf}=\sum_{\cF,\nu}D_{f_{\cF,\nu}}$. Cela définit une application linéaire
\[
D: \boldsymbol{\cD}_{\cusp}(G)\to \rI(G)_{\comp}^*.
\]

\begin{rem*}
Quand un élément de $\boldsymbol{\cD}_{\cusp}(G)$ sera noté $\bmf$, on~notera son image $D_{\bmf}$ ou $D^G_{\bmf}$ comme ci-dessus. Quand l'élément sera noté plus symboliquement $d$, on~notera son image $D[d]$ ou $D^G[d]$.
\end{rem*}

Pour $\nu\in \cN$, notons $\boldsymbol{\cD}^{\nu}_{\cusp}(G)$ le sous-espace des
\[
\bmf=(f_{\cF,\nu'})_{(\cF,\nu')\in \Fac^*_{\max}(G)}\in \boldsymbol{\cD}_{\cusp}(G)
\]
tels que $f_{\cF,\nu'}=0$ pour $\nu'\neq\nu$. On a $\boldsymbol{\cD}_{\cusp}=\prod_{\nu\in \cN}\boldsymbol{\cD}^{\nu}_{\cusp}(G)$.
Le groupe $G(F)$ agit naturellement sur $\boldsymbol{\cD}_{\cusp}(G)$ en respectant les sous-espaces $\boldsymbol{\cD}_{\cusp}^{\nu}(G)$. Il est clair que, pour $\bmf\in \boldsymbol{\cD}_{\cusp}(G)$ et $g\in G(F)$, on~a $D_{{\gbmf}-\bmf}=0$. Pour $\nu\in \cN$, notons $\cD^{\nu}_{\cusp}(G)$ l'espace des coinvariants, c'est-à-dire le quotient de $\boldsymbol{\cD}^{\nu}_{\cusp}(G)$ par le sous-espace engendré par les ${\gbmf}-\bmf$ pour $\bmf\in \boldsymbol{\cD}^{\nu}_{\cusp}(G)$ et $g\in G(F)$. On pose $\cD_{\cusp}(G)=\prod_{\nu\in \cN}\cD^{\nu}_{\cusp}(G)$. Il y a un homomorphisme naturel $\boldsymbol{\cD}_{\cusp}(G)\to\cD_{\cusp}(G)$.

\begin{rem*}
Si $\cN$ est fini, $\cD_{\cusp}(G)$ est le quotient des coinvariants de $\boldsymbol{\cD}_{\cusp}(G)$. Si $\cN$ est infini, c'est un quotient de cet espace de coinvariants.
\end{rem*}

L'application $D$ se quotiente en une application linéaire
\[
D:\cD_{\cusp}(G)\to \rI(G)_{\comp}^*.
\]

Notons $\cD'(G)$ la somme directe des $\cD_{\cusp}(M)$ où $M$ parcourt les Levi de $G$. On définit une application linéaire
\[
D^G: \cD'(G)\to \rI(G)^*
\]
qui, à $\sum_{M}d^M$, associe $\sum_{M}\Ind_{M}^G(D^M[d^M])$.
Le groupe $G(F)$ agit naturellement sur l'espace $\cD'(G)$. En effet, un élément $g\in G(F)$ envoie un Levi $M$ sur le Levi $gMg^{-1}$, une facette $(\cF_{M},\nu)\in \Fac^*_{\max}(M)$ sur une facette $(g\cF_{M},g\nu)\in \Fac^*_{\max}(gMg^{-1})$ et une fonction $f_{\cF_{M},\nu}\in C_{\cusp}(\bM_{\cF_{M}}^{\nu})$ sur une fonction $^g(f_{\cF_{M},\nu})\in C_{\cusp}(({g\bM g^{-1}})_{g\cF_{M}}^{g\nu})$. Alors $g$ envoie $\boldsymbol{\cD}_{\cusp}(M)$ sur $\boldsymbol{\cD}_{\cusp}(gMg^{-1})$ et cette application se quotiente en une application de $\cD_{\cusp}(M)$ sur $\cD_{\cusp}(gMg^{-1})$.
Il est clair que, pour $g\in G(F)$ et~\hbox{$\bmf\in \cD'(G)$}, on~a $D^G_{{\gbmf}-\bmf}=0$. On note $\cD(G)$ le quotient des coinvariants pour cette action de $G(F)$ dans $\cD'(G)$. L'application précédente se quotiente en une application linéaire
\[
D^G:\cD(G)\to \rI(G)^*.
\]

\subsection{Variantes de \texorpdfstring{$\cD(G)$}{DG}}\label{variantes}
Soit $M$ un Levi de $G$. On note $\boldsymbol{\cD}_{\cusp,G\tcomp}(M)$ le sous-espace des éléments
\[
\bmf=(f_{\cF_{M},\nu})_{(\cF_{M},\nu)\in \Fac^*_{\max}(M)}
\]
de $ \boldsymbol{\cD}_{\cusp}(M)$ tels que $f_{\cF_{M},\nu}=0$ si $\nu\not\in \cN^M_{G\tcomp}$. On note $\cD_{\cusp,G\tcomp}(M)$ son image dans $\cD_{\cusp}(M)$. On note $\cD'_{G\tcomp}(G)$ la somme directe des $\cD_{\cusp,G\tcomp}(M)$ où $M$ parcourt les Levi de $G$. On définit comme dans la section précédente le quotient des coinvariants $\cD_{G\tcomp}(G)$, qui s'identifie à un sous-espace de $\cD(G)$. L'application $D^G$ se restreint en une application linéaire
\[
D^G:\cD_{G\tcomp}(G)\to \rI(G)^*_{\comp}.
\]

Il y a une projection naturelle $\boldsymbol{\cD}_{\cusp}(M)\to \boldsymbol{\cD}_{\cusp,G\tcomp}(M)$ qui, à un élément
\[
(f_{\cF_{M},\nu})_{(\cF_{M},\nu)\in \Fac^*_{\max}(M)}\in \boldsymbol{\cD}_{\cusp}(M)
\]
associe l'élément
\[
(f_{\cF_{M},\nu})_{(\cF_{M},\nu)\in \Fac^*_{\max,G\tcomp}(M)}
\]
de $ \boldsymbol{\cD}_{\cusp,G\tcomp}(M)$. Cette projection se quotiente en une projection de $\cD_{\cusp}(M)$ sur $\cD_{\cusp,G\tcomp}(M)$.

Fixons $\nu\in \cN$. Notons $\rI(G)^{*\nu}_{\comp}$ le sous-espace des distributions invariantes à support contenu dans $G_{\comp}(F)\cap w_{G}^{-1}(\nu)$. Pour un Levi $M$ de $G$, notons $\boldsymbol{\cD}^{\nu}_{\cusp,G\tcomp}(M)$ le sous-espace des
\[
\bmf=(f_{\cF_{M},\nu'})_{(\cF_{M},\nu')\in \Fac^*_{\max}(M)}\in \boldsymbol{\cD}_{\cusp,G\tcomp}(M)
\]
tels que $f_{\cF_{M},\nu'}=0$ si $\nu'\neq\nu$. À l'aide de ces espaces, on définit comme ci-dessus un espace $\cD^{\nu}_{G\tcomp}(G)$. On a des isomorphismes naturels
\begin{align*}
\cD_{G\tcomp}(G)&\simeq \prod_{\nu\in \cN}\cD_{G\tcomp}^{\nu}(G),
\\
\rI(G)^*_{\comp}&\simeq \prod_{\nu\in \cN}\rI(G)^{*\nu}_{\comp},
\end{align*}
et l'application $D^G$ s'identifie au produit de ses restrictions
\[
D^G:\cD_{G\tcomp}^{\nu}(G)\to \rI(G)^{*\nu}_{G\tcomp}.
\]

\subsection{Calcul de \texorpdfstring{$D_{f'}^G(f_{\cF})$}{DGf}}\label{calcul}

Soit $(\cF,\nu)\!\in\!\Fac^*(G,A)$. On a défini le Levi \hbox{$M_{\cF,\nu}\!\in\!\cL_{\min}$} au paragraphe \ref{description}. Posons simplement $M=M_{\cF,\nu}$. On associe à $\cF$ la facette $\cF^{M}\in \Fac(M,A)$. L'élément $\nu$ appartient à $\cN_{G\tcomp}^M(\cF^{M})$, $(\cF^M,\nu)$ appartient à $\Fac^*_{\max}(M,A)$, $\bG_{\cF}$ s'identifie à $\bM_{\cF^M}$ et $\bG_{\cF}^{\nu}$ s'identifie à $\bM_{\cF^M}^{\nu}$, \cf \cite[Lem.\,6]{W2}. Soit $f\in C_{\cusp}(\bM^{\nu}_{\cF^M})$. On identifie $f$ à une fonction sur $\bG_{\cF}^{\nu}(k_{F})$, puis à une fonction sur $K_{\cF}^{\nu}/K_{\cF}^+$, que l'on relève en une fonction sur $K_{\cF}^{\nu}$. On prolonge cette fonction en une fonction sur $G(F)$, nulle hors de $K_{\cF}^{\nu}$. On note $f_{\cF}$ la fonction sur $G(F)$ obtenue ainsi. Elle est par construction invariante par conjugaison par $K_{\cF}^0$.

Soient $L\in \cL_{\min}$, $(\cF',\nu')\in \Fac^*_{\max,G\tcomp}(L,A)$ et $f'\in C_{\cusp}(\bL_{\cF'}^{\nu'})$. On a défini la distribution $D^G_{f'}=\Ind_{L}^G(D^L_{f'})$ sur $G(F)$. Notons $N(M,\cF^{M},L,\cF')$ l'ensemble des $n\in \Norm_{G}(A)(F)$ tels que $nMn^{-1}=L$ et $n\cF^{M}=\cF'$. Cet ensemble (qui peut être vide) est invariant à gauche par $A_{L}(F)(K_{\cF'}^0\cap \Norm_{G}(A)(F))$. On fixe un ensemble de représentants $\underline{N}(M,\cF^{M},L,\cF')$ du quotient
\[
A_{L}(F)(K_{\cF'}^0\cap \Norm_{G}(A)(F))\backslash N(M,\cF^{M},L,\cF').
\]

\skpt
\begin{prop*}
\begin{enumeratei}
\item\label{calculi}
$D^G_{f'}(f_{\cF})\neq0$ seulement si $\nu=\nu'$ et $N(M,\cF^{M},L,\cF')\neq\emptyset$.

\item\label{calculii}
Supposons $\nu=\nu'$ et $N(M,\cF^{M},L,\cF')\neq\emptyset$. Alors
\[
D^G_{f'}(f_{\cF})= \mes(K_{\cF}^0)\mes(K_{\cF'}^0) \mes(A_{L}(F)_{\rc})^{-1}\sum_{n\in \underline{N}(M,\cF^{M},L,\cF')}\<\overline{\nf},f'\>.
\]
\end{enumeratei}
\end{prop*}

\skpt
\begin{rems*}
\begin{enumerate}
\item\label{calcul1}
La fonction $\nf$ est celle définie en \ref{facettesfonctions}.

\item\label{calcul2}
Supposons $M=L$, $\cF^M=\cF'$ et $\nu=\nu'$. Si $f'$ est invariante par $K_{\cF'}^{\dag,G}$, la formule du \eqref{calculii} se simplifie et devient
\[
D_{f'}^G(f_{\cF})=\mes(K_{\cF}^0)\mes(K_{\cF'}^0)\mes(A_{L}(F)_{\rc})^{-1}[K_{\cF'}^{\dag,G}:A_{L}(F)K_{\cF'}^0]\<\bar{f},f'\>.
\]
\end{enumerate}
\end{rems*}

\begin{proof}
Puisque $f_{\cF}$ est à support dans $w_{G}^{-1}(\nu)$ et que $D^G_{f'}$ est à support dans $w_{G}^{-1}(\nu')$ il est clair que $D^G_{f'}(f_{\cF})=0$ si $\nu\neq\nu'$. Supposons désormais $\nu=\nu'$.

Fixons un sous-groupe parabolique $Q\in \cP(L)$. Commençons par calculer $f_{\cF,U_{Q}}$. On~sait définir la facette $\cF^L\subset \Imm(L_{\AD})$. Si $M\subset L$, on~a évidemment $(\cF^L)^M=\cF^M$ et la fonction $f_{\cF^L}$ sur $L(F)$ est bien définie.

\begin{lem}\label{lem1}
On a les égalités
\[
f_{\cF,U_{Q}}=\begin{cases}0,&\text{si } M\not\subset L,\\
\mes(U_{Q}(F)\cap K_{\cF}^+)f_{\cF^L}&\text{si } M\subset L. \end{cases}
\]
\end{lem}

\begin{proof}
Rappelons que, pour $l\in L(F)$, on~a $f_{\cF,U_{Q}}(l)=\int_{U_{Q}(F)}f_{\cF}(lu)\,du$. Le groupe $Q$, \resp $L$, détermine un sous-groupe parabolique $\bQ$ de $\bG_{\cF}$, \resp une composante de Levi $\bL$ de $\bQ$, de sorte que l'image de $Q(F)\cap K_{\cF}^0$, \resp $L(F)\cap K_{\cF}^0$, dans~$\bG_{\cF}(k_{F})$ soit $\bQ(k_{F})$, \resp $\bL(k_{F})$. Supposons que $f_{\cF,U_{Q}}\!\neq\!0$. Alors $Q(F)\cap K_{\cF}^{\nu}\!\neq\!\emptyset$ et il existe un sous-espace parabolique $\bQ^{\nu}$ de $\bG^{\nu}_{\cF}$, associé au parabolique $\bQ$, de sorte que l'image dans $\bG^{\nu}_{\cF}(k_{F})$ de $Q(F)\cap K_{\cF}^{\nu}$ soit $\bQ^{\nu}(k_{F})$. Notons $\bL^{\nu}$ le normalisateur de~$\bL$ dans~$\bQ^{\nu}$. Posons $V=U_{Q}(F)\cap K_{\cF}^0$. Fixons $l\in L(F)$ tel que $f_{\cF,U_{Q}}(l)\neq0$. On~écrit
\[
f_{\cF,U_{Q}}(l)=\int_{U_{Q}(F)/V}\int_{V}f_{\cF}(luv)\,dv\,du.
\]
Fixons $u\in U_{Q}(F)$ tel que l'intégrale intérieure soit non nulle. L'élément $lu$ appartient à $Q(F)\cap K_{\cF}^{\nu}$. Décomposons son image dans $\bQ^{\nu}(k_{F})$ en $\bl\bmu$ avec $\bl\in \bL^{\nu}(k_{F})$ et $\bmu\in \bU_{\bQ}(k_{F})$. On voit alors que
\[
\int_{V}f_{\cF}(luv)\,dv=c\sum_{\bmv\in \bU_{\bQ}(k_{F})}f(\bl\bmv),
\]
où $c>0$ est une constante provenant des mesures et où on~a identifié $f$ à une fonction sur $\bG^{\nu}_{\cF}(k_{F})$. Cette fonction est cuspidale. La somme ci-dessus est donc nulle si $\bQ^{\nu}\neq\bG_{\cF}^{\nu}$. Notre hypothèse de non-nullité implique donc $\bQ^{\nu}=\bG_{\cF}^{\nu}$, d'où aussi $\bQ=\bL=\bG_{\cF}$. Le tore $A$ détermine un sous-tore déployé maximal $\bA$ de $\bG_{\cF}$. On a $X_{*}(A)\simeq X_{*}(\bA)$. Les sous-tores $A_{M_{\cF}}$ et $A_{L}$ déterminent des sous-tores $\bA_{M_{\cF}}$ et $\bA_{L}$. Par définition de $M_{\cF}$, $\bA_{M_{\cF}}$ est le plus grand sous-tore déployé contenu dans le centre de $\bG_{\cF}$. Par définition de $\bL$, ce groupe est le commutant de $\bA_{L}$ dans $\bG_{\cF}$. L'égalité $\bL=\bG_{\cF}$ entraîne donc $\bA_{L}\subset \bA_{M_\cF}$. D'où aussi $A_{L}\subset A_{M_{\cF}}$. Puisque $\bQ^{\nu}=\bG_{\cF}^{\nu}$, on~a $K_{\cF}^{\nu}=(Q(F)\cap K_{\cF}^{\nu})K_{\cF}^+$. On a aussi
\[
K_{\cF}^+=(K_{\cF}^+\cap Q(F))(K_{\cF}^+\cap U_{\bar{Q}}(F)),
\]
où $\bar{Q}$ est le parabolique de composante de Levi $L$ opposé à $Q$. Donc
\[
K_{\cF}^{\nu}=(Q(F)\cap K_{\cF}^{\nu})(K_{\cF}^+\cap U_{\bar{Q}}(F)).
\]
On sait que l'ensemble $K_{\cF}^{\nu}\cap \Norm_{G}(A)(F)$ est non vide, \cf \cite[7.4.4]{BT}. Mais un élément de $\Norm_{G}(A)(F)$ qui appartient à $Q(F)U_{\bar{Q}}(F)$ appartient forcément à $\Norm_{L}(A)(F)$. Soit donc $w\in K_{\cF}^{\nu}\cap \Norm_{L}(A)(F)$. Il agit naturellement dans $A$ et conserve $A_{M_{\cF}}$. Par définition de $M=M_{\cF,\nu}$, $A_{M}$ est le plus grand sous-tore de $A_{M_{\cF}}$ contenu dans l'ensemble des points fixes de l'action de $w$. Puisque $w\in L(F)$, ce sous-tore contient~$A_{L}$. Donc $A_{L}\subset A_{M}$ et $M\subset L$. Cela démontre la première assertion du lemme.

Supposons maintenant $M\subset L$. Alors, pour $l\in L(F)$ et $u\in U_{Q}(F)$, on~a $lu\in K_{\cF}^{\nu}$ si et seulement si $l\in K_{\cF^L}^{\nu}$ et $u\in U_{Q}(F)\cap K_{\cF}^+$ (cela résulte de \cite[Lem.\,6]{W2}). On voit que la fonction $f_{\cF,U_{Q}}$ est à support dans $K_{\cF^L}^{\nu}$ et que, pour un élément $l$ de ce groupe, on~a $f_{\cF,U_{Q}}(l)=\mes(U_{Q}(F)\cap K_{\cF}^+)f(\bar{l})$, où $\bar{l}$ est l'image de $l$ dans $\bL_{\cF^L}^{\nu}(k_{F})\simeq \bM_{\cF^M}^{\nu}(k_{F})$. Par définition de la fonction $f_{\cF^L}$, on obtient la deuxième assertion du lemme, ce qui achève sa preuve.
\end{proof}

Supposons $M\subset L$. Posons
\[
I_{\cF^L,\cF'}(f,f')=\int_{L(F)}f_{\cF^L}(l)f'_{\cF'}(l)\, dl.
\]

\skpt
\begin{lem}\label{lem2}
\begin{enumeratei}
\item\label{lem2i}
Si
$\cF^L\neq\cF'$, alors $I_{\cF^L,\cF'}(f,f')=0$.

\item\label{lem2ii}
Si $\cF^L=\cF'$, alors $M=L$, $\cF^M=\cF'$ et $I_{\cF^L,\cF'}(f,f')=\mes(K_{\cF'}^0)\<\bar{f},f'\>$.
\end{enumeratei}
\end{lem}

\begin{proof}
Ici, tout se passe dans $L$, on peut aussi bien supposer $L=G$ pour simplifier les notations. On a donc $\cF^L=\cF$. Supposons $\cF\neq\cF'$. Dans l'appartement $\App(A)$, fixons
un segment $[x,x']$ joignant un point $x\in \cF^{\nu}$ à un point $x'\in \cF^{\prime\nu}$. L'hypothèse $\cF\neq\cF'$ entraîne que $x\neq x'$. Il y a une facette $\cF''\in \Fac(G,A)$ et un élément $x''\in{}]x,x'[$ tel que le segment $[x'',x'[$ soit contenu dans $\cF''$. D'après ce que sont les supports de $f_{\cF}$ et $f'_{\cF'}$, on a
\[
I_{\cF,\cF'}(f,f')=\int_{K_{\cF}^{\nu}\cap K_{\cF'}^{\nu}}f_{\cF}(g)f'_{\cF'}(g)\, dg.
\]
Si $K_{\cF}^{\nu}\cap K_{\cF'}^{\nu}=\emptyset$, cette intégrale est nulle et on~a démontré \eqref{lem2i}. Sinon, considérons un élément $g\in K_{\cF}^{\nu}\cap K_{\cF'}^{\nu}$. L'action de $g$ sur l'immeuble fixe $x$ et $x'$, donc aussi tout le segment $[x,x']$. En particulier, il fixe $[x'',x'[$, donc $g\in K_{\cF''}^{\dag}$. On a $w_{G}(g)=\nu$, donc $g\in K_{\cF''}^{\nu}$. Alors $[x'',x'[\subset \cF^{\prime\prime\nu}$. Cela entraîne $\cF''\neq\cF'$ car l'hypothèse
\[
(\cF',\nu)\in \Fac^*_{\max}(G,A)
\]
(rappelons que l'on~a supposé $L=G$) signifie que $\cF^{\prime\nu}$ est réduit à un unique point, qui est donc $x'$. Le point $x'$ est adhérent à $\cF''$, donc toute la facette $\cF'$ est contenue dans l'adhérence de $\cF''$. Cela entraîne $K_{\cF''}^0\subset K_{\cF'}^0$. En général, il n'y a pas de relation d'inclusion entre les ensembles $K_{\cF''}^{\nu}$ et $K_{\cF'}^{\nu}$ mais, si l'on suppose que l'intersection de ces deux ensembles n'est pas vide , alors $K_{\cF''}^{\nu}\subset K_{\cF'}^{\nu}$. Ici, l'intersection des deux ensembles contient $g$.
On sait qu'alors il existe un sous-groupe parabolique propre~$\bP$ de $\bG_{\cF'}$ et un sous-espace parabolique $\bP^{\nu}$ de $\bG_{\cF'}^{\nu}$ associé à $\bP$, de sorte que l'image de~$K_{\cF''}^{0}$ dans $\bG_{\cF'}(k_{F})$ soit $\bP(k_{F})$ et que l'image de $K_{\cF''}^{\nu} $ dans $\bG_{\cF'}^{\nu}(k_{F})$ soit $\bP^{\nu}(k_{F})$. Il est facile d'identifier l'ensemble $\bU_{\bP}(k_{F})$ en décrivant les facettes $\cF'$ et $\cF''$ comme en \ref{description}. On a forcément
\[
\Sigma_{\cF''}\subset \Sigma_{\cF'},\quad c_{\alpha,\cF''}=c_{\alpha,\cF'}\text{ si }\alpha\!\in\!\Sigma_{\cF''}\quadet c_{\alpha,\cF''}=c_{\alpha,\cF'}\text{ ou }c_{\alpha,\cF'}^-\text{ si }\alpha\!\in\!\Sigma_{\cF'}-\Sigma_{\cF''}.
\]
Notons $\Xi$ l'ensemble des $\alpha\in \Sigma_{\cF'}-\Sigma_{\cF''}$ tels que $c_{\alpha,\cF''}=c_{\alpha,\cF'}$. Alors $\bU_{\bP}(k_{F})$ est l'image dans $\bG_{\cF'}(k_{F})$ de $\prod_{\alpha\in \Xi}U_{\alpha,c_{\alpha,\cF'}}$. Notons $V$ ce dernier groupe. Il est inclus dans $K_{\cF'}^0$. Montrons qu'il est inclus dans $K_{\cF}^+$. En effet, soit $\alpha\in \Xi$. Les définitions entraînent que $\alpha(x'')> \alpha(x')=c_{\alpha,\cF'}$. A fortiori, $\alpha(x)> \alpha(x'')> c_{\alpha,\cF'}$. Si $\alpha\in \Sigma_{\cF}$, on~a $c_{\alpha,\cF}=\alpha(x)> c_{\alpha,\cF'}$ donc
\[
c_{\alpha,\cF}^-\geq c_{\alpha,\cF'}\quadet K_{\cF}^+\cap U_{\alpha}(F)=U_{\alpha,c_{\alpha,\cF}^-}\supset U_{\alpha,c_{\alpha,\cF'}}.
\]
Si $\alpha\not\in \Sigma_{\cF}$, $c_{\alpha,\cF}$ est le plus grand élément de $\Gamma_{\alpha}$ qui soit strictement inférieur à $\alpha(x)$, donc $c_{\alpha,\cF}\geq c_{\alpha,\cF'}$. On a encore $K_{\cF}^+\cap U_{\alpha}(F)=U_{\alpha,c_{\alpha,\cF}}\supset U_{\alpha,c_{\alpha,\cF'}}$. Cela démontre l'assertion. On peut alors écrire
\[
I_{\cF,\cF'}(f,f')=\int_{(K_{\cF}^{\nu}\cap K_{\cF'}^{\nu})/V}\int_{V}f_{\cF}(gv)f'_{\cF'}(gv)\, dv\, dg.
\]
On vient de voir que $V\subset K_{\cF}^+$ et $f_{\cF}$ est invariante par ce groupe donc l'intégrale se récrit
\[
I_{\cF,\cF'}(f,f')=\int_{(K_{\cF}^{\nu}\cap K_{\cF'}^{\nu})/V}f_{\cF}(g)\int_{V}f'_{\cF'}(gv)\, dv\, dg.
\]
Pour $g$ intervenant dans cette intégrale, notons $\bar{g}$ son image dans $\bG^{\nu}_{\cF'}(k_{F})$. On a en fait $\bar{g}\in \bP^{\nu}(k_{F})$. À une constante positive près provenant des mesures, l'intégrale intérieure n'est autre que
\[
\sum_{\bar{u}\in \bU_{\bP}(k_{F})}f'(\bar{g}\bar{ u}).
\]
Ceci est nul car $f'$ est cuspidale. Donc $I_{\cF,\cF'}(f,f')=0$, ce qui démontre le~\eqref{lem2i} de l'énoncé.

Supposons maintenant $\cF=\cF'$. Alors $M=M_{\cF,\nu}=M_{\cF',\nu}=G$ puisque $(\cF',\nu)\in \Fac^*_{\max}(G,A)$. D'où $\cF^M=\cF=\cF'$. Un calcul immédiat donne
la formule~\eqref{lem2ii} de l'énoncé, ce qui démontre celui-ci.
\end{proof}

Après ces préliminaires, démontrons la proposition. Notons $N$ un ensemble de représentants du quotient
\[
\Norm_{L}(A)(F)\backslash \Norm_{G}(A)(F)/(\Norm_{G}(A)(F)\cap K_{\cF}^0).
\]
Fixons un sous-groupe parabolique $Q \in \cP(L)$. On définit comme en \ref{lesespaces} une pseudo-mesure sur $Q(F)\backslash G(F)$. On sait que $G(F)$ est union disjointe des ensembles $Q(F)nK_{\cF}^0$ quand $n$ décrit $N$. Pour tout $n\in N$, introduisons la fonction $\varphi_{n}$ sur $G(F)$ à support dans $Q(F)nK_{\cF}^0$ telle que
\[
\varphi_{n}(lunk)=\delta_{Q}(l)\quad\text{pour tous } l\in L(F),\ u\in U_{Q}(F),\ k\in K_{\cF}^0.
\]
On note $m_{n}$ la valeur de son intégrale contre la pseudo-mesure sur $Q(F)\backslash G(F)$.
Puisque $f_{\cF}$ est invariante par $K_{\cF}^0$, la définition de $D^G_{f'}(f_{\cF})$ se récrit
\begin{equation}\tag{3}\label{calcul3}
D^G_{f'}(f_{\cF})= \sum_{n\in N}m_{n}D_{f'}^L((^n(f_{\cF}))_{U_{Q}}).
\end{equation}
Pour $n\in N$, l'action de $n$ transporte $\cF$ en une facette $n\cF$, le Levi $M=M_{\cF,\nu}$ en le Levi $nMn^{-1}=M_{n\cF,\nu}$, la facette $\cF^M$ en une facette $n\cF^M\in \Fac(nMn^{-1},A)$ et~$f$ en une fonction $\nf\in C_{\cusp}(({n\bM n^{-1}})_{n\cF^M}^{\nu})$. On a $^n(f_{\cF})=(\nf)_{n\cF}$. On applique le lemme \ref{lem1} en y remplaçant $\cF$ par $n\cF$ et $f_{\cF}$ par $(\nf)_{n\cF}$. Cette assertion nous dit que la contribution $D_{f'}^L((\nf_{\cF})_{U_{Q}})$ de $n$ à la formule~\eqref{calcul3} est nulle si $nMn^{-1}\not\subset L$. On note $N^0$ le sous-ensemble des $n\in N$ tels que $nMn^{-1}\subset L$. Pour $n\in N^0$, cette contribution est $\mes(U_{Q}(F)\cap K_{n\cF}^+)D_{f'}^L((\nf)_{(n\cF)^L})$. Par définition, on a
\begin{equation}\tag{4}\label{calcul4}
D_{f'}^L((\nf)_{(n\cF)^L})=\int_{A_{L}(F)\backslash L(F)} \int_{L(F)}(\nf)_{(n\cF)^L}(x^{-1}lx)f'_{\cF'}(l)\,dl\,dx.
\end{equation}
Notons $ N'_{n}$ un ensemble de représentants du quotient
\[
A_{L}(F)(K_{\cF'}^0\cap \Norm_{L}(A)(F))\backslash \Norm_{L}(A)(F)/(K_{(n\cF)^L}^0\cap \Norm_{L}(A)(F)).
\]
On utilise maintenant la décomposition en union disjointe
\[
L(F)=\bigsqcup_{n'\in N'_{n}} A_{L}(F)K_{\cF'}^0n'K_{(n\cF)^L}^0.
\]
Pour $n'\in N'_{n}$, notons $m'_{n'}$ la mesure de
\[
A_{L}(F)\backslash A_{L}(F)K_{\cF'}^0n'K_{(n\cF)^L}^0.
\]
Puisque les fonctions $(\nf)_{(n\cF)^L}$, \resp $f'_{\cF'}$, sont invariantes par conjugaison par $K_{(n\cF)^L}^0$, \resp $K_{\cF'}^0$, la formule~\eqref{calcul4} se récrit
\[
D_{f'}^L((\nf)_{(n\cF)^L})=\sum_{n'\in N'_{n}}m'_{n'}\int_{L(F)}{^{n'}(\nf)_{(n\cF)^L}}(l)f'_{\cF'}(l)\,dl.
\]
Pour $n'$ apparaissant ci-dessus, on~a $^{n'}(\nf)_{(n\cF)^L}=(^{n'n}f)_{(n'n\cF)^L}$ avec des définitions similaires aux précédentes. Remarquons que $M_{n'n\cF,\nu}=n'nM(n'n)^{-1}$. La formule ci-dessus se récrit
\[
D_{f'}^L((\nf)_{(n\cF)^L})=\sum_{n'\in N'_{n}}m'_{n'}I_{(n'n\cF)^L,\cF'}(^{n'n}f,f').
\]
Pour tout $n'$, on applique le lemme \ref{lem2} en y remplaçant $\cF$ par $n'n\cF$ et $f$ par $^{n'n}f$. Le~\eqref{lem2i} de ce lemme nous permet de nous limiter aux $n'$ tels que
\[
n'nM(n'n)^{-1}=L,\quad (n'n\cF)^L=n'n\cF^M=\cF'.
\]
Notons $N_{n}^{\prime0}$ l'ensemble des $n'\in N'_{n}$ satisfaisant à ces conditions. En utilisant le~\eqref{lem2ii} du lemme \ref{lem2}, on obtient
\[
D_{f'}^L((\nf)_{(n\cF)^L})=\sum_{n'\in N^{\prime0}_{n}}m'_{n'}\mes(K_{\cF'}^0)\<\overline{^{n'n}f},f'\>.
\]
En rassemblant ces calculs, on obtient
\begin{equation}\tag{5}\label{calcul5}
D^G_{f'}(f_{\cF})= \sum_{n\in N^0}\sum_{n'\in N_{n}^{\prime0}}m_{n}\mes(U_{Q}(F)\cap K_{n\cF}^+) m'_{n'}\mes(K_{\cF'}^0)\<\overline{^{n'n}f},f'\>.
\end{equation}
On vérifie que de l'application
\begin{align*}
\{(n,n')\sep n\in N^0,\, n'\in N_{n}^{\prime0}\}&\to \Norm_{G}(A)(F)\\ (n,n')&\mto n'n
\end{align*}
se déduit une bijection de l'ensemble de départ sur le quotient
\[
A_{L}(F)(K_{\cF'}^0\cap \Norm_{G}(A)(F))\backslash N(M,\cF^{M},L,\cF') .
\]
Si $N(M,\cF^{M},L,\cF')$ est vide, on~a donc $D^G_{f'}(f_{\cF})=0$, ce qui démontre le~\eqref{calculi} de la proposition. Supposons $N(M,\cF^M,L,\cF')\neq\emptyset$. On peut supposer que $\underline{N}(M,\cF^M,L,\cF')$ est égal à l'ensemble des $n'n$ pour $(n,n')$ comme ci-dessus.
Pour un tel couple, on~vérifie les propriétés suivantes
\begin{align*}
&
\begin{aligned}
K_{n\cF}^0\cap Q(F)&=(K_{n\cF}^0\cap L(F))(K_{n\cF}^0\cap U_{Q}(F)),\\
K_{n\cF}^0\cap L(F)&=K_{(n\cF)^L}^0,\\
K_{n\cF}^0\cap U_{Q}(F)&=K_{n\cF}^+\cap U_{Q}(F),
\end{aligned}
\qquad\text{\cf \cite[Lem.\,6]{W2};}
\\
&m_{n}=\mes(K_{\cF}^0)\mes(K_{(n\cF)^L}^0)^{-1}\mes(K_{n\cF}^+\cap U_{Q}(F))^{-1};
\\
&\mes(K_{(n\cF)^L}^0)=\mes(K_{\cF'}^0);
\\
&m'_{n'}=\mes(A_{L}(F)\backslash A_{L}(F)K_{\cF'}^0)=\mes(K_{\cF'}^0)\mes(A_{L}(F)_{\rc})^{-1}.
\end{align*}
En utilisant ces propriétés, la formule~\eqref{calcul5} devient celle du \eqref{calculii} de la proposition.
\end{proof}

\subsection{Relèvement de fonctions à \texorpdfstring{$G(F)$}{GF}}\label{relevement}
Pour $(\cF,\nu)\in \Fac^*(G)$, notons $C(\bG_{\cF}^{\nu})$ l'espace des fonctions à valeurs complexes sur $\bG_{\cF}^{\nu}(k_{F})$ et notons $\cE_{\cF}^{\nu}$ l'espace des fonctions sur $G(F)$, à support dans $K_{\cF}^{\nu}$ et invariantes par multiplication à droite ou à gauche par $K_{\cF}^+$. Ces espaces s'identifient comme en \ref{facettesfonctions}: pour une fonction $f\in C(\bG_{\cF}^{\nu})$, on identifie $f$ à une fonction sur $K_{\cF}^{\nu}/K_{\cF}^+$, on la relève en une fonction sur $K_{\cF}^{\nu}$ puis on l'étend en une fonction sur $G(F)$ par $0$ hors de $K_{\cF}^{\nu}$. On note $f_{\cF}$ la fonction obtenue, qui appartient à $\cE_{\cF}^{\nu}$.

Soit $M$ un Levi de $G$ et soit $(\cF_{M},\nu)\!\in\!\Fac^*_{G\tcomp}(M)$. Fixons une facette $\cF\!\in\!\Fac(G)$ telle que
\[
(\cF,\nu)\in \Fac^*(G),\quad M_{\cF,\nu}=M_{\cF_{M},\nu}\quadet \cF^M=\cF_{M}.
\]
C'est possible d'après \cite[Lem.\,6\,\&\,7]{W2}. On a $\bM_{\cF_{M}}^{\nu}=\bG_{\cF}^{\nu}$ donc $C(\bM_{\cF_{M}}^{\nu})=C(\bG_{\cF}^{\nu})$. Ainsi, pour $f\in C(\bM_{\cF_{M}}^{\nu})$, on définit la fonction $f_{\cF}\in \cE_{\cF}^{\nu}$.

\begin{lem*}
L'image de $f_{\cF}$ dans $\rI(G)$ ne dépend pas du choix de $\cF$.\end{lem*}

\begin{proof}
Soit $\cF'\in \Fac(G)$ satisfaisant aux mêmes conditions que $\cF$. On doit prouver que les images dans $\rI(G)$ de $f_{\cF}$ et $f_{\cF'}$ sont égales. On peut supposer $\cF\neq\cF'$ sinon c'est évident. On ne perd rien à supposer que $(\cF_{M},\nu)\in \Fac^*_{G\tcomp}(M,A)$. Notons $X$ l'image réciproque de $\cF^{\nu}_{M}$ dans $\Imm(G_{\AD})$. C'est un sous-ensemble de $\App(A)$. Les ensembles $\cF^{\nu}$ et $\cF^{\prime\nu}$ sont contenus dans $X$ et sont des ouverts de cet ensemble, \cf \cite[Lem.\,7]{W2}. Fixons des points $x\in \cF^{\nu}$ et $x'\in \cF^{\prime\nu}$. Le découpage de $\App(A)$ en facettes induit un découpage du segment $[x,x']$ en points et segments ouverts. C'est-à-dire que l'on~a des points $x_{i}$ pour $i=1,\dots,n$, des facettes~$\cF_{i}$ pour $i=1,\dots,n-1$ et des facettes $\cF''_{i}$ pour $i=1,\dots,n$ de sorte que
\begin{align*}
[x,x_{1}[{}&=[x,x']\cap \cF,\quad ]x_{n},x']=[x,x']\cap \cF'\,;\\
]x_{i},x_{i+1}[{}&=[x,x']\cap \cF_{i}\quad\text{pour }i=1,\dots,n-1\,;\\
\{x_{i}\}&=[x,x']\cap \cF''_{i}\quad\text{pour }i=1,\dots,n.
\end{align*}

Un élément $k\in K_{\cF_{M}}^{\nu}$ appartient à la fois à $K_{\cF}^{\nu}$ et $K_{\cF'}^{\nu}$. Donc son action sur $\Imm(G)$ fixe tout le segment $[x,x']$. Il en résulte que $\nu$ appartient à $\cN(\cF_{i})$ et à $\cN(\cF''_{i})$ pour tout $i$ et que l'on peut aussi bien ajouter des exposants $\nu$ dans les égalités ci-dessus, par exemple $]x_{i},x_{i+1}[=[x,x']\cap \cF_{i}^{\nu}$. On décrit les facettes comme en \ref{description}. Parce que $M\supset M_{\cF,\nu}\supset M_{\cF}$, on~a $\Sigma_{\cF}=\Sigma_{\cF^{M}}$ et, parce que $\cF^M=\cF_{M}$, on~a $c_{\alpha,\cF}=c_{\alpha,\cF_{M}}$ pour tout $\alpha\in \Sigma^M$. Il en est de même en remplaçant $\cF$ par $\cF'$. Fixons $i=1,\dots,n-1$. L'ensemble $\Sigma_{\cF_{i}}$ est celui des $\alpha$ tels que, pour un point $y\in \cF_{i}$, ou pour tout point $y\in \cF_{i}$, $\alpha(y)$ appartient à $\Gamma_{\alpha}$. Puisque $[x,x']\cap \cF_{i}$ est ouvert dans $[x,x']$, il revient au même de dire que $\alpha\in \Sigma_{\cF}\cap \Sigma_{\cF'}$ et $c_{\alpha,\cF}=c_{\alpha,\cF'}$. La description ci-dessus entraîne alors $\Sigma_{\cF_{i}}=\Sigma_{\cF}=\Sigma_{\cF'}=\Sigma_{\cF_{M}}$. Donc $M_{\cF_{i}}=M_{\cF}$. L'espace $\cA_{M_{\cF_{i},\nu}}$ est le sous-espace des points fixes de l'action de l'élément $k$ ci-dessus dans $\cA_{M_{\cF_{i}}}$. Il en est de même en remplaçant $\cF_{i}$ par $\cF$. Puisque $M_{\cF_{i}}=M_{\cF}$, cela entraîne $M_{\cF_{i},\nu}=M_{\cF,\nu}=M_{\cF_{M},\nu}$. Enfin, on~a évidemment $\cF_{i}^M=\cF_{M}$. Cela démontre que chaque facette $\cF_{i}$ satisfait aux mêmes conditions que $\cF$ et $\cF'$. On peut définir une fonction $f_{\cF_{i}}$ pour tout $i$ et il nous suffit de démontrer successivement que les fonctions $f_{\cF}, f_{\cF_{1}},\dots,f_{\cF_{n-1}}$ et $f_{\cF'}$ ont même image dans $\rI(G)$. On est ramené au cas de facettes consécutives, autrement dit au cas $n=1$ dans notre construction. Notons simplement $\cF''=\cF''_{1}$. Les facettes~$\cF$ et~$\cF'$ déterminent des espaces paraboliques $\bP^{\nu}$ et $\bP^{\prime\nu}$ de $\bG_{\cF''}^{\nu}$: $K_{\cF}^{\nu}$~est inclus dans $ K_{\cF''}^{\nu}$ et $\bP^{\nu}(k_{F})$ est l'image de $K_{\cF}^{\nu}$ dans $\bG_{\cF''}^{\nu}(k_{F})$. Ils ont un espace de Levi commun $\bM^{\nu}$, qui est isomorphe à $\bG_{\cF}^{\nu}$ et à $\bG_{\cF'}^{\nu}$, ou encore à $\bM_{\cF_{M}}^{\nu}$: c'est l'image de $K_{\cF}^{\nu}\cap K_{\cF'}^{\nu}$. La~fonction~$f$ s'identifie à une fonction sur $\bM^{\nu}(k_{F})$. Notons~$f^{\bP}$ la fonction sur $\bG_{\cF''}^{\nu}(k_{F})$, à support dans $\bP^{\nu}(k_{F})$, telle que $f^{\bP}(mu)=f(m)$ pour tous $m\in \bM^{\nu}(k_{F})$ et $u\in \bU_{\bP}(k_{F})$. Il résulte des constructions que $f_{\cF}=(f^{\bP})_{\cF''}$. Définissons une fonction $f''$ sur $\bG^{\nu}_{\cF''}(k_{F})$ par l'égalité
\[
f''(g)=\vert \bP(k_{F})\vert ^{-1}\sum_{x\in \bG_{\cF''}(k_{F})}f^{\bP}(x^{-1}gx)
\]
pour tout $g\in \bG_{\cF''}(k_{F})$. En posant $c=\vert \bG_{\cF''}(k_{F})\vert \vert \bP(k_{F})\vert ^{-1}$, on voit que les \hbox{fonctions~$f''_{\cF''}$} et~$c(f^{\bP})_{\cF''}$ ont même image dans $\rI(G)$. Il en est donc de même des fonctions~$f''_{\cF''}$ et~$cf_{\cF}$. On a utilisé l'espace $\bP^{\nu}$ pour construire $f''$. On peut refaire la construction en utilisant l'espace $\bP^{\prime\nu}$. Le point est que la fonction $f''$ obtenue est la même, c'est la fonction \og induite\fg de $f$. Donc les fonctions $f''_{\cF''}$ et $cf_{\cF'}$ ont même image dans $\rI(G)$ et il en est de même des fonctions $f_{\cF}$ et $f_{\cF'}$. Cela achève la démonstration.
\end{proof}

\begin{rem*}
Le groupe $ K_{\cF_{M}}^{\dag,G}$ agit naturellement sur $\bM_{\cF_{M}}^{\nu}(k_{F})$ et sur $C(\bM_{\cF_{M}}^{\nu})$. Pour $g\in K_{\cF_{M}}^{\dag,G}$ et $f\in C(\bM_{\cF_{M}}^{\nu})$, notons $\gf$ l'image de $f$ par l'action de $g$. Il résulte du lemme que les fonctions $(\gf)_{\cF}$ et $f_{\cF}$ ont même image dans $\rI(G)$.
\end{rem*}

\subsection{Un espace de fonctions}\label{unespacedefonctions}

On note $\cE(G)$ le sous-espace de $C_{\rc}^{\infty}(G(F))$ engendré par les $\cE_{\cF}^{\nu}$ quand $(\cF,\nu)$ décrit $\Fac^*(G)$. On note $\rIE(G)$ l'image de $\cE(G)$ dans $\rI(G)$.

Pour $M\in \underline{\cL}_{\min}$, le groupe $\Norm_{G}(A)(F)\cap \Norm_{G}(M)(F)$ agit sur l'ensemble $\Fac^*_{\max,G\tcomp}(M,A)$. Fixons un ensemble de représentants $ \underline{\Fac}^*_{\max,G\tcomp}(M,A)$ des orbites. Pour chaque $(\cF_{M},\nu)\in \underline{\Fac}^*_{\max,G\tcomp}(M,A)$, l'action du groupe $K_{\cF_{M}}^{\dag,G}$ sur $C(\bM_{\cF_{M}}^{\nu})$ conserve l'espace $C_{\cusp}(\bM^{\nu}_{\cF_{M}})$ et on~note $C_{\cusp}(\bM^{\nu}_{\cF_{M}})^{K_{\cF_{M}}^{\dag,G}}$ le sous-espace des invariants.
On pose
\[
\cE(G,M)=\bigoplus_{(\cF_{M},\nu)\in \underline{\Fac}^*_{\max,G\tcomp}(M,A)}C_{\cusp}(\bM^{\nu}_{\cF_{M}})^{K_{\cF_{M}}^{\dag,G}}.
\]
Pour tout $(\cF_{M},\nu)\in \underline{\Fac}^*_{\max,G\tcomp}(M,A)$, fixons comme dans la section précédente une facette $\cF_{M}^G\in \Fac(G,A)$ telle que
\[
(\cF_{M}^G,\nu)\in \Fac^*(G,A),\quad M_{\cF_{M}^G,\nu}=M\quadet (\cF_{M}^G)^M=\cF_{M}.
\]
Pour $f\in C_{\cusp}(\bM^{\nu}_{\cF_{M}})^{K_{\cF_{M}}^{\dag,G}}$, on~a défini la fonction $f_{\cF_{M}^G}$ sur $G(F)$. L'application $f\mto f_{\cF_{M}^G}$ se prolonge par linéarité en une application de $\cE(G,M)$ dans $\cE(G)$.
On~note $e(G,M):\cE(G,M)\to \rIE(G)$ la composée de cette application et de la projection $\cE(G)\to \rIE(G)$. On note $ \Im(e(G,M))$ l'image de $e(G,M)$.

\begin{lem*}
On a l'égalité $\rIE(G)=\sum_{M\in \underline{\cL}_{\min}}\Im(e(G,M))$.
\end{lem*}

\begin{proof}
Soit $(\cF,\nu)\in \Fac^*(G)$ et $f\in C(\bG_{\cF}^{\nu})$. On veut prouver que l'image de $f_{\cF}$ dans $\rI(G)$ appartient à la somme des $\Im(e(G,M))$. On ne perd rien à supposer~$f$ invariante par conjugaison par $\bG_{\cF}(k_{F})$. On sait que l'on peut écrire $f$ comme somme finie de fonctions \og induites\fg $\Ind_{\bM^{\nu}}^{\bG^{\nu}}(f_{\bM^{\nu}})$ (la définition~a été rappelée dans la preuve précédente), où $\bM^{\nu}$ est un espace de Levi de $\bG^{\nu}$ et $f_{\bM^{\nu}}$ est une fonction cuspidale sur $\bM^{\nu}(k_{F})$ invariante par conjugaison par $\bM(k_{F})$ (cette décomposition résulte par exemple de \cite[Th.\,2.11]{A}). On peut aussi bien supposer que $f$ est l'une de ces fonctions induites, disons $f=\Ind_{\bM^{\nu}}^{\bG^{\nu}}(f')$. Fixons une telle fonction et un espace parabolique~$\bP^{\nu}$ de composante de Levi $\bM^{\nu}$. Au sous-groupe parabolique $\bP$ correspond une facette $\cF'\in \Fac(G)$ dont l'adhérence contient $\cF$. On a encore $(\cF',\nu)\in \Fac^*(G)$ d'après l'existence d'un espace $\bP^{\nu}$ de parabolique associé $\bP$. On a vu dans la preuve précédente que les images dans $\rI(G)$ de $f_{\cF}$ et de $f'_{\cF'}$ étaient proportionnelles.
Cela nous ramène au problème de départ où l'on~a remplacé $\cF$ et $f$ par $\cF'$ et $f'$. En oubliant cette construction, on peut supposer $f$ cuspidale.

On peut conjuguer $\cF$ et $f$ par un élément de $G(F)$, cela ne change pas l'image de $f_{\cF}$ dans $\rI(G)$. On peut donc supposer $\cF\subset \App(A)$. Posons $M=M_{\cF,\nu}$. On a défini la facette $\cF^M\subset \Imm(M_{\AD})$ associée à $\cF$. On a $\nu\in \cN_{G\tcomp}^M(\cF^M)$ d'après \cite[Lem.\,6(i)]{W2}. La dernière assertion de ce lemme et l'égalité $M=M_{\cF,\nu}$ entraînent que $\cF^{M,\nu}$ est réduit à un point, c'est-à-dire que $(\cF^M,\nu)\in \Fac^*_{\max}(M,A)$. De nouveau, par conjugaison, on peut supposer
\[
M\in \underline{\cL}_{\min}\quadet (\cF^M,\nu)\in \underline{\Fac}^*_{\max,G\tcomp}(M,A).
\]
Les espaces $\bG_{\cF}^{\nu}$ et $\bM_{\cF^M}^{\nu}$ s'identifient, on peut considérer $f$ comme un élément de $C_{\cusp}(\bM_{\cF^M}^{\nu})$. Remarquons que $\cF$ satisfait aux les mêmes propriétés que la facette~$\cF_{M}^G$ que l'on~a fixée plus haut. D'après le lemme \ref{relevement}, les fonctions $f_{\cF}$ et $f_{\cF_{M}^G}$ ont même image dans $\rI(G)$. La fonction $f$ n'est pas forcément invariante par
conjugaison par $K_{\cF^M}^{\dag,G}$. Mais, d'après la remarque \ref{relevement}, on peut moyenner $f$ par l'action de ce groupe sans changer l'image de $f_{\cF}$ dans $\rI(G)$. On peut donc supposer $f$ invariante par $K_{\cF^M}^{\dag,G}$. On a alors $f\in \cE(G,M)$ et l'image de $f_{\cF}$ dans $\rI(G)$ est égale à l'image de $f$ par l'application $e(G,M)$.
\end{proof}

\subsection{Deux espaces en dualité}\label{deuxespaces}

Pour $M\in \underline{\cL}_{\min}$, notons $\cD_{G\tcomp}(G,M)$ l'image dans $\cD_{G\tcomp}(G)$ de
l'espace $\cD_{\cusp,G\tcomp}(M)$. Si $M=G$, $\cD_{G\tcomp}(G;G)$ est simplement l'espace $\cD_{\cusp}(G)$ déjà défini. On a
\[
\cD_{G\tcomp}(G)=\sum_{M\in \underline{\cL}_{\min}}\cD_{G\tcomp}(G,M).
\]
Pour tout $M\in \underline{\cL}_{\min}$,
on vérifie que le sous-espace
\begin{equation}\tag{1}\label{deuxespaces1}
\prod_{(\cF_{M},\nu)\in \underline{\Fac}^*_{\max,G\tcomp}(M,A)}C_{\cusp}(\bM^{\nu}_{\cF_{M}})^{K_{\cF_{M}}^{\dag,G}}\subset \boldsymbol{\cD}_{\cusp,G\tcomp}(M)
\end{equation}
s'envoie bijectivement sur $\cD_{G\tcomp}(G,M)$.

\skpt
\begin{prop*}
\begin{enumeratei}
\item\label{deuxespacesi}
Pour tout $M\in \underline{\cL}_{\min}$, l'application $e(G,M)$ est injective et la restriction de l'application $D^G$ à $\cD_{G\tcomp}(G,M)$ est injective.

\item\label{deuxespacesii}
$\rIE(G)$, \resp $D^G(\cD_{G\tcomp}(G))$, est somme directe des $\Im(e(G,M))$, \resp des $D^G(\cD_{G\tcomp}(G,M))$, quand $M$ décrit $\underline{\cL}_{\min}$.

\item\label{deuxespacesiii}
Soient $M,L\in \underline{\cL}_{\min}$. L'application
\begin{align*}
\cD_{G\tcomp}(G,M)&\to \Im(e(G;L))^*\\
\bmf&\mto (D^G_{\bmf})_{\vert \Im(e(G;L))}
\end{align*}
est nulle si $M\neq L$ et est un isomorphisme si $M=L$.
\end{enumeratei}
\end{prop*}

\begin{proof}
Tous ces espaces sont sommes ou produits d'espaces indexés par des couples $(\cF_{M},\nu)$. Pour des raisons de support, on peut fixer $\nu\in \cN$ et remplacer tous les espaces par les sous-espaces analogues où on se limite au couples $(\cF_{M},\nu')$ tels que $\nu'=\nu$, \cf \ref{variantes}. On gagne que ces espaces sont de dimension finie. En fait, d'après \eqref{deuxespaces1}, les espaces $\cD_{G\tcomp}(G,M)$ et $\cE(G,M)$ deviennent isomorphes. Fixons une base $(\bmf_{i,M})_{i=1,\dots,n_{M}}$ de $\cE(G,M)$, réunion de bases orthogonales des espaces $C_{\cusp}(\bM^{\nu}_{\cF_{M}})^{K_{\cF_{M}}^{\dag,G}}$ intervenant. En appliquant la proposition \ref{calcul} et la remarque \eqref{calcul2} qui la suit, on voit
que la matrice
\[
\big(D^G_{\bmf_{i,M}}(e(G,L)(\overline{\bmf_{j,L}}))\big)_{\substack{M,L\in \underline{\cL}_{\min},\\
i=1,\dots,n_{M},\\
j=1,\dots,n_{L}}}
\]
est diagonale, de coefficients diagonaux non nuls. Compte tenu du lemme précédent, cela entraîne toutes les assertions de l'énoncé.
\end{proof}

\Subsection{Un corollaire}\label{uncorollaire}

\skpt
\begin{cor*}
\begin{enumeratei}
\item\label{uncorollairei}
L'application linéaire $D^G:\cD_{G\tcomp}(G)\to \rI(G)^*$ est injective.

\item\label{uncorollaireii}
L'application linéaire composée $\cD_{\cusp}(G)\To{D^G}\rI(G)^*\to \rI_{\cusp}(G)^*$ est injective.
\end{enumeratei}
\end{cor*}

\begin{proof}
Cela résulte de la proposition précédente, en se rappelant que $\cE(G;G)$ est formé d'après sa définition de fonctions cuspidales.
\end{proof}

\subsection{Injectivité de \texorpdfstring{$D^G$}{DG}}\label{injectivite}
Pour $M\in \underline{\cL}_{\min}$, notons $\cD(G,M)$ l'image dans $\cD(G)$ de
l'espace $\cD_{\cusp}(M)$. On a $\cD(G)=\sum_{M\in \underline{\cL}_{\min}}\cD(G,M)$.

\skpt
\begin{prop*}
\begin{enumeratei}
\item\label{injectivitei}
L'application $D^G:\cD(G)\to \rI(G)^*$ est injective.

\item\label{injectiviteii}
Pour tout $n\in \ZZ$, l'espace $D^G(\cD(G))\cap \Ann^n\rI(G)^*$ est égal à la somme des $D^G(\cD(G,M))$ où $M$ parcourt les éléments de $\underline{\cL}_{\min}$ tels que $a_{M}\geq n$.
\end{enumeratei}
\end{prop*}

\begin{proof}
Notons $\cD^n(G)$ la somme des $\cD(G,M)$ où $M$ parcourt l'ensemble de Levi indiqué ci-dessus. Par construction, on a
\[
\cD^n(G)=\cD^{n+1}(G)+\sum_{M\in \underline{\cL}^n_{\min}}\cD(G,M).
\]
Il résulte des définitions que, pour tout $M\in \underline{\cL}^n_{\min}$, $\cD(G,M)$ est l'image naturelle dans $\cD(G)$ de $\cD_{\cusp}(M)^{W^G(M)}$.
On a aussi $D^G(\cD^n(G))\subset D^G(\cD(G))\cap \Ann^n\rI(G)^*$. On en déduit une suite d'applications
\begin{multline*}
\bigoplus_{M\in \underline{\cL}^n_{\min}} \cD_{\cusp}(M)^{W^G(M)}\to \bigoplus_{M\in \underline{\cL}^n_{\min}}\cD(G,M)\to \cD^n(G)/\cD^{n+1}(G)\\[-5pt]
\To{\delta^n} (D^G(\cD(G))\cap \Ann^n\rI(G)^*)/(D^G(\cD(G))\cap \Ann^{n+1}\rI(G)^*)\\
\to \Gr^n\rI(G)^*\simeq \bigoplus_{M\in \underline{\cL}^n_{\min}}\rI_{\cusp}(M)^{*W^G(M)}.
\end{multline*}
Le corollaire précédent, appliqué aux Levi $M\in\underline{\cL}_{\min}^n$, implique que l'application composée est injective. Les deux premières applications de la suite sont surjectives. Il en résulte que $\delta^{n}$ est injective. Pour $n=a_{G}$, on a
\[
D^G(\cD(G))\cap \Ann^{a_{G}}\rI(G)^*=D^G(\cD(G))=D^G(\cD^n(G)).
\]
L'injectivité de $\delta^{n}$ pour tout $n$ entraîne alors par récurrence que, pour tout $n$, on a
\[
D^G(\cD(G))\cap \Ann^n\rI(G)^*=D^G(\cD^n(G))\quadet \Ker(D^G)\subset \cD^n(G).
\]
Pour $n=a_{M_{\min}}+1$, cette dernière relation implique l'injectivité de $D^G$.
\end{proof}

\subsection{Variantes avec caractère central}\label{varianteaveccaracterecentral}
Pour un groupe topologique abélien et localement compact $X$, nous appelons caractère de~$X$ un homomorphisme continu de~$X$ dans $\CC^{\times}$.
Soit $\xi$ un caractère de $A_{G}(F)$. On note $C_{c,\xi}^{\infty}(G(F))$ l'espace des fonctions sur $G(F)$, à valeurs complexes, localement constantes, telles que
\[
f(ag)=\xi(a)^{-1}f(g)\quad \text{pour tous $a\in A_{G}(F)$ et tout $g\in G(F)$}
\]
et telles que l'image dans $A_{G}(F)\backslash G(F)$ du support de $f$ soit compacte. Pour $f\in C_{\rc}^{\infty}(G(F))$, notons $f_{\xi}$ la fonction définie par
\[
f_{\xi}(g)=\int_{A_{G}(F)}f(ag)\xi(a)\,da.
\]
L'application linéaire $f\mto f_{\xi}$ est une surjection de $C_{\rc}^{\infty}(G(F))$ sur $ C_{c,\xi}^{\infty}(G(F))$. De même que l'on~a défini $\rI(G)$, on définit l'espace $I_{\xi}(G)$. L'action de $A_{G}(F)$ sur $C_{\rc}^{\infty}(G(F))$ se descend en une action sur $\rI(G)$. Le groupe $A_{G}(F)$ agit dualement sur $\rI(G)^*$. L'espace $I_{\xi}(G)^*$ dual de $I_{\xi}(G)$ s'identifie au sous-espace des éléments de $\rI(G)^*$ qui se transforment selon $\xi$. Précisément, un élément $d$ de ce sous-espace s'identifie à l'élément de $I_{\xi}(G)^*$ qui envoie $f_{\xi}$ sur $d(f)$ pour tout $f\in C_{\rc}^{\infty}(G(F))$.

Soit $\cF\in \Fac(G)$. Le groupe $A_{G}(F)$ est contenu dans $K_{\cF}^{\dag}$. Il agit par multiplication sur ce groupe. La multiplication par $a\in A_{G}(F)$ envoie un sous-ensemble $K_{\cF}^{\nu}$ sur $K_{\cF}^{\nu+\nu_{a}}$, où on~a posé $\nu_{a}=w_{G}(a)$. Si $(\cF,\nu)\in \Fac^*_{\max}(G)$, on~a aussi $(\cF,\nu+\nu_{a})\in \Fac^*_{\max}(G)$. Dans ce cas, la multiplication par $a$ se descend en un isomorphisme encore noté
\[
a:C_{\cusp}(\bG_{\cF}^{\nu})\to C_{\cusp}(\bG_{\cF}^{\nu+\nu_{a}}).
\]
Ces isomorphismes définissent une action de $A_{G}(F)$ sur l'espace $\boldsymbol{\cD}_{\cusp}(G)$ (on prendra soin de la distinguer de l'action par conjugaison, qui est triviale). L'action du sous-groupe $A_{G,\tu}(F)$ est triviale. L'action se descend en une action sur $\cD_{\cusp}(G)$ que l'on~note $(a,\bmf)\mto \bmf^{a}$. Soit $\xi$ un caractère modérément ramifié de $A_{G}(F)$, c'est-à-dire trivial sur $A_{G,\tu}(F)$. On note $\cD_{\cusp,\xi}(G)$ le sous-espace des éléments $\bmf\in \cD_{\cusp}(G)$ tels que $\bmf^{a}=\xi(a)\bmf$ pour tout $a\in A_{G}(F)$.

Plus généralement, on définit de même les variantes \og à caractère central\fg de beaucoup d'objets déjà définis (au sens ci-dessus: il s'agit d'un caractère $\xi$ de $A_{G}(F)$). On les note en ajoutant un indice $\xi$ dans les notations. Notons en particulier l'égalité%
\begin{equation}\tag{1}\label{varianteaveccaracterecentral1}
D^G(\cD(G))\cap \rI(G)^*_{\xi}=D^G(\cD_{\xi}(G)),
\end{equation}
qui résulte aisément de l'injectivité de $D^G$.
Plus généralement, si $\xi_{1},\dots,\xi_{n}$ sont des caractères de $A_{G}(F)$,
\begin{equation}\tag{2}\label{varianteaveccaracterecentral2}
D^G(\cD(G))\cap \Bigl(\sum_{i=1,\dots,n}\rI(G)^*_{\xi_{i}}\Bigr)=D^G\Bigl(\sum_{i=1,\dots,n}\cD_{\xi_{i}}(G)\Bigr).
\end{equation}

\begin{proof}
On peut supposer les $\xi_{i}$ distincts. Par interpolation, si $d$ appartient au membre de gauche, les composantes $d_{i}$ de $d$ dans chaque $\rI(G)^*_{\xi_{i}}$ sont combinaisons linéaires finies de translatés $d^{a}$ pour des $a\in A_{G}(F)$. Or ces $d^{a}$ appartiennent tous à $D^G(\cD(G))$. Donc $d_{i}\in D^G(\cD(G))\cap \rI(G)^*_{\xi_{i}}$ et il reste à appliquer l'égalité \eqref{varianteaveccaracterecentral1}.
\end{proof}

\section{Éléments compacts, \texorpdfstring{$p'$}{p'}-éléments}
\addtocontents{toc}{\protect\enlargethispage{-\protect\baselineskip}}

\subsection{Retour sur les éléments topologiquement unipotents}\label{retour}
Pour $x\in G_{\tu}(F)$, l'homomorphisme $n\mto x^n$ de $\ZZ$ dans $G(F)$ se prolonge en un homomorphisme continu $z\mto x^z$ de $\ZZ_{p}$ dans $G(F)$. Il en résulte que,
\begin{enumerate}
\item\label{retour1}
si $J\subset G(F)$ est un sous-groupe fermé et s'il existe un entier $c\geq1$ premier à $p$ tel que $x^c\in J$, alors $x\in J$.

\noindent
On a:
\item\label{retour2}
l'application naturelle $G_{\tu}(F)\to G_{\AD,\tu}(F)$ est surjective; ses fibres sont les orbites de l'action par multiplication de $(Z(G)^0)_{\tu}(F)$ dans $G_{\tu}(F)$.
\end{enumerate}

\begin{proof}
Notons $\pi:G\to G_{\AD}$ l'homomorphisme naturel. L'hypothèse \ref{HypG} entraîne que $G_{\AD}(F)/\pi(G(F))$ est d'ordre premier à $p$. Pour $y\in G_{\AD,\tu}(F)$, il y a donc un entier $c\geq1$ premier à $p$ tel que $y^c\in \pi(G(F))$. D'après \eqref{retour1}, $y$ appartient à $\pi(G(F))$. Soit $x\in G(F)$ tel que $\pi(x)=y$. On décompose $x=x_{\rss}x_{u}$, où~$x_{\rss}$ est semi-simple, $x_{u}$ est unipotent et $x_{\rss}$ et $x_{u}$ commutent. Fixons un sous-tore maximal~$T$ de $G$ tel que $x_{\rss}\in T(F)$ et fixons une extension galoisienne finie $F'$ de~$F$, de degré~$d$ premier à $p$, de sorte que $T$ soit déployé sur $F'$. Posons $T_{\ad}=T/Z(G)$. L'élément $\pi(x_{\rss})$ est topologiquement unipotent. La projection $T_{\tu}(F')\to T_{\ad,\tu}(F')$ s'identifie~à
\[
X_{*}(T)\otimes_{\ZZ}(1+\frp_{F'})\to X_{*}(T_{\ad})\otimes_{\ZZ}(1+\frp_{F'}).
\]
Elle est surjective car l'hypothèse \ref{HypG} implique que l'image de $X_{*}(T)$ dans $X_{*}(T_{\ad})$ est un sous-groupe d'indice fini premier à $p$.
On peut donc choisir \hbox{$x_{1}\!\in\!T_{\tu}(F')$} tel que $\pi(x_{1})=\pi(x_{\rss})$. L'élément $x_{2}=\Norme_{F'/F}(x_{1})$ appartient à $T_{\tu}(F)$ et on~a $\pi(x_{2})=\pi(x_{\rss}^d)$. D'après \eqref{retour1}, cela implique qu'il existe $x_{3}\in T_{\tu}(F)$ tel que $\pi(x_{3})=\nobreak\pi(x_{\rss})$. L'élément $x_{3}x_{u}$ appartient à $G_{\tu}(F)$ et satisfait à $\pi(x_{3}x_{u})=y$. Cela démontre la première assertion.

L'hypothèse \ref{HypG} entraîne que $Z(G)/Z(G)^0$ est d'ordre premier à $p$. La seconde assertion s'en déduit par le même argument qui prouvait ci-dessus que $y$ appartient à $\pi(G(F))$.
\end{proof}

On a:
\begin{enumerate}\setcounter{enumi}{2}
\item\label{retour3}
pour $x\in G_{\tu}(F)$, l'image réciproque de $Z_{G_{\AD}}(x_{\ad})$ dans $G$ est égale à $Z_{G}(x)$.
\end{enumerate}

\begin{proof}
Notons simplement $Z_{G}(x_{\ad})$ cette image réciproque. Il est bien connu que $Z_{G}(x)$ est un sous-groupe distingué d'indice fini dans $Z_{G}(x_{\ad})$ et l'hypothèse \ref{HypG} implique que cet indice est premier à $p$. Soit $g\in Z_{G}(x_{\ad})$. Il existe donc un entier $c\geq1$ premier à $p$ tel que $g^c$ commute à $x$. Puisque $g\in Z_{G}(x_{\ad})$ il~existe $z\in Z(G)$ tel que $ g^{-1}xg=zx$. On a alors $z^c=1$. Or, puisque $x$ et $g^{-1}xg$ sont topologiquement unipotents, $z$ est lui aussi topologiquement unipotent. L'égalité $z^c=1$ entraîne alors que $z=1$, donc $g\in Z_{G}(x)$.
\end{proof}

On utilisera souvent la propriété suivante:

\begin{enumerate}\setcounter{enumi}{3}
\item\label{retour4}
soit $x$ un élément semi-simple de $G(F)$ et soit $y\in Z_{G}(x)(F)\cap G_{\tu}(F)$; alors $y\in G_{x,\tu}(F)$.
\end{enumerate}

En effet, l'ordre du groupe $Z_{G}(x)/G_{x}$ est borné par le nombre d'éléments du centre du revêtement simplement connexe de $G_{\AD}$. Il est donc premier à $p$ d'après \ref{HypG}.

\subsection{Éléments topologiquement nilpotents et exponentielle}\label{elementstopologiquement}
On note $\frg$ l'algèbre de Lie de $G$. On appelle \og conjugaison\fg par $G$ l'action adjointe de $G$ dans~$\frg$. Pour $g\in G$, on~note cette action de $g$, soit $\ad(g)$, soit $X\mto gXg^{-1}$. On note $\frg_{\reg}$ l'ensemble des éléments semi-simples réguliers de $\frg$.

Soit $X\in \frg(F)$, fixons un sous-tore maximal $T$ de $G$ tel que $X_{\rss}\in \frt(F)$ et fixons une extension finie $F'$ de $F$ telle que $T$ soit déployé sur $F'$. L'élément $X$ est dit topologiquement nilpotent si et seulement si $ \chi(X_{\rss})\in \frp_{F'}$ pour tout $\chi\in X^*(T)$ (cela ne dépend pas du choix de $T$). Une autre caractérisation est la suivante. Fixons une sous-algèbre d'Iwahori $\frb$ de $\frg$, c'est-à-dire $\frb=\frk_{\cF}$ pour une facette $\cF\in \Fac(G)$ de dimension maximale. Notons $\fru$ son radical pro-$p$-nilpotent, c'est-à-dire $\fru=\frk_{\cF}^+$ pour la même facette. Alors $X$ est topologiquement nilpotent si et seulement si $X$ est conjugué par un élément de $G(F)$ à un élément de $\fru(F)$. On note $\frg_{\tn}(F)$ l'ensemble des éléments topologiquement nilpotents de $\frg(F)$.

On peut définir une application exponentielle $\exp$ qui envoie un voisinage de $0$ dans $\frg(F)$ sur un voisinage de $1$ dans $G(F)$ et qui est équivariante pour les actions par conjugaison de $G(F)$. Sous l'hypothèse \ref{HypG}, ces voisinages sont les plus grands possibles, c'est-à-dire qu'on dispose de l'application exponentielle
\[
\exp:\frg_{\tn}(F)\to G_{\tu}(F)
\]
équivariante pour les actions par conjugaison de $G(F)$ et qui
un homéomorphisme entre les ensembles de départ et d'arrivée, \cf \cite[App.\,B]{DR} pour cette propriété et celles ci-dessous.

Pour $\cF\in \Fac(G)$, l'exponentielle se restreint en des homéomorphismes de \hbox{$\frg_{\tn}(F)\cap \frk_{\cF}$} sur $G_{\tu}(F)\cap K_{\cF}^0$
et de $\frk_{\cF}^+$ sur $K_{\cF}^+$. Il s'en déduit une bijection
\[
(\frg_{\tn}(F)\cap \frk_{\cF})/\frk_{\cF}^+\to (G_{\tu}(F)\cap K_{\cF}^0)/K_{\cF}^+.
\]

On peut prolonger les paires $K_{\cF}^0\supset K_{\cF}^+$ et $\frk_{\cF}\supset \frk_{\cF}^+$
en des suites $(K_{\cF,n})_{n\in \NN}$ et $(\frk_{\cF,n})_{n\in \NN}$ de sorte que
\begin{itemize}
\item
$K_{\cF}^0=K_{\cF,0}$, $K_{\cF}^+ =K_{\cF,1}$, $K_{\cF,n}\supset K_{\cF,n+1}$, $\cap_{n\in \NN}K_{\cF,n}=\{1\}$;

\item
$\frk_{\cF}=\frk_{\cF,0}$, $\frk_{\cF}^+ =\frk_{\cF,1}$, $\frk_{\cF,n}\supset \frk_{\cF,n+1}$, $\cap_{n\in \NN}\frk_{\cF,n}=\{0\}$;

\item
pour $n\!\geq\!1$, $K_{\cF,n}$ est un sous-groupe distingué de $K_{\cF}^{\dag}$ et $\frk_{\cF,n}$ est un $\fro_{F}$-idéal de~$\frk_{\cF}$;

\item
pour $n\geq1$, $K_{\cF,n}=\exp(\frk_{\cF,n})$ et l'exponentielle se réduit en un isomorphisme de groupes de $\frk_{\cF,n}/\frk_{\cF,n+1}$ sur $K_{\cF,n}/K_{\cF,n+1}$.
\end{itemize}

\subsection{Éléments \texorpdfstring{$p'$}{p'}-compacts}\label{elementsp}

Soit $x\in G(F)$. On dit que $x$ est
\begin{itemize}
\item
$p'$-compact si et seulement s'il existe un entier $c\geq1$ premier à $p$ tel que $x^c=1$;

\item
$p'$-compact $\rmod Z(G)$ si et seulement si l'image $x_{\ad}$ de $x$ dans $G_{\AD}(F)$ est $p'$-compacte.
\end{itemize}

Pour $x\in G(F)$, fixons un sous-tore maximal $T$ de $G$ contenant $x_{\rss}$. Alors:

\begin{itemize}
\item
$x$ est $p'$-compact si et seulement $x=x_{\rss}$ et il existe un entier $c\geq1$ premier à $p$ tel que $\chi(x)^c=1$ pour tout $\chi\in X^*(T)$;

\item
$x$ est $p'$-compact $\rmod Z(G)$ si et seulement $x=x_{\rss}$ et il existe un entier $c\geq1$ premier à $p$ tel que $\chi(x)^c=1$ pour tout $\chi\in X^*(T)$ dont la restriction à $Z(G)$ est triviale.
\end{itemize}

Puisque les classes de conjugaison de sous-tores maximaux de $G$ sont en nombre fini, on peut choisir une extension $F'/F$ d'ordre premier à $p$ telle que tout sous-tore maximal de $G$ soit déployé sur $F'$. Notons $c$ le nombre d'éléments non nuls de son corps résiduel. C'est un entier premier à $p$. On voit
qu'il n'y a qu'un nombre fini de classes de conjugaison par $G(F)$ d'éléments $p'$-compacts et que l'on~a $x^c=1$ pour tout élément $p'$-compact de $G(F)$. Plus généralement, on~a $x^c\in G_{\tu}(F)$ pour tout élément compact $x\in G(F)$. On a:
\begin{enumerate}
\item\label{elementsp1}
tout élément compact $x\in G(F)$ s'écrit de façon unique $x=x_{p'}x_{\tu}$, où $x_{p'}$ est $p'$-compact, $x_{\tu}\in G_{\tu}(F)$ et $x_{p'}$ et $x_{\tu}$ commutent; de plus, $x_{p'}$ et $x_{\tu}$ appartiennent à $\overline{x^{\ZZ}}$.
\end{enumerate}

\begin{proof}
Comme on vient de le dire, $x^c$ est topologiquement unipotent. Soit~$c'$ l'inverse de $c$ dans $\ZZ_{p}$. On pose $x_{\tu}=(x^c)^{c'}$ et $x_{p'}=xx_{\tu}^{-1}$. Ces termes satisfont aux conditions requises et on voit que ce sont les seules solutions possibles.
\end{proof}

On déduit de \eqref{elementsp1} et de \ref{retour}\,\eqref{retour2} et \eqref{retour3} que
\begin{enumerate}\setcounter{enumi}{1}
\item\label{elementsp2}
tout élément $x\in G(F)$ compact $\rmod Z(G)$ s'écrit $x=x_{p'}x_{\tu}$, où $x_{p'}$ est $p'$\nobreakdash-compact $\rmod Z(G)$, $x_{\tu}\in G_{\tu}(F)$ et $x_{p'}$ et $x_{\tu}$ commutent; le groupe $(Z(G)^0)_{\tu}(F)$ agit sur l'ensemble des solutions: un élément $z\in (Z(G)^0)_{\tu}(F)$ envoie le couple $(x_{p'},x_{\tu})$ sur $(x_{p'}z^{-1},x_{\tu}z)$; les solutions forment une unique orbite pour cette action; de plus, pour toute solution $(x_{p'},x_{\tu})$, les deux éléments $x_{p'}$ et $x_{\tu}$ appartiennent à l'adhérence du groupe $Z(G)(F)x^\ZZ$.
\end{enumerate}

On utilisera aussi la variante suivante. Soit $\underline{Z}$ un sous-groupe algébrique de $Z(G)$. Pour $x\in G(F)$, on dit que $x$ est compact $\rmod \underline{Z}$, \resp $p'$-compact $\rmod \underline{Z}$, si et seulement si l'image de $x$ dans $(G/\underline{Z})(F)$ est compacte, \resp $p'$-compacte. Supposons que le quotient $Z(G)(F)/\underline{Z}(F)$ soit compact. Alors $x$ est compact $\rmod \underline{Z}$ si et seulement s'il est compact $\rmod Z(G)$ (la propriété analogue pour \og $p'$-compact\fg est évidemment fausse). On a

\begin{enumerate}\setcounter{enumi}{2}
\item\label{elementsp3}
tout élément $x\in G(F)$ compact $\rmod Z(G)$ s'écrit $x=x_{p'}x_{\tu}$, où $x_{p'}$ est $p'$\nobreakdash-compact $\rmod \underline{Z}$, $x_{\tu}\in G_{\tu}(F)$ et $x_{p'}$ et $x_{\tu}$ commutent; le groupe $\underline{Z}_{\tu}(F)$ agit sur l'ensemble des solutions comme en \eqref{elementsp2}; les solutions forment une unique orbite pour cette action.
\end{enumerate}

\subsection{\texorpdfstring{$p'$}{p'}-éléments}\label{pelements}

Soit $x\in G(F)$. On lui associe un sous-groupe parabolique $Q[x]$ de $G$ et une composante de Levi $L[x]$ de $Q[x]$ de la façon suivante, \cf \cite{C}. L'élément~$x_{\rss}$ agit par conjugaison sur $\frg$. On fixe une extension galoisienne finie $F'$ de $F$ telle que toutes les valeurs propres appartiennent à $F^{\prime\times}$. On note $\Sigma$ l'ensemble de ces valeurs propres et, pour $\sigma\in \Sigma$, $\frg_{\sigma}$ l'espace propre. Alors l'algèbre de Lie $\mfrq[x]$ est la somme des $\frg_{\sigma}$ sur les $\sigma\in \Sigma$ telles que $\vert \sigma\vert _{F'}\leq 1$ et $\frl[x]$ est la somme des $\frg_{\sigma}$ sur les $\sigma\in \Sigma$ telles que $\vert \sigma\vert _{F'}= 1$.
Le couple $(Q[x], L[x])$ étant uniquement défini et ne dépendant que de $x_{\rss}$, il est défini sur $F$ et conservé par $Z_{G}(x_{\rss})$. Cela entraîne
\begin{enumerate}
\item\label{pelements1}
$Z_{G}(x_{\rss})\subset L[x]$, a fortiori $Z_{G}(x)\subset L[x]$ et $x\in L[x]$.
\end{enumerate}

Par construction, $x$, vu comme élément de $L[x]$, est compact $\rmod Z(L[x])$ (on dira que $x$ est compact dans $L[x]\bmod Z(L[x])$). Le Levi $L[x]$ est le plus grand Levi $L$ de~$G$ tel que $x$ appartienne à $L$ et que $x$ soit compact dans $L\bmod Z(L)$.

Nous dirons que $x$ est un $p'$-élément si et seulement si $x$ est $p'$-compact dans $L[x]\bmod Z(L[x])$. Avec les notations ci-dessus, cela équivaut à ce que $x=x_{\rss}$ soit semi-simple et, pour tout $\sigma\in \Sigma$ telle que $\vert \sigma\vert _{F}=1$, $\sigma$ soit une racine de l'unité d'ordre premier à $p$. On note $G(F)_{p'}$ l'ensemble des $p'$-éléments de $G(F)$.

\skpt
\begin{lem*}
\begin{enumeratei}
\item\label{pelementsi}
Tout élément $x\in G(F)$ s'écrit $x=x_{p'}x_{\tu}$ où $x_{p'}$ est un $p'$-élément, $x_{\tu}$ est topologiquement unipotent et $x_{p'}$ et $x_{\tu}$ commutent.

\item\label{pelementsii}
Pour une telle décomposition, les éléments $x_{p'}$ et $x_{\tu}$ appartiennent à $L[x](F)$ et on~a $L[x_{p'}]=L[x]$.

\item\label{pelementsiii}
Pour une telle décomposition, on a
\[
Z_{G}(x)=Z_{G}(x_{p'})\cap Z_{G}(x_{\tu})\quadet G_{x}=(G_{x_{p'}})_{x_{\tu}}.
\]

\item\label{pelementsiv}
La décomposition est unique modulo l'action de $(Z(L[x])^0)_{\tu}(F)$ similaire à celle de \ref{elementsp}\,\eqref{elementsp2}.
\end{enumeratei}
\end{lem*}

\begin{proof}
Considérons une décomposition $x=yx_{\tu}$ où $x_{\tu}$ est topologiquement unipotent et $y$ et $x_{\tu}$ commutent. Montrons que
\begin{equation}\tag{2}\label{pelements2}
y,x_{\tu}\in L[x]\quadet L[x]=L[y].
\end{equation}
Les éléments $y$ et $x_{\tu}$ commutent à $x$ et la première assertion résulte de \eqref{pelements1}. Notons ici $x_{\ad}$, $y_{\ad}$ et $x_{\tu,\ad}$ les images de $x$, $y$ et $x_{\tu}$ dans $L[x]_{\AD}(F)$. L'adhérence $\overline{y_{\ad}^{\ZZ}}$ du groupe engendré par $y_{\ad}$ est contenue dans le produit des deux groupes $\overline{x_{\ad}^{\ZZ}}$ et $\overline{x_{\tu,\ad}^{\ZZ}}$. Ce dernier groupe est compact puisque $x_{\tu}$ est topologiquement unipotent. Le premier est compact par définition de $L[x]$. Donc $\overline{y_{\ad}^{\ZZ}}$ est compacte, c'est-à-dire que~$y$ est compact $\rmod Z(L[x])$. Or $L[y]$ est le plus grand Levi $L$ tel que $y\in L$ et $y$ soit compact $\rmod Z(L)$. Donc $L[x]\subset L[y]$. On obtient l'inclusion opposée en échangeant les rôles de $x$ et $y$ (on~a $y=xx_{\tu}^{-1}$). Cela prouve \eqref{pelements2}.

On a vu que $x$ appartenait à $ L[x]$ et était compact $\rmod Z(L[x])$. En appliquant \ref{elementsp}\,\eqref{elementsp2} dans le groupe $L[x]$, on obtient une décomposition $x=x_{p'}x_{\tu}$ où $x_{p'}, x_{\tu}\in L[x](F)$, $x_{p'}$ est $p'$-compact $\rmod Z(L[x])$, $x_{\tu}$ est topologiquement unipotent et $x_{p'}$ et $x_{\tu}$ commutent. D'après \eqref{pelements2}, on~a $L[x_{p'}]=L[x]$, donc $x_{p'}$ est $p'$-compact $\rmod Z(L[x_{p'}])$, donc c'est un $p'$-élément. Cela démontre \eqref{pelementsi}.

Pour une décomposition $x=x_{p'}x_{\tu}$ comme en \eqref{pelementsi}, les éléments $x_{p'}$ et $x_{\tu}$ commutent à $x$ donc appartiennent à $L[x]$ d'après \eqref{pelements1}. Cela démontre la première assertion de \eqref{pelementsii} et la seconde résulte de \eqref{pelements2}. En définitive, les décompositions $x=x_{p'}x_{\tu}$ satisfaisant à \eqref{pelementsi} sont exactement les décompositions dans $L[x](F)$ où l'on impose que $x_{p'}$ est $p'$-compact $\rmod Z(L[x])$. Le \eqref{pelementsiv} résulte donc de \ref{elementsp}\,\eqref{elementsp2} appliqué dans le groupe $L[x]$.

Pour \eqref{pelementsiii}, les commutants $Z_{G}(x)$ et $Z_{G}(x_{p'})$ sont contenus dans $L[x]$ d'après \eqref{pelements1} et la dernière assertion de \eqref{pelementsii}. On ne perd rien à supposer $L[x]=G$. Il est clair que $Z_{G}(x_{p'})\cap Z_{G}(x_{\tu})\subset Z_{G}(x)$. L'image $x_{\tu,\ad}$ de $x_{\tu}$ dans $G_{\AD}$ appartient à l'adhérence du groupe engendré par $x_{\ad}$. L'image $g_{\ad}$ d'un élément $g\in Z_{G}(x)$ commute donc à $x_{\tu,\ad}$. D'après \ref{retour}\,\eqref{retour2}, cela implique que $g\in Z_{G}(x_{\tu})$. Alors $g$ appartient forcément aussi à $Z_{G}(x_{p'})$. Cela démontre la première égalité de \eqref{pelementsiii}. La composante neutre de $Z_{G}(x_{p'})\cap Z_{G}(x_{\tu})$ est clairement $(G_{x_{p'}})_{x_{\tu}}$, d'où la seconde égalité.
\end{proof}

Pour $x\in G(F)$, appelons $p'$-décomposition de $x$ une décomposition $x=x_{p'}x_{\tu}$ satisfaisant aux conditions du \eqref{pelementsi} de l'énoncé.

\subsection{Décomposition de \texorpdfstring{$G(F)$}{GF} associée aux \texorpdfstring{$p'$}{p'}-éléments}\label{p'decomposition}

Pour tout $\epsilon\in G(F)_{p'}$, posons $C(\epsilon)=\epsilon G_{\epsilon,\tu}(F)=\epsilon\,\exp(\frg_{\epsilon,\tn}(F))$. Évidemment, l'action de $G(F)$ par conjugaison conserve $G(F)_{p'}$ et, pour $g\in G(F)$, on~a $gC(\epsilon)g^{-1}=C(g\epsilon g^{-1})$. Posons
\[
C_{G}(\epsilon)=\bigcup_{g\in G(F)}C(g^{-1}\epsilon g)=\bigcup_{g\in G(F)}g^{-1}C(\epsilon )g.
\]
On a déjà utilisé le discriminant de Weyl $D^G(x)$ d'un élément $x\in G(F)$. Il y a de même un discriminant de Weyl $D^G(X)$ pour un élément $X\in \frg(F)$.

\skpt
\begin{lem*}
\begin{enumeratei}
\item\label{p'decompositioni}
Pour $\epsilon\in G(F)_{p'}$, l'ensemble $C_{G}(\epsilon) $ est ouvert et fermé.

\item\label{p'decompositionii}
Pour $\epsilon\in G(F)_{p'}$, il existe un nombre réel $d^G(\epsilon)>0$ tel que, pour tout $X\in \frg_{\epsilon,\tn}(F)$, on ait
\[
D^G(\epsilon \exp(X))=d^G(\epsilon)D^{G_{\epsilon}}(X).
\]

\item\label{p'decompositioniii}
On a les égalités
\[
G(F)=\bigcup_{\epsilon\in G(F)_{p'}}C(\epsilon)=\bigcup_{\epsilon\in G(F)_{p'}}C_{G}(\epsilon).
\]

\item\label{p'decompositioniv}
Soient $\epsilon ,\epsilon '\in G(F)_{p'}$. Alors $C(\epsilon )=C(\epsilon ')$ ou $C(\epsilon )\cap C(\epsilon ')=\emptyset$. L'égalité a lieu si et seulement s'il existe $z\in (Z(L[\epsilon ])^0)_{\tu}(F)$ tel que $\epsilon '=\epsilon z$. Dans ce cas, on a
\[
L[\epsilon ]=L[\epsilon '],\quad G_{\epsilon}=G_{\epsilon'}\quadet z\in G_{\epsilon,\tu}(F).
\]
\end{enumeratei}
\end{lem*}

\begin{proof}
La propriété suivante est bien connue (\cf \cite[Cor.\,2.4]{HC}): soit \hbox{$x\in G(F)$} un élément semi-simple et $V_{x}'$ un voisinage de $1$ dans $G_{x}(F)$; alors il existe un voisinage $V_{x}$ de $1$ dans $G_{x}(F)$ contenu dans $V_{x}'$ tel qu'en posant
\[
B(x,V_{x})=\{ g^{-1} xyg\sep y\in V_{x}, g\in G(F)\},
\]
cet ensemble $B(x,V_{x})$ soit ouvert et fermé. Appliquons ceci à un élément semi-simple $x \in C(\epsilon)$. On a $G_{x}\subset G_{\epsilon}$ d'après le~\eqref{pelementsiii} du lemme \ref{pelements}. Prenons
\[
V'_{x}=\{y\in G_{x}(F)\sep xy\in C(\epsilon)\},
\]
dont on déduit un voisinage $V_{x}$. D'après le choix de $V'_{x}$, $B(x,V_{x})$ est contenu dans $C_{G}(\epsilon)$. La réunion des $B(x,V_{x})$ quand $x$ parcourt les éléments semi-simples de $C(\epsilon)$ contient tous ces éléments semi-simples. Étant invariante par conjugaison, elle contient tous les éléments semi-simples de $C_{G}(\epsilon)$ et contient en fait $C_{G}(\epsilon)$ tout entier: un élément quelconque est conjugué à un élément arbitrairement voisin de sa partie semi-simple. Mais $G_{\epsilon,\tu}(F)$ est compact modulo conjugaison par $G_{\epsilon}(F)$ et donc $C_{G}(\epsilon)$ est compact modulo conjugaison par $G(F)$. On peut extraire un sous-recouvrement de $C_{G}(\epsilon)$ par un nombre fini d'ensembles $B(x,V_{x})$. Il s'agit en fait d'une réunion finie puisque ces ensembles sont contenus dans $C_{G}(\epsilon)$. Puisqu'ils sont ouverts et fermés, on obtient \eqref{p'decompositioni}.

Les deux parties de la formule~\eqref{p'decompositionii} sont insensibles au remplacement de $X$ par sa partie semi-simple. On peut donc supposer $X$ semi-simple. On fixe un sous-tore maximal $T$ de $G_{\epsilon}$ tel que $X\in \frt(F)$. On a aussi $\epsilon\in T(F)$ puisque $T$ commute à~$\epsilon$. Fixons une extension finie $F'$ de $F$ telle que $T$ soit déployé sur $F'$. On note $\Sigma$ l'ensemble des racines de $T$ dans $G$, $\Sigma^{L[\epsilon]}$, \resp $\Sigma^{G_{\epsilon}}$, le sous-ensemble des racines dans $L[\epsilon]$, \resp $G_{\epsilon}$. On a $\Sigma^{G_{\epsilon}}\subset \Sigma^{L[\epsilon]}$. Par définition,
\begin{align*}
D^G(\epsilon)&=\bigl\vert \textstyle\prod_{\alpha\in \Sigma}(1-\alpha(\epsilon \exp(X)))\bigr\vert _{F}
\\
D^{G_{\epsilon}}(X)&=\bigl\vert \textstyle\prod_{\alpha\in \Sigma^{G_{\epsilon}}}\alpha(X)\bigr\vert _{F}.
\end{align*}
Puisque $\exp(X)$ est topologiquement unipotent, $\alpha(\exp(X))$ appartient à $1+\frp_{F'}$ pour tout $\alpha\in \Sigma$. A fortiori $\vert \alpha(\exp(X))\vert _{F'}=1$.
Si $\alpha\in \Sigma-\Sigma^{L[\epsilon]}$, on~a $\vert \alpha(\epsilon)\vert _{F'}\neq1$ par définition de $L[\epsilon]$. Puisque $\vert \alpha(\exp(X))\vert _{F'}=1$, on a
\[
\vert 1-\alpha(\epsilon \exp(X))\vert _{F'}=\vert 1-\alpha(\epsilon)\vert _{F'}
\]
qui est non nul. Si $\alpha\in \Sigma^{L[\epsilon]}$, $\alpha(\epsilon)$ est une racine de l'unité d'ordre premier à $p$ puisque~$\epsilon$ est $p'$-compact $\rmod Z(L[\epsilon])$. Cette racine est égale à $1$ si et seulement si $\alpha\in \Sigma^{G_{\epsilon}}$. Si $\alpha\in \Sigma^{L[\epsilon]}-\Sigma^{G_{\epsilon}}$, on a
\[
\alpha(\epsilon)\in \fro_{F'}^{\times}-(1+\frp_{F'})\quadet\alpha(\exp(X))\in 1+\frp_{F'}
\]
donc encore
\[
\vert 1-\alpha(\epsilon \exp(X))\vert _{F'}=\vert 1-\alpha(\epsilon)\vert _{F'}=1.
\]
Enfin, si $\alpha\in \Sigma^{G_{\epsilon}}$, on~a $\alpha(\epsilon)=1$ donc
\[
\vert 1-\alpha (\epsilon \exp(X))\vert _{F'}=\vert 1-\alpha(\exp(X))\vert _{F'}=\vert 1-\exp(\alpha(X))\vert _{F'}
\]
avec $\alpha(X)\in \frp_{F'}$, donc
\[
\vert 1-\alpha (\epsilon \exp(X))\vert _{F'}=\vert \alpha(X)\vert _{F'}.
\]
L'assertion \eqref{p'decompositioni} en résulte.

Le \eqref{pelementsi} du lemme \ref{pelements} implique la première égalité du \eqref{p'decompositioni} d'où trivialement la deuxième.

Soient $\epsilon,\epsilon'\in G(F)_{p'}$. Supposons $C(\epsilon )\cap C(\epsilon ')\neq\emptyset$ et fixons $x$ dans cette intersection. On peut écrire
\[
x=\epsilon u=\epsilon 'u'\quad\text{avec }u\in G_{\epsilon ,\tu}(F)\text{ et } u'\in G_{\epsilon ',\tu}(F).
\]
Ces deux décompositions sont des $p'$-décompositions. D'après le~\eqref{pelementsiv} du lemme \ref{pelements}, il existe $z\in (Z(L[\epsilon ])^0)_{\tu}(F)$ tel que $\epsilon '=\epsilon z$. Inversement, s'il existe un tel $z$, l'assertion~\eqref{pelements2} de \ref{pelements} appliquée à $x=\epsilon'=\epsilon z$ nous dit que $L[\epsilon ]=L[\epsilon ']$. D'après l'assertion~\eqref{pelements1} de \ref{pelements}, on a
\[
G_{\epsilon }=L[\epsilon ]_{\epsilon }\quadet G_{\epsilon '}=L[\epsilon ']_{\epsilon '}=L[\epsilon ]_{\epsilon '}.
\]
Puisque $z\in Z(L[\epsilon ])^0(F)$, l'égalité $\epsilon '=\epsilon z$ entraîne que $L[\epsilon ]_{\epsilon '}=L[\epsilon ]_{\epsilon }$ donc aussi $G_{\epsilon '}=G_{\epsilon }$. On a aussi $z\in G_{\epsilon ,\tu}(F)$ et on en déduit que
\[
C(\epsilon ')=\epsilon 'G_{\epsilon ',\tu}(F)=\epsilon zG_{\epsilon ,\tu}(F)=\epsilon G_{\epsilon ,\tu}(F)=C(\epsilon ).
\]
Cela démontre \eqref{p'decompositioniv}.
\end{proof}

\subsection{Le cas d'un Levi}\label{lecasdunlevi}

Soit $M$ un Levi de $G$. On a dans $M$ la même propriété que dans $G$, à savoir
\[
M(F)=\bigcup_{\epsilon\in M(F)_{p'}}C^M(\epsilon),
\]
où $C^M(\epsilon)=\epsilon \exp(\mfrm_{\epsilon,\tn}(F))$. Le lemme ci-dessous énonce une propriété plus fine.

\begin{lem*}
Soit $M$ un Levi de $G$. Alors
\begin{enumeratei}
\item\label{lecasdunlevii}
$G(F)_{p'}\cap M(F)\subset M(F)_{p'}$;

\item\label{lecasdunleviii}
$M(F)=\bigcup_{\epsilon\in G(F)_{p'}\cap M(F)}C^M(\epsilon)$.
\end{enumeratei}
\end{lem*}

\begin{proof}
Soit $\epsilon\in G(F)_{p'}\cap M(F)$. Pour démontrer que $\epsilon\in M(F)_{p'}$, on doit prouver que toute valeur propre de $\ad(\epsilon)$ dans $\mfrm(F)$ qui est de valeur absolue $1$ dans une extension convenable de $F$, est une racine de l'unité d'ordre premier à $p$. Mais une valeur propre de $\ad(\epsilon)$ dans $\mfrm(F)$ est aussi une valeur propre de $\ad(\epsilon)$ dans $\frg(F)$. La propriété voulue résulte du fait que $\epsilon\in G(F)_{p'}$.

Soit $x\in M(F)$.
Écrivons une $p'$-décomposition de $x$ dans $G(F)$: $x=\epsilon \exp(X)$, où $X\in \frg_{\epsilon,\tn}(F)$.
On a $A_{M}\subset G_{x} $ puisque $x\in M(F)$ et $G_{x}\subset G_{\epsilon}$ d'après le~\eqref{pelementsiii} du lemme~\ref{pelements}. Donc $A_{M}\subset G_{\epsilon}$, ce qui implique $\epsilon\in M(F)$. Alors $\exp(X)=\epsilon^{-1}x$ appartient aussi à $M(F)$. On en déduit que $X$ appartient à $\mfrm_{\epsilon,\tn}(F)$ donc $x$ appartient à $C^M(\epsilon)$, où~$\epsilon$ appartient à $G(F)_{p'}\cap M(F)$.
\end{proof}

\Subsection{Un lemme sur les classes de conjugaison et les éléments \texorpdfstring{$p'$-compacts mod\,$Z(G)$}{pZG}}\label{commutant}

\begin{lem*}
Soit $\epsilon\in G(F)$ un élément $p'$-compact $\rmod Z(G)$ et soit $\cF\in \Fac(G)$. Supposons $\epsilon\in K_{\cF}^{\dag}$. Soit $X\in \frg_{\epsilon,\tn}(F)\cap \frk_{\cF}$. Alors, pour tout élément $x\in \epsilon \exp(X) K_{\cF}^+$, il existe $Y\in \frg_{\epsilon}(F)\cap \frk_{\cF}^+$ tel que $x$ soit conjugué à $\epsilon \exp(X+Y)$ par un élément de~$K_{\cF}^+$.
\end{lem*}

La preuve est standard, on la rappelle pour être complet.

\begin{proof}
Fixons un entier $c\geq1$ premier à $p$ tel que $\epsilon^c\in Z(G)(F)$. On~intro\-duit les deux polynômes
\begin{align*}
P(T)&=c^{-1}(T^{c-1}+\cdots+T+1),\\
Q(T)&=c^{-1}(1-T)(c-1+(c-2)T+\cdots+T^{c-2})
\end{align*}
à coefficients dans $\ZZ_{p}$. On a $P(T)+Q(T)=1$.
L'opérateur $\ad(\epsilon)$ dans $\frg(F)$ est semi-simple et ses valeurs propres dans $\bar{F}$ sont des racines $c$-ièmes de l'unité. Alors l'espace~$\frg$ se décompose en somme de l'espace propre associé à la racine $1$, qui est~$\frg_{\epsilon}$, et de l'espace somme des espaces propres associés aux racines différentes de $1$, notons-le $\frg_{\neq1}$. L'opérateur $P(\ad(\epsilon))$, \resp $Q(\ad(\epsilon))$, est le projecteur sur l'espace $\frg_{\epsilon}$, \resp $\frg_{\neq1}$, relativement à cette décomposition.
On introduit des suites $(K_{\cF,n})_{n\in \NN}$ et $(\frk_{\cF,n})_{n\in \NN}$ comme en \ref{elementstopologiquement}. L'action $\ad(\epsilon)$ par conjugaison conserve $\frk_{\cF,n}$ pour tout $n$. On a donc
\begin{align*}
\frk_{\cF,n}&=\frk_{\cF,n,\epsilon}\oplus \frk_{\cF,n,\neq1},
\\
\tag*{où}
\frk_{\cF,n,\epsilon}&=\frk_{\cF,n}\cap \frg_{\epsilon}(F)=P(\ad(\epsilon))(\frk_{\cF,n})\\
\tag*{et}
\frk_{\cF,n,\neq1}&=\frk_{\cF,n}\cap \frg_{\neq1}(F)=Q(\ad(\epsilon))(\frk_{\cF,n}).
\end{align*}
Posons $\bar{\frk}_{\cF,n}=\frk_{\cF,n}/\frk_{\cF,n+1}$. L'opérateur $\ad(\epsilon)$ se réduit en un opérateur de cet espace, que l'on~note encore $\ad(\epsilon)$. On a encore
\begin{align*}
\tag{1}\label{commutant1}
\bar{\frk}_{\cF,n}&=\bar{\frk}_{\cF,n,\epsilon}\oplus \bar{\frk}_{\cF,n,\neq1},
\\
\tag*{où}
\bar{\frk}_{\cF,n,\epsilon}&=\frk_{\cF,n,\epsilon}/\frk_{\cF,n+1,\epsilon}=P(\ad(\epsilon))(\bar{\frk}_{\cF,n})\\
\tag*{et}
\bar{\frk}_{\cF,n,\neq1}&=\frk_{\cF,n,\neq1}/\frk_{\cF,n+1,\neq1}=Q(\ad(\epsilon))(\bar{\frk}_{\cF,n}).
\end{align*}

On va prouver par récurrence sur $n\geq1$ qu'il existe
\[
Y_{n}\in \frg_{\epsilon}(F)\cap \frk_{\cF}^+,\quad Z_{n}\in \frk_{\cF,n}\quadet k_{n}\in K_{\cF}^+
\]
tels que
\[
k_{n}^{-1}xk_{n}=\epsilon \exp(X+Y_{n})\exp(Z_{n}).
\]
Pour $n=1$, il suffit de prendre $Y_{1}=0$, $k_{1}=1$ et pour $Z_{1}$ l'élément de $\frk_{\cF}^+$ tel que $x=\epsilon \exp(X)\exp(Z_{1})$. Supposons ces termes définis au rang $n$. Posons
\[
x_{n}=\epsilon \exp(X+Y_{n}).
\]
Notons $\bar{Z}_{n}$ la réduction de $Z_{n}$ dans $\bar{\frk}_{\cF,n}$ et décomposons $\bar{Z}_{n}$ en $\bar{Z}_{n,\epsilon}+\bar{Z}_{n,\neq1}$ conformément à la décomposition \eqref{commutant1}. Relevons $\bar{Z}_{n,\epsilon}$ en un élément $Y'_{n}\in \frk_{\cF,n,\epsilon}$. Les opérateurs $\ad(\exp(X+Y_{n}))$ et $\ad(x_{n})$ conservent eux aussi les espaces $\frk_{\cF,n'}$ pour tout $n'$ et ils se descendent en des opérateurs sur $\bar{\frk}_{\cF,n}$. Parce que $\exp(X+Y_{n})$ commute à $\epsilon$, ces opérateurs commutent à $\ad(\epsilon)$ donc préservent la décomposition \eqref{commutant1}. Sur l'espace $\bar{\frk}_{\cF,\neq1}$, les valeurs propres de $\ad(\epsilon)$ sont toutes différentes de $1$ tandis que l'opérateur $\ad(X+Y_{n})$ est nilpotent. Il en résulte que $1-\ad(x_{n}^{-1})$ est inversible sur cet espace. On peut donc fixer un élément $Y''_{n}\in \frk_{\cF,n}$ tel que $\bar{Z}_{n,\neq1}$ soit la réduction de $(1-\ad(x_{n}^{-1}))(Y''_{n})$. Posons $h_{n}=\exp(Y''_{n})$. C'est un élément de $K_{\cF,n}$. On a
\[
h_{n}x_{n}\exp(Z_{n})h_{n}^{-1}=x_{n}\exp(\ad(x_{n}^{-1})(Y''_{n}))\exp(Z_{n})\exp(-Y''_{n}).
\]
Mais on voit que
\[
\exp(\ad(x_{n}^{-1})(Y''_{n}))\exp(Z_{n})\exp(-Y''_{n})\in \exp(Y'_{n})K_{\cF,n+1}.
\]
Il existe $Y_{n+1}\in \frg_{\epsilon}(F)\cap \frk_{\cF}^+ $ tel que $\exp(X+Y_{n})\exp(Y'_{n})=\exp(X+Y_{n+1})$. En posant $k_{n+1}=h_{n}k_{n}$, on obtient
$k_{n+1}^{-1}xk_{n+1}=\epsilon \exp(X+Y_{n+1})z_{n+1}$ avec $z_{n+1}\in K_{\cF,n+1}$ et il reste à prendre pour $Z_{n+1}$ l'élément de $\frk_{\cF,n+1}$ tel que $z_{n+1}=\exp(Z_{n+1})$ pour obtenir les éléments cherchés au rang $n+1$.

On peut extraire des suites $(Y_{n})_{n\geq1}$ et $(k_{n})_{n\geq1}$ des sous-suites convergentes (en fait, on n'en a même pas besoin, les suites définies ci-dessus convergent). À la limite, on obtient des éléments $Y\in \frg_{\epsilon}(F)\cap \frk_{\cF}^+$ et $k\in K_{\cF}^+$ tels que $k^{-1}xk=\epsilon \exp(X+Y)$. Cela achève la preuve.
\end{proof}

\Subsection{Un corollaire}\label{corollaire}
\begin{cor*}
Soit $\epsilon\in G(F)$ un élément $p'$-compact $\rmod Z(G)$, soit $\cF\in \Fac(G)$ et soit $g\in G(F)$. Supposons que $\epsilon\in K_{\cF}^{\dag}$ et que $g^{-1}\epsilon g$ appartienne à $\epsilon K_{\cF}^+$. Alors $g$ appartient à $Z_{G}(\epsilon)(F)K_{\cF}^+$.
\end{cor*}

\begin{proof}
D'après le lemme précédent, quitte à multiplier $g$ à droite par un élément de $K_{\cF}^+$, on peut supposer $g^{-1}\epsilon g=\epsilon \exp(Y)$ avec $Y\in \frg_{\epsilon,\tn}(F)$. Puisque~$\epsilon$ et $g^{-1}\epsilon g$ sont tous deux $p'$-compacts $\rmod Z(G)$, l'assertion \ref{elementsp}\,\eqref{elementsp2} entraîne que \hbox{$Y\in \frz(G)(F)$}. Donc $g$ appartient à l'image réciproque dans $G$ de $Z_{G_{\AD}}(\epsilon_{\ad})$. On a déjà dit que le groupe $Z_{G}(\epsilon)$ est d'indice fini premier à $p$ dans cette image réciproque. Il existe donc un entier $c\geq1$ premier à $p$ tel que $g^c\in Z_{G}(\epsilon)$. Mais, puisque~$Y$ est central, on~a $g^{-c}\epsilon g^c=\epsilon \exp(cY)$. Donc $Y=0$ et $g^{-1}\epsilon g=\epsilon$, c'est-à-dire $g\in Z_{G}(\epsilon)(F)$.
\end{proof}

\Subsection{Un lemme sur les éléments \texorpdfstring{$p'$}{p'}-compacts et les espaces de Levi}\label{p'unlemme}
\begin{lem*}
Soit $\epsilon\in G(F)$ un élément $p'$-compact $\rmod Z(G)$ et soit $(\cF,\nu)\in \Fac^*(G)$. Supposons $\epsilon\in K_{\cF}^{\nu}$. Notons $\bar{\epsilon}$ la réduction de $\epsilon$ dans $\bG_{\cF}^{\nu}(k_{F})$. Soit $\bP^{\nu}$ un espace parabolique de $\bG_{\cF}^{\nu}$. Supposons $\bar{\epsilon}\in \bP^{\nu}(k_{F})$. Alors il existe une composante de Levi~$\bM^{\nu}$ de $\bP^{\nu}$ telle que $\bar{\epsilon}\in \bM^{\nu}(k_{F})$.
\end{lem*}

La preuve s'inspire de celle de \cite[Lem.\,9]{P}.

\begin{proof}
On note $\bP$ le sous-groupe parabolique de $\bG_{\cF}$ associé à $\bP^{\nu}$. L'espace $\boldsymbol{\fru}_{\bP}$ possède une filtration finie $(\boldsymbol{\fru}_{i})_{i=1,\dots,n}$ définie par récurrence par $\boldsymbol{\fru}_{1}=\boldsymbol{\fru}_{\bP}$ et, pour $i\geq1$, $\boldsymbol{\fru}_{i+1}$ est l'espace engendré par les $[U_{1},U_{i}]$ pour $U_{1}\in \boldsymbol{\fru}_{1}$ et $U_{i}\in \boldsymbol{\fru}_{i}$. La filtration se termine par $\boldsymbol{\fru}_{n}=\{0\}$. On note $\bU_{i}=\exp(\boldsymbol{\fru}_{i})$. Les $\bU_{i}$ sont des sous-groupes distingués de $\bU_{\bP}$, les quotients $\bU_{i}/\bU_{i+1}$ sont abéliens et isomorphes aux $\boldsymbol{\fru}_{i}/\boldsymbol{\fru}_{i+1}$. On a $\bU_{n}=\{1\}$. Nous allons prouver par récurrence sur $i$ que
\begin{enumerate}
\item\label{p'unlemme1}
il existe une composante de Levi $\bM_{i}$ et un élément $u_{i}\in \bU_{i}(k_{F})$ tels que $\ad(\bar{\epsilon})(\bM_{i})=\ad(u_{i})(\bM_{i})$.
\end{enumerate}

Pour $i=1$, on choisit pour $\bM_{1}$ n'importe quelle composante de Levi de $\bP$. Puisque $\bar{\epsilon}\in \bP^{\nu}(k_{F})$, $\ad(\bar{\epsilon})(\bM_{1})$ est encore une telle composante de Levi de $\bP$ et on sait que deux telles composantes sont conjuguées par un élément de $\bU_{\bP}(k_{F})=\bU_{1}(k_{F})$. On~choisit pour $u_{1}$ l'élément qui conjugue $\bM_{1}$ en $\ad(\bar{\epsilon})(\bM_{1})$. Supposons le problème résolu au rang $i\geq1$.
Fixons un entier $c\geq1$ premier à $p$ tel que $\epsilon^c\in Z(G)$. On voit par récurrence sur un entier $m\geq1$ que
\[
\ad(\bar{\epsilon}^m)(\bM_{i})=\ad( u_{i,m})(\bM_{i}),
\]
où $u_{i,m}=\ad(\bar{\epsilon}^{m-1})(u_{i})\cdots\ad(\bar{\epsilon})(u_{i})u_{i}$. Pour $m=c$, on obtient $\bM_{i}=\ad(u_{i,c})(\bM_{i})$, donc $u_{i,c}=1$. Remarquons que $\ad(\bar{\epsilon})$ conserve la filtration $(\bU_{j})_{j=1,\dots,n}$. En notant~$X_{i}$ la réduction de $u_{i}$ dans $\bU_{i}(k_{F})/\bU_{i+1}(k_{F})$ et en écrivant ce groupe additivement, on obtient
\[
\ad(\bar{\epsilon}^{c-1})(X_{i})+\cdots+\ad(\bar{\epsilon})(X_{i})+X_{i}=0.
\]
Par le même argument que dans la preuve de \ref{commutant}, cela entraîne l'existence d'un élément $Y_{i}\in \bU_{i}(k_{F})/\bU_{i+1}(k_{F})$ tel que $X_{i}=Y_{i}-\ad(\bar{\epsilon})(Y_{i})$. On relève $Y_{i}$ en $v_{i}\in \bU_{i}(k_{F})$ et on pose $\bM_{i+1}=\ad(v_{i})(\bM_{i})$. On calcule
\[
\ad(\bar{\epsilon})(\bM_{i+1})=\ad(u_{i+1})(\bM_{i+1}),
\]
où $u_{i+1}=\ad(\bar{\epsilon})(v_{i})u_{i}v_{i}^{-1}$. Par construction de $v_{i}$, on voit que $u_{i+1}\in \bU_{i+1}(k_{F})$, ce~qui résout le problème en $i+1$.

Pour $i=n$, le Levi $\bM=\bM_{n}$ est conservé par $\ad(\bar{\epsilon})$. Alors $\bM^{\nu}=\bar{\epsilon}\bM$ est une composante de Levi de $\bP^{\nu}$ contenant $\bar{\epsilon}$.
\end{proof}

\Subsection{Points fixes dans \texorpdfstring{$\Imm(G_{\AD})$ d'un élément $p'$-compact mod\,$Z(G)$}{IGZG}}\label{pointsfixes}
Soit $\epsilon \in G(F)$ un élément $p'$-compact $\rmod Z(G)$. Notons $M$ le commutant dans~$G$ du tore déployé $A_{G_{\epsilon}}$. C'est un Levi de $G$. On a $\epsilon\in M$ et $G_{\epsilon}\subset M$ ($\epsilon$ et $G_{\epsilon}$ commutent à~$A_{G_{\epsilon}}$). On a aussi $A_{M}=A_{G_{\epsilon}}$. En effet, puisque $G_{\epsilon}\subset M$, on~a $ A_{G_{\epsilon}}\subset M$~et, puisque~$M$ commute à $A_{G_{\epsilon}}$ par construction de $M$, on~a $A_{G_{\epsilon}}\subset A_{M}$. Inversement, puisque $\epsilon\in M$, on~a $A_{M}\subset G_{\epsilon}$ et, puisque~$A_{M}$ commute à $M$, donc aussi au sous-ensemble $G_{\epsilon}$, on~a $A_{M}\subset A_{G_{\epsilon}}$.

Comme en \ref{groupesdelevi}, notons $\Imm(G_{\AD},M)$ la réunion dans $\Imm(G_{\AD})$ des appartements $\App(A_{M'})$ associés aux Levi minimaux $M'$ de $M$. L'action de $M(F)$ sur $\Imm(G_{\AD})$ conserve le sous-ensemble $\Imm(G_{\AD},M)$, il en est donc de même des actions de $\epsilon$ et de $G_{\epsilon}(F)$. L'espace vectoriel $\cA_{M} /\cA_{G}$ agit naturellement sur $\Imm(G_{\AD},M)$ (il agit sur chaque $\App(A_{M'})$ et ces actions se recollent). Cette action commute à celles de $\epsilon$ et de $G_{\epsilon}(F)$. Notons $\Imm(G_{\AD})^{\epsilon}$ et $\Imm(G_{\AD},M)^{\epsilon}$ les sous-ensembles de points fixes de l'action de $\epsilon$ dans $\Imm(G_{\AD})$, \resp $\Imm(G_{\AD},M)$. L'ensemble $\Imm(G_{\AD},M)^{\epsilon}$ est stable par l'action de $\cA_{M}/\cA_{G}$. On note $\Imm(G_{\AD},M)^{\epsilon}/(\cA_{M}/\cA_{G})$ l'ensemble quotient. L'action de $G_{\epsilon}(F)$ se descend en une action sur ce quotient.

\skpt
\begin{prop*}
\begin{enumeratei}
\item\label{pointsfixesi}
On a l'égalité $\Imm(G_{\AD})^{\epsilon}=\Imm(G_{\AD},M)^{\epsilon}$.

\item\label{pointsfixesii}
L'ensemble $\Imm(G_{\AD},M)^{\epsilon}/(\cA_{M}/\cA_{G})$ muni de son action de $G_{\epsilon}(F)$ s'identifie canoniquement à $\Imm(G_{\epsilon,\AD})$.

\item\label{pointsfixesiii}
Soit $\cF\in \Fac(G)$ telle que
\[
\cF\cap \Imm(G_{\AD})^{\epsilon}\neq\emptyset.
\]
Alors l'image de $\cF\cap \Imm(G_{\AD})^{\epsilon}$ est contenue dans une facette $\cF'\in \Fac(G_{\epsilon})$. L'action naturelle de $\epsilon$ sur $K_{\cF}^0$ se descend en une action algébrique sur $\bG_{\cF}$ et le groupe $\bG_{\epsilon,\cF'}$ s'identifie à la composante neutre $\bG_{\cF}^{\epsilon,0}$ du sous-groupe des points fixes $\bG_{\cF}^{\epsilon}$ par cette action. Le~groupe $K_{\cF'}^0$ est le groupe des $g\in G_{\epsilon}(F)\cap K_{\cF}^0$ tels que $w_{G_{\epsilon}}(g)=0$. On a
\[
K_{\cF'}^+=G_{\epsilon}(F)\cap K_{\cF}^+,\quad \frk_{\cF'}=\frg_{\epsilon}(F)\cap \frk_{\cF},\quad \frk_{\cF'}^+=\frg_{\epsilon}(F)\cap \frk_{\cF}^+.
\]
\end{enumeratei}
\end{prop*}

\begin{proof}
Dans la seconde preuve de \cite[Th.\,1.9]{PY}, G.\,Prasad et J.-K.\,Yu démon\-trent que $\Imm(G_{\AD})^{\epsilon}=\Imm(G_{\AD},M)^{\epsilon}$ et que cet ensemble, muni de son action de $G_{\epsilon}(F)$, s'identifie à l'immeuble étendu du groupe $G_{\epsilon,\ad}=G_{\epsilon}/Z(G)$. L'identification est canonique à translations près par les éléments de $\cA_{M}/\cA_{G}$. Cela équivaut aux assertions \eqref{pointsfixesi} et \eqref{pointsfixesii}. Soit $\cF\in \Fac(G)$ telle que $\cF\cap \Imm(G_{\AD})^{\epsilon}\neq\emptyset$. Soit $x\in \cF\cap \Imm(G_{\AD})^{\epsilon}\neq\emptyset$, notons $y$ son image dans $\Imm(G_{\epsilon,\AD})$. Soit $\cF'\in \Fac(G_{\epsilon})$ la facette à laquelle appartient $y$. Plongeons les immeubles pour les groupes adjoints dans les immeubles étendus. Comme on le sait, on peut définir un schéma en groupes~$\cG_{x}$ défini sur $\fro_{F}$ satisfaisant entre autres que $\cG_{x}(\fro_{F})$ est le sous-groupe des éléments de $G(F)$ dont l'action sur l'immeuble étendu $\Imm(G)$ fixe $x$. On sait que la partie réductive de la composante neutre de sa fibre spéciale n'est autre que $\bG_{\cF}$. Dans \cite[\S3]{P}, \cf en particulier les définitions de 3.1 de cette référence, Prasad montre que~$\epsilon$ agit naturellement sur $\cG_{x}$ et il définit la composante neutre $\cG_{x}^{\epsilon,0}$ du sous-schéma des points fixes $\cG_{x}^{\epsilon}$. Il prouve que la composante neutre $\cG_{\epsilon,y}^{0}$ de $\cG_{\epsilon,y}$ s'identifie à ce schéma $\cG_{x}^{\epsilon,0}$. En passant aux parties réductives des fibres spéciales, on obtient que~$\bG_{\epsilon,\cF'}$ s'identifie à $\bG_{\cF}^{\epsilon,0}$. Cela entraîne
\begin{equation}\tag{1}\label{pointsfixes1}
K_{\cF'}^0\subset K_{\cF}^0\cap G_{\epsilon}(F).
\end{equation}
En passant aux algèbres de Lie, l'égalité $\cG_{\epsilon,y}^{0}=\cG_{x}^{\epsilon,0}$ entraîne aussi
\begin{equation}\tag{2}\label{pointsfixes2}
\frk_{\cF'}=\frk_{\cF}\cap \frg_{\epsilon}(F),\quad \frk^+_{\cF'}=\frk^+_{\cF}\cap \frg_{\epsilon}(F).
\end{equation}

D'après \ref{description}\,\eqref{description5}, ces égalités entraînent que $\cF'$ ne dépend pas de $x$ et est uniquement déterminée, ce qui est la première assertion de \eqref{pointsfixesiii}.
L'égalité
\begin{equation}\tag{3}\label{pointsfixes3}
K_{\cF'}^+=K_{\cF}^+\cap G_{\epsilon}(F)
\end{equation}
se déduit par l'exponentielle de la seconde égalité de \eqref{pointsfixes2}.

On peut préciser \eqref{pointsfixes1}. D'après la définition de $K_{\cF'}^0$, $K_{\cF'}^0$ est contenu dans l'ensemble des $g\in K_{\cF}^0\cap G_{\epsilon}(F) $ tels que $w_{G_{\epsilon}}(g)=0$. Inversement, un tel $g$ fixe $x$ donc aussi~$y$ (la~projection de $\Imm(G_{\AD})^{\epsilon}$ sur $\Imm(G_{\epsilon,\AD})$ étant compatible avec l'action de $G_{\epsilon}(F)$). Donc $g\in K_{\cF'}^{\dag}$. Puisque $w_{G_{\epsilon}}(g)=0$, on~a $g\in K_{\cF'}^0$. Cela démontre l'assertion de l'énoncé concernant $K_{\cF'}^0$.
\end{proof}

\section{Quasi-caractères}
\subsection{Transformées de Fourier et intégrales orbitales}\label{transformees}
Plusieurs notions que l'on~a introduites sur le groupe $G(F)$ ont des analogues sur l'algèbre de Lie $\frg(F)$. Pour $f\in C_{\rc}^{\infty}(\frg(F))$ et $X\in \frg_{\reg}(F)$, on définit l'intégrale orbitale $I^G(X,f)$. On définit aussi l'espace $\rI(\frg)$ quotient de $C_{\rc}^{\infty}(\frg(F))$ par le sous-espace des fonctions dont toutes les intégrales orbitales sont nulles.

On fixe un caractère continu $\psi$ de $F$ de conducteur $\frp_{F}$ et une forme bilinéaire symétrique non dégénérée $\<{\cdot},{\cdot}\>$ sur $\frg(F)$, invariante par conjugaison par $G(F)$. On sait que l'on peut la choisir telle que, pour toute $\cF\in \Fac(G)$, $\frk_{\cF}^+$ soit le \og dual\fg de~$\frk_{\cF}$, c'est-à-dire que $\frk_{\cF}^+$ est l'ensemble des $X\in \frg(F)$ tels que $\<X,Y\>\in \frp_{F}$ pour tout $Y\in \frk_{\cF}$. On suppose qu'il en est ainsi. On définit la transformation de Fourier $f\mto \hat{f}$ dans $C_{\rc}^{\infty}(\frg(F))$ par
\[
\hat{f}(X)=\int_{\frg(F)}f(Y)\psi(\<X,Y\>)\,dY,
\]
où $dY$ est la mesure auto-duale. La transformation $f\mto \hat{f}$ de $C_{\rc}^{\infty}(\frg(F))$ se descend en une transformation de $\rI(\frg)$.

Notons $\Nil(\frg)$ l'ensemble des orbites nilpotentes dans $\frg(F)$. Pour tout $\cO\in \Nil(\frg)$, on fixe une mesure sur $\cO$ invariante par conjugaison. Cela permet de définir l'inté\-grale orbitale $I_{\cO}$ sur $C_{\rc}^{\infty}(\frg(F))$, puis sa transformée de Fourier $f\mto I_{\cO}(\hat{f})$. \hbox{Notons}~$\frg_{\reg}$ l'ensemble des éléments semi-simples et réguliers de $\frg$. D'après Harish-Chandra (\hbox{\cite[Th.\,4.4]{HC}}), il existe une fonction $\hat{j}(\cO)$ localement intégrable sur $\frg(F)$ et localement constante sur $\frg_{\reg}(F)$, de sorte que
\[
I_{\cO}(\hat{f})=\int_{\frg(F)}f(Y)\hat{j}(\cO,Y)\,dY
\]
pour toute $f\in C_{\rc}^{\infty}(\frg(F))$.

Notons $E$ l'espace des fonctions sur $\fro_{F}$ engendré par les fonctions $\lambda\mto \vert \lambda\vert ^n$ pour $n\in \ZZ$. Il est immédiat que, si deux éléments de $E$ coïncident sur $\frp_{F}^n$ pour un entier $n\geq0$, alors, ils sont égaux.

On sait que, pour tout $\cO\in \Nil(\frg)$ et pour tout $Y\in \frg_{\reg}(F)$, la fonction $\lambda\mto \hat{j}(\cO,\lambda^2 Y)$, définie sur $\fro_{F}$, appartient à $E$.

\subsection{Quasi-caractères et quasi-caractères de niveau~\texorpdfstring{$0$}{0}}\label{quasicaracteres}

Si $\theta$ est une fonction localement intégrable sur $G(F)$ et invariante par conjugaison, il lui est associée la distribution invariante $D$ sur $G(F)$ définie par
\[
D(f)=\int_{G(F)}f(g)\theta(g)\,dg.
\]

On appelle quasi-caractère de $G(F)$ une distribution invariante $D$ sur $G(F)$ associée à une fonction $\theta_{D}$ sur $G(F)$ localement intégrable et invariante par conjugaison, qui satisfait à la propriété suivante:
\begin{enumerate}
\item\label{quasicaracteres1}
pour tout $x\in G(F)$ semi-simple, il existe un voisinage $\frV$ de $0$ dans $\frg_{x}(F)$ et, pour tout $\cO\in \Nil(\frg_{x})$, il existe un nombre complexe $c_{D,\cO}$ de sorte que, pour presque tout $Y\in \frV$, on ait l'égalité
\[
\theta_{D}(x \exp(Y))=\sum_{\cO\in \Nil(\frg_{x})}c_{D,\cO}\hat{j}(\cO,Y).
\]
\end{enumerate}

\skpt
\begin{rems*}
\begin{enumerate}\setcounter{enumi}{1}
\item\label{quasicaracteres2}
Le voisinage $\frV$ n'est évidemment pas uniquement déterminé, par contre les constantes $c_{D,\cO}$ le sont car les fonctions $\hat{j}(\cO)$ sont linéairement indépendantes dans tout voisinage de $0$.

\item\label{quasicaracteres3}
Appliquée à $x$ fortement régulier, cette propriété implique que, quitte à modifier $\theta_{D}$ sur un ensemble de mesure nulle, on peut supposer $\theta_{D}$ définie et localement constante sur $G_{\reg}(F)$.

\item\label{quasicaracteres4}
La condition que $\theta_{D}$ est localement intégrable est redondante car une fonction invariante par conjugaison et satisfaisant à la condition \eqref{quasicaracteres1} est automatiquement localement intégrable.

\item\label{quasicaracteres5}
Soit $D\in \rI(G)^*$ et soit $\alpha$ une fonction sur $G(F)$, localement constante et invariante par conjugaison. On définit la distribution $\alpha D$ par $(\alpha D)(f)=D(\alpha f)$. Si $D$ est un quasi-caractère, alors $\alpha D$ l'est aussi.
\end{enumerate}
\end{rems*}

Soit $f\in C_{\rc}^{\infty}(G(F))$ une fonction très cuspidale. On a défini la distribution $D_{f}$ en~\ref{lesespaces}. On a:
\begin{enumerate}\setcounter{enumi}{5}
\item\label{quasicaracteres6}
$D_{f}$ est un quasi-caractère.
\end{enumerate}
\Cf \cite[5.9]{W1}. La fonction $\theta_{D_{f}}$ se calcule de la façon suivante. Soit $x\in G_{\reg}(F)$. Notons~$M^x$ le commutant de $A_{G_{x}}$ dans $G$. C'est le plus petit Levi $M'$ tel que $x\in M'(F)$. De plus $x$ est elliptique dans $M^{x}(F)$. On sait définir l'intégrale orbitale pondérée
\[
J^G_{M^{x}}(x,f)=D^G(x)^{1/2}\int_{A_{M^{x}}(F)\backslash G(F)}f(g^{-1}xg)v_{M^{x}}(g)\,dg.
\]
Le poids $v_{M^{x}}$ est calculé relativement à un sous-groupe compact spécial de $G(F)$ fixé, \cf \cite{A1}, mais, parce que $f$ est très cuspidale, l'intégrale ci-dessus ne dépend pas de ce choix, \cf \cite[Lem.\,5.2]{W1}. Pour tout tore déployé $A'$, posons $m(A')=\mes(A'(F)_{\rc})$. D'après \cite[9(1)]{W2}, on~a alors
\begin{equation}\tag{7}\label{quasicaracteres7}
\theta_{D_{f}}(x)=(-1)^{a_{M^{x}}-a_{G}}D^G(x)^{-1/2}m(A_{M^{x}})m(A_{G})^{-1}J^G_{M^{x}}(x,f).
\end{equation}

On appelle quasi-caractère de niveau~$0$ une distribution invariante $D$ associée à une fonction $\theta_{D}$ sur $G(F)$ localement intégrable et invariante par conjugaison, qui satisfait à la condition
\begin{enumerate}\setcounter{enumi}{7}
\item\label{quasicaracteres8}
pour tout $\epsilon\in G(F)_{p'}$, pour tout $\cO\in \Nil(\frg_{\epsilon})$, il existe un nombre complexe $c_{D,\cO}$ de sorte que, pour presque tout $Y\in \frg_{\epsilon,\tn}(F)$, on ait l'égalité
\[
\theta_{D}(\epsilon \exp(Y))=\sum_{\cO\in \Nil(\frg_{\epsilon})}c_{D,\cO}\hat{j}(\cO,Y).
\]
\end{enumerate}

On a des variantes de ces définitions pour l'algèbre de Lie. Soit $D$ une distribution invariante sur $\frg(F)$ associée à une fonction $\theta_{D}$ sur $\frg(F)$ localement intégrable et invariante par conjugaison. On dit que $D$ est un quasi-caractère si $\theta_{D}$ satisfait à la condition
\begin{enumerate}\setcounter{enumi}{8}
\item\label{quasicaracteres9}
pour tout $X\in \frg(F)$ semi-simple, il existe un voisinage $\frV$ de $0$ dans $\frg_{X}(F)$ et, pour tout $\cO\in \Nil(\frg_{X})$, il existe un nombre complexe $c_{D,\cO}$ de sorte que, pour presque tout $Y\in \frV$, on ait l'égalité
\[
\theta_{D}(X+Y)=\sum_{\cO\in \Nil(\frg_{X})}c_{D,\cO}\hat{j}(\cO,Y).
\]
\end{enumerate}

On dit que c'est un quasi-caractère de niveau~$0$ si $\theta_{D}$ est à support dans $\frg_{\tn}(F)$ et que, pour tout $\cO\in \Nil(\frg)$, il existe un nombre complexe $c_{D,\cO}$ de sorte que, pour presque tout $Y\in \frg_{\tn}(F)$, on ait l'égalité
\[
\theta_{D}(Y)=\sum_{\cO\in \Nil(\frg)}c_{D,\cO}\hat{j}(\cO,Y).
\]

Si $f\in C_{\rc}^{\infty}(\frg(F))$ est une fonction très cuspidale, on définit la distribution $D_{f}$ sur $\frg(F)$ comme on l'a fait sur le groupe. Cette distribution est un quasi-caractère.

\skpt
\begin{lem*}
\begin{enumeratei}
\item\label{quasicaracteresi}
Tout quasi-caractère sur $G(F)$, \resp $\frg(F)$, de niveau~$0$ est un quasi-caractère.

\item\label{quasicaracteresii}
soit $D$ un quasi-caractère sur $G(F)$; alors $D$ est de niveau~$0$ si et seulement si, pour tout $\epsilon\in G(F)_{p'}$ et presque tout $X\in \frg_{\epsilon,\tn}(F)$, la fonction $\lambda\mto \theta_{D}(\epsilon \exp(\lambda^2X))$ définie sur $\fro_{F}$ appartient à $E$, \cf \ref{transformees} .

\item\label{quasicaracteresiii}
Fixons une décomposition $G(F)=\bigcup_{\epsilon\in B}C_{G}(\epsilon)$, où $B$ est un sous-ensemble de $G(F)_{p'}$. Soit $D$ une distribution invariante associée à une fonction $\theta_{D}$ qui satisfait à la condition \eqref{quasicaracteres8} restreinte aux $\epsilon\in B$. Alors $D$ est un quasi-caractère de niveau~$0$.
\end{enumeratei}
\end{lem*}

\begin{proof}
Il résulte de \cite[Lem.\,6.3(iii)]{W1} que, pour tout quasi-caractère $D$ de niveau~$0$ sur $\frg(F)$, il existe une fonction très cuspidale $f$ telle que $D$ coïncide avec $D_{f}$ sur $\frg_{\tn}(F)$. Puisque $D_{f}$ est un quasi-caractère, $D$ l'est aussi, \cf remarque \eqref{quasicaracteres5}.

Soit maintenant $D$ un quasi-caractère de niveau~$0$ sur $G(F)$. Soit $x\in G(F)$ un élément semi-simple. Écrivons une $p'$-décomposition $x=\epsilon \exp(X)$ avec $X\in \frg_{\epsilon,\tn}(F)$. Soit $Y\in \frg_{x}(F)$. D'après \ref{pelements}\,\eqref{pelements1}, on~a $Y\in \frg_{\epsilon}(F)$ et $Y$ commute à $X$. Donc $x\exp(Y)=\epsilon \exp(X+Y)$. Si $Y$ est assez petit, $X+Y$ est topologiquement nilpotent. Par hypothèse sur $D$, on~a donc
\[
\theta_{D}(x \exp(Y))=\theta_{D}(\epsilon \exp(X+Y))=\sum_{\cO\in \Nil(\frg_{\epsilon})}c(D,\cO)\hat{j}(\cO,X+Y).
\]
Notons $D'$ le quasi-caractère de niveau~$0$ sur $\frg_{\epsilon}(F)$ dont la fonction $\theta_{D'}$ associée est définie par
\[
\theta_{D'}(Z)= \sum_{\cO\in \Nil(\frg_{\epsilon})}c(D,\cO)\hat{j}(\cO,Z)
\]
pour $Z\in \frg_{\epsilon,\tn}(F)$. On vient de voir que c'est un quasi-caractère. L'hypothèse que $x$ est semi-simple implique que $X$ l'est aussi. Il existe donc un voisinage $\frV$ de $0$ dans $(\frg_{\epsilon})_{X}(F)$ tel que, pour presque tout $Y\in \frV$, on ait l'égalité
\[
\theta_{D'}(X+Y)=\sum_{\cO\in \Nil((\frg_{\epsilon})_{X})}c_{D',\cO}\hat{j}(\cO,Y).
\]
Autrement dit, $\theta_{D}(x \exp(Y))$ est égal à l'expression de droite pour presque tout \hbox{$Y\!\in\!\frV$}. Puisque $(\frg_{\epsilon})_{X}=\frg_{x}$, c'est exactement la condition \eqref{quasicaracteres1} requise. Donc $D$ est un quasi-caractère, ce qui démontre \eqref{quasicaracteresi}.

Soit $D$ un quasi-caractère sur $G(F)$. Si $D$ est de niveau~$0$, la propriété énoncée au~\eqref{quasicaracteresii} résulte de ce que les fonctions $\lambda\mto \hat{j}(\cO,\lambda^2X)$ appartiennent à $E$. Inversement, supposons la propriété en question satisfaite. Soit $\epsilon\in G(F)_{p'}$, fixons un voisinage $\frV$ et des constantes $c_{D,\cO}$, de sorte que \eqref{quasicaracteres1} soit satisfaite. Soit $X\in \frg_{\epsilon,\tn}(F)$ et supposons $\epsilon \exp(\lambda X)\in G_{\reg}$ pour tout $\lambda\in \fro_{F}$ (ceci est vérifié pour presque tout $X$). Les fonctions
\begin{align*}
\lambda&\mto \theta_{D}(\epsilon \exp(\lambda^2X))
\\
\tag*{et}
\lambda&\mto \sum_{\cO\in \Nil(\frg_{\epsilon})}c_{D,\cO}\hat{j}(\cO,\lambda^2X)
\end{align*}
appartiennent toutes deux à $E$ (la première par hypothèse). Pour un entier $n$ assez grand, $\lambda^2 X$ appartient à $\frV$ pour tout $\lambda\in \frp_{F}^n$, donc les deux fonctions coïncident sur~$\frp_{F}^n$. Elles sont alors égales. Pour $\lambda=1$, cela démontre l'égalité
\[
\theta_{D}(\epsilon \exp(X))=\sum_{\cO\in \Nil(\frg_{\epsilon})}c_{D,\cO}\hat{j}(\cO,X)
\]
et cela pour presque tout $X\in \frg_{\epsilon,\tn}(F)$. Donc $D$ est un quasi-caractère de niveau~$0$, ce qui démontre \eqref{quasicaracteresii}.

Pour \eqref{quasicaracteresiii}, il s'agit de voir que \eqref{quasicaracteres8} est aussi satisfaite pour $\epsilon\in G(F)_{p'}-B$. Ce problème étant insensible à la conjugaison par $G(F)$, on peut supposer que $\epsilon\in C(\epsilon')$ pour un $\epsilon'\in B$. Écrivons $\epsilon=\epsilon'\exp(Z)$, avec $Z\in \frg_{\epsilon',\tn}(F)$. D'après le~\eqref{p'decompositioniv} du lemme~\ref{p'decomposition}, on~a $G_{\epsilon}=G_{\epsilon'}$ et $Z$ est un élément central dans $\frg_{\epsilon}(F)$. Pour $Y\in \frg_{\epsilon,\tn}(F)$, on~a alors $\epsilon \exp(Y)=\epsilon' \exp(Z+Y)$ et, en appliquant l'hypothèse \eqref{quasicaracteres8} pour $\epsilon'$, on a
\[
\theta_{D}(\epsilon \exp(Y))=\sum_{\cO\in \Nil(\frg_{\epsilon})}c_{D,\cO}\hat{j}(\cO,Z+Y).
\]
Mais les fonctions $\hat{j}(\cO)$ sont invariantes par translations par tout élément central dans $\frg_{\epsilon}(F)$. Le $Z$ disparaît de l'expression ci-dessus et on aboutit à une expression comme en \eqref{quasicaracteres8} de $\theta_{D}(\epsilon \exp(Y))$.
\end{proof}

\skpt
\begin{rem*}
\begin{enumerate}\setcounter{enumi}{9}
\item
Cette preuve et la remarque \eqref{quasicaracteres4} entraînent qu'une fonction $\theta_{D}$ invariante par conjugaison par $G(F)$ et satisfaisant à la condition \eqref{quasicaracteres8} est forcément localement intégrable.
\end{enumerate}
\end{rem*}

\subsection{Variante}\label{variante1}
Soit $\underline{Z}$ un sous-groupe algébrique de $Z(G)$ tel que $Z(G)(F)/ \underline{Z}(F)$ soit compact.

\begin{lem*}
Soit $\theta$ une fonction définie presque partout sur $G_{\comp}(F)$, invariante par conjugaison et satisfaisant à la condition suivante:
\begin{itemize}
\item
pour tout tout élément $\epsilon\in G(F)$ qui est $p'$-compact $\rmod\underline{Z}$, pour tout $\cO\in \Nil(\frg_{\epsilon})$, il existe un nombre complexe $c_{\cO}$ de sorte que, pour presque tout $Y\in \frg_{\epsilon,\tn}(F)$, on ait l'égalité
\[
\theta(\epsilon \exp(Y))=\sum_{\cO\in \Nil(\frg_{\epsilon})}c_{\cO}\hat{j}(\cO,Y).
\]

\end{itemize}
Alors $\theta$ est la fonction associée à un quasi-caractère de niveau~$0$.
\end{lem*}

\begin{proof}
C'est un cas particulier du \eqref{quasicaracteresiii} du lemme \ref{quasicaracteres}: on prend pour ensemble $B$ la réunion de celui des éléments de $G(F)_{p'}$ qui ne sont pas $p'$-compacts $\rmod Z(G)$ et de l'ensemble des éléments $p'$-compacts $\rmod \underline{Z}$. On a bien $G(F)=\bigcup_{\epsilon\in B}C_{G}(\epsilon)$ d'après \ref{elementsp}\,\eqref{elementsp3}. Puisque $\theta$ est à support dans $G_{\comp}(F)$, cette fonction est nulle sur $C_{G}(\epsilon)$ si $\epsilon$ n'est pas $p'$-compact $\rmod Z(G)$. Elle satisfait par hypothèse le développement requis sur $C_{G}(\epsilon)$ si $\epsilon$ est $p'$-compact $\rmod \underline{Z}$.
\end{proof}

\Subsection{Induction de quasi-caractères}\label{induction}

\begin{lem*}
Soient $M$ un Levi de $G$ et $D^M$ un quasi-caractère de $M(F)$. Posons $D= \Ind_{M}^G(D^M)$ .
\begin{enumeratei}
\item\label{inductioni}
La distribution $D$ est un quasi-caractère.

\item\label{inductionii}
Si $D^M$ est de niveau~$0$, $D$ est de niveau~$0$.
\end{enumeratei}
\end{lem*}

Preuve. Le \eqref{inductioni} est démontré dans \cite[Lem.\,2.3]{W3}. On note $\theta_{D^M}$ la fonction associée à~$D^M$. Pour $x\in G_{\reg}(F)$, l'ensemble $\{g^{-1}xg\sep g\in G(F)\cap M(F)\}$ se décompose en un nombre fini de classes de conjugaison par $M(F)$. Fixons un ensemble $X_{M}(x)$ de représentants de ces classes. Par un calcul d'intégration facile, $D$ est associé à la fonction $\theta_{D}$ définie sur $G_{\reg}(F)$ par la formule\vspace*{-2pt}\enlargethispage{.5\baselineskip}%
\begin{equation}\tag{1}\label{induction1}
\theta_{D}(x)=\sum_{y\in X_{M}(x)}D^G(x)^{-1/2}D^M(y)^{1/2}\theta_{D^M}(y).
\end{equation}
Soient $\epsilon\in G(F)_{p'}$ et $X\in \frg_{\epsilon,\tn}(F)$. Posons $x=\epsilon \exp(X)$ et supposons $x\in G_{\reg}(F)$. Pour tout $y\in X_{M}(x)$, soit $g_{y}\in G(F)$ tel que $g_{y}^{-1}xg_{y}=y$. Posons $\epsilon_{y}=g_{y}^{-1}\epsilon g_{y}$ et $X_{y}=g_{y}^{-1}X g_{y}$. On a $y=\epsilon_{y}\exp(X_{y})$ et ceci est une $p'$-décomposition de $y$. Le~grou\-pe~$A_{M}$ commute à $y$ puisque $y\in M(F)$. Il commute à $\epsilon_{y}$ et $X_{y}$ d'après le~\eqref{pelementsiii} du lemme \ref{pelements}. Donc $\epsilon_{y}\in M(F)$ et $X_{y}\in \mfrm_{\epsilon_{y}}(F)$. De plus $\epsilon_{y}$ est un $p'$-élément dans $M(F)$ d'après le~\eqref{lecasdunlevii} du lemme \ref{lecasdunlevi}. Soit $\lambda\in F^{\times}$ tel que $\lambda X$ soit encore topologiquement nilpotent. Posons $x'=\epsilon \exp(\lambda X)$. L'élément $x'$ appartient encore à $G_{\reg}(F)$. Les éléments $\epsilon_{y}\exp(\lambda X_{y})=g_{y}^{-1}x'g_{y}$ sont tous des éléments de $M(F)$ conjugués à $x'$ par un élément de $G(F)$. Montrons que
\begin{enumerate}\setcounter{enumi}{1}
\item\label{induction2}
si $y,y'\in X_{M}(x)$, avec $y\neq y'$, alors $\epsilon_{y}\exp(\lambda X_{y})$ et $\epsilon_{y'}\exp(\lambda X_{y'})$ ne sont pas conjugués par un élément de $M(F)$.
\end{enumerate}

\begin{proof}
Supposons qu'il existe $m\in M(F)$ tel que $m^{-1}\epsilon_{y}\exp(\lambda X_{y})m=\epsilon_{y'}\exp(\lambda X_{y'})$. En posant $h=g_{y}mg_{y'}^{-1}$, on~a alors $h^{-1}x'h=x'$. Puisque $x'\in G_{\reg}(F)$, on~a $h\in G_{x'}(F)$, d'où $h\in G_{\epsilon}(F)$ d'après le~\eqref{pelementsiii} du lemme \ref{pelements}. On voit alors que $m^{-1}\epsilon_{y}m=\epsilon_{y'}$, d'où aussi $m^{-1}\lambda X_{y}m=\lambda X_{y'}$. Il en résulte que $m^{-1}X_{y}m =X_{y'}$, puis que $m^{-1}ym=y'$. Mais cela est contradictoire avec la définition de $X_{M}(x)$, ce qui démontre \eqref{induction2}.

En conséquence de \eqref{induction2}, l'ensemble $X_{M}(x')$ a au moins autant d'éléments que $X_{M}(x)$. La situation étant symétrique en $x$ et $x'$, ces deux ensembles ont même nombre d'éléments. Alors \eqref{induction2} nous dit que l'on peut choisir pour $X_{M}(x')$ l'ensemble $\{\epsilon_{y}\exp(\lambda X_{y})\sep y\in X_{M}(x)\}$. On remplace maintenant $\lambda$ par $\lambda^2$, avec $\lambda\in \fro_{F}$. L'égalité~\eqref{induction1} pour $\epsilon \exp(\lambda^2X)$ devient\vspace*{-5pt}
\begin{multline*}
\theta_{D}(\epsilon \exp(\lambda^2X))\\[-3pt]
=\sum_{y\in X_{M}(x)}D^G(\epsilon \exp(\lambda^2X))^{-1/2}D^M(\epsilon_{y}\exp(\lambda^2X_{y}))^{1/2}\theta_{D^M}(\epsilon_{y}\exp(\lambda^2X_{y})).
\end{multline*}
D'après le~\eqref{p'decompositionii} du lemme \ref{p'decomposition}, cela se récrit\vspace*{-5pt}
\begin{multline*}
\theta_{D}(\epsilon \exp(\lambda^2X))\\[-3pt]
=\sum_{y\in X_{M}(x)}d^G(\epsilon)^{-1/2}D^{G_{\epsilon}}(\lambda^2X)^{-1/2} d^M(\epsilon_{y})^{1/2}D^{M_{\epsilon_{y}}}(\lambda^2X_{y})^{1/2}\theta_{D^M}(\epsilon_{y}\exp(\lambda^2X_{y})).
\end{multline*}
Comme fonction de $\lambda$, le terme $D^{G_{\epsilon_{y}}}(\lambda^2X)^{-1/2}$, \resp $D^{M_{\epsilon_{y}}}(\lambda^2X_{y})^{1/2}$, est produit d'une constante et d'une puissance entière de $\vert \lambda\vert _{F}$. Si $D^M$ est de niveau~$0$, le terme $ \theta_{D^M}(\epsilon_{y}\exp(\lambda^2X_{y}))$ appartient à $E$ d'après le~\eqref{quasicaracteresii} du lemme \ref{quasicaracteres}. Donc la fonction $\lambda\mto \theta_{D}(\epsilon \exp(\lambda^2X))$ appartient à $E$, ce qui prouve que $D$ est un quasi-caractère de niveau~$0$ d'après le même lemme.
\end{proof}

Un cas particulier de la formule~\eqref{induction1} nous servira plus tard. Supposons
\[
x\in G_{\reg}(F)\cap M_{\rell}(F).
\]
Fixons un ensemble de représentants $\underline{N}(M)$ du quotient $\Norm_{G}(M)(F)/ M(F)$. On~véri\-fie que l'on peut choisir pour ensemble $X_{M}(x)$ l'ensemble
\[
X_{M}(x)=\{n^{-1}xn\sep n\in \underline{N}(M)\}.
\]
On a aussi $D^M(n^{-1}xn)=D^M(x)$ pour tout $n$. La formule devient
\begin{equation}\tag{3}\label{induction3}
\theta_{D}(x)=D^G(x)^{-1/2}D^M(x)^{1/2}\sum_{n\in \underline{N}(M)}\theta_{D^M}(n^{-1}xn).
\end{equation}

\subsection{L'espace \texorpdfstring{$\cD(\frg)$}{Dg}}\label{lespace}

Soit $\cF\in \Fac(G)$. On note $\boldsymbol{\frg}_{\cF}$ l'algèbre de Lie du groupe $\bG_{\cF}$ et $\boldsymbol{\frg}_{\cF,\nil}$ son sous-ensemble des éléments nilpotents. On note $C_{\cusp}(\boldsymbol{\frg}_{\cF,\nil})$ l'espace des fonctions sur $\boldsymbol{\frg}_{\cF}(k_{F})$, à valeurs complexes, à support nilpotent, qui sont invariantes par conjugaison par $\bG_{\cF}(k_{F})$ et cuspidales. On note $\Fac_{\max}(G)$ l'ensemble des facettes réduites à un point (c'est-à-dire les sommets de l'immeuble).

À l'aide de ces objets, on définit des espaces $\boldsymbol{\cD}_{\cusp}(\frg)$, $\cD_{\cusp}(\frg)$, $\boldsymbol{\cD}(\frg)$ et $\cD(\frg)$ de façon similaire à ceux de \ref{facettesfonctions}. On pose\vspace*{-3pt}
\begin{align*}
\boldsymbol{\cD}_{\cusp}(\frg)&=\bigoplus_{\cF\in \Fac_{\max}(G)}C_{\cusp}(\boldsymbol{\frg}_{\cF,\nil}),
\\
\boldsymbol{\cD}(\frg)&=\bigoplus_{M}\boldsymbol{\cD}_{\cusp}(\mfrm),
\end{align*}
où $M$ parcourt les Levi de $G$. Le groupe $G(F)$ agit naturellement sur ces espaces et on~note $\cD_{\cusp}(\frg)$ et $\cD(\frg)$ les espaces de coinvariants. On définit une application linéaire $D^G:\boldsymbol{\cD}(\frg)\to \rI(\frg)^*$: pour un Levi $M$, une facette $\cF_{M}\in \Fac_{\max}(M)$ et une fonction $f\in C_{\cusp}(\boldsymbol{\mfrm}_{\cF_{M},\nil})$, on relève $f$ en une fonction $f_{\cF_{M}}$ sur $\mfrm(F)$ qui est très cuspidale et on pose $D^G_{f}=\Ind_{M}^G(D^M_{f_{\cF_{M}}})$ (en notant encore $\Ind_{M}^G:\rI(\mfrm)^*\to \rI(\frg)^*$ l'homomorphisme d'induction). L'application $D^G$ se quotiente en une application définie sur $\cD(\frg)$.

\skpt
\begin{prop*}
\begin{enumeratei}
\item\label{lespacei}
L'application $D^G:\cD(\frg)\to \rI(\frg)^*$ est injective.

\item\label{lespaceii}
Son image est l'espace des quasi-caractères de niveau~$0$ sur $\frg(F)$.
\end{enumeratei}
\end{prop*}

\begin{proof}
Pour $\cF\in \Fac(G)$, notons $ \cE_{\cF,\nil}$ l'espace des fonctions sur $\frg(F)$ à valeurs complexes à support dans $\frk_{\cF}\cap \frg_{\tn}(F)$ et invariantes par translations par $\frk_{\cF}^+$. Notons $\cE(\frg)$ le sous-espace de $C_{\rc}^{\infty}(\frg(F))$ engendré par les $\cE_{\cF,\nil}$ pour $\cF\in \Fac(G)$. On note $\rIE(\frg)$ son image dans $\rI(\frg)$. Les démonstrations des paragraphes \ref{calcul} à \ref{uncorollaire} qui concernent des fonctions sur $G(F)$ s'adaptent à l'algèbre de Lie et conduisent aux mêmes conclusions:
\begin{enumerate}
\item\label{lespace1}
la composée de $D^G$ et de la restriction $\rI(\frg)^*\to \rIE(\frg)^*$ est un isomorphisme de $\cD(\frg)$ sur cet espace;

\item\label{lespace2}
l'application linéaire composée $\cD_{\cusp}(\frg)\To{D^G}\rI(\frg)^*\to \rI_{\cusp}(\frg)^*$ est injective.

\smallskip\noindent
L'assertion \eqref{lespace1} entraîne que $D^G$ est injective.
\end{enumerate}

Notons $\cH$ le sous-espace de $ C_{\rc}^{\infty}(\frg(F))$ engendré par les fonctions $f$ pour lesquelles il existe une sous-algèbre d'Iwahori $\frb$ de $\frg$ de sorte que $f$ soit invariante par $\frb$. En notant $\fru$ le radical pro-$p$-nilpotent de $\frb$, cette invariance équivaut à ce que $\hat{f}$ soit à support dans $\fru$. En se rappelant que $\frg_{\tn}(F)$ est l'ensemble des $X\in \frg(F)$ tels qu'il existe une telle algèbre $\fru$ de sorte que $X\in \fru$, on voit facilement que $\cH$ est l'ensemble des fonctions $f$ tels que le support de $\hat{f}$ soit contenu dans $\frg_{\tn}(F)$. Remarquons que, pour $f\in \cE(\frg)$, le support de $f$ est contenu dans $\frg_{\tn}(F)$, donc $\hat{f}\in \cH$.
Notons $\rIH$ l'image de $\cH$ dans $\rI(\frg)$ et $\hat{\rI}\cE(\frg)$ l'image de $\rIE(\frg)$ par transformation de Fourier.

Rappelons qu'un élément $X\in \frg(F)$ est dit entier s'il existe une facette $\cF\in \Fac(G)$ de sorte que $X\in \frk_{\cF}$. Notons $\rI(\frg)^*_{\ent}$ l'espace des distributions invariantes sur $\frg(F)$ à support entier.

Soient $M$ un Levi de $G$, $\cF_{M}\in \Fac_{\max}(M)$ et $f\in C_{\cusp}(\boldsymbol{\mfrm}_{\cF_{M},\nil})$. On a $D^G_{f}=\Ind_{M}^G(D^M_{f_{\cF_{M}}})$. On vérifie que la transformée de Fourier $\hat{D}^G_{f}$ de la distribution $D^G_{f}$ est égale à $ \Ind_{M}^G(D^M_{\hat{f}_{\cF_{M}}})$. La fonction $\hat{f}_{\cF_{M}}$ étant à support dans $\frk_{\cF_{M}}$ donc entier, on~a $\hat{D}^G_{f}\in \rI(\frg)^*_{\ent}$. On note $\hat{D}^G(\cD(\frg))$ l'image de $D^G(\cD(\frg))$ par transformée de Fourier.

Notons $i:\CC[\Nil(\frg)]\to \rI(\frg)^*$ l'application qui, à une orbite nilpotente $\cO$, associe l'intégrale orbitale $I_{\cO}$. Son image est bien sûr contenue dans $\rI(\frg)^*_{\ent}$. On a un diagramme commutatif
\[
\xymatrix@R=.5cm{
&&\rIH^*\ar[dr]^-r&\\
\hat{D}^G(\cD(\frg))\ar[r]^{j}&\rI(\frg)^*_{\ent}\ar[ru]^-{\res}&&\hat{\rI}\cE(\frg)^*\\
&&\CC[\Nil(\frg)]\ar[ul]^-i\ar[uu]^-{p_{1}}\ar[ur]_-{p_2}&
}
\]
L'application $j$ est l'injection naturelle. Les applications $\res$ et $r$ sont les restrictions. Les applications $p_{1}$ et $p_{2}$ sont celles qui rendent le diagramme commutatif. Dans \hbox{\cite[Th.\,2.1.5]{D}}, DeBacker a prouvé les assertions suivantes:
\begin{enumerate}\setcounter{enumi}{2}
\item\label{lespace3}
les applications $\res$ et $p_{1}$ ont même image;

\item\label{lespace4}
$p_{1}$ et $p_{2}$ sont injectives.
\end{enumerate}

Soit $\bmf\in \cD(\frg)$. Par construction, $D^G_{\bmf}$ est à support dans $\frg_{\tn}(F)$. Pour prouver que~$D^G_{\bmf}$ est un quasi-caractère de niveau~$0$, il suffit de prouver que $D^G_{\bmf}$ coïncide sur cet ensemble avec une combinaison linéaire de transformées de Fourier d'intégrales orbitales nilpotentes. D'après ce que l'on~a dit ci-dessus, il suffit de prouver que~$\hat{D}^G_{\bmf}$ coïncide sur $\cH$ avec une combinaison linéaire d'intégrales orbitales nilpotentes. Autrement dit, il suffit de prouver que $\res\circ j(\hat{D}^G_{\bmf})$ appartient à l'image de $p_{1}$. Cela résulte de \eqref{lespace3}.

Inversement, soit $d$ un quasi-caractère de niveau~$0$. Notons $\hat{d}_{1}\in \rIH^*$ la restriction de $\hat{d}$ à $\rIH$. Le quasi-caractère $d$ coïncide sur $\frg_{\tn}(F)$ avec une combinaison linéaire de transformées de Fourier d'intégrales orbitales nilpotentes. Donc $\hat{d}$ coïncide sur $\cH$ avec une combinaison linéaire d'intégrales orbitales nilpotentes. C'est-à-dire que $\hat{d}_{1}$ appartient à l'image de $p_{1}$. Soit $N_{1}\in \CC[\Nil(\frg)]$ tel que $\hat{d}_{1}=p_{1}(N_{1})$. Par transformation de~Fourier, l'assertion \eqref{lespace1} dit que l'application $r\circ \res\circ j$ est un isomorphisme. Il existe donc $\bmf\in \cD(\frg)$ tel que $r(\hat{d}_{1})=r\circ \res\circ j(\hat{D}^G_{\bmf})$. D'après \eqref{lespace3}, il existe $N_{2}\in \CC[\Nil(\frg)]$ tel que $\res\circ j(\hat{D}^G_{\bmf})=p_{1}(N_{2})$. On a alors
\[
p_{2}(N_{2})=r\circ p_{1}(N_{2})=r\circ \res\circ j(\hat{D}^G_{\bmf})=r(\hat{d}_{1})=r\circ p_{1}(N_{1})=p_{2}(N_{1}).
\]
D'après \eqref{lespace4}, cela entraîne $N_{1}=N_{2}$. D'où
\[
\hat{d}_{1}=p_{1}(N_{1})=p_{1}(N_{2})=\res\circ j(\hat{D}^G_{\bmf}).
\]
Autrement dit, les distributions $\hat{d}$ et $\hat{D}^G_{\bmf}$ coïncident sur $\cH$. Donc $d$ et $D^G_{\bmf}$ coïncident sur $\frg_{\tn}(F)$. Mais elles sont toutes deux à support dans cet ensemble. Donc $d=D^G_{\bmf}$, ce qui achève la démonstration.
\end{proof}

\subsection{Filtrations sur l'algèbre de Lie}\label{filtrationssurlalgebredelie}
Pour tout $n \in \ZZ$, notons $\Fil^n\rI(\frg)$ l'image dans $\rI(\frg)$ du sous-espace des $f\in C_{\rc}^{\infty}(\frg(F))$ qui satisfont à la condition: pour tout Levi $M$ tel que $\dim(A_{M})>n$ et tout $X\in \frg_{\reg}(F)\cap \mfrm(F)$, on~a $I^G(X,f)=0$. On a
\[
\Fil^{a_{G}-1}\rI(\frg)=\{0\}\subset \Fil^{a_{G}}\rI(\frg)=\rI_{\cusp}(\frg)\subset\cdots\subset \Fil^{a_{M_{\min}}}\rI(\frg)=\rI(\frg),
\]
et, en posant $\Gr^n\rI(\frg)=\Fil^n\rI(\frg)/\Fil^{n-1}\rI(\frg)$, on a
\[
\Gr^n\rI(\frg)\simeq \bigoplus_{M\in \underline{\cL}_{\min}^n}\rI_{\cusp}(\mfrm)^{W^G(M)}.
\]

Notons $\Ann^n\rI(\frg)^*$ l'annulateur de $\Fil^{n-1}\rI(\frg)$ dans $\rI(\frg)^*$. On a
\[
\Ann^{a_{M_{\min}}+1}\rI(\frg)^*=\{0\}\subset \Ann^{a_{M_{\min}}}\rI(\frg)^*\subset\cdots\subset \Ann^{a_{G}}\rI(\frg)^*=\rI(\frg)^*,
\]
et, en posant $\Gr^n\rI(\frg)^*=\Ann^n\rI(\frg)^*/\Ann^{n+1}\rI(\frg)^*$,
\begin{equation}\tag{1}\label{filtrationssurlalgebredelie1}
\Gr^n\rI(\frg)^*\simeq \bigoplus_{M\in \underline{\cL}_{\min}^n}\rI_{\cusp}(\mfrm)^{*W^G(M)}.
\end{equation}

Notons $Qc_{0}(\frg)$ l'espace des quasi-caractères de niveau~$0$ sur $\frg(F)$. Notons $Qc_{0}^n(\frg)$ la somme des $\Ind_{M}^G(Qc_{0}(\mfrm))$ sur les Levi $M$ tels que $\dim(A_{M})\geq n$. Notons aussi $\cD^n(\frg)$ la somme des images dans $\cD(\frg)$ des $\cD_{\cusp}(\mfrm)$ sur les mêmes Levi $M$.

\begin{lem*}
Pour tout $n\in \ZZ$, on~a l'égalité
\[
Qc_{0}^n(\frg)=Qc_{0}(\frg)\cap \Ann^n\rI(\frg)^*=D^G(\cD^n(\frg)),
\]
et l'isomorphisme
\[
Qc_{0}^n(\frg)/Qc_{0}^{n+1}(\frg)\simeq \bigoplus_{M\in \underline{\cL}^n}D^M(\cD_{\cusp}(\mfrm)^{W^G(M)}).
\]
\end{lem*}

\begin{proof}
L'égalité $Qc_{0}^n(\frg)=D^G(\cD^n(\frg))$ résulte des définitions et de la proposition \ref{lespace} appliquée aux Levi $M$ tels que $\dim(A_{M})\geq n$. Notons $\rI(\frg)^*_{\nil}$ le sous-espace des distributions sur $\frg(F)$ à support nilpotent. Autrement dit $\rI(\frg)^*_{\nil}=i(\CC[\Nil])$ avec les notations de la preuve précédente. On a prouvé en \cite[Prop.\,I.5.3]{MW} que $\rI(\frg)^*_{\nil}\cap \Ann^n\rI(\frg)^*$ était la somme des $\Ind_{M}^G(\rI(\mfrm)^*_{\nil})$ où $M$ parcourt les Levi tels que $\dim(A_{M})\geq n$. Les filtrations sont invariantes par transformées de Fourier. Cela résulte de ce que, pour toute $f\in C_{\rc}^{\infty}(\frg(F))$, tout Levi $M$ et tout $X\in \mfrm(F)\cap \frg_{\reg}(F)$, $I^G(X,\hat{f})$ ne dépend que des $I^G(Y,f)$ pour $Y\in \mfrm(F)\cap \frg_{\reg}(F)$. Notons $\hat{\rI}(\frg)^*_{\nil}$ l'image de $\rI(\frg)^*_{\nil}$ par transformée de Fourier. On obtient que $\hat{\rI}(\frg)^*_{\nil}\cap \Ann^n\rI(\frg)^*$ est la somme des $\Ind_{M}^G(\hat{\rI}(\mfrm)^*_{\nil})$ où $M$ parcourt les Levi tels que $\dim(A_{M})\geq n$. Soit $d\in Qc_{0}(\frg)\cap \Ann^n\rI(\frg)^* $. Par définition d'un quasi-caractère de niveau~$0$, $d$ est la restriction à $\frg_{\tn}(F)$ d'un élément $d_{0}\in \hat{\rI}(\frg)^*_{\nil}$. Puisque $d\in \Ann^n\rI(\frg)^*$, $d_{0}$ annule tout élément de $\Fil^{n-1}\rI(\frg)$ à support topologiquement nilpotent. L'espace $\Fil^{n-1}\rI(\frg)$ est invariant par l'action de $F^{\times}$ par homothétie. Puisque $\hat{\rI}(\frg)^*_{\nil}$ est engendré par des éléments homogènes, sinon par l'action de $F^{\times}$, du moins par l'action du sous-groupe des carrés, dire que $d_{0}$ annule tout élément de $\Fil^{n-1}\rI(\frg)$ à support topologiquement nilpotent équivaut à dire qu'il annule tout $\Fil^{n-1}\rI(\frg)$, c'est-à-dire que $d_{0}\in \hat{\rI}(\frg)^*_{\nil}\cap \Ann^n\rI(\frg)^*$. Donc $d_{0}=\sum_{M}\Ind_{M}^G(d_{M,0})$, où l'on somme sur les Levi $M$ tels que $\dim(A_{M})\geq n$ et où $d_{M,0}$ est un certain élément de $\hat{\rI}(\mfrm)^*_{\nil}$. Donc $d= \sum_{M}\Ind_{M}^G(d_{M})$, où $d_{M}$ est la restriction de $d_{M,0}$ à $\mfrm_{\tn}(F)$. Cet élément $d_{M}$ appartient à $Qc_{0}(\mfrm)$. Donc $d\in Qc_{0}^n(\frg)$. Cela démontre l'inclusion $Qc_{0}(\frg)\cap \Ann^n\rI(\frg)^*\subset Qc_{0}^n(\frg)$ et l'inclusion opposée est évidente.

On a
\[
Qc_{0}^n(\frg)/Qc_{0}^{n+1}(\frg)=D^G(\cD^n(\frg))/D^G(\cD^{n+1}(\frg))
\]
et, d'après ce que l'on vient de prouver, cet espace s'envoie injectivement dans $\Gr^n\rI(\frg)^*$.
Par définition, $D^G(\cD^n(\frg))$ est la somme
\[\textstyle
\text{de }D^G(\cD^{n+1}(\frg))\quadet \text{de }D^G(\bigoplus_{M\in \underline{\cL}_{\min}^n}\cD_{\cusp}(\mfrm)).
\]
Ce dernier espace a donc même image dans $\Gr^n\rI(\frg)^*$ que $Qc_{0}^n(\frg)/Qc_{0}^{n+1}(\frg)^*$. Soit
\[
M\in \underline{\cL}_{\min}^n\quadet \bmf\in \cD_{\cusp}(\mfrm).
\]
L'image de $D^G_{\bmf}$ dans $\Gr^n\rI(\frg)^*$ est nulle dans les composantes $\rI_{\cusp}(\mfrm')^{*W^G(M')}$ de la décomposition \eqref{filtrationssurlalgebredelie1} pour $M'\neq M$ et est l'image naturelle de $D^M_{\bmf}$ dans $\rI_{\cusp}(\mfrm)^{* W^G(M)}$ pour $M'=M$ (c'est-à-dire l'image de $D^M_{\bmf}$ dans $\rI_{\cusp}(\mfrm)^*$ que l'on restreint au sous-espace $\rI_{\cusp}(\mfrm)^{W^G(M)}$). En moyennant sur le groupe $ W^G(M)$, on obtient que $Qc_{0}^n(\frg)/Qc_{0}^{n+1}(\frg)$ s'identifie à la somme des images de $\cD_{\cusp}(\mfrm)^{W^G(M)}$ dans $\rI_{\cusp}(\mfrm)^{*W^G(M)}$.
D'après le~\eqref{lespace2} de \ref{lespace}, $\cD_{\cusp}(\mfrm)^{W^G(M)}$ s'envoie injectivement dans $\rI_{\cusp}(\mfrm)^{*W^G(M)}$, ce qui achève la démonstration.
\end{proof}

\Subsection{ Un premier théorème}\label{unpremiertheoreme}

\begin{thm*}
L'image de l'application $D^G:\cD(G)\to \rI(G)^*$ est formée de quasi-caractères de niveau~$0$ sur $G(F)$.
\end{thm*}

\begin{proof}
En vertu du \eqref{inductionii} du lemme \ref{induction}, il suffit de prouver que, pour tout élément $\bmf\in \cD_{\cusp}(G)$, $D^G_{\bmf}$ est un quasi-caractère de niveau~$0$ sur $G(F)$. La famille~$\bmf$ peut être infinie mais la définition des quasi-caractères de niveau~$0$ est locale et, loca\-lement, seuls un nombre fini de fonctions de la famille $\bmf$ interviennent. On peut donc aussi bien supposer qu'il y a une seule fonction. C'est-à-dire que l'on peut fixer $(\cF,\nu)\in \Fac^*_{\max}(G)$ et $f\in C_{\cusp}(\bG_{\cF}^{\nu})$ et supposer que $\bmf$ est réduit à l'unique fonction~$f$. On a alors $D^G_{\bmf}=D^G_{f_{\cF}}$. C'est un quasi-caractère d'après \ref{quasicaracteres}\,\eqref{quasicaracteres6}. On note $\theta$ sa fonction localement intégrable associée. Soit $\epsilon$ un élément de $G(F)$ qui est $p'$-compact $\rmod Z(G)$. On doit calculer $\theta(\epsilon \exp(X))$ pour tout $X\in \frg_{\epsilon,\tn}(F)$ tel que $\epsilon \exp(X)$ soit fortement régulier. On peut évidemment supposer $\nu=w_{G}(\epsilon)$, sinon cette fonction est nulle. Fixons $X$. Notons $T$ le commutant de $X$ dans $G_{\epsilon}$, qui est aussi le commutant de $\epsilon \exp(X)$ dans $G$. Notons $M$ le commutant de $A_{T}$ dans $G$. C'est un Levi de $G$ contenant $\epsilon$ et $\exp(X)$ par construction. On pose $M_{\epsilon}=G_{\epsilon}\cap M$. C'est un Levi de $G_{\epsilon}$ et on~a $X\in \mfrm_{\epsilon}(F)$. L'élément $X$ est elliptique dans $\mfrm_{\epsilon}(F)$ et $\epsilon \exp(X)$ est elliptique dans $M(F)$. Notons que $A_{T}=A_{M_{\epsilon}}=A_{M}$. D'après \ref{quasicaracteres}\,\eqref{quasicaracteres7}, on a
\begin{multline}\tag{1}\label{unpremiertheoreme1}
\theta(\epsilon \exp(X))=(-1)^{a_{M}-a_{G}}D^G(\epsilon \exp(X))^{-1/2}m(A_{M})m(A_{G})^{-1}J^G_{M}(\epsilon \exp(X),f_{\cF})
\\
=(-1)^{a_{M}-a_{G}}m(A_{M})m(A_{G})^{-1} \int_{A_{M}(F)\backslash G(F)}f_{\cF}(g^{-1} \epsilon \exp(X)g)v_{M}(g)\, dg.
\end{multline}
Pour $g\in G(F)$, on a
\begin{enumerate}\setcounter{enumi}{1}
\item\label{unpremiertheoreme2}
$g^{-1}\epsilon \exp(X)g\in K_{\cF}^{\nu}$ si et seulement si $g^{-1}\epsilon g\in K_{\cF}^{\nu}$ et $g^{-1} \exp(X)g\in K_{\cF}^0$.
\end{enumerate}

En effet, soit $c\geq1$ un entier premier à $p$ tel que $\epsilon^c\in Z(G)(F)$. Supposons $g^{-1}\epsilon \exp(X)g\in K_{\cF}^{\nu}$. Alors $g^{-1} \epsilon^c\exp(cX)g\in K_{\cF}^{\dag}$. On a aussi
\[
g^{-1}\epsilon^cg\in Z(G)(F)\subset K_{\cF}^{\dag}.
\]
Donc $g^{-1}\exp(cX)g\in K_{\cF}^{\dag}$. Il en résulte que $g^{-1}\exp(X)g\in K_{\cF}^{\dag}$ et, puisque c'est un élément topologiquement unipotent, on~a forcément $g^{-1}\exp(X)g\in K_{\cF}^0$. Puisque $g^{-1}\epsilon \exp(X)g\in K_{\cF}^{\nu}$, cela entraîne $g^{-1}\epsilon g\in K_{\cF}^{\nu}$. La réciproque est évidente. D'où \eqref{unpremiertheoreme2}.

On a aussi:
\begin{enumerate}\setcounter{enumi}{2}
\item\label{unpremiertheoreme3}
la classe de conjugaison par $G(F)$ de $\epsilon$ coupe $K_{\cF}^{\nu}$ en un nombre fini de classes de conjugaison par $K_{\cF}^0$.
\end{enumerate}

En effet, notons $\Cl(\epsilon)$ cette classe de conjugaison par $G(F)$. C'est un sous-ensemble fermé de $G(F)$, donc son intersection avec $K_{\cF}^{\nu}$ est compacte. L'application
\begin{align*}
Z_{G}(\epsilon)(F)\backslash G(F)&\to \Cl(\epsilon)\\
g&\mto g^{-1}\epsilon g
\end{align*}
est un homéomorphisme. Puisque les orbites de l'action de $K_{\cF}^{0}$ sur l'ensemble de départ sont ouvertes, il en est de même de celles sur l'ensemble d'arrivée. Puisque $\Cl(\epsilon)\cap K_{\cF}^{\nu}$ est compact, il n'y a donc qu'un nombre fini de telles orbites, d'où \eqref{unpremiertheoreme3}.

Notons $\Gamma_{\epsilon}$ l'ensemble des $g\in G(F)$ tels que $g^{-1}\epsilon g\in K_{\cF}^{\nu}$. En conséquence de \eqref{unpremiertheoreme3}, l'ensemble $G_{\epsilon}(F)\backslash \Gamma_{\epsilon}/K_{\cF}^0$ est fini. Fixons-en un ensemble de représentants $\boldsymbol{\gamma}$. Pour tout $\gamma\in \boldsymbol{\gamma}$, définissons comme en \ref{facettesfonctions} la facette $\gamma\cF$ et posons
\[
m_{\gamma}=\mes(G_{\epsilon}(F)\cap K_{\gamma\cF}^0)^{-1},
\]
la mesure étant calculée relativement à la mesure implicitement fixée sur $G_{\epsilon}(F)$. On~vérifie la formule d'intégration\vspace*{-3pt}
\[
\int_{A_{M}(F)\backslash\Gamma_{\epsilon}}\varphi(g)\,dg=\sum_{\gamma\in \boldsymbol{\gamma}}m_{\gamma}\int_{A_{M}(F)\backslash G_{\epsilon}(F)}\int_{K_{\cF}^0}\varphi(g\gamma k)\,dk\,dg
\]
pour toute fonction intégrable $\varphi$ sur $A_{M}(F)\backslash \Gamma_{\epsilon}$. On applique cela à la fonction\vspace*{-3pt}\enlargethispage{.5\baselineskip}%
\[
\varphi(g)=f_{\cF}(g^{-1}\epsilon \exp(X) g)v_{M}(g).
\]
La fonction $f_{\cF} $ est invariante par conjugaison par~$K_{\cF}^0$. On~voit que\vspace*{-3pt}
\[
f_{\cF}(k^{-1}\gamma^{-1}g^{-1}\epsilon \exp(X)g \gamma k)=(^{\gamma}f)_{\gamma\cF}(\epsilon \exp(g^{-1}Xg))
\]
pour tous $\gamma\in \boldsymbol{\gamma}$, $g\in G_{\epsilon}(F)$ et $k\in K_{\cF}^0$. Pour de mêmes
$\gamma$ et $g$, posons\vspace*{-3pt}
\[
v_{M,\gamma}(g)=\int_{K_{\cF}^0}v_{M}(g\gamma k)\,dk.
\]
Grâce à \eqref{unpremiertheoreme1} et \eqref{unpremiertheoreme2} et à l'égalité $A_{M_{\epsilon}}=A_{M}$, on obtient\vspace*{-3pt}
\begin{multline*}\tag{4}\label{unpremiertheoreme4}
\theta(\epsilon \exp(X))=(-1)^{a_{M_{\epsilon}}-a_{G}}m(A_{M_{\epsilon}})m(A_{G})^{-1}\\
\cdot\sum_{\gamma\in \boldsymbol{\gamma}}m_{\gamma}
\int_{A_{M_{\epsilon}}(F)\backslash G_{\epsilon}(F)}(^{\gamma}f)_{\gamma\cF}(\epsilon \exp(g^{-1}X g))v_{M,\gamma} (g)\,dg.
\end{multline*}
Fixons $\gamma\in \boldsymbol{\gamma}$.
Par définition de cet ensemble, on~a $\epsilon\in K_{\gamma\cF}^{\nu}$. Introduisons la facette $\cF'_{\gamma}\in \Imm(G_{\epsilon,\AD})$ associée à $\gamma\cF$, \cf proposition \ref{pointsfixes}\,\eqref{pointsfixesiii}. Le groupe $\bG_{\epsilon,\cF_{\gamma}'}$ est la composante neutre de $(\bG_{\gamma\cF})^{\epsilon}$.
Notons $\bar{\epsilon}$ la réduction de $\epsilon$ dans $\bG^{\nu}_{\gamma\cF}(k_{F})$ et $(^{\gamma}f)'$ la fonction sur $\bG_{\epsilon,\cF'_{\gamma}}(k_{F})$, à support unipotent, définie par $(^{\gamma}f)'(x)=(^{\gamma}f)(\bar{\epsilon}x)$, pour tout élément unipotent $x\in \bG_{\epsilon,\cF'_{\gamma}}(k_{F})$. De $(^{\gamma}f)'$ se déduit une fonction $(^{\gamma}f)'_{\cF'_{\gamma}}$ sur~$G_{\epsilon}(F)$.

\skpt
\begin{lem*}
\begin{enumeratei}
\item\label{unpremiertheoremei}
Pour tout $g\in G_{\epsilon}(F)$, on~a l'égalité
\[
(^{\gamma}f)_{\gamma\cF}(\epsilon \exp(g^{-1} Xg))=(^{\gamma}f)'_{\cF_{\gamma}'}(\exp(g^{-1}Xg)).
\]

\item\label{unpremiertheoremeii}
Si $\dim(A_{G_{\epsilon}})> \dim(A_{G})$ ou si $\cF'_{\gamma}$ n'est pas réduit à un point, $ (^{\gamma}f)'=0$.

\item\label{unpremiertheoremeiii}
Supposons que $A_{G_{\epsilon}}=A_{G}$ et que $\cF'_{\gamma}$ soit réduit à un point, c'est-à-dire $\cF_{\gamma}'\in \Fac_{\max}(G_{\epsilon})$. Alors $(^{\gamma}f)'$ est cuspidale.
\end{enumeratei}
\end{lem*}

\begin{proof}
Pour simplifier la notation, on suppose $\gamma=1$, on pose $\cF'=\cF'_{\gamma}$ et $f'=(^{\gamma}f)'$.

Pour $y\in K_{\cF}^0\cap G_{\epsilon}(F)$, la réduction $\bar{y}$ de $y$ dans $\bG_{\cF}(k_{F})$ appartient à $(\bG_{\cF})^{\epsilon}(k_{F})$. Mais, si $y$ est topologiquement unipotent, on peut remplacer dans cette relation le groupe $(\bG_{\cF})^{\epsilon}$ par sa composante neutre $\bG_{\epsilon,\cF'}$. De plus cette réduction $\bar{y}$ est unipotente. En appliquant cela à $y=\exp(g^{-1}Xg)$, on en déduit le~\eqref{unpremiertheoremei} de l'énoncé.

Notons $\bA_{\cF'}$ le plus grand sous-tore central et déployé de $\bG_{\epsilon,\cF'}$. Notons $\bM'$ son commutant dans $\bG_{\cF}$, qui est un Levi de ce groupe. Fixons un élément $x_{*}\in X_{*}(\bA_{\cF'})$ en position générale, notons $\bP'$ le sous-groupe de $\bG_{\cF}$ engendré par $\bM'$ et par les groupes radiciels relatifs à l'action de $\bA_{\cF'}$ associés aux racines $\alpha$ telles que $\<\alpha,x_{*}\>>0$. C'est un élément de $\cP(\bM')$.
Le tore $\bA_{\cF'}$ est contenu dans l'ensemble des points fixes de l'action de $\bar{\epsilon}$ sur $\bG_{\cF}$. Il en résulte que cette action conserve $\bM'$ et $\bP'$. Alors $\bP^{\prime\nu}=\bar{\epsilon}\bP'$ est un espace parabolique de $\bG^{\nu}_{\cF}$ d'espace de Levi $\bM^{\prime\nu}=\bar{\epsilon}\bM'$. Sur la clôture algébrique $\bar{k}_{F}$, les valeurs propres de l'action de $\bar{\epsilon}$ dans $\boldsymbol{\fru}_{\bP'}$ sont des racines de l'unité, puisque $\epsilon$ est $p'$-compact $\rmod Z(G)$. L'espace des points fixes de l'action de $\bar{\epsilon}$ dans $\boldsymbol{\frg}_{\cF}$ est $\boldsymbol{\frg}_{\epsilon,\cF'}$, qui est inclus dans $\boldsymbol{\mfrm}'$ par définition de $\bM'$. Donc les valeurs propres de l'action de $\bar{\epsilon}$ dans $\boldsymbol{\fru}_{\bP'}$ sont toutes différentes de $1$. Il en résulte aisément que, pour tout élément unipotent $\bar{n}\in \bM'(k_{F})$, l'action $1-\ad(\bar{\epsilon}\bar{n})$ dans $\boldsymbol{\fru}_{\bP'}(k_{F})$ est un isomorphisme. Pour un tel élément $\bar{n}$, il s'en déduit l'égalité\vspace*{-3pt}\enlargethispage{.2\baselineskip}%
\[
\sum_{\bar{u}\in \bU_{\bP'}(k_{F})}f(\bar{\epsilon}\bar{n}\bar{u})=\sum_{\bar{u}\in \bU_{\bP'}(k_{F})}f(\bar{u}^{-1}\bar{\epsilon}\bar{n}\bar{u}).
\]
Puisque $f$ est invariante par conjugaison, le membre de droite n'est autre que $\vert \bU_{\bP'}(k_{F})\vert f(\bar{\epsilon}\bar{n})$. Supposons que $\bP^{\prime\nu}$ soit un espace parabolique propre de $\bG_{\cF}^{\nu}$. Alors le membre de gauche est nul puisque $f$ est cuspidale et invariante par conjugaison. Dans ce cas $f(\bar{\epsilon}\bar{n})=0$. A fortiori, cela est vrai pour tout élément unipotent $\bar{n}\in \bG_{\epsilon,\cF'}(k_{F})$, auquel cas $f(\bar{\epsilon}\bar{n})=f'(\bar{n})$. D'où $f'=0$. Pour démontrer le~\eqref{unpremiertheoremeii} de l'énoncé, il suffit de prouver que, sous les hypothèses de cette assertion, l'espace parabolique $\bP^{\prime\nu}$ est propre, ou encore que $\bM'\neq\bG_{\cF}$. On démontre la contraposée de cette assertion. Supposons donc $\bM'=\bG_{\cF}$. Notons $\bA_{\cF}$ le plus grand sous-tore central déployé dans $\bG_{\cF}$ et notons $\bA_{\cF}^{\nu}$ le plus grand sous-tore contenu dans l'ensemble des points fixes de l'action de $\bar{\epsilon}$ dans $\bA_{\cF}$. L'hypothèse $\bM'=\bG_{\cF}$ signifie que $\bA_{\cF'}\subset \bA_{\cF}$. Puisque $\bar{\epsilon}$ agit trivialement sur $\bA_{\cF'}$, cela entraîne $\bA_{\cF'}\subset \bA_{\cF}^{\nu}$ d'où l'égalité $\bA_{\cF'}= \bA_{\cF}^{\nu}$ car l'inclusion opposée est immédiate. Puisque $(\cF,\nu)\in \Fac_{\max}(G)$, on~a $\dim( \bA_{\cF}^{\nu})=\dim(A_{G})$. D'où $\dim(\bA_{\cF'})=\dim(A_{G})$. Or on~a évidemment $\dim(\bA_{\cF'})\geq \dim(A_{G_{\epsilon}})\geq \dim(A_{G})$. Ces trois dimensions sont donc égales. L'égalité $\dim(\bA_{\cF'})= \dim(A_{G_{\epsilon}})$ équivaut à ce que $\cF'$ soit réduit à un point. L'égalité des trois dimensions précédentes est donc le contraire de l'hypothèse de \eqref{unpremiertheoremeii}, ce qui achève la démonstration de cette assertion.

Supposons maintenant que $\dim(A_{G_{\epsilon}})=\dim(A_{G})$ et que $\cF'$ soit réduit à un point, prouvons que $f'$ est cuspidale. Fixons un sous-groupe parabolique propre $\bP'_{\epsilon}$ de $\bG_{\epsilon,\cF'}$ de composante de Levi $\bM'_{\epsilon}$. Notons $\bA'$ le plus grand sous-tore central déployé de~$\bM'_{\epsilon}$ et~$\bM'$ le commutant de $\bA'$ dans $\bG_{\cF}$. Fixons un élément $x_{*}\in X_{*}(\bA')$ tel que \hbox{$\<\alpha,x_{*}\>\!>\!0$} pour toute racine $\alpha$ de $\bA'$ dans $\boldsymbol{\fru}_{\bP'_{\epsilon}}$. On définit un sous-groupe parabolique $\bP'$ de~$\bG_{\cF}$ par la condition: $\boldsymbol{\fru}_{\bP'}$ est la somme des espaces radiciels associés aux racines $\alpha$ de~$\bA'$ dans $\boldsymbol{\frg}_{\cF}$ telles que $\<\alpha,x_{*}\>>0$. En supposant $x_{*}$ en position générale, $\bP'$ a pour composante de Levi $\bM'$ et on~a $\bP'\cap \bG_{\epsilon,\cF'}=\bP'_{\epsilon}$. De plus $\bP'$ est propre puisque~$\bP'_{\epsilon}$ l'est. L'argument est maintenant similaire à celui de la preuve de \eqref{unpremiertheoremeii}. L'ensemble $\bP^{\prime\nu}=\bar{\epsilon}\bP'$ est un espace parabolique de $\bG_{\cF}^{\nu}$ d'espace de Levi $\bM^{\prime\nu}=\bar{\epsilon}\bM'$. Pour un élément unipotent $\bar{n}\in \bM'_{\epsilon}(k_{F})$, on~a les égalités\vspace*{-3pt}
\begin{align*}
\sum_{\bar{u}\in \bU_{\bP'_{\epsilon}}(k_{F})}\hspace*{-2mm}f'(\bar{n}\bar{u})=\sum_{\bar{u}\in \bU_{\bP'_{\epsilon}}(k_{F})}\hspace*{-2mm}f(\bar{\epsilon}\bar{n}\bar{u})
&=\vert \bU_{\bP'}(k_{F})\vert ^{-1}\sum_{\bar{x}\in \bU_{\bP'}(k_{F})}\sum_{\bar{u}\in \bU_{\bP'_{\epsilon}}(k_{F})}\hspace*{-2mm}f(\bar{x}\bar{\epsilon}\bar{n}\bar{u}\bar{x}^{-1})\\
&=\vert \bU_{\bP'}(k_{F})\vert ^{-1}\vert \bU_{\bP'_{\epsilon}}(k_{F})\vert \sum_{\bar{x}\in \bU_{\bP'}(k_{F})}f(\bar{\epsilon}\bar{n}\bar{x})
\end{align*}
et cette dernière expression est nulle car $f$ est cuspidale. La nullité de la première expression signifie que $f'$ l'est. Cela démontre le lemme.
\end{proof}

Notons $\boldsymbol{\gamma}_{\max}$ le sous-ensemble des $\gamma\in \boldsymbol{\gamma}$ tels que $\cF'_{\gamma}$ soit réduit à un point. Si $\dim(A_{G_{\epsilon}})\!>\!\dim(A_{G})$, le~\eqref{unpremiertheoremeii} du lemme et la formule~\eqref{unpremiertheoreme4} impliquent que \hbox{$\theta(\epsilon \exp(X))\!=\!0$} et on~a terminé. Supposons désormais $\dim(A_{G_{\epsilon}})=\dim(A_{G})$. Alors les assertions \eqref{unpremiertheoremei} et \eqref{unpremiertheoremeii} du lemme entraînent que la formule~\eqref{unpremiertheoreme4} se transforme en
\begin{multline*}\tag{5}\label{unpremiertheoreme5}
\theta(\epsilon \exp(X))=(-1)^{a_{M_{\epsilon}}-a_{G}}m(A_{M_{\epsilon}})m(A_{G})^{-1}\\
\cdot\sum_{\gamma\in \boldsymbol{\gamma}_{\max}}m_{\gamma}
\int_{A_{M_{\epsilon}}(F)\backslash G_{\epsilon}(F)}(^{\gamma}f)'_{\cF'_{\gamma}}(\exp(g^{-1}X g))v_{M,\gamma} (g)\,dg.
\end{multline*}
Fixons $\gamma\in \boldsymbol{\gamma}_{\max}$ et calculons $v_{M,\gamma}(g)$ pour $g\in G_{\epsilon}(F)$. Le poids $v_{M}$ est calculé grâce au choix d'un sous-groupe compact spécial de $G(F)$. Fixons aussi un tel sous-groupe compact $K_{\epsilon}$ du groupe $G_{\epsilon}(F)$, qui permet de définir un poids $v_{M_{\epsilon}}$. Une formule d'Arthur, que l'on~a reprise en \cite[Lem.\,3.3]{W1}, dit que, pour $g\in G_{\epsilon}(F)$, $v_{M,\gamma}(g)-v_{M_{\epsilon}}(g)$ est une somme finie de termes $v'_{Q_{\epsilon}}(g)$, où $Q_{\epsilon}$ est un sous-groupe parabolique propre de~$G_{\epsilon}$ contenant $M_{\epsilon}$ et $v'_{Q_{\epsilon}}$ est une fonction localement intégrable sur $G_{\epsilon}(F)$ invariante à gauche par $M_{\epsilon}(F)$ et par $U_{Q_{\epsilon}}(F)$. Pour de telles données, notons $L_{\epsilon}$ la composante de Levi de $Q_{\epsilon}$ contenant $M_{\epsilon}$. En utilisant la décomposition d'Iwasawa, on calcule
\begin{multline*}
\int_{A_{M_{\epsilon}}(F)\backslash G_{\epsilon}(F)}(^{\gamma}f)'_{\cF'_{\gamma}}(\exp(g^{-1}X g))v_{Q_{\epsilon}}' (g)\,dg\\
=c\int_{K_{\epsilon}}\int_{A_{M_{\epsilon}}(F)\backslash L_{\epsilon}(F)}\int_{U_{Q_{\epsilon}}(F)}(^{\gamma}f)'_{\cF'_{\gamma}}(k^{-1}u^{-1}l^{-1}\exp(X)luk)v'_{Q_{\epsilon}}(lk)\,du\,dl\,dk,
\end{multline*}
où $c>0$ ne dépend que des mesures de Haar. Puisque $X$ est régulier, l'intégrale intérieure en $u$ se transforme en
\[
c(X)\int_{U_{Q_{\epsilon}}(F)}(^{\gamma}f)'_{\cF'_{\gamma}}(k^{-1}l^{-1}\exp(X)luk)\, du,
\]
où $c(X)>0$ est un certain déterminant. Or cette intégrale est nulle car $(^{\gamma}f)'_{\cF'_{\gamma}}$ est très cuspidale d'après le~\eqref{unpremiertheoremeiii} du lemme. Les fonctions $v'_{Q_{\epsilon}}$ ne contribuent donc pas à la formule~\eqref{unpremiertheoreme5} et on peut récrire cette formule
\begin{multline*}
\theta(\epsilon \exp(X))=(-1)^{a_{M_{\epsilon}}-a_{G}}m(A_{M_{\epsilon}})m(A_{G})^{-1}\\
\cdot\sum_{\gamma\in \boldsymbol{\gamma}_{\max}}m_{\gamma}
\int_{A_{M_{\epsilon}}(F)\backslash G_{\epsilon}(F)}(^{\gamma}f)'_{\cF'_{\gamma}}(\exp(g^{-1}X g))v_{M_{\epsilon}} (g)\,dg.
\end{multline*}
En comparant avec \ref{quasicaracteres}\,\eqref{quasicaracteres7}, on obtient
\[
\theta(\epsilon \exp(X))=(-1)^{a_{G_{\epsilon}}-a_{G}}m(A_{G_{\epsilon}})m(A_{G})^{-1}\sum_{\gamma\in \boldsymbol{\gamma}_{\max}}m_{\gamma}\theta_{\gamma}(\exp(X)),
\]
où on~a noté $\theta_{\gamma}$ la fonction associée au quasi-caractère $D^{G_{\epsilon}}_{(^{\gamma}f)'}$. Pour tout $\gamma\in \boldsymbol{\gamma}_{\max}$, notons $\varphi_{\gamma}$ la fonction sur $\boldsymbol{\frg}_{\epsilon,\cF'_{\gamma}}(k_{F})$, à support nilpotent, telle que $\varphi_{\gamma}(Y)=(^{\gamma}f)'(\exp(Y))$ pour tout élément nilpotent $Y\in \boldsymbol{\frg}_{\epsilon,\cF'_{\gamma}}(k_{F})$. Il est clair que la fonction $Y\mto \theta_{\gamma}(\exp(Y))$, définie sur $\frg_{\epsilon,\tn}(F)$, n'est autre que la fonction associée à la distribution $D^{G_{\epsilon}}_{\varphi_{\gamma}}$ sur $\frg_{\epsilon}(F)$. Notons $\theta_{\varphi_{\gamma}}$ cette fonction. On obtient la formule finale
\begin{equation}\tag{6}\label{unpremiertheoreme6}
\theta(\epsilon \exp(X))=(-1)^{a_{G_{\epsilon}}-a_{G}}m(A_{G_{\epsilon}})m(A_{G})^{-1}\sum_{\gamma\in \boldsymbol{\gamma}_{\max}}m_{\gamma}\theta_{\varphi_{\gamma}}(X).
\end{equation}
On applique le~\eqref{lespaceii} de la proposition \ref{lespace}. Il entraîne que chaque fonction $\theta_{\varphi_{\gamma}}$ est combinaison linéaire de fonctions $Y\mto \hat{j}(\cO,Y)$ où $\cO$ parcourt les orbites nilpotentes de $\frg_{\epsilon}(F)$. Il en est donc de même
de la fonction $X\mto \theta(\epsilon \exp(X))$. Cela signifie que~$D^G_{f}$ est un quasi-caractère de niveau~$0$ sur $G(F)$.
\end{proof}

\Subsection{Restriction aux éléments elliptiques d'un quasi-caractère sur \texorpdfstring{$G(F)$ de niveau~$0$}{GF0}}\label{restriction}

Soit $D$ un quasi-caractère sur $G(F)$, notons $\theta_{D}$ sa fonction associée. Disons que $D$ est de niveau~$0$ sur les elliptiques si et seulement si $\theta_{D}$ satisfait à la condition suivante:
\begin{enumerate}
\item\label{restriction1}
pour tout élément $\epsilon\!\in\!G(F)$ qui est $p'$-compact $\rmod Z(G)$ et pour tout $X\!\in\!\frg_{\epsilon,\tn}(F)$ tel que $\epsilon \exp(X)\in G_{\rell}(F)$, la fonction $\lambda\mto \theta_{D}(\epsilon \exp(\lambda^2X))$ sur $\fro_{F}$ appartient à $E$.
\end{enumerate}

\skpt
\begin{rems*}
\begin{enumerate}\setcounter{enumi}{1}
\item\label{restriction2}
Pour $\epsilon\!\in\!G(F)_{p'}$ et $X\!\in\!\frg_{\epsilon,\tn}(F)$, $\epsilon \exp(X)$ ne peut être elliptique que si $A_{G_{\epsilon}}\!=\!A_{G}$. A fortiori, $\epsilon$ est $p'$-compact $\rmod Z(G)$.

\item\label{restriction3}
La même preuve qu'au \eqref{quasicaracteresii} du lemme \ref{quasicaracteres} montre que \eqref{restriction1} équivaut à la condition:
\begin{itemize}
\item
pour tout élément $\epsilon\in G(F)$ qui est $p'$-compact $\rmod Z(G)$ et pour tout $\cO\in \Nil(\frg_{\epsilon})$, il existe $c_{D,\cO}\in \CC$ de sorte que, pour tout $X\in \frg_{\epsilon,\tn}(F)$ tel que $\epsilon \exp(X)\in G_{\rell}(F)$, on ait l'égalité
\[
\theta_{D}(\epsilon \exp(X))=\sum_{\cO\in \Nil(\frg_{\epsilon})}c_{D,\cO}\hat{j}(\cO,X).
\]
\end{itemize}

\item\label{restriction4}
Compte tenu de la remarque précédente, la même preuve qu'au \eqref{quasicaracteresiii} du lemme \ref{quasicaracteres} montre que, si $B$ est un sous-ensemble de $G(F)_{p'}$ tel que $G(F)=\bigcup_{\epsilon\in B} C_{G}(\epsilon)$, alors~$D$ est un quasi-caractère de niveau~$0$ sur les elliptiques si et seulement si la condition~\eqref{restriction1} est satisfaite pour tout $\epsilon\in B$.
\end{enumerate}
\end{rems*}

\begin{prop*}
Soit $D$ un quasi-caractère sur $G(F)$ de niveau~$0$ sur les elliptiques. Alors il existe $\bmf\in \cD_{\cusp}(G)$ tel que $D$ coïncide avec $D^G_{\bmf}$ sur $G_{\rell}(F)$.
\end{prop*}

\begin{proof}
Comme on l'a vu dans \ref{variantes}, l'espace $\cD_{\cusp}(G)$ se décompose en produit d'espaces indexés par $\cN$. Il en est de même de l'espace des quasi-caractères de niveau~$0$ sur les elliptiques. On peut donc fixer $\nu\in \cN$ et supposer que
$D$ est à support dans $w_{G}^{-1}(\nu)$. L'intersection $w_{G}^{-1}(\nu)\cap G_{\rell}(F)$ est contenue dans une réunion finie d'ensembles $ C_{G}(\epsilon )$, où $\epsilon$ est $p'$-compact $\rmod Z(G)$. Ces ensembles sont disjoints ou confondus et ils sont ouverts et fermés. On peut fixer un élément $p'$-compact $\rmod Z(G)$, supposer que $D$ est à support dans $ C_{G}(\epsilon)$ et prouver:
\begin{enumerate}\setcounter{enumi}{4}
\item\label{restriction5}
il existe $\bmf\in C_{\cusp}(G)$ tel que $D^G_{\bmf}$ soit à support dans $C_{G}(\epsilon)$ et coïncide avec $D$ sur $C_{G}(\epsilon)\cap G_{\rell}(F)$.
\end{enumerate}

Comme on l'a dit dans la remarque \eqref{restriction2}, on peut supposer $A_{G_{\epsilon}}=A_{G}$, sinon la solution de \eqref{restriction5} est triviale.

Écrivons comme dans la remarque \eqref{restriction3}
\[
\theta_{D}(\epsilon \exp(X))=\sum_{\cO\in \Nil(\frg_{\epsilon})}c_{D,\cO}\hat{j}(\cO,X)
\]
pour tout $X\in \frg_{\epsilon,\tn}(F)$ tel que $\epsilon \exp(X)\in G_{\rell}(F)$.
Notons $\tau$ la fonction sur $\frg_{\epsilon}(F)$, à support topologiquement nilpotent, qui est égale sur cet ensemble au membre de droite ci-dessus. En la moyennant sur le groupe $Z_{G}(\epsilon)(F)/G_{\epsilon}(F)$, on peut la supposer invariante par $Z_{G}(\epsilon)(F)$. La fonction $\tau$ est
la fonction associée à un quasi-caractère de niveau~$0$ sur $\frg_{\epsilon}(F)$. Celui-ci est de la forme $D^{G_{\epsilon}}_{\boldsymbol{\varphi}}$ pour un élément $\boldsymbol{\varphi}\in \cD(\frg_{\epsilon})$ d'après le~\eqref{lespaceii} de la proposition \ref{lespace}. Si on se restreint aux éléments elliptiques, les distributions induites disparaissent. Donc il existe $\boldsymbol{\varphi}\in \cD_{\cusp}(\frg_{\epsilon})$ tel que $\tau$ coïncide avec $\theta_{D^{G_{\epsilon}}_{\boldsymbol{\varphi}}}$ sur $\frg_{\epsilon,\rell}(F)$. Fixons un tel $\boldsymbol{\varphi}=(\varphi_{b})_{b\in B}$, où $B$ est un sous-ensemble fini de $\Fac_{\max}(G_{\epsilon})$. Parce que $\tau$ est invariante par $Z_{G}(\epsilon)(F)$, on peut supposer que $\boldsymbol{\varphi}$ l'est aussi (pour l'action naturelle de ce groupe sur $\cD_{\cusp}(\frg_{\epsilon})$).
Puisque $A_{G_{\epsilon}}=A_{G}$, $\Imm(G_{\epsilon,\AD})$ s'identifie à $\Imm(G_{\AD})^{\epsilon}$, \cf \ref{pointsfixes}. Pour tout $b\in B$, il y a donc une unique facette $\cF_{b}\in \Imm(G_{\AD})$ telle que $\cF_{b}\cap \Imm(G_{\AD})^{\epsilon}=b$. On a $\epsilon\in K_{\cF_{b}}^{\nu}$ et $\cF_{b}^{\nu}$ est réduit à un point, autrement dit $(\cF_{b},\nu)\in \Fac^*_{\max}(G)$. De la fonction $\varphi_{b}$, on déduit une fonction $f_{b}$ sur $\bG_{\cF_{b}}^{\nu}(k_{F})$ de la façon suivante. Par l'exponentielle, on identifie~$\varphi_{b}$ à une fonction $f_{b,1}$ sur $\bG_{\epsilon,b}(k_{F})$ à support unipotent, telle que $f_{b,1}(\exp(X))=\varphi_{b}(X)$ pour tout élément nilpotent $X\in \boldsymbol{\frg}_{\epsilon,b}(k_{F})$. En notant $\bar{\epsilon}$ la réduction de $\epsilon$ dans $\bG_{\cF_{b}}^{\nu}(k_{F})$, on définit une fonction $f_{b,2}$ sur $\bG_{\cF_{b}}^{\nu}(k_{F})$, à support dans $\bar{\epsilon}\bG_{\epsilon,b}(k_{F})$, telle que $f_{b,2}(\bar{\epsilon}x)=f_{b,1}(x)$ pour tout $x\in \bG_{\epsilon,b}(k_{F})$. Enfin, on pose
\[
f_{b}(x)=\vert Z_{\bG_{\cF_{b}}}(\bar{\epsilon})(k_{F})\vert ^{-1}
\sum_{y\in \bG_{\cF_{b}}(k_{F})}f_{b,2}(y^{-1}xy)
\]
pour tout $x\in \bG_{\cF_{b}}^{\nu}(k_{F})$.

\skpt
\begin{lem*}
\begin{enumeratei}
\item\label{restrictioni}
On a l'égalité $f_{b}(\bar{\epsilon}\exp(X))=\varphi_{b}(X)$ pour tout élément nilpotent $X\in \boldsymbol{\frg}_{\epsilon,b}(k_{F})$.

\item\label{restrictionii}
la fonction $f_{b}$ appartient à $ C_{\cusp}(\bG_{\cF_{b}}^{\nu})$.
\end{enumeratei}
\end{lem*}

\begin{proof}
Rappelons que, d'après la proposition \ref{pointsfixes}, $\bG_{\epsilon,b}$ est la composante neutre du groupe $Z_{\bG_{\cF_{b}}}(\bar{\epsilon})$. Pour un élément nilpotent $X\in \boldsymbol{\frg}_{\epsilon,b}(k_{F})$, l'élément $\bar{\epsilon}\exp(X)$ a pour composante semi-simple $\bar{\epsilon}$ et pour composante unipotente $\exp(X)$. Le support de la fonction $f_{b,2}$ est formé de tel éléments. En conséquence, pour un tel élément $x=\bar{\epsilon}\exp(X)$ et pour $y\in \bG_{\cF_{b}}(k_{F})$, on n'a $f_{b,2}(y^{-1} xy)\neq0$ que si $y\in Z_{\bG_{\cF_{b}}}(\bar{\epsilon})(k_{F})$. On obtient
\[
f_{b}(x)=\vert Z_{\bG_{\cF_{b}}}(\bar{\epsilon})(k_{F})\vert ^{-1}\sum_{y\in Z_{\bG_{\cF_{b}}}(\bar{\epsilon})(k_{F})}f_{b,2}(\bar{\epsilon}\exp(y^{-1}Xy)).
\]
D'après la définition de $f_{b,2}$, le~\eqref{restrictioni} de l'énoncé résulte de l'assertion:
\begin{itemize}
\item
la fonction $\varphi_{b}$ est invariante par l'action de $Z_{\bG_{\cF_{b}}}(\bar{\epsilon})(k_{F})$.
\end{itemize}
D'après le corollaire \ref{corollaire}, ce groupe est l'image naturelle dans $\bG_{\cF_{b}}(k_{F})$ du groupe $Z_{G}(\epsilon)(F)\cap K_{\cF_{b}}^0$. Or $\boldsymbol{\varphi}$ est invariant par $Z_{G}(\epsilon)(F)$, donc $\varphi_{b}$ est invariante par $Z_{G}(\epsilon)(F)\cap K_{\cF_{b}}^0$. Cela démontre \eqref{restrictioni}.

Pour prouver \eqref{restrictionii}, on doit montrer que, pour tout espace parabolique propre $\bP^{\nu}$ de $\bG_{\cF_{b}}^{\nu}$ et pour tout $x\in \bP^{\nu}(k_{F})$, on a
\[
\sum_{u\in \bU_{\bP}(k_{F})}f_{b}(xu)=0.
\]
Il suffit de prouver que, pour tout espace parabolique propre $\bP^{\nu}$ de $\bG_{\cF_{b}}^{\nu}$ et pour tout $x\in \bP^{\nu}(k_{F})$, on a
\begin{equation}\tag{6}\label{restriction6}
\sum_{u\in \bU_{\bP}(k_{F})}f_{b,2}(xu)=0.
\end{equation}
En effet, la première assertion pour $\bP^{\nu}$ résulte de la seconde appliquée à chaque espace $y^{-1}\bP^{\nu}y$ pour $y\in \bG_{\cF_{b}}(k_{F})$. Fixons donc $\bP^{\nu}$. Si $\bP^{\nu}(k_{F})$ ne coupe pas le support de $f_{b,2}$, l'assertion \eqref{restriction6} est claire. Supposons que $\bP^{\nu}(k_{F})$ coupe ce support. Il existe donc un élément nilpotent $X\in \boldsymbol{\frg}_{\epsilon,b}(k_{F})$ tel que $\bar{\epsilon}\exp(X)\in \bP^{\nu}(k_{F})$. Puisque $\bar{\epsilon}$ et $\exp(X)$ commutent, puisque $\ad(\bar{\epsilon})$ est d'ordre premier à $p$ tandis que $\ad(\exp(X))$ est d'ordre une puissance de $p$, $\ad(\bar{\epsilon})$ appartient au groupe engendré par $\ad(\bar{\epsilon}\exp(X))$. Ce dernier opérateur conserve $\bP$, donc $\ad(\bar{\epsilon})$ conserve lui aussi $\bP$, c'est-à-dire $\bar{\epsilon}\in\nobreak \bP^{\nu}(k_{F})$. D'après le lemme \ref{p'unlemme}, on peut fixer une composante de Levi $\bM^{\nu}$ de $\bP^{\nu}$ qui contient $\bar{\epsilon}$. Notons $\bM_{\epsilon}$ et $\bP_{\epsilon}$ les composantes neutres des groupes des points fixes de l'opérateur $\ad(\bar{\epsilon})$ dans $\bM$ et $\bP$. Le groupe $\bP_{\epsilon}$ est un sous-groupe parabolique de $\bG_{\epsilon,b}$ et $\bM_{\epsilon}$ en est une composante de Levi. Montrons que
\begin{equation}\tag{7}\label{restriction7}
\bP_{\epsilon}\neq\bG_{\epsilon,b}.
\end{equation}

Notons $\bA_{\bM}$, \resp $\bA_{\cF_{b}}$ le plus grand sous-tore central déployé de $\bM$, \resp $\bG_{\cF_{b}}$. Notons $\bA_{\bM,\epsilon}$, \resp $\bA_{\cF_{b},\epsilon}$, la composante neutre du sous-groupe des points fixes de l'opérateur $\ad(\bar{\epsilon})$ dans $\bA_{\bM}$, \resp $\bA_{\cF_{b}}$. Puisque $\bP^{\nu}$ est propre, $\bA_{\bM,\epsilon}$ contient strictement $\bA_{\cF_{b},\epsilon}$. Donc $\dim(\bA_{\bM,\epsilon})>\dim(\bA_{\cF_{b},\epsilon})\geq \dim(A_{G})$. Le tore $\bA_{\bM,\epsilon}$ est contenu dans le plus grand sous-tore central déployé $\bA_{\bM_{\epsilon}}$ de $\bM_{\epsilon}$. Donc $\dim(\bA_{\bM_{\epsilon}})> \dim(A_{G})$. On a supposé $\dim(A_{G})=\dim(A_{G_{\epsilon}})$ et on~a $\dim(A_{G_{\epsilon}})=\dim(\bA_{\bG_{\epsilon,b}})$ puisque $b$ est un sommet de $\Imm(G_{\epsilon,\AD})$. Donc $\dim(\bA_{\bM_{\epsilon}})> \dim(\bA_{\bG_{\epsilon,b}})$, ce qui signifie que $\bM_{\epsilon}$ est un Levi propre de $\bG_{\epsilon,b}$. Cela démontre \eqref{restriction7}.

On peut dans \eqref{restriction6} se restreindre au cas où $x\in \bM^{\nu}(k_{F})$, c'est-à-dire $x=\bar{\epsilon}m$ avec $m\in \bM(k_{F})$. Montrons que
\begin{enumerate}\setcounter{enumi}{7}
\item\label{restriction8}
si $m\not\in \exp(\boldsymbol{\mfrm}_{\epsilon,\nil})(k_{F})$, l'ensemble des $u\in \bU_{\bP}(k_{F})$ tels que $xu$ appartient au support de $f_{b,2}$ est vide; si $m=\exp(Y)$ avec $Y\in \boldsymbol{\mfrm}_{\epsilon,\nil}(k_{F})$, alors cet ensemble est contenu dans celui des $\exp(N)$ pour $N\in \boldsymbol{\fru}_{\bP_{\epsilon}}(k_{F})$.
\end{enumerate}

Dire que $xu$ appartient au support de $f_{b,2}$ implique que $xu=\bar{\epsilon} \exp(X)$ avec $X\in \boldsymbol{\frg}_{\epsilon,b,\nil}(k_{F})$, ou encore $mu=\exp(X)$. Si on applique $\ad(\bar{\epsilon})$ à cette égalité, puisque $X$ commute à $\bar{\epsilon}$ et que $\ad(\bar{\epsilon})$ conserve $\bM$ et $\bU_{\bP}$, on obtient $mu=\ad(\bar{\epsilon})(m)\ad(\bar{\epsilon})(u)$, ce qui entraîne $m=\ad(\bar{\epsilon})(m)$ et $u=\ad(\bar{\epsilon})(u)$. Cette dernière relation implique que~$u=\exp(N)$ pour un $N\in \boldsymbol{\fru}_{\bP_{\epsilon}}(k_{F})$. La première, jointe au fait que $m$ est forcément unipotent (puisque $\exp(X)$ l'est), entraîne que $m=\exp(Y)$ pour un $Y\in \boldsymbol{\mfrm}_{\epsilon,\nil}(k_{F})$. Cela démontre \eqref{restriction8}.

On peut donc supposer $m=\exp(Y)$ avec $Y\in \boldsymbol{\mfrm}_{\epsilon,\nil}(k_{F})$. Il résulte de \eqref{restriction8} et de la définition de $f_{b,2}$ que
\[
\sum_{u\in \bU_{\bP}(k_{F})}f_{b,2}(xu)=\sum_{N\in \boldsymbol{\fru}_{\bP_{\epsilon}}(k_{F}) }\varphi_{b}(Y+N).
\]
Puisque $\bP_{\epsilon}$ est un sous-groupe parabolique propre de $\bG_{\epsilon,b}$, ceci est nul puisque $\varphi_{b}$ est cuspidale. Cela achève la preuve du lemme.
\end{proof}

Reprenons la preuve de la proposition. On définit la distribution $D^G_{f_{b}}$. Notons $\theta_{b}$ sa fonction associée. Il est clair que son support est contenu dans la réunion des conjugués du support de $f_{b}$. D'après la construction de cette fonction et le lemme~\ref{commutant}, ce support est contenu dans l'ensemble des $k^{-1}C(\epsilon)k$ pour $k\in K_{\cF_{b}}^0$. Donc le support de $\theta_{b}$ est contenu dans $C_{G}(\epsilon)$ et il nous suffit de calculer $\theta_{b}(\epsilon \exp(X))$ pour $X\in \frg_{\epsilon,\tn}(F)$.
C'est le calcul que l'on~a fait dans la section précédente. L'ensemble $\Gamma_{\epsilon}$ contient $Z_{G}(\epsilon)(F)K_{\cF_{b}}^0$. Montrons que
\begin{enumerate}\setcounter{enumi}{8}
\item\label{restriction9}
on peut remplacer l'ensemble $\Gamma_{\epsilon}$ par $Z_{G}(\epsilon)(F)K_{\cF_{b}}^0$.
\end{enumerate}

Il suffit de montrer que, si $\gamma\in \Gamma_{\epsilon}-Z_{G}(\epsilon)(F)K_{\cF_{b}}^0$ et si $g\in G_{\epsilon}(F)\gamma K_{\cF_{b}}^0$, $g^{-1}\epsilon \exp(X) g$ n'appartient pas au support de $f_{b}$. Supposons que $g^{-1}\epsilon \exp(X) g$ appartienne à ce support. Comme on l'a dit ci-dessus, il existe donc $k\in K_{\cF_{b}}^0$ et $Y\in \frg_{\epsilon,\tn}(F)$ tels que $g^{-1}\epsilon \exp(X)g=k^{-1}\epsilon \exp(Y)k$. En fixant un entier $c\geq1$ premier à $p$ tel que $\epsilon^c\in Z(G)(F)$, et en élevant l'égalité précédente à la puissance $c$, on obtient
$\exp(cg^{-1}Xg)=\exp(ck^{-1}Yk)$, d'où $\exp(g^{-1}Xg)=\exp(k^{-1}Yk)$, d'où aussi $g^{-1}\epsilon g=k^{-1} \epsilon k$. Mais alors $gk^{-1}\in Z_{G}(\epsilon)(F)$ et $\gamma\in Z_{G}(\epsilon)(F)K_{\cF_{b}}^0 $ contrairement à l'hypothèse. Cela démontre \eqref{restriction9}.

On peut donc supposer $\boldsymbol{\gamma}\subset Z_{G}(\epsilon)(F)$. On vérifie qu'alors, les constantes $m_{\gamma}$ intervenant dans la section précédente sont toutes égales. Le \eqref{restrictioni} du lemme ci-dessus assure que les fonctions $\varphi_{\gamma}$ intervenant dans l'assertion \ref{unpremiertheoreme}\,\eqref{unpremiertheoreme6} sont les images $^{\gamma}\varphi_{b}$ de~$\varphi_{b}$ par l'action de $\gamma$ sur $\cD(\frg_{\epsilon})$. Cette assertion
devient
\[
\theta_{b}(\epsilon \exp(X))=c_{b} \sum_{\gamma\in \boldsymbol{\gamma}}
\theta_{^{\gamma}\varphi_{b}}(X),
\]
où $c_{b}$ est une certaine constante non nulle.
On pose $\bmf=\sum_{b\in B}c_{b}^{-1}\vert \boldsymbol{\gamma}\vert^{-1} f_{b}$. C'est un élément de $\cD_{\cusp}(G)$.
La distribution $D^G_{\bmf}$ est à support dans $C_{G}(\epsilon)$ et, en notant $\theta'$ sa fonction associée, on a
\[
\theta'(\epsilon \exp(X))= \vert \boldsymbol{\gamma}\vert ^{-1}\sum_{\gamma\in \boldsymbol{\gamma}}\sum_{b\in B}\theta_{^{\gamma}\varphi_{b}}(X)=\vert \boldsymbol{\gamma}\vert ^{-1}\sum_{\gamma\in \boldsymbol{\gamma}}\theta_{^{\gamma}\boldsymbol{\varphi}}(X).
\]
Mais $\boldsymbol{\varphi}$ est invariante par $Z_{G}(\epsilon)(F)$ donc par $\boldsymbol{\gamma}$. D'où
\[
\theta'(\epsilon \exp(X))=\theta_{\boldsymbol{\varphi}}(X)=\tau(X),
\]
puis $\theta'(\epsilon \exp(X))\!=\!\theta_{D}(\epsilon \exp(X))$ si $X\!\in\!\frg_{\epsilon,\tn}(F)$ et $\epsilon \exp(X)\!\in\!G_{\rell}(F)$. Cela prouve~\eqref{restriction5} et la proposition.
\end{proof}

\subsection{Filtration sur le groupe et quasi-caractères}\label{filtrationsurlegroupe}

Notons $Qc(G)$ l'espace des quasi-caractères sur $G(F)$. Notons $Qc^n(G)$ la somme des $\Ind_{M}^G(Qc(M))$ sur les Levi~$M$ tels que $\dim(A_{M})\geq n$.

\begin{prop*}
Pour tout $n\in \ZZ$, on~a l'égalité
$Qc^n(G)=Qc(G)\cap \Ann^n\rI(G)^*$.
\end{prop*}

\begin{proof}
L'inclusion $Qc^n(G)\subset Qc(G)\cap \Ann^n\rI(G)^*$ est évidente. Soit $D\in Qc(G)\cap \Ann^n\rI(G)^*$, notons $\theta$ sa fonction associée. L'hypothèse que $D$ appartient à $\Ann^n\rI(G)^*$ signifie que, pour tout Levi $M$ tel que $\dim(A_{M})<n$ et tout $m\in M_{\rell}(F)\cap G_{\reg}(F)$, on~a $\theta(m)=0$. Pour $M\in \underline{\cL}_{\min}^n$, notons $\theta_{M}$ la fonction définie presque partout sur $M(F)$ par $\theta_{M}(x)=D^G(x)^{1/2}D^M(x)^{-1/2}\theta(x)$ pour $x\in M(F)$. Notons $D_{M}$ la distribution qui s'en déduit.
Par l'isomorphisme \eqref{filtrations1} de \ref{filtrations}, $D$ s'envoie sur la somme des restrictions des $D_{M}$ à $\rI_{\cusp}(M)^{W^G(M)}$. Montrons que
\begin{enumerate}
\item\label{filtrationsurlegroupe1}
il existe un quasi-caractère $D'_{M}\in Qc(M)$ tel que $D_{M}$ et $D'_{M}$ aient même restriction à $\rI_{\cusp}(M)$.
\end{enumerate}

Commençons par fixer un élément semi-simple $x\in M(F)$ et prouvons que
\begin{enumerate}\setcounter{enumi}{1}
\item\label{filtrationsurlegroupe2}
il existe un voisinage $\frV$ de $0$ dans $\mfrm_{x}(F)$ et, pour tout $\cO\in \Nil(\mfrm_{x})$, il existe un nombre complexe $c_{D,\cO}$ de sorte que, pour presque tout $Y\in \frV$ tel que $x \exp(Y)$ soit elliptique dans $M(F)$, on ait l'égalité
\[
\theta_{M}(x \exp(Y))=\sum_{\cO\in \Nil(\mfrm_{x})}c_{D,\cO}\hat{j}(\cO,Y).
\]
\end{enumerate}

Si $\dim(A_{M_{x}})> \dim(A_{M})$, un tel élément $x \exp(Y)$ n'est jamais elliptique dans $M(F)$ et l'assertion est triviale. Supposons $\dim(A_{M_{x}})=\dim(A_{M})=n$. Puisque $D$ est un quasi-caractère, on peut en tout cas fixer un voisinage $\frV'$ de $0$ dans $\frg_{x}(F)$ et, pour tout $\cO\in \Nil(\frg_{x})$, un nombre complexe $c'_{D,\cO}$ de sorte que, pour presque tout $Y\in \frV'$, on ait l'égalité
\begin{equation}\tag{3}\label{filtrationsurlegroupe3}
\theta(x \exp(Y))=\sum_{\cO\in \Nil(\frg_{x})}c'_{D,\cO}\hat{j}(\cO,Y).
\end{equation}

Notons $\tau$ la fonction du membre de droite, qui est définie presque partout sur $\frg_{x}(F)$.
Soit $L^x$ un Levi de $G_{x}$ tel que $\dim(A_{L^x})<n$ et soit $Y\in \frl^x_{\rell}(F)\cap \frV'$. On note~$L$ le commutant de $A_{L^x}$ dans $G$. Alors $x\in L(F)$, $L_{x}=L^x$, $A_{L}=A_{L^x}$ et $x \exp(Y)$ est elliptique dans $L(F)$. Puisque $\dim(A_{L})<n$, $\theta(x \exp(Y))=0$. Donc $\tau(Y)=0$. Puisque $\tau$ est combinaison linéaire de fonctions homogènes pour l'action de $F^{\times,2}$ par homothéties, on peut supprimer ci-dessus la restriction: $Y\in \frV'$. La restriction de $\tau$ à $\frg_{x,\tn}(F)$ est donc la fonction associée à un élément de $Qc_{0}(\frg_{x})\cap \Ann^{n}\rI(\frg_{x})^*$, avec les notations de \ref{filtrationssurlalgebredelie}. D'après le lemme de ce paragraphe, on peut écrire cet élément comme somme sur les $L^x\in \underline{\cL}_{\min}^{G_{x},n}$ de distributions $\Ind_{L^x}^{G_{x}}(d_{L^x})$, où $d_{L^x}\in Qc(\frl^x)$. On~peut supposer que $M_{x}\in \underline{\cL}_{\min}^{G_{x},n}$ et que $d_{M_{x}}$ est invariante par $ W^{G_{x}}(M_{x})$. Notons~$\tau_{x}$ la fonction associée à $d_{M_{x}}$. En utilisant \ref{induction}\,\eqref{induction3}, on vérifie que l'on~a l'égalité
\[
\tau(Y)=cD^{G_{x}}(Y)^{-1/2}D^{M_{x}}(Y)^{1/2}\tau_{x}(Y)
\]
pour presque tout $Y\in
\mfrm_{x,\rell}(F)\cap \mfrm_{x,\tn}(F)$, où $c=\vert W^{G_{x}}(M_{x})\vert $. D'autre part, si $\frV'$ est assez petit, il existe $c'>0$ tel que
\[
D^G(x \exp(Y))^{1/2}D^M(x \exp(Y))^{-1/2}=c'D^{G_{x}}(Y)^{1/2}D^{M_{x}}(Y)^{-1/2}
\]
pour tout $Y\in \frV'\cap \mfrm_{x}(F)$ tel que $x \exp(Y)\in G_{\reg}(F)$. L'égalité \eqref{filtrationsurlegroupe3} entraîne
\[
\theta_{M}(x \exp(Y))=cc'\tau_{x}(Y)
\]
pour presque tout $Y\in
\frV'\cap\mfrm_{x,\rell}(F)\cap \mfrm_{x,\tn}(F)$.
Cette fonction est de la forme du membre de droite de l'égalité de l'assertion \eqref{filtrationsurlegroupe2}. En prenant $\frV=\frV'\cap \mfrm_{x}(F)$, cela résout cette assertion \eqref{filtrationsurlegroupe2}.

On peut supposer que le membre de droite de l'égalité de cette assertion est invariant par l'action du groupe $\Norm_{G_{x}}(M_{x})(F) $. Posons
\[
V_{x}=\{m^{-1}x \exp(Y)m\sep m\in M(F), Y\in \frV\}.
\]
Quitte à restreindre $\frV$, cet ensemble est ouvert et fermé dans $M(F)$. On définit un quasi-caractère $d_{x}$ de $M(F)$, à support dans $V_{x}$, dont la fonction associée satisfait à
\[
\theta_{d_{x}}(m^{-1}x \exp(Y)m)=\sum_{\cO\in \Nil(\mfrm_{x})}c_{D,\cO}\hat{j}(\cO,Y)
\]
pour tout $m\in M(F)$ et $Y\in \frV$.
Cette fonction coïncide avec $\theta_{M}$ sur $V_{x}\cap M_{\rell}(F)$. On construit ces ensembles $V_{x}$ et ces quasi-caractères $d_{x}$ pour tout élément semi-simple $x\in M(F)$. La réunion des $V_{x}$ est $M(F)$ tout entier, on peut en extraire un recouvrement localement fini $M(F)=\bigcup_{j\in J}V_{x_{j}}$. Il existe une famille $(\alpha_{j})_{j\in J}$, où~$\alpha_{j}$ est une fonction localement constante sur $M(F)$, invariante par conjugaison et à support dans $V_{x_{j}}$, de sorte que $\sum_{j\in J}\alpha_{j}(m)=1$ pour tout $m\in M(F)$. En utilisant la remarque \eqref{quasicaracteres5} de \ref{quasicaracteres}, la distribution
\[
D'_{M}=\sum_{j\in J}\alpha_{j}d_{x_{j}}
\]
est bien définie puisque le recouvrement est localement fini, c'est un quasi-caractère et elle coïncide avec $D_{M}$ sur les éléments elliptiques de $M(F)$, donc elle a même restriction que $D_{M}$ à $\rI_{\cusp}(M)$. Cela démontre \eqref{filtrationsurlegroupe1}.

La distribution
\[
\sum_{M\in \underline{\cL}^n} \vert W^G(M)\vert ^{-1}\Ind_{M}^G(D'_{M})
\]
appartient à $Qc^n(G)$ et, d'après \ref{filtrations}, son image dans $\Gr^n\rI(G)^*$ est égale à celle de $D$. Cela démontre l'inclusion
\[
Qc(G)\cap \Ann^{n}\rI(G)^*\subset (Qc(G)\cap \Ann^{n+1}\rI(G)^*)+Qc^n(G).
\]
L'inclusion $Qc(G)\cap \Ann^{n}\rI(G)^*\subset Qc^n(G)$ s'en déduit par récurrence descendante sur~$n$. Cela démontre la proposition.
\end{proof}

\subsection{Un deuxième théorème}\label{undeuxiemetheoreme}

Notons $Qc_{0}(G)$ l'espace des quasi-caractères de niveau~$0$ sur $G(F)$.

\begin{thm*}
L'application $D^G:\cD(G)\to \rI(G)^*$ est un isomorphisme de $\cD(G)$ sur $Qc_{0}(G)$.
\end{thm*}

\begin{proof}
Compte tenu du théorème \ref{unpremiertheoreme}, il suffit de prouver que $Qc_{0}(G)$ est contenu dans l'image de l'application $D^G$. Soit $n\in \ZZ$. Prouvons que
\begin{enumerate}
\item\label{undeuxiemetheoreme1}
$Qc_{0}(G)\cap \Ann^n\rI(G)^*$ est contenu dans la somme de $D^G(\cD^n(G))$ et de $Qc_{0}(G)\cap \Ann^{n+1}\rI(G)^*$.
\end{enumerate}

Remarquons que $D^G(\cD^n(G))$ est contenu dans $Qc_{0}(G)\cap \Ann^n\rI(G)^*$ d'après \ref{unpremiertheoreme}. Soit $D\in Qc_{0}(G)\cap \Ann^n\rI(G)^*$. Pour $M\in \underline{\cL}_{\min}^n$, on~a défini dans le paragraphe précédent la distribution $D_{M}$ sur $M(F)$, de fonction associée $\theta_{M}$, et on~a prouvé l'existence d'un quasi-caractère $D'_{M}$ sur $M(F)$ qui coïncidait avec $D_{M}$ sur $M_{\rell}(F)$. Montrons que
\begin{enumerate}\setcounter{enumi}{1}
\item\label{undeuxiemetheoreme2}
$D'_{M}$ est de niveau~$0$ sur les elliptiques.
\end{enumerate}

D'après les remarques \eqref{restriction2} et \eqref{restriction4} de \ref{restriction} et le lemme \ref{lecasdunlevi}, on peut fixer $\epsilon\in G(F)_{p'}\cap M(F)$ et démontrer que, pour tout $X\in \mfrm_{\epsilon,\tn}(F)$ tel que $\epsilon \exp(X)\in M_{\rell}(F)$, la fonction $\lambda\mto \theta_{D'_{M}}(\epsilon \exp(\lambda^2X))$ appartient à $E$. Pour un tel $X$, on~a par définition
\begin{multline*}
\theta_{D'_{M}}(\epsilon \exp(\lambda^2X))\\
=\theta_{M}(\epsilon \exp(\lambda^2X))=D^G(\epsilon \exp(\lambda^2X))^{1/2}D^M(\epsilon \exp(\lambda^2X))^{-1/2}\theta_{D}(\epsilon \exp(\lambda^2X)).
\end{multline*}
D'après le~\eqref{p'decompositionii} du lemme \ref{p'decomposition}, la fonction
\[
\lambda\mto D^G(\epsilon \exp(\lambda^2X))^{1/2}D^M(\epsilon \exp(\lambda^2X))^{-1/2}
\]
est produit d'une constante et d'une puissance entière de $\vert \lambda\vert $. Puisque $D\in Qc_{0}(G)$, la fonction $\lambda\mto \theta_{D}(\epsilon \exp(\lambda^2X))$ appartient à $E$. Il en résulte que la fonction $\lambda\mto \theta_{D'_{M}}(\epsilon \exp(\lambda^2X))$ appartient à $E$, d'où \eqref{undeuxiemetheoreme2}.

D'après la proposition \ref{restriction}, il existe donc $D''_{M}\in D^M(\cD_{\cusp}(M))$ tel que $D''_{M}$ coïncide avec $D'_{M}$ sur $M_{\rell}(F)$, donc aussi avec $D_{M}$. La distribution
\[
\sum_{M\in \underline{\cL}_{\min}^n} \vert W^G(M)\vert ^{-1}\Ind_{M}^G(D''_{M})
\]
appartient alors à $D^G(\cD^n(G))$ et son image dans $\Gr^n\rI(G)^*$ est égale à celle de $D$. Cela prouve \eqref{undeuxiemetheoreme1}.

Pour $n=a_{M_{\min}}+1$, on~a $Qc_{0}(G)\cap \Ann^n\rI(G)^*=\{0\}=D^G(\cD^n(G))$. Grâce à \eqref{undeuxiemetheoreme1}, on démontre par récurrence descendante que $Qc_{0}(G)\cap \Ann^n\rI(G)^*=D^G(\cD^n(G))$ pour tout $n$. Pour $n=a_{G}$, cela démontre l'égalité $Qc_{0}(G)=D^G(\cD(G))$.
\end{proof}

\section{Représentations et représentations de niveau \texorpdfstring{$0$}{0}}\label{representations}

\subsection{Représentations de niveau~\texorpdfstring{$0$}{0}}
Notons $\Irr(G)$ l'ensemble des classes d'isomorphismes de représentations lisses irréductibles de $G(F)$ (dans des espaces complexes). Une telle représentation $\pi$, dans un espace $V$, est dite de niveau~$0$ si et seulement s'il existe une facette $\cF\in \Fac(G)$ telle que l'espace des invariants $V^{K_{\cF}^+}$ soit non nul. On note $\Irr(G)^0$ le sous-ensemble des (classes d'isomorphismes de) représentations de niveau~$0$ et $\Irr(G)^{>0}$ celui des représentations qui ne sont pas de niveau~$0$. On~note~$p^0 $ la projection $\CC[\Irr(G)]\to \CC[\Irr(G)^0]$ qui annule $\CC[\Irr(G)^{>0}]$. On sait que~$p^0$ est associée à l'action d'un idempotent du centre de Bernstein, \cf \cite[Cor.\,3.9(i)]{B}. Il se déduit de cet idempotent une action sur divers objets, par exemple $C_{\rc}^{\infty}(G(F))$, que l'on~note encore $p^0$. On sait que $p^0$ commute, en un sens compréhensible, aux opérations d'induction et de passage au module de Jacquet: par exemple, si $P$ est un sous-groupe parabolique de $G$ de composante de Levi $M$ et si $\pi^M\in \Irr(M)$, alors $p^0(\Ind_{P}^G(\pi^M))=\Ind_{P}^G(p^0(\pi^M))$, \cf \cite[Th.\,5.2]{MP}.

On note $\Temp(G)$ le sous-ensemble de $\Irr(G)$ des représentations tempérées et on pose $\Temp(G)^0=\Temp(G)\cap \Irr(G)^0$.

Si $\xi$ est un caractère de $A_{G}(F)$, on~note $\Irr(G)_{\xi}$ le sous-ensemble de $\Irr(G)$ des représentations dont le caractère central coïncide avec $\xi$ sur $A_{G}(F)$. On définit de même $\Irr(G)_{\xi}^0$, $\Temp(G)_{\xi}$ etc. Évidemment, $\Irr(G)_{\xi}^0$ est vide si $\xi$ n'est pas modérément ramifié et $\Temp(G)_{\xi}$ est vide si $\xi$ n'est pas unitaire.

Pour toute représentation irréductible $\pi$, on~note $\Theta_{\pi}$
son caractère-distribution. D'après Harish-Chandra, c'est un quasi-caractère, \cf \cite[Th.\,16.2]{HC}. On note $\theta_{\pi}$ sa fonction associée. Par linéarité, l'application $\pi\mto \Theta_{\pi}$ se prolonge en une application
\[
\Theta:\CC[\Irr(G)]\to \rI(G)^*
\]
qui est injective.

\subsection{Représentations elliptiques}\label{representationselliptiques}
Arthur a défini un ensemble de \og représentations elliptiques\fg, \cf \cite[\S3]{A2}. Une telle représentation est une combinaison linéaire finie, à coefficients dans $\CC$, de représentations irréductibles tempérées. Sa définition dépend de divers choix, en particulier de normalisations d'opérateurs d'entrelacement. Nous supposons ces choix fixés et nous notons $\Ell(G)$ l'ensemble des représentations ellip\-ti\-ques. C'est un ensemble linéairement indépendant. L'espace $\CC[\Ell(G)]$ ne dépend pas des choix. Pour un Levi $M\in \cL_{\min}$, le sous-espace $\CC[\Ell(M)]\subset \CC[\Temp(M)]$ est invariant par l'action naturelle du groupe $ W^G(M)$ sur $\CC[\Temp(M)]$. On note $\CC[\Ell(M)]^{W^G(M)}$ le sous-espace des invariants. On a alors la décomposition en somme directe
\[
\CC[\Temp(G)]=\bigoplus_{M\in\underline{\cL}_{\min}}\Ind_{M}^G(\CC[\Ell(M)]^{W^G(M)}).
\]

Introduisons l'application d'Harish-Chandra $H_{G}:G(F)\to \cA_{G}$, notons $\cA_{G}^*$ l'espace vectoriel réel dual de $\cA_{G}$ et $\cA_{G,\CC}^*$ son complexifié. Un élément $\chi\in \cA_{G,\CC}^*$ définit un caractère de $G(F)$, que l'on~note encore $\chi$, défini par $\chi(g)=\exp(\<\chi,H_{G}(g)\>)$. Si $\pi\in \Irr(G)$, on~note $\pi_{\chi}$ la représentation $\pi_{\chi}(g)=\chi(g)\pi(g)$. Si $\chi\in i\cA_{G}^*$, cette opération $\pi\mto \pi_{\chi}$ conserve l'espace $\CC[\Ell(G)]$. On note $\Ell(G)_{\RR}$ l'ensemble des $\pi_{\chi}$ pour $\pi\in \Ell(G)$ et $\chi\in \cA_{G}^*$. On a alors la décomposition plus générale
\begin{equation}\tag{1}\label{representationselliptiques1}
\CC[\Irr(G)]=\bigoplus_{M\in\underline{\cL}_{\min}}\Ind_{M}^G(\CC[\Ell(M)_{\RR}]^{W^G(M)}).
\end{equation}

Via l'injection $\Theta$, cette décomposition est compatible avec la filtration de $\rI(G)^*$ définie en \ref{filtrations}. C'est-à-dire que, pour $n\in \ZZ$, l'intersection de $\Ann^n\rI(G)^*$ et de l'image de $\Theta$ est égale à l'image par cette injection de la sous-somme de l'expression ci-dessus, où on limite la sommation aux $M$ tels que $\dim(A_{M})\geq n$. En particulier, l'application composée
\begin{equation}\tag{2}\label{representationselliptiques2}
\CC[\Ell(G)_{\RR}]\To{\Theta}\rI(G)^*\to \rI_{\cusp}(G)^*
\end{equation}
est injective.

En fait, chaque représentation elliptique est combinaison linéaire de sous-repré\-sen\-tations d'une même induite d'une série discrète d'un Levi de $G$. Il en résulte d'une part que toute représentation elliptique admet un caractère central. D'autre part que ou bien toutes ces composantes sont de niveau~$0$, ou bien aucune ne l'est. On~note $\Ell(G)^0$, \resp $\Ell(G)^{>0}$, les représentations elliptiques du premier cas, \resp du second. Puisque $p^0$ \og commute\fg à l'induction, on obtient des variantes des égalités précédentes telles que
\begin{align*}\tag{3}\label{representationselliptiques3}
\CC[\Temp(G)^0]&=\bigoplus_{M\in\underline{\cL}_{\min}}\Ind_{M}^G(\CC[\Ell(M)^0]^{W^G(M)});\\\CC[\Irr(G)^0]&=\bigoplus_{M\in\underline{\cL}_{\min}}\Ind_{M}^G(\CC[\Ell(M)^0_{\RR}]^{W^G(M)});\\
\CC[\Irr(G)^{>0}]&=\bigoplus_{M\in\underline{\cL}_{\min}}\Ind_{M}^G(\CC[\Ell(M)_{\RR}^{>0}]^{W^G(M)}).
\end{align*}
Il y a aussi des variantes de ces égalités en imposant que toutes les représentations se transforment selon un caractère fixé $\xi$ de $A_{G}(F)$.

\subsection{Le résultat de \texorpdfstring{\cite{W2}}{W2}}\label{leresultat}

Soit $\pi\in \Irr(G)^0$. En \cite{W2}, on~a associé à $\pi$ un élément de $ \cD_{\cusp}(G)$ que l'on avait noté $\Theta^G_{\pi,\cusp}$. On modifie légèrement la définition en y incluant les mesures $\mes(A_{G}(F)\backslash K_{\cF}^{\dag})^{-1}$ qui intervenaient dans les formules de cette référence. On note $\Delta^G_{\pi,\cusp}$ ou simplement $\Delta_{\pi,\cusp}$ l'élément de $\cD_{\cusp}(G)$ ainsi obtenu.

Soit $M$ un Levi de $G$. Choisissons un sous-groupe parabolique $P\in \cP(M)$. On associe au module de Jacquet $\pi_{P}$ un élément $\Delta^M_{\pi_{P},\cusp}$. Notons $\Delta^M_{\pi_{M},\cusp,G\tcomp}$ sa projection dans $\cD_{\cusp,G\tcomp}(M)$, autrement dit sa restriction au sous-ensemble des éléments de $M(F)$ qui sont compacts $\rmod Z(G)$. Elle ne dépend pas du choix de $P$.

Le résultat principal de \cite{W2} est que
\begin{enumerate}
\item\label{leresultat1}
$\Theta_{\pi}$ coïncide sur $G_{\comp}(F)$ avec la distribution
\[
\sum_{M\in \cL_{\min}}\vert W^M\vert \vert W^G\vert ^{-1}D^G[\Delta^M_{\pi_{M},\cusp,G\tcomp}].
\]
\end{enumerate}

Autrement dit, notons $\Delta_{\pi}$ l'image naturelle de
\[
\sum_{M\in \cL_{\min}}\vert W^M\vert \vert W^G\vert ^{-1}\Delta^M_{\pi_{M},\cusp,G\tcomp}
\]
dans $\cD_{G\tcomp}(G)$. Alors $\Theta_{\pi}$ coïncide avec $D^G[\Delta_{\pi}]$ sur $G_{\comp}(F)$. On prolonge $\pi\mto \Delta_{\pi}$ par linéarité en une application
\[
\Delta:\CC[\Irr(G)^0]\to \cD_{G\tcomp}(G).
\]

\begin{rem*}
La somme intervenant dans \eqref{leresultat2} peut se récrire
\begin{equation}\tag{2}\label{leresultat2}
\sum_{M\in \underline{\cL}_{\min}}\vert W^G(M)\vert ^{-1}D^G[\Delta^M_{\pi_{M},\cusp,G\tcomp}].
\end{equation}
\end{rem*}

\subsection{Surjectivité}\label{surjectivite}

Soit $\xi$ un caractère unitaire et modérément ramifié de $A_{G}(F)$. L'application $\Delta$ se restreint en une application linéaire $\CC[\Temp(G)_{\xi}^0]\to \cD_{G\tcomp,\xi}(G)$.

\skpt
\begin{lem*}
\begin{enumeratei}
\item\label{surjectivitei}
L'application $\Delta:\CC[\Temp(G)_{\xi}^0]\to \cD_{G\tcomp,\xi}(G)$ est surjective.

\item\label{surjectiviteii}
L'application $\Delta_{\cusp}:\CC[\Ell(G)_{\xi}^0]\to \cD_{\cusp,\xi}(G)$ est bijective.
\end{enumeratei}
\end{lem*}

\begin{proof}
Notons ici $q_{\xi}$ la projection de $C_{\rc}^{\infty}(G(F))$ dans $C_{c,\xi}^{\infty}(G(F))$ ou de $\rI(G)$ dans $I_{\xi}(G)$. On a le diagramme naturel
\[
\xymatrix@C=.5cm@R=.3cm{
&\cD_{G\tcomp,\xi}(G)\ar@{=}[r]&\cD_{G\tcomp,\xi}(G)\ar[dd]_-{D^G}\ar[dr]^-{p_1}&\\
\CC[\Temp(G)_{\xi}^0]\ar[ur]^-\Delta\ar[dr]_-\Theta&&&(q_{\xi}(\rIE(G)))^*\\
&I_{\xi}(G)^*\ar[r]^-{p_{3}}&I_{\xi}(G)_{\comp}^*\ar[ur]_-{p_{2}}
}
\]
Les applications $p_{2}$ et $p_{3}$ sont les restrictions et $p_{1}\!=\!p_{2}\circ D^G$. Le triangle de droite est donc commutatif. La partie gauche du diagramme l'est d'après \ref{leresultat}\,\eqref{leresultat1}. Montrons que
\begin{enumerate}
\item\label{surjectivite1}
$p_{2}\circ p_{3}\circ \Theta$ est surjective.
\end{enumerate}

Notons que l'espace $(q_{\xi}(\rIE(G)))^*$ est de dimension finie. Chaque élément de $\cE(G)$ est invariant par un groupe $K_{\cF}^+$ pour une certaine facette $\cF\in \Fac(G)$. Donc $\cE(G)$ est annulé par toute représentation qui n'est pas de niveau~$0$. D'après \cite[Th.\,0(d)]{Kh} (plus exactement une variante avec caractère central), l'application $\CC[\Temp(G)_{\xi}]\to (q_{\xi}(\rIE(G)))^*$ similaire à $p_{2}\circ p_{3}\circ \Theta$ est surjective. D'après la propriété précédente, on peut aussi bien remplacer $\Temp(G)_{\xi}$ par $\Temp(G)_{\xi}^0$, d'où \eqref{surjectivite1}.

Une variante avec caractère central de la proposition \ref{deuxespaces} nous dit que l'application~$p_{1}$ est un isomorphisme. Alors le~\eqref{surjectivitei} de l'énoncé se déduit de \eqref{surjectivite1}.

En utilisant \ref{leresultat}\,\eqref{leresultat1}, l'injectivité de l'application du \eqref{surjectiviteii} se déduit de celle de l'application \eqref{representationselliptiques2} de \ref{representationselliptiques}.
D'après le~\eqref{surjectivitei} déjà démontré et la relation \eqref{representationselliptiques3} de \ref{representationselliptiques}, pour démontrer la surjectivité de l'application du \eqref{surjectiviteii}, il reste à prouver que,
si $M$ est un Levi propre de $G$, si $\pi^M\in \CC[\Ell(M)^0_{\xi}]$ et si l'on~note $\pi=\Ind_{M}^G(\pi^M)$, alors $\Delta_{\pi,\cusp}=0$. La projection $\rI(G)^*\to \rI_{\cusp}(G)^*$ annule toutes les distributions dont le support ne coupe pas $G(F)_{\comp}$, car tout élément elliptique régulier est compact $\rmod Z(G)$. La projection de $\Theta_{\pi}$ est donc égale à celle de $D^G[\Delta_{\pi}]$. Cette projection $\rI(G)^*\to \rI_{\cusp}(G)^*$ annule aussi toutes les distributions induites à partir d'un Levi propre. Il en résulte que la projection de $\Theta_{\pi}$ est nulle et que la projection de $D^G[\Delta_{\pi}]$ est égale à celle de $D^G[\Delta_{\pi,\cusp}]$. Donc la projection dans $\rI_{\cusp}(G)^*$ de $D^G[\Delta_{\pi,\cusp}]$ est nulle. D'après le~\eqref{uncorollaireii} du corollaire \ref{uncorollaire}, on~a donc $\Delta_{\pi,\cusp}=0$, ce qui achève la démonstration.
\end{proof}

\subsection{Produit scalaire elliptique}\label{produitscalaireelliptique}
Fixons un caractère unitaire $\xi$ de $A_{G}(F)$. L'espace
\[
\CC[\Temp(G)_{\xi}]=\bigoplus_{M\in\underline{\cL}_{\min}}\Ind_{M}^G(\CC[\Ell(M)_{\xi}]^{W^G(M)})
\]
est muni d'un produit scalaire elliptique noté $\<.,\>_{\rell}$, \cf \cite[\S6]{A2}. Les composantes ci-dessus indexées par un Levi propre sont dans le noyau de ce produit. Par contre, le produit se restreint en un produit hermitien défini positif sur $\CC[\Ell(G)_{\xi}]$, pour lequel $\Ell(G)_{\xi}$ est une base orthogonale. Rappelons la définition du produit scalaire. Fixons un ensemble $\cT_{\rell}$ de représentants des classes de conjugaison de sous-tores elliptiques maximaux de $G$. Pour $T\in \cT_{\rell}$, on pose $W(T)=\Norm_{G}(T)(F)/T(F)$ et on munit $A_{G}(F)\backslash T(F)$ de la mesure de Haar de masse totale $1$. Pour $\pi,\pi'\in \CC[\Temp(G)_{\xi}]$, on~a alors
\begin{equation}\tag{1}\label{produitscalaireelliptique1}
\<\pi,\pi'\>_{\rell}=\sum_{T\in \cT_{\rell}}\vert W(T)\vert ^{-1}\int_{A_{G}(F)\backslash T(F)}\overline{\theta_{\pi}(t)}\theta_{\pi'}(t)D^G(t)\,dt.
\end{equation}
Remarquons que ceci ne dépend d'aucune mesure.

On va munir $\cD_{\cusp,\xi}(G)$ d'un produit scalaire que l'on appelle aussi elliptique. Pour $\cF,\cF'\in \Fac(G)$,
notons $N(\cF,\cF')$ l'ensemble des $g\in G(F)$ tels que $g\cF=\cF'$ et fixons un ensemble de représentants $\underline{N}(\cF,\cF')$ du quotient
\[
A_{G}(F)K_{\cF'}^0\backslash N(\cF,\cF').
\]
On a établi en \ref{variantes} une décomposition
\begin{equation}\tag{2}\label{produitscalaireelliptique2}
\boldsymbol{\cD}_{\cusp}(G)=\prod_{\nu\in \cN}
\boldsymbol{\cD}_{\cusp}(G)^{\nu}.
\end{equation}
On écrit $\bmf=\prod_{\nu\in \cF}\bmf^{\nu}$ la décomposition d'un élément $\bmf\in \boldsymbol{\cD}_{\cusp}(G)$.
Fixons $\nu\in \cN$. Soient $(\cF,\nu),(\cF',\nu)\in \Fac^*_{\max}(G)$ et soient
$f\in C_{\cusp}(\bG_{\cF}^{\nu})$ et $f'\in C_{\cusp}(\bG_{\cF'}^{\nu})$. On pose
\[
\<f,f'\>_{\rell}= \mes(A_{G}(F)\backslash A_{G}(F)K_{\cF}^0)^2\sum_{g\in \underline{N}(\cF,\cF')}\<{\gf},f'\>.
\]
Ce produit se prolonge par linéarité en un produit sur $\boldsymbol{\cD}_{\cusp}(G)^{\nu}$.
Soient $\bmf, \bmf'\in \boldsymbol{\cD}_{\cusp,\xi}(G)$. On voit que $\<\bmf^{\nu+\nu'},\bmf^{\prime\nu+\nu'}\>=\<\bmf^{\nu},\bmf^{\prime\nu}\>$ pour tout $\nu\in \cN$ et tout $\nu'\in w_{G}(A_{G}(F))$. On pose
\[
\<\bmf,\bmf'\>_{\rell}=\sum_{\nu\in \cN/w_{G}(A_{G}(F))}\<\bmf^{\nu},\bmf^{\prime\nu}\>_{\rell}.
\]
On vérifie que, pour $g\in G(F)$, $\<{\gbmf},\bmf'\>_{\rell}=\<\bmf,\bmf'\>_{\rell}$. C'est-à-dire que $\gbmf-\bmf$ appartient au noyau du produit scalaire elliptique. Celui-ci se descend donc en un produit sur $\cD_{\cusp,\xi}(G)$.

\begin{prop*}
L'application
\begin{align*}
\CC[\Ell(G)_{\xi}^0]&\to \cD_{\cusp,\xi}(G)\\
\pi&\mto \Delta_{\pi,\cusp}
\end{align*}
est un isomorphisme isométrique pour les produits scalaires elliptiques.
\end{prop*}

\begin{proof}
On peut supposer que l'ensemble de représentants $\underline{\Fac}^*_{\max}(G,A)$ est stable par l'opération $(\cF,\nu)\mto (\cF,\nu+\nu')$ pour tout $\nu'\in w_{G}(A_{G}(F))$. Représentons $\Delta_{\pi,\cusp}$ comme un élément de
\[
\prod_{(\cF,\nu)\in \underline{\Fac}^*_{\max}(G,A)}C_{\cusp}(\bG_{\cF}^{\nu})^{K_{\cF}^{\dag}},
\]
\cf \ref{deuxespaces}. On note $(f_{\pi,\cF,\nu})_{(\cF,\nu)\in \underline{\Fac}^*_{\max}(G,A)}$ cet élément.
Soit
\[
T\in \cT_{\rell}\quadet t\in T(F)\cap G_{\reg}(F).
\]
D'après \cite[18(3)]{W2}, on~a l'égalité
\[
\theta_{\pi}(t)=D^G(t)^{-1/2}\sum_{(\cF,\nu)\in \underline{\Fac}^*_{\max}(G,A)} I^G(t,f_{\pi,\cF,\nu}).
\]
On a identifié ici l'élément $f_{\pi,\cF,\nu}$ à la fonction $(f_{\pi,\cF,\nu})_{\cF}$. La somme est finie car les termes ne sont non nuls que si $\nu=w_{G}(t)$.

\begin{rem*}
La formule de la référence contient un terme $\mes(A_{G}(F)\backslash K_{\cF}^{\dag})^{-1}$ que l'on~a incorporé à la définition de $f_{\pi,\cF,\nu}$.
\end{rem*}

Fixons un ensemble de représentants $\underline{\cN}$ des orbites dans $\cN$ pour l'action de $w_{G}(A_{G}(F))$ par addition. Notons $ \underline{\Fac}^*_{\max}(G,A;\underline\cN)$ le sous-ensemble des éléments $(\cF,\nu)\in \underline{\Fac}^*_{\max}(G,A)$ tels que $\nu\in \underline{\cN}$.
La formule ci-dessus se récrit
\[
\theta_{\pi}(t)=D^G(t)^{-1/2}\sum_{(\cF,\nu)\in \underline{\Fac}^*_{\max}(G,A;\underline{\cN})} \sum_{\nu'\in w_{G}(A_{G}(F))}I^G(t,f_{\pi,\cF,\nu+\nu'}).
\]
Parce que la restriction à $A_{G}(F)$ du caractère central de $\pi$ est $\xi$, on voit que la somme intérieure n'est autre que
\[
\mes(A_{G}(F)_{\rc})^{-1}\int_{A_{G}(F)}I^G(ta,f_{\pi,\cF,\nu})\xi(a)^{-1}\,da.
\]
D'où
\[
\theta_{\pi}(t)=\mes(A_{G}(F)_{\rc})^{-1}\int_{A_{G}(F)}\hspace*{-2mm} D^G(ta)^{-1/2}\hspace*{-2mm}\sum_{(\cF,\nu)\in \underline{\Fac}^*_{\max}(G,A;\underline{\cN})}\hspace*{-2mm}I^G(ta,f_{\pi,\cF,\nu})\xi(a)^{-1}\,da.
\]
La formule~\eqref{produitscalaireelliptique1} devient
\begin{multline*}
\<\pi,\pi'\>_{\rell}=\mes(A_{G}(F)_{\rc})^{-1}\\[-5pt]
\cdot\sum_{(\cF,\nu)\in \underline{\Fac}^*_{\max}(G,A;\underline{\cN})} \sum_{T\in \cT_{\rell}}\vert W(T)\vert ^{-1}\int_{T(F)}D^G(t)^{1/2}I^G(t,\overline{f_{\pi,\cF,\nu}})\theta_{\pi'}(t)\,dt.
\end{multline*}
Les fonctions $f_{\pi,\cF,\nu}$ sont cuspidales, leurs intégrales orbitales sont nulles hors des éléments elliptiques. La formule de Weyl conduit donc à l'égalité
\[
\<\pi,\pi'\>_{\rell}=\mes(A_{G}(F)_{\rc})^{-1}\sum_{(\cF,\nu)\in \underline{\Fac}^*_{\max}(G,A;\underline{\cN})} \Theta_{\pi'}(\overline{f_{\pi,\cF,\nu}}).
\]
Fixons $(\cF,\nu)$ intervenant ci-dessus. Puisque la fonction $f_{\pi,\cF,\nu}$ est à support compact $\rmod Z(G)$, on peut exprimer $\Theta_{\pi'}(\overline{f_{\pi,\cF,\nu}})$ à l'aide de $\Delta_{\pi'}$. Puisque la fonction $f_{\pi,\cF,\nu}$ est cuspidale, les distributions induites ne comptent pas, d'où
\[
\Theta_{\pi'}(\overline{f_{\pi,\cF,\nu}})=D^G[\Delta_{\pi',\cusp}](\overline{f_{\pi,\cF,\nu}}).
\]
Ceci est calculé par la proposition \ref{calcul} et la remarque \eqref{calcul2} qui la suit. On obtient
\[
\Theta_{\pi'}(\overline{f_{\pi,\cF,\nu}})=c_{\cF}\<f_{\pi,\cF,\nu},f_{\pi',\cF,\nu}\>,
\]
où
\[
c_{\cF}=\mes(K_{\cF}^0)^2\mes(A_{G}(F)_{\rc})^{-1}[K_{\cF}^{\dag}:A_{G}(F)K_{\cF}^0].
\]
De la même façon, on a
\[
\<f_{\pi,\cF,\nu},f_{\pi',\cF,\nu}\>_{\rell}=d_{\cF}\<f_{\pi,\cF,\nu},f_{\pi',\cF,\nu}\>,
\]
où
\[
d_{\cF}=\mes(A_{G}(F)\backslash A_{G}(F)K_{\cF}^0)^2[K_{\cF}^{\dag}:A_{G}(F)K_{\cF}^0].
\]
D'où
\[
\Theta_{\pi'}(\overline{f_{\pi,\cF,\nu}})= c_{\cF}d_{\cF}^{-1}\<f_{\pi,\cF,\nu},f_{\pi',\cF,\nu}\>_{\rell},
\]
puis
\[
\<\pi,\pi'\>_{\rell}=\mes(A_{G}(F)_{\rc})^{-1}\sum_{(\cF,\nu)\in \underline{\Fac}^*_{\max}(G,A;\underline{\cN})} c_{\cF}d_{\cF}^{-1}\<f_{\pi,\cF,\nu},f_{\pi',\cF,\nu}\>_{\rell}.
\]
On vérifie que
\[
c_{\cF}=\mes(A_{G}(F)_{\rc}) d_{\cF}.
\]
D'où
\[
\<\pi,\pi'\>_{\rell}=\sum_{(\cF,\nu)\in \underline{\Fac}^*_{\max}(G,A;\underline{\cN})}\<f_{\pi,\cF,\nu},f_{\pi',\cF,\nu}\>_{\rell}.
\]
Par définition, la membre de droite n'est autre que $\<\Delta_{\pi,\cusp},\Delta_{\pi',\cusp}\>_{\rell}$. Cela démontre que l'application de l'énoncé est isométrique. Elle est bijective d'après le~\eqref{surjectiviteii} du lemme \ref{surjectivite}. Cela achève la démonstration.
\end{proof}

\subsection{Un lemme préliminaire}\label{unlemme}

Rappelons que l'on~a défini en \ref{restriction} la notion de quasi-caractère sur $G(F)$ de niveau~$0$ sur les elliptiques.

\begin{lem*}
Soit $\pi\in \CC[\Ell(G)_{\RR}]$. Supposons que $\Theta_{\pi}$ soit un quasi-caractère de niveau~$0$ sur les elliptiques. Alors $\pi$ appartient à $\CC[\Ell(G)^0_{\RR}]$.
\end{lem*}

\begin{proof}
On peut décomposer la représentation $\pi$ en $\pi=\sum_{i=1,\dots,h}\pi_{i}$ où chaque $\pi_{i}$ est combinaison linéaire de représentations irréductibles dont le caractère central se restreint à $A_{G}(F)$ en un même caractère $\xi_{i}$. On suppose les $\xi_{i}$ tous distincts. Par interpolation, on peut trouver une famille de nombres complexes $(c_{j})_{j=1,\dots,h}$ et une famille $(a_{j})_{j=1,\dots,h}$ d'éléments de $A_{G}(F)$ de sorte que
\[
\Theta_{\pi_{i}}(g)=\sum_{j=1,\dots,h}c_{j}\Theta_{\pi}(a_{j}g)
\]
pour tout $g\in G(F)$. Puisque $\Theta_{\pi}$ est un quasi-caractère de niveau~$0$ sur les elliptiques, cette formule entraîne que $\Theta_{\pi_{i}}$ l'est aussi. Cela nous ramène au cas où $\pi$ elle-même est combinaison linéaire de représentations irréductibles dont le caractère central se restreint à $A_{G}(F)$ en un même caractère que l'on~note $\xi$. On peut aussi tordre le problème par un élément $\chi\in \cA_{G}^*$: $\Theta_{\pi_{\chi}}$ est encore un quasi-caractère de niveau~$0$ sur les elliptiques et, si l'on prouve que $\pi_{\chi}$ appartient à
$\CC[\Ell(G)^0_{\RR}]$, la même assertion s'ensuit pour $\pi$. On peut donc supposer $\xi$ unitaire. D'après la proposition \ref{restriction}, il existe $d'\in \cD_{\cusp}(G)$ tel que $\Theta_{\pi}$ coïncide avec $D^G[d']$ sur $G_{\rell}(F)$. Rappelons que $A_{G}(F)$ agit naturellement sur $\cD_{\cusp}(G)$. On note $(a,\delta)\mto \delta^{a}$ cette action.
L'élément $d'$ n'a pas forcément de caractère central pour l'action de $A_{G}(F)$ mais on peut le remplacer par un élément qui admet $\xi$ pour tel caractère central. En effet, on peut d'abord remplacer $d'$ par l'intégrale $\int_{A_{G}(F)_{\rc}}\xi(a)^{-1}d^{\prime a} \,da$. Cela ne modifie pas la propriété ci-dessus puisque $\pi$ se transforme par $A_{G}(F)$ selon $\xi$. Mais cela assure que $A_{G}(F)_{\rc}$ agit sur $d'$ selon le caractère $\xi$.
Utilisons la décomposition \eqref{produitscalaireelliptique2} de \ref{produitscalaireelliptique} et écrivons $d'=\prod_{\nu\in \cN}d^{\prime\nu}$. Pour tout $\nu\in \cN$, $D^G(d^{\prime\nu})$ est à support dans $w_{G}^{-1}(\nu)$ et coïncide avec~$\Theta_{\pi}$ sur $G_{\rell}(F)\cap w_{G}^{-1}(\nu)$.
Fixons un ensemble de représentants $\underline{\cN}$ des orbites de l'action par addition de $w_{G}(A_{G}(F))$ dans $\cN$. On voit qu'il existe un unique $ d\in \cD_{\cusp,\xi}(G)$ tel que $ d^{\nu}=d^{\prime\nu}$ pour tout $\nu\in \underline{\cN}$. Puisque $\pi$ se transforme par $A_{G}(F)$ selon $\xi$, $\Theta_{\pi}$ coïncide avec $D^G(d)$ sur $G_{\rell}(F)$. D'après le~\eqref{surjectiviteii} du lemme \ref{surjectivite}, il existe $\pi^0\in \CC[\Ell(G)_{\xi}^0]$ tel que $\Delta_{\cusp}(\pi^0)=d$. Alors, d'après l'assertion \eqref{leresultat1} de \ref{leresultat}, $\Theta_{\pi^0}$ coïncide sur $G_{\rell}(F)$ avec $D^G(d)$, donc aussi avec $\Theta_{\pi}$. D'après l'injectivité de l'application \eqref{representationselliptiques2} de \ref{representationselliptiques}, on~a $\pi=\pi^0$, ce qui achève la démonstration.
\end{proof}

\Subsection{Caractérisation des représentations de niveau~\texorpdfstring{$0$}{0}}\label{caracterisation}

\begin{thm*}
Soit $\pi\in \CC[\Irr(G)]$. Alors $\pi\in \CC[\Irr(G)^0]$ si et seulement si $\Theta_{\pi}$ est un quasi-caractère de niveau~$0$ sur $G(F)$.
\end{thm*}

\begin{proof}
Supposons $\pi\in \CC[\Irr(G)^0]$. D'après \ref{leresultat}\,\eqref{leresultat1} et le théorème \ref{unpremiertheoreme}, $\Theta_{\pi}$ coïncide avec un quasi-caractère de niveau~$0$ sur les éléments de $G(F)$ qui sont compacts $\rmod Z(G)$. Soit $\epsilon\in G(F)_{p'}$, supposons que $\epsilon$ n'est pas compact $\rmod Z(G)$. Alors $L[\epsilon]$ est un Levi propre et on~a
\[
L[\epsilon \exp(X)]=L[\epsilon]\quadet Q[\epsilon \exp(X)]=Q[\epsilon]\quad \text{pour tout }X\in \frg_{\epsilon,\tn}(F).
\]
D'après \cite[Th.\,5.2]{C}, on~a $\theta_{\pi}(\epsilon \exp(X))=\delta_{Q[\epsilon]}(\epsilon \exp(X))^{1/2}\theta_{\pi_{Q[\epsilon]}}(\epsilon \exp(X))$ pour tout $X\in \frg_{\epsilon,\tn}(F)$. Un même calcul qu'au \eqref{p'decompositionii} du lemme \ref{p'decomposition} montre que $\delta_{Q[\epsilon]}(\epsilon \exp(X))=\delta_{Q[\epsilon]}(\epsilon)$. La représentation $\pi_{Q[\epsilon]}$ est de niveau~$0$. Puisque $L[\epsilon]\subsetneq G$, on peut raisonner par récurrence sur le rang semi-simple de $G$ et supposer prouvé que $\Theta_{\pi_{Q[\epsilon]}}$ est un quasi-caractère de niveau~$0$ sur $L[\epsilon](F)$. Autrement dit que, pour des constantes $c_{\cO}$ convenables, on a
\[
\theta_{\pi_{Q[\epsilon]}}(\epsilon \exp(X))=\sum_{\cO\in \Nil(\frl[\epsilon]_{\epsilon})}c_{\cO}\hat{j}(\cO,X)
\]
pour tout $X\in \frl[\epsilon]_{\epsilon,\tn}(F)$. Puisque $\frl[\epsilon]_{\epsilon}=\frg_{\epsilon}$, on en déduit un développement analogue de $\theta_{\pi}(\epsilon \exp(X))$. Donc $\Theta_{\pi}$ coïncide avec un quasi-caractère de niveau~$0$ sur $C_{G}(\epsilon)$. Ceci est donc vrai pour tout $\epsilon\in G(F)_{p'}$ et cela signifie que $\Theta_{\pi}$ est un quasi-caractère de niveau~$0$.

Pour $n\in \ZZ$, notons $\CC[\Irr(G)]^n$ la sous-somme de l'expression \ref{representationselliptiques}\,\eqref{representationselliptiques1} où l'on somme sur les Levi $M\in \underline{\cL}^m_{\min}$ pour $m\geq n$. On va prouver
\begin{enumerate}
\item\label{caracterisation1}
soit $\pi\in \CC[\Irr(G)]^n$; supposons que $\Theta_{\pi}$ est un quasi-caractère de niveau~$0$ sur $G(F)$; alors $\pi$ appartient à la somme de $\CC[\Irr(G)]^{n+1}$ et de
\[
\bigoplus_{M\in \underline{\cL}^n_{\min}}\Ind_{M}^G(\CC[\Ell(M)_{\RR}^0])^{W^G(M)}.
\]
\end{enumerate}

Écrivons $\pi\!=\!\sum_{M\in \underline{\cL}_{\min}}\Ind_{M}^G(\pi^M)$, où $\pi^M\!\in\!\CC[\Ell(M)_{\RR}]^{W^G(M)}$ et $\pi^M\!=\!0$ si $a_{M}\!<\!n$. Le caractère $\Theta_{\pi}$ appartient à $Qc_{0}(G)$ par hypothèse donc à $Qc_{0}(G)\cap \Ann^n\rI(G)^*$. D'après le~\eqref{injectiviteii} de la proposition \ref{injectivite} et le théorème \ref{undeuxiemetheoreme}, cet espace a même image dans $\Gr^n\rI(G)^*$ que
\[
\bigoplus_{M\in \underline{\cL}^n_{\min}}\Ind_{M}^G(D^M(\cD_{\cusp}(M)^{W^G(M)})).
\]
Autrement dit, pour tout $M\in \underline{\cL}^n_{\min}$, on peut fixer $d^M\in \cD_{\cusp}(M)^{W^G(M)}$ tel que $\Theta_{\pi^M}$ coïncide avec $D^M(d^M)$ sur les éléments elliptiques de $M(F)$. A fortiori, $\Theta_{\pi^M}$ est un quasi-caractère sur $M(F)$ de niveau~$0$ sur les elliptiques. D'après le lemme~\ref{unlemme}, $\pi^M$ appartient à $\CC[\Ell(M)_{\RR}^0]$ et donc aussi à $\CC[\Ell(M)_{\RR}^0]^{W^G(M)}$. Posons
$\sigma=\sum_{M\in \underline{\cL}^n_{\min}}\Ind_{M}^G(\pi^M)$. Alors $\Theta_{\pi}$ et $\Theta_{\sigma}$ ont même image dans $\Gr^n\rI(G)^*$. Donc $\pi-\sigma$ appartient à $\CC[\Irr(G)]^{n+1}$. Puisque $\sigma$ appartient au dernier espace indiqué dans l'assertion \eqref{caracterisation1}, cela démontre cette assertion.

D'après le sens \og seulement si\fg déjà démontré du théorème, l'assertion \eqref{caracterisation1} implique le sens \og si\fg par récurrence descendante sur $n$.
\end{proof}

\section{Endoscopie}

\subsection{Données endoscopiques}\label{donneesendoscopiques}
On note $W_{F}$ le groupe de Weil de $F$. Il contient le groupe d'inertie $I_{F}$. On note $I_{F}^s$ le groupe d'inertie sauvage de $F$, c'est-à-dire le $p$-Sylow du groupe $I_{F}$.

La théorie de l'endoscopie a été développée par Langlands et ses successeurs. Nous ne la reprendrons pas et nous adopterons sa formulation telle qu'elle figure dans \cite{MW}. Une donnée endoscopique de $G$ est un triplet $\bG'=(G',\cG',s)$, où $G'$ est un groupe réductif connexe défini et quasi-déployé sur $F$, $\cG'$ est un sous-groupe du $L$-groupe $^LG=\hat{G}\rtimes W_{F}$ et $s$ est un élément semi-simple du groupe complexe $\hat{G}$. Le triplet satisfait aux conditions décrites en \cite[I.1.5]{MW}. En particulier, le groupe $\hat{G}'$ s'identifie à la composante neutre $Z_{\hat{G}}(s)^0$ du commutant $Z_{\hat{G}}(s)$. On dit que la donnée est elliptique si $Z(\hat{G})^{\Gamma}$ est un sous-groupe d'indice fini de $Z(\hat{G}')^{\Gamma}$. On définit comme en \cite[I.1.10]{MW} une correspondance entre classes de conjugaison stable d'éléments semi-simples dans $G'(F)$ et $G(F)$ (on parlera aussi bien d'une correspondance entre éléments semi-simples dans $G'(F)$ et $G(F)$). On dit que la donnée est relevante s'il existe un couple d'éléments $(x',x)\in G'_{\reg}(F)\times G_{\reg}(F)$ qui se correspondent. Nous ne considérerons que des données relevantes. Une donnée elliptique est toujours relevante.

L'hypothèse \ref{HypG} a les conséquences suivantes:
\begin{enumerate}
\item\label{donneesendoscopiques1}
il existe une extension de $F$ de degré premier à $p$ telle que le groupe $G'$ soit déployé sur cette extension;

\item\label{donneesendoscopiques2}
$W_{F}\cap \cG'$ est un sous-groupe de $W_{F}$ d'indice fini premier à $p$.
\end{enumerate}

\begin{proof}
Le \eqref{donneesendoscopiques1} résulte de ce que le rang de $G'$ est le même que celui de~$G$ et que l'on~a déjà remarqué que, pour tout tore sur $F$ de dimension inférieure ou égale à ce rang, il existe une extension $F'$ de $F$ de degré premier à $p$ telle que ce tore soit déployé sur $F'$.

Pour \eqref{donneesendoscopiques2}, fixons une telle extension galoisienne $F'$ de $F$ de degré premier à $p$, telle que $G'$ soit déployé sur $F'$. L'action de~$\Gamma_{F'}$ sur $\hat{G}'$ est triviale, a fortiori l'action de~$\Gamma_{F'}$ sur $\hat{G}$ est triviale. Le noyau de l'homomorphisme naturel $W_{F}\to \Gal(F'/F)$ est $W_{F'}$.
Par définition, la projection $\cG'\to W_{F}$ est scindée, c'est-à-dire que l'on peut fixer un homomorphisme continu $i:W_{F}\to \cG'$ qui est une section de cette projection. Pour $w\in W_{F}$, écrivons $i(w)=(j(w),w)$. La restriction de $j$ à $W_{F'}$ est un caractère de ce groupe à valeurs dans $Z_{\hat{G}}(s)$. Notons $W^1$ le sous-groupe des $w\in W_{F'}$ tels que $j(w)\in Z_{\hat{G}}(s)^0$. L'hypothèse \ref{HypG} implique que $Z_{\hat{G}}(s)^0$ est un sous-groupe de $Z_{\hat{G}}(s)$ d'indice premier à $p$. Donc $W^1$ est d'indice fini premier à $p$ dans $W_{F'}$, donc aussi dans $W_{F}$. Pour $w\in W^1$, on~a $(1,w)=j(w)^{-1}i(w)$ et ce terme appartient à $ \cG'$ puisque $j(w)\in Z_{\hat{G}}(s)^0=\hat{G}'\subset \cG'$. Donc $W^1\subset W_{F}\cap \cG'$ et \eqref{donneesendoscopiques2} est démontré.
\end{proof}

On a besoin de munir toute donnée endoscopique $\bG'$ de données auxiliaires $G'_{1},C_{1},\hat{\xi}_{1}$, \cf \cite[I.2.1]{MW}. Le tore $C_{1}(F)$ est naturellement muni d'un caractère $\lambda_{1}:C_{1}(F)\to \CC^{\times}$. Il y a une notion d'unitarité pour de telles données, \cf loc.\,cit.\,I.7.1, qui implique que $\lambda_{1}$ est unitaire.
Nous dirons que les données auxiliaires sont modérément ramifiées si et seulement si elles satisfont aux conditions suivantes:
\begin{enumerate}\setcounter{enumi}{2}
\item\label{donneesendoscopiques3}
$G'_{1}$ et $C_{1}$ sont déployés sur une extension de $F$ de degré premier à $p$;

\item\label{donneesendoscopiques4}
on~a $\hat{\xi}_{1}(1,w)=(1,w)$ pour tout $w\in I_{F}^s$.
\end{enumerate}

\begin{lem*}
Sous notre hypothèse \ref{HypG}, il existe des données auxiliaires modérément ramifiées et unitaires.
\end{lem*}

\begin{proof}
Choisir une suite
\[
1\to C_{1}\to G'_{1}\to G'\to 1
\]
équivaut à choisir une suite duale
\[
1\to \hat{G'}\to \hat{G}'_{1}\to \hat{C}_{1}\to 1
\]
Pour celle-ci, il y a le choix standard suivant. Fixons une paire de Borel $(\hat{B}',\hat{T})$ de~$\hat{G}'$ préservée par l'action galoisienne. On pose $\hat{G}'_{1}=(\hat{G}'\times \hat{T})/\diag(Z(\hat{G}'))$, où $\diag$ est le plongement diagonal. Le centre de $\hat{G}'_{1}$ est isomorphe à $\hat{T}$, donc est connexe. Le tore~$\hat{C}_{1}$ est $\hat{T}/Z(\hat{G})$ et on sait bien que le groupe des caractères de ce tore possède une base conservée par $\Gamma_{F}$. Donc $C_{1}$ est un tore induit. Si $F'$ est une extension finie de~$F$ de degré à $p$ sur laquelle $G'$ est déployé, les actions galoisiennes sur $\hat{G}'_{1}$ et $\hat{C}_{1}$ sont triviales sur $\Gamma_{F'}$. Dualement, $C_{1}$ est déployé sur $F'$ et, puisque $G'$ l'est aussi, $G'_{1}$ l'est également.

Pour construire $\hat{\xi}_{1}$ satisfaisant à \eqref{donneesendoscopiques4}, le mieux est de reprendre la preuve de Langlands \cite[Lem.\,4]{L}. On s'aperçoit qu'il y a un unique argument à préciser. La~clé de la preuve est que, quand $S_{2}\to S_{1}$ est une injection de tores définis sur $F$, tout caractère $\chi_{2}:S_{2}(F)\to \CC^{\times}$ se prolonge en un caractère $\chi_{1}:S_{1}(F)\to \CC^{\times}$ (rappelons que, pour nous, caractère signifie homomorphisme continu). On doit préciser que tout caractère modé\-rément ramifié $\chi_{2}$ se prolonge en un caractère modérément ramifié $\chi_{1}$. Le sous-groupe $S_{1,\tu}(F)$ de $S_{1}(F)$ est compact, donc le produit~$S_{2}(F)S_{1,\tu}(F)$ est fermé dans $S_{1}(F)$. Puisque $\chi_{2}$ est modérément ramifié, il est trivial sur $S_{2,\tu}(F)=S_{2}(F)\cap S_{1,\tu}(F)$. Donc on peut prolonger $\chi_{2}$ en un caractère $\chi':S_{2}(F)S_{1,\tu}(F)\to \CC^{\times}$ qui est trivial sur $S_{1,\tu}(F)$. Puisque $S_{2}(F)S_{1,\tu}(F)$ est fermé dans $S_{1}(F)$, on peut ensuite prolonger $\chi'$ en un caractère $\chi_{1}:S_{1}(F)\to \CC^{\times}$. Il est encore trivial sur $S_{1,\tu}(F)$, c'est-à-dire qu'il est modérément ramifié.

Avec ce complément, la preuve de \cite{L} permet de construire $\hat{\xi}_{1}$ satisfaisant à \eqref{donneesendoscopiques4}. Le même argument qu'en \cite[I.7.1(3)]{MW} permet de le modifier afin d'assurer que les données sont unitaires.
\end{proof}

On fixe un ensemble de représentants $\cE(G)$ des classes d'équivalence de données endoscopiques elliptiques de $G$.
Pour tout $\bG'=(G',\cG',s)\in \cE(G)$, on fixe des données auxiliaires $G'_{1},C_{1},\hat{\xi}_{1}$.
On suppose que ces données auxiliaires sont modérément ramifiées et unitaires. On peut alors fixer et on fixe un facteur de transfert unitaire~$\Delta_{1}$. Signalons que, pour nous, ce facteur ne contient pas le terme $\Delta_{\textit{IV}}$ de \cite{LS}, que nous avons incorporé à la définition des intégrales orbitales en suivant Arthur.

Toute donnée endoscopique $\bG'$ de $G$ apparaît comme une \og donnée de Levi\fg d'une donnée endoscopique elliptique $\bG''$. Celle-ci n'est pas unique mais choisissons-la. Alors les données auxiliaires que l'on~a fixées pour $\bG''$ se restreignent en des données auxiliaires pour $\bG'$. De même, si $M$ est un Levi de $G$, toute donnée endoscopique $\bM'$ de $G$ apparaît comme une \og donnée de Levi\fg d'une donnée endoscopique elliptique~$\bG'$ de~$G$. De nouveau, elle n'est pas unique mais on en choisit une. Alors les données auxiliaires que l'on~a fixées pour $\bG'$ se restreignent en des données auxiliaires pour~$\bM'$.

\subsection{Décomposition de \texorpdfstring{$\CC[\Ell(G)]$}{CLG}}\label{decomposition}

Il y a une décomposition
\[
\rI_{\cusp}(G)=\bigoplus_{\bG'\in\cE(G)}\rI_{\cusp}(G)_{\bG'}
\]
où chaque facteur $\rI_{\cusp}(G)_{\bG'}$ est le sous-espace des $f\in \rI_{\cusp}(G)$ dont les transferts à tout $\bG''\neq\bG'$ sont nuls, \cf \cite[Prop.\,I.4.11]{MW}. Fixons un caractère unitaire $\xi$ de $A_{G}(F)$. Il y a une variante de la décomposition ci-dessus où l'on impose que les fonctions se transforment selon ce caractère $\xi$. Dualement, il y a une décomposition
\begin{equation}\tag{1}\label{decomposition1}
\CC[\Ell(G)_{\xi}]=\bigoplus_{\bG'\in\cE(G)}\CC[\Ell(G)_{\xi}]_{\bG'}.
\end{equation}
Elle est orthogonale pour le produit scalaire elliptique. Pour $\pi\in \CC[\Ell(G)_{\xi}]$, notons $\pi=\sum_{\bG'\in \cE(G)}\pi_{\bG'}$ l'écriture de $\pi$ selon cette décomposition.

\skpt
\begin{prop*}
\begin{enumeratei}
\item\label{decompositioni}
Si $\pi\in \CC[\Ell(G)^0_{\xi}]$, alors $\pi_{\bG'}\in \CC[\Ell(G)^0_{\xi}]$ pour tout $\bG'\in \cE(G)$.

\item\label{decompositionii}
Si $\pi\in \CC[\Ell(G)^{>0}_{\xi}]$, alors $\pi_{\bG'}\in \CC[\Ell(G)^{>0}_{\xi}]$ pour tout $\bG'\in \cE(G)$.
\end{enumeratei}
\end{prop*}

\begin{proof}
Soient $\pi\in \CC[\Ell(G)_{\xi}]$ et $\bG'\in \cE(G)$. Soit $x\in G_{\rell}(F)$. On sait que les classes de conjugaison par $G(F)$ contenues dans la classe de conjugaison stable de~$x$ forment un espace principal homogène sous l'action d'un groupe abélien fini~$K(x)$. Pour $k\in K(x)$, notons $kx$ un représentant de l'image par $k$ de la classe de conjugaison de $x$ (on choisit $1x=x$). Notons $K(x)^{\vee}$ le groupe des caractères de~$K(x)$. Chaque $\bG'\in \cE(G)$ détermine un sous-ensemble $K(x)^{\vee}_{\bG'}$ de $K(x)^{\vee}$ et $K(x)^{\vee}$ est union disjointe des $K(x)^{\vee}_{\bG'}$ quand $\bG'$ décrit $\cE(G)$. Pour tout $\bG'\in \cE(G)$, on~a alors l'égalité
\begin{equation}\tag{2}\label{decomposition2}
\theta_{\pi_{\bG'}}(x)=\vert K(x)\vert ^{-1}\sum_{\kappa\in K(x)^{\vee}_{\bG'}}\sum_{k\in K(x)}\kappa(k)\theta_{\pi}(kx).
\end{equation}
Soient $\epsilon\in G(F)$ un élément $p'$-compact $\rmod Z(G)$ et $X\in \frg_{\epsilon,\tn}(F)$. Posons $x=\epsilon \exp(X)$ et supposons $x\in G_{\rell}(F)$. Pour tout $k\in K(x)$, fixons $g_{k}\in G(\bar{F})$ tel que $kx=g_{k}^{-1}xg_{k}$. Posons $\epsilon_{k}=g_{k}^{-1}\epsilon g_{k}$ et $X_{k}=g_{k}^{-1} X g_{k}$. Alors $x_{k}=\epsilon_{k}\exp(X_{k})$ est une $p'$-décomposition de $x_{k}$. Soit $\lambda\in \fro_{F}\setminus\{0\}$, posons $y=\epsilon \exp(\lambda^2 X)$. On a $K(x)=K(y)$ et on voit comme dans la preuve du lemme \ref{induction} que l'on peut choisir $y_{k}=\epsilon_{k}\exp(\lambda^2 X_{k})$ pour tout $k\in K(y)$. L'égalité \eqref{decomposition2} pour $y$ devient
\[
\theta_{\pi_{\bG'}}(\epsilon \exp(\lambda^2 X))=\vert K(x)\vert ^{-1}\sum_{\kappa\in K(x)^{\vee}_{\bG'}}\sum_{k\in K(x)}\kappa(k)\theta_{\pi}(\epsilon_{k}\exp(\lambda^2 X_{k})).
\]
Supposons $\pi\in \CC[\Ell(G)^0_{\xi}]$. Comme fonctions de $\lambda$, les termes du membre de droite appartiennent à $E$. Donc aussi la fonction $\lambda\mto \theta_{\pi_{\bG'}}(\epsilon \exp(\lambda^2 X))$. C'est-à-dire que~$\theta_{\pi_{\bG'}}$ est un quasi-caractère de niveau~$0$ sur les elliptiques. D'après le lemme \ref{unlemme}, $\pi_{\bG'}$ appartient à $\CC[\Ell(G)_{\xi}^0]$. Cela démontre le~\eqref{decompositioni} de l'énoncé.

Supposons $\pi\in \CC[\Ell(G)^{>0}_{\xi}]$, soit $\sigma\in \CC[\Ell(G)^0_{\xi}]$. Parce que la décomposition \eqref{decomposition1} est orthogonale, on a
\[
\<\pi_{\bG'},\sigma\>_{\rell}=\<\pi_{\bG'},\sigma_{\bG'}\>_{\rell}=\<\pi,\sigma_{\bG'}\>_{\rell}.
\]
On vient de prouver que $\sigma_{\bG'}$ était de niveau~$0$. Donc le produit de droite est nul. Il~en est de même de celui de gauche. Cela étant vrai pour tout $\sigma\in \CC[\Ell(G)^0_{\xi}]$, cela entraîne $\pi_{\bG'}\in \CC[\Ell(G)^{>0}_{\xi}]$.
\end{proof}

\subsection{Le cas quasi-déployé}\label{lecasquasideploye}
Supposons $G$ quasi-déployé. On note $\SI(G)$ le quotient de $C_{\rc}^{\infty}(G(F))$ par le sous-espace des fonctions dont toutes les intégrales orbitales stables sont nulles. En identifiant les intégrales orbitales à des formes linéaires sur $\rI(G)$, c'est aussi le quotient de $\rI(G)$ par le sous-espace des éléments dont toutes les intégrales orbitales stables sont nulles. Notons
$\SI(G)^*$ l'espace des distributions stables sur $G(F)$, c'est-à-dire le dual de $\SI(G)$. On note $\CC[\Irr(G)]^{\st}$ le sous-espace des $\pi\in \CC[\Irr(G)]$ telles que $\Theta_{\pi}\in \SI(G)^*$. On définit de même des espaces $\CC[\Ell(G)]^{\st}$ etc.

On définit une projection linéaire $p^{\st}:\CC[\Irr(G)]\to \CC[\Irr(G)]^{\st}$ de la façon suivante.
Il y a une donnée endoscopique elliptique maximale $\bG=(G,{^LG},1)$ que l'on suppose appartenir à $\cE(G)$. Dans la décomposition \eqref{decomposition1} de \ref{decomposition}, on~a l'égalité $\CC[\Ell(G)_{\xi}]^{\st} =\CC[\Ell(G)_{\xi}]_{\bG}$. Alors $p^{\st}:\CC[\Ell(G)_{\xi}]\to \CC[\Ell(G)_{\xi}]^{\st}$ est la projection associée à cette décomposition \eqref{decomposition1} de \ref{decomposition}, c'est-à-dire qu'elle annule les composantes $\CC[\Ell(G)_{\xi}]_{\bG'}$ pour $\bG'\neq\bG$.
Plus généralement, le sous-espace $\CC[\Ell(G)_{\RR}]^{\st}$ est engendré par les $\pi_{\chi}$, où $\pi$ parcourt $\CC[\Ell(G)_{\xi}]^{\st}$, $\xi$ parcourt le groupe des caractères unitaires de $A_{G}(F)$ et $\chi$ parcourt $\cA_{G}^*$. On définit $p^{\st}:\CC[\Ell(G)_{\RR}]\to \CC[\Ell(G)_{\RR}]^{\st}$ de sorte que, si $\pi\in \CC[\Ell(G)_{\xi}]$ et $\chi\in \cA_{G}^*$, on ait $p^{\st}(\pi_{\chi})=p^{\st}(\pi)_{\chi}$.
Le corollaire XI.3.1 de \cite{MW} équivaut à l'égalité suivante, similaire à \ref{representationselliptiques}\,\eqref{representationselliptiques1}:\vspace*{-3pt}\enlargethispage{.6\baselineskip}%
\begin{equation}\tag{1}\label{lecasquasideploye1}
\CC[\Irr(G)]^{\st}=\bigoplus_{M\in \underline{\cL}_{\min}}\Ind_{M}^G(\CC[\Ell(M)_{\RR}]^{\st,W^G(M)}).
\end{equation}
Des projections que l'on vient de définir pour chaque Levi se déduit alors la projection~$p^{\st}$ cherchée.

\skpt
\begin{cor*}
\begin{enumeratei}
\item\label{lecasquasideployei}
Pour $\pi\in \CC[\Irr(G)^0]$, on~a $p^{\st}(\pi)\in \CC[\Irr(G)^0]^{\st}$. Pour $\pi\in \CC[\Irr(G)^{>0}]$, on~a $p^{\st}(\pi)\in \CC[\Irr(G)^{>0}]^{\st}$.

\item\label{lecasquasideployeii}
On a l'égalité $p^0\circ p^{\st}=p^{\st}\circ p^0$.
\end{enumeratei}
\end{cor*}

\begin{proof}
Le \eqref{lecasquasideployei} résulte immédiatement de la proposition \ref{decomposition} et le~\eqref{lecasquasideployei} est équivalent au \eqref{lecasquasideployei}.
\end{proof}

\subsection{L'hypothèse \texorpdfstring{\ref{HypendoG}}{Hyp}}\label{Hypendo}

Il n'est pas clair que l'hypothèse \ref{HypG} que nous avons posée entraîne $(\Hyp)(G')$ pour tout groupe endoscopique $G'$ de $G$, et encore moins qu'elle entraîne $(\Hyp)(G'_{1})$ pour un groupe auxiliaire $G'_{1}$. Considérons l'hypothèse\vspace*{-6pt}
\begin{equation}\tag*{$\textup{(Hyp)}_\rendo(G)$}\label{HypendoG} \begin{cases}
\text{l'hypothèse \ref{HypG} est satisfaite;}\\
\text{pour tout $\bG'\in \cE(G)$, les hypothèses $(\Hyp)(G')$ et $(\Hyp)(G'_{1})$}\\[-3pt]
\hspace*{7cm}\text{sont satisfaites.}
\end{cases}
\end{equation}

Puisque l'hypothèse \ref{HypG} implique $(\Hyp)(M)$ pour tout Levi $M$ de $G$ et d'après ce que l'on~a dit à la fin de la section \ref{donneesendoscopiques}, l'hypothèse \ref{HypendoG} a les conséquences suivantes:
\begin{itemize}
\item
pour toute donnée endoscopique $\bG'$ de $G$ (pas forcément elliptique), les hypothèses $(\Hyp)(G')$ et $(\Hyp)(G'_{1})$ sont satisfaites;

\item
pour tout Levi $M$ de $G$, l'hypothèse $(\Hyp)_{\rendo}(M)$ est satisfaite.

\smallskip
\noindent
Dans la suite de la section, nous imposons l'hypothèse renforcée \ref{HypendoG}.
\end{itemize}


\subsection{Représentations de niveau~\texorpdfstring{$0$}{0} et données auxiliaires}\label{representationsetdonnees}

Soit $\bG'=(G',\cG',s)$ une donnée endoscopique relevante de $G$. On s'intéresse exclusivement aux représentations irréductibles de $G'_{1}(F)$ dont le caractère central coïncide avec $\lambda_{1}$ sur $C_{1}(F)$. On note $\Irr(G'_{1})_{\lambda_{1}}$ l'ensemble de ces représentations, avec les variantes $\Temp(G'_{1})_{\lambda_{1}}$, $\Irr(G'_{1})^0_{\lambda_{1}}$ etc. Ces ensembles dépendent du choix des données auxiliaires. Pour que les résultats qui suivent aient vraiment un sens, il vaut mieux qu'ils n'en dépendent pas trop. Selon le formalisme que l'on~a développé en \cite{MW}, pour deux choix de données auxiliaires, il y a une bijection canonique entre les ensembles associés à chacune des données. Expliquons cela.
Choisissons d'autres données auxiliaires modérément ramifiées $(G'_{2},C_{2},\hat{\xi}_{2})$. Notons $G'_{12}$ le produit fibré de $G'_{1}$ et $G'_{2}$ au-dessus de $G'$. On a deux suites exactes
\[
1\to C_{2}\to G'_{12}\to G'_{1}\to 1,\,\,1\to C_{1}\to G'_{12}\to G'_{2}\to 1.
\]
On a introduit en \cite[I.2.5]{MW} un homomorphisme $\lambda_{12}:G'_{12}(F)\to \CC^{\times}$. Sa restriction à $C_{1}(F)\times C_{2}(F)$ est le produit des caractères $\lambda_{1}$ et $\lambda_{2}^{-1}$. Soit $\pi_{1}\in \Irr(G'_{1})_{\lambda_{1}}$, réalisée dans un espace complexe $V$. Par la première suite ci-dessus, $\pi_{1}$ se relève en une représentation $\pi_{12}$ de $G'_{12}(F)$. La restriction à $C_{1}(F)$, \resp $C_{2}(F)$, de son carac\-tère central est $\lambda_{1}$, \resp $1$. Posons $\pi_{21}=\lambda_{12}^{-1}\otimes \pi_{12}$. Alors la restriction à $C_{1}(F)$, \resp $C_{2}(F)$, du caractère central de $\pi_{21}$ est $1$, \resp $\lambda_{2}$. La représentation $\pi_{21}$ se descend par la deuxième suite ci-dessus en une représentation $\pi_{2}$ de $G'_{2}(F)$ dans $V$ dont le caractère central coïncide avec $\lambda_{2}$ sur $C_{2}(F)$. L'application $\pi_{1}\mto \pi_{2}$ est une bijection de $\Irr(G'_{1})_{\lambda_{1}}$ sur $\Irr(G'_{2})_{\lambda_{2}}$. Parce que $\lambda_{1}$ et $\lambda_{2}$ sont unitaires, le caractère $\lambda_{12}$ l'est aussi et la bijection envoie $\Temp(G'_{1})_{\lambda_{1}}$ sur $\Temp(G'_{2})_{\lambda_{2}}$. Parce que les données auxiliaires sont modérément ramifiées, le caractère $\lambda_{12}$ est \og modérément ramifié\fg au sens que sa restriction à l'ensemble $G'_{12,\tu}(F)$ est triviale. Remarquons de plus que $\Imm(G'_{1,\AD})=\Imm(G'_{\AD})=\Imm(G'_{2,\AD})$. Une facette $\cF\in \Fac(G')$ détermine un sous-groupe $K_{\cF}^+$ de $G'(F)$ et des sous-groupes similaires dans $G'_{1}(F)$ et $G'_{2}(F)$, que l'on~note $K_{1,\cF}^+$ et $K_{2,\cF}^+$. On voit que si $\pi_{1}$ admet des invariants non nuls par le sous-groupe $K_{1,\cF}^{+}$, alors $\pi_{2}$ admet des invariants non nuls par le sous-groupe $K_{2,\cF}^{+}$. La bijection envoie donc $\Irr(G'_{1})^0_{\lambda_{1}}$ sur $\Irr(G'_{2})^0_{\lambda_{2}}$.

\subsection{Correspondance entre éléments semi-simples}\label{correspondance}
Soit $\bG'=(G',\cG',s)$ une donnée endoscopique de $G$. Soient $x'\in G'_{\reg}(F)$ et $x\in G_{\reg}(F)$ deux éléments qui se correspondent. Notons $T'$ et $T$ les commutants de $x'$ dans $G'$ et de $x$ dans $G$. Ce sont des tores et il y a un unique isomorphisme $\xi_{T,T'}:T\to T'$ défini sur $F$ de sorte que $\xi_{T,T'}(x)=x'$, \cf \cite[I.1.10]{MW}. D'où un isomorphisme encore noté $\xi_{T,T'}:\frt\to \frt'$.

\skpt
\begin{cor*}
\begin{enumeratei}
\item\label{correspondancei}
Si $x=\epsilon \exp(X)$ est une $p'$-décomposition de $x$, alors $x'=\xi_{T,T'}(\epsilon)\exp(\xi_{T,T'}(X))$ est une $p'$-décomposition de $x'$.

\item\label{correspondanceii}
Soit $x=\epsilon \exp(X)$ une $p'$-décomposition de $x$. Pour tout $\lambda\in F^{\times}$ tel que $ \lambda X\in \frg_{\epsilon,\tn}(F)$, $\xi_{T,T'}(\epsilon)\exp(\xi_{T,T'}(\lambda X))$ est un élément de $G'_{\reg}(F)$ qui correspond à $\epsilon \exp(\lambda X)$.

\item\label{correspondanceiii}
Supposons $\bG'$ elliptique. Alors $x$ est compact $\rmod Z(G)$ si et seulement si $x'$ est compact $\rmod Z(G')$.
\end{enumeratei}
\end{cor*}

\begin{proof}
Puisque $\xi_{T,T'}$ est un isomorphisme, l'élément $\xi_{T,T'}(X)$ est topologiquement nilpotent. Posons $\epsilon'=\xi_{T,T'}(\epsilon)$, introduisons le Levi $L[\epsilon]$ de $G$ et le Levi similaire $L'[\epsilon']$ de $G'$. Notons $\Sigma$ l'ensemble de racines de $T$ dans $G$, $\Sigma'$ celui des racines de $T'$ dans $G'$ et $\Sigma^{L[\epsilon]}$ et $\Sigma^{L'[\epsilon']}$ les sous-ensembles des racines dans $L[\epsilon]$, \resp $L'[\epsilon']$. L'isomorphisme $\xi_{T,T'}$ transporte $\Sigma$ en un sous-ensemble de $X^*(T')$ et on sait que $\Sigma'\subset \xi_{T,T'}(\Sigma)$. Il résulte alors des définitions que $\Sigma^{L'[\epsilon']}=\Sigma'\cap \xi_{T,T'}(\Sigma^{L[\epsilon]})$. D'où $\xi_{T,T'}(Z(L[\epsilon]))\subset Z(L'[\epsilon'])$. Il existe un entier $c\geq1$ premier à $p$ tel que $\epsilon^c\in Z(L[\epsilon])$. Cela entraîne $(\epsilon')^c\in Z(L'[\epsilon'])$. Donc $\epsilon'$ est un $p'$-élément dans $G'(F)$. Cela démontre le~\eqref{correspondancei}.

Le \eqref{correspondanceii} est immédiat. Supposons $\bG'$ elliptique. Alors $Z(G)(F)$ est un sous-groupe cocompact de $Z(G')(F)$ et \eqref{correspondanceiii} en résulte.
\end{proof}

\subsection{Facteur de transfert}\label{facteurdetransfert}
Soit $\bG'=(G',\cG',s)$ une donnée endoscopique relevante de $G$. Rappelons que le facteur de transfert est une fonction définie sur l'ensemble $\bD_{1}$ des couples $(x'_{1},x)$ tels que: $x'_{1}\in G'_{1,\reg}(F)$, $x\in G_{\reg}(F)$ et, en notant~$x'$ l'image de $x'_{1}$ dans $G'(F)$, les classes de conjugaison stable de $x'$ et de $x$ se correspondent. Considérons un tel couple et fixons une
$p'$-décomposition $x=\epsilon \exp(X)$. Par le~\eqref{correspondancei} du
lemme \ref{correspondance}, on en déduit une $p'$-décomposition $x'=\epsilon' \exp(X')$. L'application $\frg'_{1,\tn}(F)\to \frg'_{\tn}(F)$ est surjective. Fixons $X'_{1}\in \frg'_{1,\tn}(F)$ au-dessus de $X'$. Définissons $\epsilon'_{1}\in G'_{1}(F)$ par l'égalité $x'_{1}=\epsilon'_{1}\exp(X'_{1})$. On voit que $\epsilon'_{1}$ est un $p'$-élément, donc que l'égalité $x'_{1}=\epsilon'_{1}\exp(X'_{1})$ est une $p'$-décomposition. Pour $\lambda\in \fro_{F}\setminus\{0\}$, le couple $(\epsilon'_{1}\exp(\lambda X'_{1}),\epsilon \exp(\lambda X))$ appartient encore à l'ensemble de définition $\bD_{1}$ du facteur de transfert.

\begin{lem*}
Pour tout $\lambda\in \fro_{F}\setminus\{0\}$, on~a l'égalité
\[
\Delta_{1}(\epsilon'_{1}\exp(\lambda^2 X'_{1}),\epsilon \exp(\lambda^2 X))=\Delta_{1}(\epsilon'_{1}\exp( X'_{1}),\epsilon \exp( X)).
\]
\end{lem*}

Preuve. Notons $T$, $T'$ et $T'_{1}$ les commutants de $x$, $x'$ et $x'_{1}$ dans $G$, $G'$ et $G'_{1}$. Pour calculer le facteur de transfert, on peut choisir des $\chi$-data modérément ramifiées. Les termes $\Delta_{I}$ et $\Delta_{\textit{III},1}$ de \cite{LS} ne dépendent que des tores $T$ et $T'_{1}$. Le terme $\Delta_{\textit{III},2}$ est la valeur en $x'_{1}$ d'un caractère de $T'_{1}(F)$. Celui-ci est construit à l'aide des $\chi$\nobreakdash-data et est modérément ramifié. Donc $\Delta_{\textit{III},2}(\epsilon'_{1}\exp( X'_{1}),\epsilon \exp( X))$ ne dépend pas du couple $(X'_{1},X)$. Il reste le terme $\Delta_{\textit{II}}$. Un calcul similaire à celui de \cite[Lem.\, 10.3]{W4} montre que, comme fonction de $(X'_{1},X)$, $\Delta_{\textit{II}}(\epsilon'_{1}\exp( X'_{1}),\epsilon \exp( X))$ est proportionnel au terme analogue $\Delta_{\textit{II}}(X',X)$ défini sur les algèbres de Lie. D'après \cite[Lem.\,3.2.1]{F}, la fonction $\lambda\mto \Delta_{\textit{II}}(\lambda X',\lambda X)\Delta_{\textit{II}}(X',X)^{-1}$ définie sur $F^{\times}$ est un caractère quadratique. Elle est donc constante sur les carrés. Le lemme en résulte. $\square$

\subsection{Actions des centres}\label{actiondescentres}
Soit $\bG'=(G',\cG',s)\in \cE(G)$. Puisque cette donnée est elliptique, le groupe $Z(G)(F)$ s'identifie à un sous-groupe cocompact de $Z(G')(F)$. Le sous-groupe $A_{G}(F)\subset Z(G)(F)$ est égal au sous-groupe $A_{G'}(F)\subset Z(G')(F)$. Le~groupe $A_{G'_{1}}(F)$ s'envoie surjectivement sur $A_{G'}(F)$ (parce que $C_{1}$ est un tore induit). Notons $\underline{Z}'_{1}$ le groupe des $z'_{1}\in Z(G'_{1})$ dont l'image dans $Z(G')$ appartient à $Z(G)$. Ce groupe contient $C_{1}$ et $A_{G'_{1}}$ et $ \underline{Z}'_{1}(F)$ est un sous-groupe cocompact de $Z(G'_{1})(F)$. Pour tout couple $(x'_{1},x)\in \bD_{1}$ et pour tout $z'_{1}\in \underline{Z}'_{1}(F)$, le couple $(z'_{1}x'_{1},zx)$ appartient aussi à $\bD_{1}$, où $z$ est l'image de $z'_{1}$ dans $Z(G)(F)$. D'après \cite[Lem.\,4.4.A]{LS}, il existe un caractère $\zeta_{1}$ de $\underline{Z}'_{1}(F)$ de sorte que, pour $(x'_{1},x)\in \bD_{1}$, on ait l'égalité
\[
\Delta_{1}(z'_{1}x'_{1},zx)=\zeta(z'_{1})\Delta_{1}(x'_{1},x).
\]
Ce caractère coïncide avec $\lambda_{1}^{-1}$ sur $C_{1}(F)$.
Grâce à l'hypothèse \ref{HypendoG}, on a
\begin{enumerate}
\item\label{actiondescentres1}
le caractère $\zeta_{1}$ est trivial sur le sous-ensemble $\underline{Z}'_{1,\tu}(F)$ de $\underline{Z}'_{1}$.
\end{enumerate}

La preuve est similaire à celle du lemme précédent.

\subsection{Un premier résultat de transfert}\label{unpremierresultat}
Soit $\bG'=(G',\cG',s)\in \cE(G)$. On note $\SI_{\lambda_{1}}(G'_{1})$ la variante de $\SI(G'_{1})$ où l'on impose que les fonctions se transforment par translations par $C_{1}(F)$ selon le caractère $\lambda_{1}^{-1}$. On définit de façon similaire $\SI_{\lambda_{1}}(G'_{1})^*$ (c'est le dual de $\SI_{\lambda_{1}}(G'_{1})$). On note $\Aut(\bG')$ le groupe d'automorphismes de $\bG'$, \cf \cite[I.1.5]{MW}. Ce groupe agit naturellement sur $\SI_{\lambda_{1}}(G'_{1})$ et $\SI_{\lambda_{1}}(G'_{1})^*$, \cf \cite[I.2.6]{MW}.

\begin{rem*}
La définition adoptée ici et dans \cite{MW} n'est pas celle que l'on trouve dans d'autres références. Elle diffère de la définition adoptée notamment par Arthur par torsion par un caractère de sorte que le facteur $\Delta_{1}$ soit invariant par cette action. Cela explique la disparition dans nos formules des caractères que l'on trouve dans celles d'Arthur.
\end{rem*}

On sait définir un transfert $\rI(G)\to \SI_{\lambda_{1}}(G'_{1})$, dont l'image est contenue dans le sous-espace des invariants par $\Aut(\bG')$. D'où une application transposée $ \SI_{\lambda_{1}}(G'_{1})^*\to \rI(G)^*$ entre les espaces duaux. Arthur démontre dans \cite{A3} qu'il existe
une unique application linéaire dite de transfert spectral
\[
\transfert:\CC[\Irr(G'_{1})_{\lambda_{1}}]^{\st}\to\CC[\Irr(G)]
\]
qui rend le diagramme suivant commutatif
\[
\xymatrix@C=1.5cm{
\CC[\Irr(G'_{1})_{\lambda_{1}}]^{\st}\ar[r]^-{\transfert}\ar[d]_-\Theta&\CC[\Irr(G)]\ar[d]^-\Theta\\
\SI_{\lambda_{1}}(G'_{1})^*\ar[r]& \rI(G)^*
}
\]
Elle est invariante par l'action de $\Aut(\bG')$, c'est-à-dire qu'elle se factorise en
\[
\transfert:\CC[\Irr(G'_{1})_{\lambda_{1}}]^{\st}\to\CC[\Irr(G'_{1})_{\lambda_{1}}]^{\st,\Aut(\bG')}\to\CC[\Irr(G)],
\]
où la première application est la projection naturelle.

\begin{prop*}
Cette application $\transfert$ envoie $\CC[\Irr(G'_{1})^0_{\lambda_{1}}]^{\st}$ dans $\CC[\Irr(G)^0]$.
\end{prop*}

\begin{proof}
Soit $\pi'\in \CC[\Irr(G'_{1})^0_{\lambda_{1}}]^{\st}$, posons $\pi=\transfert(\pi')$. Soit $\epsilon\in G(F)_{p'}$ et soit $X\in \frg_{\epsilon,\tn}(F)$. Posons $x=\epsilon \exp(X)$ et supposons $x\in G_{\reg}(F)$. Fixons un ensemble de représentants $X^{\bG'}(x)$ des classes de conjugaison stable dans $G'(F)$ qui correspondent à celle de $x$ dans $G(F)$. Pour tout $x'\in X^{\bG'}(x)$, on fixe un relèvement~$x'_{1}$ de~$x'$ dans $G'_{1}(F)$. Par définition, on~a l'égalité
\begin{equation}\tag{1}\label{unpremierresultat1}
\theta_{\pi}(x)D^G(x)^{1/2}=\sum_{x'\in X^{\bG'}(x)}\Delta_{1}(x'_{1},x)\theta_{\pi'}(x'_{1})D^{G'}(x')^{1/2}.
\end{equation}
On note $T$ le commutant de $x$ dans $G$ et, pour tout $x'\in X^{\bG'}(x)$, on~note $T_{x'}$ celui de $x'$ dans $G'$. On a des isomorphismes $\xi_{T,T_{x'}}:T\to T_{x'}$, \cf \ref{correspondance}. On pose
\[
\epsilon_{x'}=\xi_{T,T_{x'}}(\epsilon),\quad X_{x'}=\xi_{T,T_{x'}}(X).
\]
Alors $x'\!=\!\epsilon_{x'}\exp(X_{x'})$ est une $p'$-décomposition. On fixe un relèvement $X_{1,x'}$ de~$X_{x'}$ dans $\frg'_{1,\tn}\mkern-2.5mu(F)$ et un relèvement $\epsilon_{1,x'}$ de $\epsilon_{x'}$ dans $G'_{1}\mkern-2.5mu(F)$ de sorte que~\hbox{$x'_{1}\!\!=\!\epsilon_{1,x'}\exp(X_{1,x'})$}.
Soit $\lambda\in F^{\times}$ tel que $\lambda X\in \frg_{\epsilon,\tn}(F)$, posons $y=\epsilon \exp(\lambda X)$. Alors, pour tout $x'\in X^{\bG'}(x)$, l'élément $\epsilon_{x'}\exp(\lambda X_{x'})$ correspond à $ y$. Montrons que
\begin{enumerate}\setcounter{enumi}{2}
\item\label{unpremierresultat2}
si $x',x''\in X^{\bG'}(x)$ et $x'\neq x''$, alors $\epsilon_{x'}\exp(\lambda X_{x'})$ n'est pas stablement conjugué à $\epsilon_{x''}\exp(\lambda X_{x''})$.
\end{enumerate}

Supposons qu'ils sont stablement conjugués. Il existe alors $h\in G'(\bar{F})$ tel que
\[
h^{-1}\epsilon_{x'}\exp(\lambda X_{x'})h=\epsilon_{x''}\exp(\lambda X_{x''}).
\]
Puisque les éléments en question sont fortement réguliers, l'opérateur de conjugaison par $h$ se restreint en un isomorphisme $\Int_{h}:T_{x''}\to T_{x'}$. Le composé
\[
\Xi=\xi_{T,T_{x'}}^{-1}\circ \Int_{h}\circ\xi_{T,T_{x''}}
\]
est un automorphisme de $T$ fixant $y$. De plus, par construction des isomorphismes $\xi_{T,T_{x'}}$ et $\xi_{T,T_{x''}}$, $\Xi$ est forcément de la forme $\Int_{h'}$ pour un élément $h'\in \Norm_{G}(T)(\bar{F})$. Puisqu'il fixe $y$ et que $y$ est fortement régulier, c'est l'identité. Mais alors
\[
\Int_{h}(\epsilon_{x''})=\epsilon_{x'},\quad \Int_{h}(X_{x''})=X_{x'},
\]
donc $h^{-1}x'h=x''$, ce qui contredit la définition de $X^{\bG'}(x)$. Cela démontre \eqref{unpremierresultat2}.

Cette assertion entraîne que $X^{\bG'}(y)$ a au moins autant d'éléments que $X^{\bG'}(x)$. La situation étant symétrique en $x$ et $y$, ces ensembles ont même nombre d'éléments. Alors \eqref{unpremierresultat2} entraîne que l'on peut choisir pour $X^{\bG'}(y)$ l'ensemble des $\epsilon_{x'}\exp(\lambda X_{x'})$ pour $x'\in X^{\bG'}(x)$. On applique cela en remplaçant $\lambda$ par $\lambda^2$ pour $\lambda\in \fro_{F}\setminus\{0\}$. L'égalité \eqref{unpremierresultat1} pour $y=\epsilon \exp(\lambda^2 X)$ devient
\begin{multline*}\tag{3}\label{unpremierresultat3}
\theta_{\pi}(\epsilon \exp(\lambda^2X))D^G(\epsilon \exp(\lambda^2 X))^{1/2}\\[-3pt]
=\sum_{x'\in X^{\bG'}(x)}\Delta_{1}(\epsilon_{1,x'}\exp(\lambda^2X_{1,x'}),\epsilon \exp(\lambda^2X))
\\[-10pt]
\cdot\theta_{\pi'}(\epsilon_{1,x'}\exp(\lambda^2X_{1,x'}))D^{G'}(\epsilon_{x'}\exp(\lambda^2X_{x'}))^{1/2}.
\end{multline*}
On considère ces termes comme des fonctions de $\lambda$. Comme on l'a dit plusieurs fois, les discriminants de Weyl contribuent par des puissances de $\vert \lambda\vert _{F}$. D'après le lemme~\ref{facteurdetransfert}, le facteur de transfert est constant. Puisque $\pi'$ est de niveau~$0$, le théorème \ref{caracterisation} dit que $\Theta_{\pi'}$ est un quasi-caractère de niveau~$0$. Donc les termes $\theta_{\pi'}(\epsilon_{1,x'}\exp(\lambda^2X_{1,x'}))$ appartiennent à $E$. Donc $\lambda\mto \theta_{\pi}(\epsilon \exp(\lambda^2X))$ appartient à $E$, c'est-à-dire que $\Theta_{\pi}$ est un quasi-caractère de niveau~$0$. D'après le même théorème \ref{caracterisation}, $\pi$ appartient à $\CC[\Irr(G)^0]$.
\end{proof}

\subsection{Une réciproque partielle}\label{unereciproque}
Soit $\bG'=(G',\cG',s)\in \cE(G)$.
L'application $\transfert$ se restreint en une application linéaire
\[
\transfert_{\rell}:\CC[\Ell(G'_{1})_{\lambda_{1}}]^{\st}\to \CC[\Ell(G)].
\]
\begin{prop*}
Soit $\pi\in \CC[\Ell(G)^0]$. Supposons que $\pi$ appartienne à l'image de l'application $\transfert_{\rell}$. Alors il existe $\pi'\in \CC[\Ell(G'_{1})^0_{\lambda_{1}}]^{\st}$ telle que $\transfert_{\rell}(\pi')=\pi$.
\end{prop*}

\begin{proof}[\proofname. Première étape]
Fixons $\sigma'\in\CC[\Ell(G'_{1})_{\lambda_{1}}]^{\st}$ tel que \[
\transfert_{\rell}(\sigma')\!=\!\pi.
\]
Décomposons $\sigma'$ en somme $\sum_{i=1,\dots,h}\sigma'_{i}$, où $\sigma'_{i}$ est une combinaison linéaire de représentations elliptiques de $G'_{1}(F)$ dont le caractère central se restreint en un même carac\-tère $\xi'_{i}$ de $A_{G'_{1}}(F)$. On suppose les $\xi'_{i}$ distincts. On a introduit un caractère $\zeta_{1}$ en \ref{actiondescentres}. Le caractère $(\zeta_{1})_{ \vert A_{G'_{1}}(F)}\xi'_{i}$ de $A_{G'_{1}}(F)$ se factorise par la projection $A_{G'_{1}}(F)\to A_{G}(F)$ en un caractère $\xi_{i}$ de ce dernier groupe et les caractères $\xi_{i}$ sont tous distincts. On a l'égalité $\pi=\sum_{i=1,\dots,h}\transfert_{\rell}(\sigma'_{i})$ et, puisque $\pi$ est de niveau~$0$, on en déduit par interpolation que toutes les composantes $\transfert_{\rell}(\sigma'_{i})$ sont de niveau~$0$. On peut remplacer $\pi$ par chacune de ces composantes et cela ramène le problème au cas où $A_{G'_{1}}(F)$ agit sur $\sigma'$ selon un caractère $\xi'$. Ensuite, on peut encore tordre $\sigma'$ par un élément de $\cA_{G'_{1}}^*$ et supposer que $\xi'$ est unitaire.

\subsubsection*{Deuxième étape: construction d'un quasi-caractère de niveau~$0$}

Pour tout élément $\epsilon_{1}\in G'_{1}(F)$ qui est $p'$-compact $\rmod \underline{Z}'_{1}$, fixons un voisinage $\frV_{\epsilon_{1}}$ de $0$ dans $\frg'_{1,\epsilon_{1}}(F)$ et une famille $(c_{\epsilon_{1},\cO})_{\cO\in \Nil(\frg'_{1,\epsilon_{1}})}$ de nombres complexes de sorte que, pour tout $X_{1}\in \frV_{\epsilon_{1}}$ tel que $\epsilon_{1}\exp(X_{1})\in G'_{1,\reg}(F)$, on ait l'égalité
\begin{equation}\tag{1}\label{unereciproque1}
\theta_{\sigma'}(\epsilon_{1}\exp(X_{1}))=\sum_{\cO\in \Nil(\frg'_{1,\epsilon_{1}})}c_{\epsilon_{1},\cO}\hat{j}(\cO,X_{1}).
\end{equation}
Le voisinage $\frV_{\epsilon_{1}}$ n'est pas uniquement déterminé mais les nombres complexes $c_{\epsilon_{1},\cO}$ le sont. Soit $z_{1}\in \underline{Z}'_{1,\tu}(F)$. On peut remplacer $\epsilon_{1}$ par $\epsilon_{1}z_{1}$. Cela ne change pas l'ensemble $\Nil(\frg'_{1,\epsilon_{1}})$. Montrons que:
\begin{enumerate}\setcounter{enumi}{1}
\item\label{unereciproque2}
pour tout $\cO\in \Nil(\frg'_{1,\epsilon_{1}})$, la fonction $z_{1}\mto c_{\epsilon_{1}z_{1},\cO}$ est localement constante;

\item\label{unereciproque3}
on peut choisir les voisinages $\frV_{\epsilon_{1}}$ de sorte que $\frV_{\epsilon_{1}z_{1}}\!=\!\frV_{\epsilon_{1}}$ pour tout $z_{1}\!\in\!\underline{Z}'_{1,\tu}(F)$.
\end{enumerate}

En effet, fixons d'abord un voisinage $\frV_{\epsilon'_{1}}$ pour tout $\epsilon'_{1}\in \epsilon_{1}\underline{Z}'_{1,\tu}(F)$.
Choisissons un sous-groupe $\frY_{\epsilon'_{1}}\subset \underline{\frz}'_{1,\tn}(F)$ de sorte que $\frY_{\epsilon'_{1}}+\frV_{\epsilon'_{1}}\subset \frV_{\epsilon'_{1}}$. On vérifie que, si $z_{1}\in \exp(\frY_{\epsilon'_{1}}) $, on peut choisir $\frV_{\epsilon'_{1}z_{1}}=\frV_{\epsilon'_{1}}$ et que l'on a
$c_{\epsilon'_{1}z_{1},\cO}=c_{\epsilon'_{1},\cO}$ pour tout~$\cO\in \Nil(\frg'_{1,\epsilon_{1}})$. Cela prouve déjà \eqref{unereciproque2}. L'ensemble $\epsilon_{1}\underline{Z}'_{1,\tu}(F)$ est compact et est réunion des ensembles ouverts $\epsilon'_{1}\exp(\frY_{\epsilon'_{1}})$. On extrait un recouvrement fini
\[
\epsilon_{1}\underline{Z}'_{1,\tu}(F)=\bigcup_{\epsilon'_{1}\in I}\epsilon'_{1}\exp(\frY_{\epsilon'_{1}}),
\]
où $I$ est un sous-ensemble fini de $\epsilon_{1}\underline{Z}'_{1,\tu}(F)$. On pose $\frV=\cap_{\epsilon'_{1}\in I}\frV_{\epsilon'_{1}}$. On voit alors que l'on peut remplacer nos $\frV_{\epsilon'_{1}}$ initiaux par le voisinage constant $\frV$. Cela démontre~\eqref{unereciproque3}.

On suppose désormais que la condition \eqref{unereciproque3} est satisfaite. Pour $\cO\in \Nil(\frg'_{1,\epsilon_{1}})$, posons
\[
\bar{c}_{\epsilon_{1},\cO}=\mes(\underline{Z}'_{1,\tu}(F))^{-1}\int_{\underline{Z}'_{1,\tu}(F)}c_{\epsilon_{1}z_{1},\cO}\,dz_{1}.
\]
On a $\bar{c}_{\epsilon_{1}z_{1},\cO}=\bar{c}_{\epsilon_{1},\cO}$ pour tout $z_{1}\in \underline{Z}'_{1,\tu}(F)$. On définit une fonction $\tau$ sur $G'_{1,\reg}(F)$, à support contenu dans l'ensemble $G'_{1,\comp}(F)$ des éléments de $G'_{1}(F)$ qui sont compacts $\rmod Z(G'_{1})$ de la façon suivante. Pour $x_{1}\in G'_{1,\comp}(F)\cap G'_{1,\reg}(F)$, on choisit une $p'$-décomposition $x_{1}=\epsilon_{1}\exp(X_{1})$ telle que l'élément $\epsilon_{1}$ soit $p'$-compact $\rmod \underline{Z}'_{1}$, \cf \ref{elementsp}\,\eqref{elementsp3}. On pose
\[
\tau(x_{1})=\sum_{\cO\in \Nil(\frg'_{1,\epsilon_{1}})}\bar{c}_{\epsilon_{1},\cO}\hat{j}(\cO,X_{1}).
\]
Cela ne dépend pas de la $p'$-décomposition choisie car, d'après \ref{elementsp}\,\eqref{elementsp3}, toute autre décomposition satisfaisant à la même propriété est de la forme
\[
x_{1}=(\epsilon_{1}\exp(Z_{1}))\exp(X_{1}-Z_{1}),
\]
avec $Z_{1}\in \underline{\frz}'_{1,\tn}(F)$ et on~a fait ce qu'il fallait pour que la formule ci-dessus soit insensible à un tel changement. Il résulte de cette définition que, pour tout élément $\epsilon_{1}\in G'_{1}(F)$ qui est $p'$-compact $\rmod \underline{Z}'_{1}$ et pour tout $X_{1}\in \frV_{\epsilon_{1}}$ tel que $\epsilon_{1}\exp(X_{1})\in G'_{1,\reg}(F)$, on~a l'égalité
\begin{equation}\tag{4}\label{unereciproque4}
\tau(\epsilon_{1}\exp(X_{1}))=\mes(\underline{Z}'_{1,\tu}(F))^{-1}\int_{\underline{Z}'_{1,\tu}(F)}\theta_{\sigma'}(\epsilon_{1}z_{1}\exp(X_{1}))\,dz_{1}.
\end{equation}
Montrons que
\begin{enumerate}\setcounter{enumi}{4}
\item\label{unereciproque5}
la fonction $\tau$ est constante sur les classes de conjugaison stable.
\end{enumerate}

Soient $x_{1},x_{2}\in G'_{1,\reg}(F)$ que l'on suppose stablement conjugués. On veut prouver que $\tau(x_{1})=\tau(x_{2})$. Remarquons que, ou bien $x_{1}$ et $x_{2}$ sont tous deux compacts $\rmod Z(G'_{1})$, ou bien aucun d'eux ne l'est. Dans le second cas, on~a \hbox{$\tau(x_{1})=\tau(x_{2})=0$}. Supposons que $x_{1}$ et $x_{2}$ sont compacts $\rmod Z(G'_{1})$. On fixe $h\in G'_{1,\reg}(\bar{F})$ tel que $x_{2}=h^{-1}x_{1}h$. On choisit une $p'$-décomposition $x_{1}=\epsilon_{1}\exp(X_{1})$ telle que $\epsilon_{1}$ soit $p'$-compact $\rmod \underline{Z}'_{1}$. Posons $\epsilon_{2}=h^{-1}\epsilon_{1}h$ et $X_{2}=h^{-1}X_{1}h$. Alors $x_{2}=\epsilon_{2}\exp(X_{2})$ et cette égalité est une $p'$-décomposition telle que $\epsilon_{2}$ soit $p'$-compact $\rmod \underline{Z}'_{1}$. Pour \hbox{$\lambda\in \fro_{F}\setminus\{0\}$}, les éléments $\epsilon_{1}\exp(\lambda^2X_{1})$ et $\epsilon_{2}\exp(\lambda^2X_{2})$ sont encore stablement conjugués. Si $\val_{F}(\lambda)$ est assez grand, on~a $\lambda^2X_{1}\in \frV_{\epsilon_{1}}$ et $\lambda^2X_{2}\in \frV_{\epsilon_{2}}$, donc $\tau(\epsilon_{1}\exp(\lambda^2X_{1}))$ et $\tau(\epsilon_{2}\exp(\lambda^2X_{2}))$ sont calculés par la formule~\eqref{unereciproque4}. Pour $z_{1}\in \underline{Z}'_{1,\tu}(F)$, les éléments $\epsilon_{1}z_{1}\exp(\lambda^2X_{1})$ et $\epsilon_{2}z_{1}\exp(\lambda^2X_{2})$ sont encore stablement conjugués. Puisque $\sigma'$ est stable, on~a donc
\[
\theta_{\sigma'}(\epsilon_{1}z_{1}\exp(\lambda^2X_{1}))=\theta_{\sigma'}(\epsilon_{2}z_{1}\exp(\lambda^2X_{2})).
\]
La formule~\eqref{unereciproque4} entraîne alors l'égalité $\tau(\epsilon_{1}\exp(\lambda^2X_{1}))=\tau(\epsilon_{2}\exp(\lambda^2X_{2}))$. Cela est vrai pour $\val_{F}(\lambda)$ assez grande. Mais, comme fonctions de $\lambda$, les deux termes de cette égalité appartiennent à $E$. Ils sont donc égaux pour tout $\lambda$. En $\lambda=1$, cela entraîne l'égalité cherchée $\tau(x_{1})=\tau(x_{2})$ qui prouve \eqref{unereciproque5}.

En particulier, $\tau$
est invariante par conjugaison. D'après le lemme \ref{variante1}, $\tau$ est la fonction associée à un quasi-caractère $D_{1}$ sur $G'_{1}(F)$ de niveau~$0$. Comme $\Theta_{\sigma'}$, ce quasi-caractère se transforme par $A_{G'_{1}}(F)$ selon le caractère $\xi$ et par $C_{1}(F)$ selon $\lambda_{1}$.

\subsubsection*{Troisième étape, fin de la preuve de la proposition}

On vient de construire un quasi-caractère $D_{1}$ sur $G'_{1}(F)$. D'après le théorème \ref{undeuxiemetheoreme}, il existe un élément $d_{1}\in \cD_{\comp}(G'_{1})$ tel que $D_{1}=D^{G'_{1}}[d_{1}]$. Notons $d_{1,\cusp}$ la composante cuspidale de $d_{1}$. Cet élément conserve les mêmes propriétés que $D_{1}$ de transformation par $A_{G'_{1}}(F)$ et $C_{1}(F)$. On note $\pi'$ l'élément de $\CC[\Ell(G'_{1})_{\xi'}^0]$ tel que $\Delta_{\cusp}(\pi')=d_{1,\cusp}$, \cf le~\eqref{surjectiviteii} du lemme \ref{surjectivite}. Par définition, $\theta_{\pi'}$ coïncide avec $\tau$ sur $G'_{1,\rell}(F)$ et cette fonction est invariante par conjugaison stable. Parce que $\pi'$ est elliptique, cela entraîne que $\pi'$ est stable, c'est-à-dire $\pi'\in \CC[\Ell(G'_{1})_{\xi'}^0]^{\st}$, \cf \cite[Th.\,6.1]{A3}.
Il est clair que la restriction à $C_{1}(F)$ du caractère central de $\pi'$ est $\lambda_{1}$.

Pour prouver la proposition, il suffit de prouver que $\transfert_{\rell}(\pi')=\pi$. Or $\pi$ et~$\pi'$ sont elliptiques. D'après \cite[Th.\,6.2]{A3}, il suffit de prouver que $\transfert_{\rell}(\pi')$ et~$\pi$ coïncident sur $G_{\rell}(F)$. C'est-à-dire qu'en posant $\Pi=\transfert_{\rell}(\pi')$,
il suffit de prouver
\begin{enumerate}\setcounter{enumi}{5}
\item\label{unereciproque6}
pour tout $x\in G_{\rell}(F)$, on~a l'égalité $\theta_{\Pi}(x)=\theta_{\pi}(x)$.
\end{enumerate}

On fixe une $p'$-décomposition $x=\epsilon \exp(X)$. Soit $\lambda\in \fro_{F}\setminus\{0\}$.
L'égalité \eqref{unpremierresultat3} de \ref{unpremierresultat} dit que
\begin{multline*}
\theta_{\Pi}( \epsilon \exp(\lambda^2X))D^G(\epsilon \exp(\lambda^2X))^{1/2}\\
=\sum_{x'\in X^{\bG'}(x)}\Delta_{1}(\epsilon_{1,x'}\exp(\lambda^2X_{1,x'}),\epsilon \exp(\lambda^2X))
\\[-8pt]
\cdot\theta_{\pi'}(\epsilon_{1,x'}\exp(\lambda^2X_{1,x'}))D^{G'}(\epsilon_{x'}\exp(\lambda^2X_{x'}))^{1/2}.
\end{multline*}
Les éléments intervenant à droite sont elliptiques, on peut remplacer la fonction~$\theta_{\pi'}$ par~$\tau$. Puisque $\epsilon$ est $p'$-compact $\rmod Z(G)$, les $\epsilon_{1,x'}$ sont $p'$-compacts $\rmod \underline{Z}'_{1}$. Si~$\val_{F}(\lambda)$ est assez grand, on utilise \eqref{unereciproque4} et on obtient
\begin{multline*}
\theta_{\Pi}( \epsilon \exp(\lambda^2X))D^G(\epsilon \exp(\lambda^2X))^{1/2}\\
=\mes(\underline{Z}'_{1,\tu}(F))^{-1}\int_{\underline{Z}'_{1,\tu}(F)}\sum_{x'\in X^{\bG'}(x)} \Delta_{1}(\epsilon_{1,x'}\exp(\lambda^2X_{1,x'}),\epsilon \exp(\lambda^2X))\\
\cdot\theta_{\sigma'}(\epsilon_{1,x'}z_{1}\exp(\lambda^2X_{1,x'}))D^{G'}(\epsilon_{x'}\exp(\lambda^2X_{x'}))^{1/2}\,dz_{1}.
\end{multline*}
Notons $z$ et $z'$ les images de $z_{1}$ dans $Z(G)_{\tu}(F)$ et $ Z(G')_{\tu}(F)$. D'après \ref{actiondescentres}\,\eqref{actiondescentres1}, on peut remplacer ci-dessus $\Delta_{1}(\epsilon_{1,x'}\exp(\lambda^2X_{1,x'}),\epsilon \exp(\lambda^2X))$ par
\[
\Delta_{1}(\epsilon_{1,x'}z_{1}\exp(\lambda^2X_{1,x'}),\epsilon z \exp(\lambda^2X)).
\]
On peut aussi remplacer $D^{G'}(\epsilon_{x'}\exp(\lambda^2X_{x'}))$ par $D^{G'}(\epsilon_{x'}z'\exp(\lambda^2X_{x'}))$. On obtient%
\begin{equation}\tag{7}\label{unereciproque7}
\theta_{\Pi}( \epsilon \exp(\lambda^2X))D^G(\epsilon \exp(\lambda^2X))^{1/2}=\mes(\underline{Z}'_{1,\tu}(F))^{-1}\int_{\underline{Z}'_{1,\tu}(F)}B(z_{1})\,dz_{1},
\end{equation}
où
\begin{multline*}
B(z_{1})=\sum_{x'\in X^{\bG'}(x)}\Delta_{1}(\epsilon_{1,x'}z_{1}\exp(\lambda^2X_{1,x'}),\epsilon z \exp(\lambda^2X))
\\[-8pt]
\cdot\theta_{\sigma'}(\epsilon_{1,x'}z_{1}\exp(\lambda^2X_{1,x'}))D^{G'}(\epsilon_{x'}z'\exp(\lambda^2X_{x'}))^{1/2}.
\end{multline*}
Il est clair que $B(z_{1})$ est le membre de droite de la formule~\eqref{unpremierresultat3} de \ref{unpremierresultat} appliquée à la représentation $\sigma'$, pour le point $\epsilon z \exp(\lambda^2X)$ de $G(F)$. Puisque $\pi=\transfert_{\rell}(\sigma')$, on obtient
\[
B(z_{1})=\theta_{\pi}(\epsilon z \exp(\lambda^2X))D^G(\epsilon z \exp(\lambda^2X))^{1/2}.
\]
Le $z$ disparaît du discriminant de Weyl. Il disparaît aussi du premier terme car $\pi$ est de niveau~$0$ par hypothèse. D'où
\[
B(z_{1})=\theta_{\pi}(\epsilon \exp(\lambda^2X))D^G(\epsilon \exp(\lambda^2X))^{1/2},
\]
puis, d'après \eqref{unereciproque7},
\[
\theta_{\Pi}( \epsilon \exp(\lambda^2X)) =\theta_{\pi}( \epsilon \exp(\lambda^2X)).
\]
On a supposé $\val_{F}(\lambda)$ assez grande. Mais, puisque $\pi'$ est de niveau~$0$, $\Pi$ l'est aussi d'après la proposition \ref{unpremierresultat}. Il en est de même pour $\pi$ par hypothèse. Les deux membres de la formule ci-dessus appartiennent donc à $E$. Étant égaux pour $\val_{F}(\lambda)$ assez grande, ils sont égaux pour tout $\lambda$. En $\lambda=1$, cela prouve \eqref{unereciproque6} et la proposition~\ref{unereciproque}.
\end{proof}

\subsection{Le théorème principal}\label{letheoreme}

Soit $\bG'=(G',\cG',s)\in \cE(G)$. Rappelons que l'on impose l'hypothèse \ref{HypendoG}.
\begin{thm*}
L'application
\[
\transfert:\CC[\Irr(G'_{1})_{\lambda_{1}}]^{\st}\to \CC[\Irr(G)]
\]
satisfait à l'égalité $p^0\circ \transfert=\transfert\circ p^0$.
\end{thm*}

\begin{rem*}
Les deux $p^0$ de l'égalité ne sont pas les mêmes: le premier vit sur le groupe $G$ et le second sur $G'_{1}$.
\end{rem*}

\begin{proof}
On a l'égalité
\begin{equation}\tag{1}\label{letheoreme1}
\CC[\Irr(G'_{1})_{\lambda_{1}}]^{\st}=\bigoplus_{M'\in \underline{\cL}_{\min}^{G'}}\Ind_{M'_{1}}^{G'_{1}}(\CC[\Ell(M'_{1})_{\RR,\lambda_{1}}]^{\st,W^{G'}(M')}),
\end{equation}
\cf \ref{lecasquasideploye}\,\eqref{lecasquasideploye1}, dont on~a adapté les notations de façon que l'on espère compréhensible. En particulier, chaque Levi $M'$ de $G'$ détermine un unique Levi $M'_{1}$ de $G'_{1}$ (son~image réciproque dans ce groupe). Notons $\underline{\cL}_{\min,G\trel}^{G'}$ l'ensemble des éléments de $\underline{\cL}_{\min}^{G'}$ qui sont \og relevants\fg pour $G$ (dans le cas où $G$ est quasi-déployé, tout élément de $\underline{\cL}_{\min}^{G'}$ est relevant). Il y a une application naturelle
de $\underline{\cL}_{\min,G\trel}^{G'}$ dans $\underline{\cL} _{\min}^G$, \cf \hbox{\cite[1.3.4]{MW}}. Si $M'\mto M$ par cette application, il se déduit de $\bG'$ une donnée endoscopique $\bM'=(M',\cM',s^M)$ de $M$, qui est elliptique.
Soit $M'\in \underline{\cL}_{\min,G\trel}^{G'}$, notons $M\in \underline{\cL} _{\min}^G$ le Levi qui lui correspond. On a l'égalité $\cA_{M}=\cA_{M'}$. L'application $\transfert_{\rell}: \CC[\Ell(M'_{1})_{\lambda_{1}}]^{\st}\to \CC[\Ell(M)]$ se prolonge en une application linéaire
$\CC[\Ell(M'_{1})_{\RR,\lambda_{1}}]^{\st}\to \CC[\Ell(M)_{\RR}]$ compatible aux opérateurs de torsion par des éléments $\chi\in \cA_{M}^*=\cA_{M'}^*$. On la compose avec la projection naturelle sur les invariants par $W^G(M)$ et on obtient par restriction une application linéaire
\[
\CC[\Ell(M'_{1})_{\RR,\lambda_{1}}]^{\st,W^{G'}(M')}\to \CC[\Ell(M)_{\RR}]^{W^G(M)}.
\]
Le transfert commute à l'induction et on en déduit une application linéaire de l'espace
\[
\Ind_{M'_{1}}^{G'_{1}}(\CC[\Ell(M'_{1})_{\RR,\lambda_{1}}]^{\st,W^{G'}(M')})
\]
dans le membre de droite de la relation \eqref{representationselliptiques1} de \ref{representationselliptiques}, c'est-à-dire dans $ \CC[\Irr(G)]$. Si maintenant $M'\in \underline{\cL}_{\min}^{G'}-\underline{\cL}_{\min,G\trel}^{G'}$, on envoie l'espace
\[
\Ind_{M'_{1}}^{G'_{1}}(\CC[\Ell(M'_{1})_{\RR,\lambda_{1}}]^{\st,W^{G'}(M')})
\]
dans $ \CC[\Irr(G)]$ par l'application nulle. La somme des applications ainsi définies sur tout $M'\in \underline{\cL}_{\min}^{G'}$ est une application linéaire du membre de droite de \eqref{letheoreme1} ci-dessus dans $ \CC[\Irr(G)]$. C'est
l'application $\transfert$. Il résulte de ces considérations que, pour démontrer le théorème, il suffit de prouver l'égalité analogue pour les applications $\transfert_{\rell}$ associés aux différents Levi de $G'$ qui sont relevants pour $G$. On ne perd rien à considérer seulement l'application $\transfert_{\rell}$ associée à $G'$ lui-même. C'est-à-dire qu'il suffit de prouver
\begin{enumerate}\setcounter{enumi}{1}
\item\label{letheoreme2}
on~a l'égalité $p^0\circ \transfert_{\rell}=\transfert_{\rell}\circ p^0$.
\end{enumerate}

Comme dans le corollaire \ref{lecasquasideploye}, cela équivaut aux deux assertions
\begin{enumerate}\setcounter{enumi}{2}
\item\label{letheoreme3}
si $\pi'\in \CC[\Ell(G'_{1})_{\lambda_{1}}^0]^{\st}$, on~a $\transfert_{\rell}(\pi')\in \CC[\Ell(G)^0]$;

\item\label{letheoreme4}
si $\pi'\in \CC[\Ell(G'_{1})_{\lambda_{1}}]^{\st}$ satisfait à $p^0(\pi')=0$, alors $p^0\circ \transfert_{\rell}(\pi')=0$.
\end{enumerate}

La proposition \ref{unpremierresultat} est justement l'assertion \eqref{letheoreme3}. Démontrons \eqref{letheoreme4}. Comme toujours, on peut décomposer $\pi'$ en $\sum_{i=1,\dots,h}\pi'_{i}$ où $\pi'_{i}$ est une combinaison linéaire de représentations elliptiques de $G'_{1}(F)$ dont le caractère central se restreint en un même caractère unitaire $\xi'_{i}$ de $A_{G'_{1}}(F)$. Ces représentations satisfont aux mêmes hypothèses que $\pi'$ et il suffit de prouver que $p^0\circ \transfert_{\rell}(\pi'_{i})=0$ pour tout $i$. En oubliant cette construction, on peut fixer un caractère unitaire $\xi'$ de $A_{G'_{1}}(F)$ et supposer que $\pi'\in \CC[\Ell(G'_{1})_{\xi',\lambda_{1}}]^{\st}$. Le caractère $\xi'$ détermine un caractère $\xi$ de $A_{G}(F)$ comme on l'a vu dans la preuve de \ref{unereciproque} et l'application $\transfert_{\rell}$ envoie $\CC[\Ell(G'_{1})_{\xi',\lambda_{1}}]^{\st}$ dans $\CC[\Ell(G)_{\xi}]$. L'action du groupe $\Aut(\bG)$ conserve le premier espace: elle préserve le caractère $\xi'$ car celui-ci est uniquement déterminé par $\xi$. On peut remplacer $\pi'$ par sa projection sur le sous-espace des invariants, cela ne change pas $\transfert_{\rell}(\pi')$. Cela ne modifie pas non plus la condition $p^0(\pi')=0$ car l'action de $\Aut(\bG')$ conserve le sous-espace $\CC[\Ell(G'_{1})^0_{\xi',\lambda_{1}}]^{\st}$. Bref, on peut supposer $\pi'\in \CC[\Ell(G'_{1})_{\xi',\lambda_{1}}]^{\st,\Aut(\bG')}$. On~pose $\pi=\transfert_{\rell}(\pi')$. Il faut maintenant se rappeler la décomposition \eqref{decomposition1} de~\ref{decomposition}. L'application $\transfert_{\rell}$ se restreint en un isomorphisme
\[
\CC[\Ell(G'_{1})_{\xi',\lambda_{1}}]^{\st,\Aut(\bG')}\to \CC[\Ell(G)_{\xi}]_{\bG'}
\]
qui est une similitude pour les produits scalaires elliptiques. On note $c$ le rapport de similitude. En particulier, $\pi\in \CC[\Ell(G)_{\xi}]_{\bG'}$. Il résulte de la proposition~\ref{decomposition} que $p^0(\pi)$ appartient aussi à $\CC[\Ell(G)_{\xi}]_{\bG'}$, donc appartient à l'image de $\transfert_{\rell}$. D'après la proposition \ref{unereciproque}, il existe $\pi^{\prime0}\in \CC[\Ell(G_{1})^0_{\lambda_{1}}]^{\st}$ telle que $\transfert_{\rell}(\pi^{\prime0})=p^0(\pi)$. Les mêmes arguments que ci-dessus permettent de supposer que $\pi^{\prime0}\in \CC[\Ell(G'_{1})^0_{\xi',\lambda_{1}}]^{\st,\Aut(\bG')}$. Parce que $p^0$ est une projection orthogonale, on a
$\<p^0(\pi),p^0(\pi)\>_{\rell}=\<p^0(\pi),\pi\>_{\rell}$. Donc
\begin{align*}
\<p^0(\pi),p^0(\pi)\>_{\rell}&=\<\transfert_{\rell}(\pi^{\prime0}),\transfert_{\rell}(\pi')\>_{\rell}\\
&=c^{-1}\<\pi^{\prime0},\pi'\>_{\rell}=c^{-1}\<\pi^{\prime0},p^0(\pi')\>_{\rell},
\end{align*}
toujours parce que $p^0$ est une projection orthogonale. Puisque $p^0(\pi')=0$, cela entraîne
$\<p^0(\pi),p^0(\pi)\>_{\rell}=0$, d'où $p^0(\pi)=0$. Cela achève la preuve de \eqref{letheoreme4} et du théorème.
\end{proof}

\section{Quelques conséquences}

\subsection{Le cas quasi-déployé derechef}\label{derechef}
Supposons $G$ quasi-déployé.
\begin{cor*}
La projection de Bernstein $p^0$ de $\rI(G)$ se quotiente en une projection $p^0$ de $\SI(G)$.
\end{cor*}

\begin{proof}
Notons $I^{\inst}(G)$ le noyau de la projection $\rI(G)\to \SI(G)$. On doit prouver que $p^0$ conserve ce sous-espace. C'est-à-dire, soit $f\in I^{\inst}(G)$, on doit prouver que $p^0(f)\in I^{\inst}(G)$. Il suffit pour cela de prouver que, pour tout $\pi^{\st}\in \CC[\Irr(G)]^{\st}$, on~a $\Theta_{\pi^{\st}}(p^0(f))=0$, \cf \cite[Cor.\,XI.5.2(i)]{MW}. D'après les propriétés des projecteurs de Bernstein, on a
\[
\Theta_{\pi^{\st}}(p^0(f))=\Theta_{p^0(\pi^{\st})}(f).
\]
D'après le~\eqref{lecasquasideployeii} du corollaire \ref{lecasquasideploye}, on~a $p^0(\pi^{\st})=p^0\circ p^{\st}(\pi^{\st})=p^{\st}\circ p^0(\pi^{\st})$. Donc $p^0(\pi^{\st})\in \CC[\Irr(G)]^{\st}$. Alors $\Theta_{p^0(\pi^{\st})}(f)=0$ puisque $f\in I^{\inst}(G)$.
\end{proof}

\subsection{Transfert de fonctions}\label{transfertdefonctions}
Soit $\bG'=(G',\cG',s)\in \cE(G)$.
\begin{cor*}
On suppose \ref{HypendoG} satisfaite. L'application $\transfert:\rI(G)\to \SI_{\lambda_{1}}(G'_{1})$ satisfait à l'égalité $p^0\circ \transfert=\transfert\circ p^0$.
\end{cor*}

Cela se déduit du théorème \ref{letheoreme} comme le corollaire précédent se déduisait du corollaire \ref{lecasquasideploye}.

\subsection{Stabilité dans l'espace \texorpdfstring{$\cD_{\cusp}$}{Dcusp}}\label{stabilite}
Supposons $G$ quasi-déployé. Disons qu'un élément $d\in \cD_{\cusp}(G)$ est stable si la distribution $D^G[d]$ l'est et qu'il est stable sur les elliptiques si la distribution $D^G[d]$ l'est, c'est-à-dire si $\theta_{D^G[d]}$ est constante sur les classes de conjugaison stable contenues dans $G_{\rell}(F)$. On note $\cD_{\cusp}(G)^{\st}$ le sous-espace des éléments stables dans $\cD_{\cusp}(G)$.

\begin{prop*}
Soit $d\in \cD_{\cusp}(G)$. Alors $d$ est stable si et seulement si $d$ est stable sur les elliptiques.
\end{prop*}

\begin{proof}
Utilisons la décomposition $\cD_{\cusp}(G)=\prod_{\nu\in \cN}\cD_{\cusp}(G)^{\nu}$. Il~est clair que $d$ est stable, \resp stable sur les elliptiques, si et seulement si chaque composante $d^{\nu}$ est stable, \resp stable sur les elliptiques. On peut donc fixer $\nu\in \cN$ et supposer $d\in \cD_{\cusp}(G)^{\nu}$. Notons $\cX$ l'ensemble des restrictions à $A_{G}(F)_{\rc}$ de caractères modérément ramifiés de $A_{G}(F)$. Pour $\xi\in \cX$, posons
\[
d_{\xi}=\mes(A_{G}(F)_{\rc})^{-1}\int_{A_{G}(F)_{\rc}}\xi(a)^{-1}{d^{a}}\, da.
\]
Alors $d=\sum_{\xi\in \cX}d_{\xi}$. De nouveau, $d$ est stable, \resp stable sur les elliptiques, si et seulement si chaque composante $d_{\xi}$ est stable, \resp stable sur les elliptiques. On peut fixer $\xi\in \cX$ et supposer que $d$ se transforme par $A_{G}(F)_{\rc}$ selon $\xi$. Prolongeons $\xi$ en un caractère unitaire de $A_{G}(F)$. Alors il existe un unique $\bmd\in \cD_{\cusp}(G)$ tel que:
\begin{itemize}
\item
$\bmd\in \cD_{\cusp,\xi}(G)$, c'est-à-dire $\bmd$ se transforme par $A_{G}(F)$ selon $\xi$;

\item
pour $\nu'\in \cN$, la composante $\bmd^{\nu'}$ est nulle si $\nu'\not\in \nu+w_{G}(A_{G}(F))$;

\item
$\bmd^{\nu}=d^{\nu}$.
\end{itemize}

On voit que $d$ est stable, \resp stable sur les elliptiques, si et seulement si $\bmd$ l'est. En~oubliant cette construction, on peut fixer un caractère unitaire $\xi$ de $A_{G}(F)$ et supposer $d\in \cD_{\cusp,\xi}(G)$. D'après le~\eqref{surjectiviteii} du lemme \ref{surjectivite}, il existe une unique $\pi\in \CC[\Ell(G)_{\xi}^0]$ telle que $\Delta_{\cusp}(\pi)=d$. Alors, d'après \ref{leresultat}\,\eqref{leresultat1}, $\Theta_{\pi}$ est stable sur les \hbox{elliptiques}. D'après~\cite[Th.\,6.1]{A3}, $\pi$ est stable. Soient $M\in \cL_{\min}$ et $P\in \cP(M)$. Puisque $\pi$ est stable, il en est de même du module de Jacquet $\pi_{P}$. D'après
\ref{leresultat}\,\eqref{leresultat1}, $\Delta^M_{\pi_{P},\cusp}$ est stable sur les elliptiques. A fortiori, sa projection $\Delta^M_{\pi_{M},\cusp,G\tcomp}$ dans $\cD_{\cusp,G\tcomp}(M)$ est stable sur les elliptiques (cette projection consiste à restreindre le support à un sous-ensemble invariant par conjugaison stable). Supposons que~$M$ est un Levi propre. On peut raisonner par récurrence et supposer notre proposition déjà démontrée pour $M$. Donc $\Delta^M_{\pi_{M},\cusp,G\tcomp}$ est stable. C'est-à-dire que $D^M(\Delta^M_{\pi_{M},\cusp,G\tcomp})$ est une distribution stable sur $M(F)$. La stabilité se conserve par induction donc $D^G(\Delta^M_{\pi_{M},\cusp,G\tcomp})$ est une distribution stable sur $G(F)$. La formule \ref{leresultat}\,\eqref{leresultat1} exprime alors $D^G[d]$ comme la différence entre $\Theta_{\pi}$ qui est stable et la somme sur les $M\neq G$ de distributions stables. Donc $D^G[d]$ est stable et $d$ l'est alors par définition.
\end{proof}

\subsection{Transfert et module de Jacquet}\label{transfertetmoduledeJacquet}
Soit $\bG'=(G',\cG',s)\in \cE(G)$. On a rappelé dans la preuve du théorème \ref{letheoreme} l'existence d'une application
\begin{equation}\tag{1}\label{transfertetmoduledeJacquet1}
\underline{\cL}_{\min,G\trel}^{G'}\to\underline{\cL}_{\min}^G.
\end{equation}
Notons-la $j$. Rappelons quelques-unes de ses propriétés. Soit $M'\in
\underline{\cL}_{\min,G\trel}^{G'}$, posons
$M=j(M')$. Alors $M'$ se complète en une donnée endoscopique $\bM'=(M',\cM',s^M)$ de $M$, qui est elliptique.
On sait définir une correspondance entre classes de conjugaison stable dans $G'_{\reg}(F)$ et dans $G_{\reg}(F)$. On définit de même à l'aide de~$\bM'$ une correspondance entre classes de conjugaison stable dans $M'_{\reg}(F)$ et dans $M_{\reg}(F)$. Appelons correspondance pour $\bG'$ la première et correspondance pour $\bM'$ la seconde. On a:
\begin{enumerate}\setcounter{enumi}{1}
\item\label{transfertetmoduledeJacquet2}
soient $x'\in M'(F)$ et $x\in M(F)\cap G_{\reg}(F)$; supposons que $x$ et $x'$ se correspondent pour $\bM'$; alors ils se correspondent pour $\bG'$.
\end{enumerate}

Signalons que la réciproque est fausse.
Des données auxiliaires que l'on~a fixées pour $\bG'$ se déduisent des données pour $\bM'$. En particulier $M'_{1}$ est l'image réciproque de $M'$ dans $G'_{1}$. On a fixé un facteur de transfert pour la donnée $\bG'$, notons-le plus précisément $\Delta_{1}^{\bG'}$. Il y a un unique facteur de transfert $\Delta_{1}^{\bM'}$ pour la donnée $\bM'$ de sorte que
\begin{enumerate}\setcounter{enumi}{2}
\item\label{transfertetmoduledeJacquet3}
pour $x$, $x'$ comme en \eqref{transfertetmoduledeJacquet2} et pour $x'_{1}\in M'_{1}(F)$ au-dessus de $x'$, on ait
$\Delta_{1}^{\bM'}(x_{1}',x)=\Delta_{1}^{\bG'}(x'_{1},x)$.
\end{enumerate}

À l'aide de ce facteur de transfert, on définit comme en \ref{unpremierresultat} l'application
\[
\transfert:\CC[\Irr(M'_{1})_{\lambda_{1}}]^{\st}\to \CC[\Irr(M)].
\]

Puisque $\bM'$ est une donnée endoscopique elliptique de $M$, il y a un isomorphisme $A_{M'}\simeq A_{M}$. Pour un Levi $L$ de $G$ contenant $M$, on identifie $A_{L}\subset A_{M}$ à un sous-ensemble de $A_{M'}$ et on peut définir le commutant $Z_{G'}(A_{L})$ de $A_{L}$ dans $G'$. C'est un Levi de $G'$ contenant $M'$. Notons $J(M')$ l'ensemble des Levi $L$ de $G$ contenant $M$ tels que $Z_{G'}(A_{L})=M'$.
Les racines de $A_{M'}$ dans $G'$ sont aussi des racines de $A_{M}$ dans~$G$. La condition $Z_{G'}(A_{L})=M'$ équivaut à ce que la restriction à $\cA_{L}$ de toute racine de $A_{M'}$ dans $G$ soit non nulle. Soit $L\in J(M')$ et soit $Q\in \cP(L)$. À $Q$ est associée une chambre dans $\cA_{L}$. L'hyperplan de $\cA_{M}$ annulé par une racine de $A_{M}$ dans $G$ ne coupe cette chambre que si cette racine s'annule sur $\cA_{L}$. Ainsi l'hyperplan associé à une racine de $A_{M'}$ dans $G'$ ne coupe pas cette chambre. Celle-ci est donc incluse dans la chambre de $\cA_{M'}$ associée à un unique sous-groupe parabolique $P'\in \cP(M')$, que l'on~note $P'_{Q}$.

Soit $\pi'\in \CC[ \Irr(G')_{\lambda_{1}}]^{\st}$, posons $\pi=\transfert(\pi')$. Soient $M\in \underline{\cL}_{\min}$ et
$P\in \cP(M)$. Notons $\pi_{P}$ le module de Jacquet normalisé de $\pi$ relatif à $P$ et $ \Theta_{\pi_{P}\vert\rell}$ la restriction de $\Theta_{\pi_{P}}$ à $M_{\rell}(F)$.

\begin{prop*}
Sous ces hypothèses, on~a les égalités
\begin{align*}
\Theta_{\pi_{P}}&=\sum_{M'\in \underline{\cL}_{\min, G\trel}^{G'}}\hspace*{-5mm}\vert W^{G'}(M')\vert ^{-1}\hspace*{-3mm}\sum_{\substack{w\in W^G/W^M\\w(M)\in J(M')}}\hspace*{-3mm}\Ind_{w^{-1}(j(M'))}^M(w^{-1}\circ \transfert(\Theta_{\pi'_{P'_{w(P)}}}));\\
\Theta_{\pi_{P}\vert\rell}&=\sum_{\substack{M'\in \underline{\cL}_{\min,G\trel}^{G'}\\ j(M')= M}}\hspace*{-5mm}\vert W^{G'}(M')\vert ^{-1}\hspace*{-3mm}\sum_{w\in W^{G}(M)}\hspace*{-3mm}w^{-1}\circ \transfert(\Theta_{\pi'_{P'_{w(P)}}\vert\rell}).
\end{align*}
\end{prop*}

\begin{proof}
La première égalité est une autre formulation de \cite[Th.\,5.6]{H}. Déduisons-en la seconde. Considérons une représentation induite intervenant dans le membre de droite de la première. S'il s'agit d'une induite propre, la restriction de son caractère à $M_{\rell}(F)$ est nulle. Ces induites disparaissent et la sommation en $w$ se limite aux $w$ tels que $w^{-1}(j(M'))=M$. Mais $M$ et $j(M')$ appartiennent tous deux à notre ensemble de représentants $\underline{\cL}_{\min}$ des classes de conjugaison de Levi. Ils sont conjugués si et seulement s'ils sont égaux. La somme en $w$ est donc vide si $j(M')\neq M$. Si $j(M')=M$, elle devient une somme sur les $w\in W^G/W^M$ tels que $w(M)=M$, c'est-à-dire $w\in W^G(M)$. Enfin, si $j(M')=M$ et $w\in W^G(M)$, puisque $\bM'$ est une donnée endoscopique elliptique de $M$, un élément de $M'(F)$ qui se transfère en un élément de $M_{\rell}(F)$ appartient lui-même à $M'_{\rell}(F)$. Autrement dit, on~a dans ce cas $(w^{-1}\circ \transfert(\Theta_{\pi'_{P'_{w(P)}}}))_{\vert\rell}=w^{-1}\circ \transfert(\Theta_{\pi'_{P'_{w(P)}}\vert\rell})$. On obtient alors la deuxième égalité de l'énoncé.
\end{proof}

\Subsection{ L'espace \texorpdfstring{$\cD_{\cusp}(G)$}{DG} et le transfert}\label{finale}

\begin{prop*}
On suppose \ref{HypendoG} satisfaite. Soient $d\in \cD_{\cusp}(G)$ et, pour tout $\bG'\in \cE(G)$, soit $d^{\bG'}\in \cD_{\cusp,\lambda_{1}}(G'_{1})^{\st}$. Supposons que $D^G[d]$ coïncide avec $\sum_{\bG'\in \cE(G)}\transfert(D^{G'_{1}}[d^{\bG'}])$ sur $G_{\rell}(F)$. Alors\vspace*{-3pt}
\[
D^G[d]=\sum_{\bG'\in \cE(G)}\transfert(D^{G'_{1}}[d^{\bG'}]).
\]
\end{prop*}

\begin{proof}
Soit $\bG'\in \cE(G)$. Des injections naturelles $Z(\hat{G})\to Z(\hat{G}')\to Z(\hat{G}'_{1})$ se déduisent des homomorphismes\vspace*{-3pt}
\begin{equation}\tag{1}\label{finale1}
\cN^{G'_{1}}\to \cN^{G'}\to \cN.
\end{equation}

Pour $\nu\in \cN$, on définit la composante $d^{\nu}$ de $d$ comme dans le paragraphe précédent. On définit $d^{\bG',\nu}$ comme le produit des composantes $d^{\bG',\nu'}$ sur les $\nu'\in \cN^{G'_{1}}$ qui s'envoient sur $\nu$ par la suite ci-dessus. La famille $((d^{\bG'})_{\bG'\in \cE(G)},d)$ satisfait à l'hypothèse ou la conclusion de la proposition si et seulement si chaque famille $((d^{\bG',\nu})_{\bG'\in \cE(G)},d^{\nu})$ les satisfait. On peut donc fixer $\nu\in \cN$ et supposer $d\in \cD_{\cusp}(G)^{\nu}$. Notons $\cX$ l'ensemble des restrictions à $A_{G}(F)_{\rc}$ de caractères modérément ramifiés de $A_{G}(F)$. Pour $\bG'\in \cE(G)$, cet ensemble s'identifie à celui des restrictions à $A_{G'_{1}}(F)_{\rc}$ de caractères modérément ramifiés de $A_{G'_{1}}(F)$ qui coïncident avec $\lambda_{1}$ sur $A_{G'_{1}}(F)\cap C_{1}(F)$: à $\xi\in \cX$, on associe d'abord l'image réciproque de $\xi$ par la surjection $A_{G'_{1}}(F)_{\rc}\to A_{G}(F)_{\rc}$, puis le produit $\xi^{\bG'}$ de ce caractère par la restriction de $\zeta_{1}$ à $A_{G'_{1}}(F)_{\rc}$. Ainsi, comme dans le paragraphe \ref{stabilite}, on associe à $d$, \resp $d^{\bG'}$, des composantes $d_{\xi}$, \resp $d^{\bG'}_{\xi^{\bG'}}$, pour $\xi\in \cX$. La famille $((d^{\bG'})_{\bG'\in \cE(G)},d)$ satisfait à l'hypothèse ou la conclusion de la proposition si et seulement si chaque famille $((d^{\bG'}_{\xi^{\bG'}})_{\bG'\in \cE(G)},d_{\xi})$ y satisfait. On peut fixer $\xi\in \cX$ et supposer que $d=d_{\xi}$ et $d^{\bG'}=d^{\bG'}_{\xi^{\bG'}}$ pour tout $\bG'$. Prolongeons $\xi$ en un caractère unitaire de $A_{G}(F)$, qui s'identifie comme ci-dessus pour tout $\bG'$ à un caractère $\xi^{\bG'}$ de $A_{G'_{1}}(F)$ dont la restriction à $A_{G'_{1}}(F)\cap C_{1}(F)$ coïncide avec la restriction de $\lambda_{1}$. On~construit un élément $\bmd\in \cD_{\cusp,\xi}(G)$ comme en \ref{stabilite} et de façon similaire des éléments $\bmd^{\bG'}\in \cD_{\cusp,\lambda_{1},\xi^{\bG'}}(G'_{1})$. La famille $((d^{\bG'})_{\bG'\in \cE(G)},d)$ satisfait à l'hypothèse ou la conclusion de la proposition si et seulement si la famille $((\bmd^{\bG'})_{\bG'\in \cE(G)},\bmd)$ y satisfait. En résumé, on peut fixer un caractère unitaire $\xi$ de $A_{G}(F)$, qui détermine pour tout $\bG'$ un tel caractère $\xi^{\bG'}$ de $A_{G'_{1}}(F)$, et supposer $d\in \cD_{\cusp,\xi}(G)$ et $d^{\bG'}\in \cD_{\cusp,\lambda_{1},\xi^{\bG'}}(G'_{1})$ pour tout $\bG'$.

À ce point, on peut simplifier le problème en se ramenant au cas où une unique donnée $\bG'$ intervient. Pour cela, on introduit la représentation $\pi\in \CC[\Ell(G)_{\xi}^0]$ telle que $\Delta_{\pi,\cusp}=d$. On utilise la décomposition \eqref{decomposition1} de \ref{decomposition} et on écrit conformément $\pi=\sum_{\bG'\in \cE(G)}\pi_{\bG'}$. On sait que $\pi_{\bG'}$ satisfait aux les mêmes propriétés que $\pi$ d'après la proposition \ref{decomposition}. On note $d_{\bG'}=\Delta_{\cusp}(\pi_{\bG'})$. On a $d=\sum_{\bG'}d_{\bG'}$. On voit que si $((d^{\bG'})_{\bG'\in \cE(G)},d)$ satisfait à l'hypothèse de l'énoncé, alors, pour tout $\bG'\in \cE(G)$, $D^G[d_{\bG'}]$ coïncide sur $G_{\rell}(F)$ avec $\transfert(D^{G'_{1}}[d^{\bG'}])$. Inversement, si $D^G[d_{\bG'}]=\transfert(D^{G'_{1}}[d^{\bG'}])$ pour tout $\bG'$, la conclusion de l'énoncé est satisfaite. Cela nous ramène au cas où la famille $(d^{\bG'})_{\bG'\in \cE(G)}$ a au plus une composante non nulle. On note désormais $\bG'$ l'indice de cette composante et on~note $(d',d)\in \cD_{\cusp,\lambda_{1};\xi'}(G'_{1})^{\st}\times \cD_{\cusp,\xi}(G)$ le couple auquel se réduisent les données de départ.

On introduit la représentation $\pi\in \CC[\Ell(G)_{\xi}^0]$ telle que $\Delta_{\pi,\cusp}=d$ et la représentation $\pi'\in \CC[\Ell(G'_{1})^0_{\lambda_{1},\xi'}]^{\st}$ telle que $\Delta^{G'_{1}}_{\pi',\cusp}=d'$. Alors, d'après \ref{leresultat}\,\eqref{leresultat1}, $\Theta_{\pi}$~coïncide avec $\Theta_{\transfert(\pi')}$ sur les elliptiques. D'après \cite[Th.\,6.2]{A3}, $\pi=\transfert(\pi')$. Soient $M\in \underline{\cL}_{\min}$ et $P\in \cP(M)$. D'après \ref{leresultat}\,\eqref{leresultat1} appliqué en remplaçant $G$ par $M$, la deuxième égalité de la proposition \ref{transfertetmoduledeJacquet} nous dit que $D^M[\Delta^M_{\pi_{P},\cusp}]$ coïncide sur $M_{\rell}(F)$ avec
\begin{equation}\tag{3}\label{finale3}
\sum_{\substack{M'\in \underline{\cL}_{\min,G\trel}^{G'}\\ j(M')=M}}\hspace*{-3mm}\vert W^{G'}(M')\vert ^{-1}\sum_{w\in W^{G}(M)}\hspace*{-2mm}w^{-1}\circ \transfert( D^{M'_{1}}[\Delta^{M'_{1}}_{\pi'_{P'_{w(P)}},\cusp}]).
\end{equation}
Remarquons que, pour $M'$ et $w$ intervenant ci-dessus, on peut transformer la donnée endoscopique $\bM'$ de $M$ en une donnée $w^{-1}(\bM')$. Le terme que l'on somme peut se récrire $\transfert(D^{w^{-1}(M'_{1})}[d^{w^{-1}(M'_{1})}])$ pour un élément convenable $d^{w^{-1}(M'_{1})}\in \cD_{\cusp,w^{-1}(\lambda_{1})}(w^{-1}(M'_{1}))$. C'est-à-dire que la somme \eqref{finale3} est de la même forme que celle qui intervient dans l'énoncé de la proposition.
Supposons $M\neq G$. En raisonnant par récurrence, on suppose notre proposition prouvée pour $M$. Alors $D^M[\Delta^M_{\pi_{P},\cusp}]$ est égal à \eqref{finale3}. Pour un couple $(x'_{1},x)\in G'_{1,\reg}(F)\times G_{\reg}(F)$ d'éléments qui se correspondent, $x$ est compact $\rmod Z(G)$ si et seulement si $x'_{1}$ est compact $\rmod Z(G'_{1})$: cela parce
que $Z(G)(F)$ est un sous-groupe cocompact de $Z(G')(F)$. Restreindre $D^M[\Delta^M_{\pi_{P},\cusp}]$ à $M(F)\cap G_{\comp}(F)$ équivaut à restreindre chaque terme de \eqref{finale3} indexé par $M'$ à $M_{1}'(F)\cap G'_{1,\comp}(F)$. Après restriction à $M(F)\cap G_{\comp}(F)$, on sait que le terme $\Delta^M_{\pi_{P},\cusp}$ ne dépend plus de $P$. Il en est de même des termes intervenant dans~\eqref{finale3}. On obtient que $D^M[\Delta^M_{\pi_{M},\cusp,G\tcomp}]$ est égal à
\[
\sum_{\substack{M'\in \underline{\cL}_{\min,G\trel}^{G'}\\ j(M')=M}}\hspace*{-4mm}\vert W^{G'}(M')\vert ^{-1}\sum_{w\in W^{G}(M)} \hspace*{-2mm}w^{-1}\circ \transfert(D^{M'_{1}}[\Delta^{M'_{1}}_{\pi'_{M'},\cusp,G'_{1}\tcomp}]).
\]
Induisons à $G$. Évidemment, l'induction est insensible à la torsion par un élément de~$W^G$. Les $w$ disparaissent de la formule et la somme en $w$ est remplacée par la multiplication par $\vert W^G(M)\vert $. D'autre part, l'induction commute au transfert. On obtient
\begin{multline*}\tag{4}\label{finale4}
D^G[\Delta^M_{\pi_{M},\cusp,G\tcomp}]\\
= \sum_{\substack{M'\in \underline{\cL}_{\min,G\trel}^{G'}\\ j(M')=M}}\hspace*{-5mm}\vert W^{G'}(M')\vert ^{-1}\vert W^G(M)\vert \transfert(D^{G'_{1}}[\Delta^{M'_{1}}_{\pi'_{M'},\cusp,G'_{1}\tcomp}]).
\end{multline*}
On utilise l'égalité \ref{leresultat}\,\eqref{leresultat2} qui
peut se récrire
\[
D^G[d]= \Theta_{\pi\vert G_{\comp}(F)}-\sum_{\substack{M\in \underline{\cL}_{\min}\\M\neq G}}\vert W^G(M)\vert ^{-1}D^G[\Delta^M_{\pi_{M},\cusp,G\tcomp}],
\]
où $ \Theta_{\pi\vert G_{\comp}(F)}$ est la restriction de $\Theta_{\pi}$ à $G_{\comp}(F)$.
En utilisant \eqref{finale4}, on obtient
\begin{multline*}
D^G[d]
= \Theta_{\pi\vert G_{\comp}(F)}-\hspace*{-2mm}\sum_{\substack{M\in \underline{\cL}_{\min}\\M\neq G}}\sum_{\substack{M'\in \underline{\cL}_{\min,G\trel}^{G'}\\ j(M')=M}}\hspace*{-5mm}\vert W^{G'}(M')\vert ^{-1}
\transfert(D^{G'_{1}}[\Delta^{M'_{1}}_{\pi'_{M'},\cusp,G'_{1}\tcomp}])
\\
=\Theta_{\pi\vert G_{\comp}(F)}-\hspace*{-2mm}\sum_{\substack{M'\in \underline{\cL}_{\min,G\trel}^{G'}\\ M'\neq G'}}\hspace*{-5mm}\vert W^{G'}(M')\vert ^{-1} \transfert(D^{G'_{1}}[\Delta^{M'_{1}}_{\pi'_{M'},\cusp,G'_{1}\tcomp}]).
\end{multline*}
On peut remplacer l'ensemble de sommation $ \underline{\cL}_{\min,G\trel}^{G'}$ par $ \underline{\cL}_{\min}^{G'}$: pour un Levi non relevant, le terme que l'on somme est nul.
En utilisant de nouveau \ref{leresultat}\,\eqref{leresultat2} cette fois dans $G'_{1}$, on obtient
\[
D^G[d]= \Theta_{\pi\vert G_{\comp}(F)}+\transfert(D^{G'_{1}}[d']-\Theta_{\pi'\vert G'_{1,\comp}(F)}).
\]
On a déjà dit que, grâce à Arthur, on avait $\pi=\transfert(\pi')$ donc les termes $\Theta_{\pi\vert G_{\comp}(F)}$ et $\transfert(\Theta_{\pi'\vert G'_{1,\comp}(F)})$ disparaissent. D'où
\[
D^G[d]= \transfert(D^{G'_{1}}[d']),
\]
ce qu'il fallait démontrer.
\end{proof}

\Changelback

% End literal retained input author-body.tex

\begin{thebibliography}{99}
\bibitem{A1} J. Arthur - “The local behaviour of weighted orbital integrals” , Duke Math. J. 56 (1988) no. 2, p. 223-293
\bibitem{A2} J. Arthur - “On elliptic tempered characters” , Acta Math. 171 (1993) no. 1, p. 73-138
\bibitem{A3} J. Arthur - “On local character relations” , Selecta Math. (N.S.) 2 (1996) no. 4, p. 501-579
\bibitem{A} A.-M. Aubert - “Séries de Harish-Chandra de modules et correspondance de Howe modulaire” , J. Algebra 165 (1994) no. 3, p. 576-601
\bibitem{B} J. N. Bernstein - “Le “centre” de Bernstein” , in Representations of reductive groups over a local field (J. N. Bernstein, P. Deligne, D. Kazhdan \& M.-F. Vignéras, eds.) , Travaux en Cours , Hermann, Paris, 1984, p. 1-32, rédigé par P. Deligne
\bibitem{BT} F. Bruhat \& J. Tits - “Groupes réductifs sur un corps local : I. Données radicielles valuées” , Publ. Math. Inst. Hautes Études Sci. 41 (1972), p. 5-251
\bibitem{C} W. Casselman - “Characters and Jacquet modules” , Math. Ann. 230 (1977) no. 2, p. 101-105
\bibitem{D} S. DeBacker - “Homogeneity results for invariant distributions of a reductive $p$-adic group” , Ann. Sci. École Norm. Sup. (4) 35 (2002) no. 3, p. 391-422
\bibitem{DR} S. DeBacker \& M. Reeder - “Depth-zero supercuspidal $L$-packets and their stability” , Ann. of Math. (2) 169 (2009) no. 3, p. 795-901
\bibitem{F} A. Ferrari - “Théorème de l’indice et formule des traces” , Manuscripta Math. 124 (2007) no. 3, p. 363-390
\bibitem{HR} T. Haines \& M. Rapoport - “On parahoric subgroups” , appendice à [28]
\bibitem{HC} Harish-Chandra - Admissible invariant distributions on reductive $p$-adic groups , University Lecture Series , vol. 16 , American Mathematical Society, Providence, RI, 1999, rédigé par S. DeBacker et P. Sally
\bibitem{HV} G. Henniart \& M.-F. Vignéras - “A Satake isomorphism for representations modulo $p$ of reductive groups over local fields” , J. reine angew. Math. 701 (2015), p. 33-75
\bibitem{H} K. Hiraga - “On functoriality of Zelevinski involutions” , Compositio Math. 140 (2004) no. 6, p. 1625-1656
\bibitem{Kh} D. Kazhdan - “Cuspidal geometry of $p$-adic groups” , J. Analyse Math. 47 (1986), p. 1-36
\bibitem{KV} D. Kazhdan \& Y. Varshavsky - “Endoscopic decomposition of certain depth zero representations” , in Studies in Lie theory , Progress in Math. , vol. 243 , Birkhäuser Boston, Boston, MA, 2006, p. 223-301
\bibitem{KM} J.-L. Kim \& F. Murnaghan - “Character expansions and unrefined minimal $K$-types” , Amer. J. Math. 125 (2003) no. 6, p. 1199-1234
\bibitem{K} R. E. Kottwitz - “Isocrystals with additional structure. II” , Compositio Math. 109 (1997) no. 3, p. 255-339
\bibitem{Lanard1} T. Lanard - “Sur les $\ell $-blocs de niveau zéro des groupes $p$-adiques” , Compositio Math. 154 (2018) no. 7, p. 1473-1507
\bibitem{Lanard2} T. Lanard - “Sur les $\ell $-blocs de niveau zéro des groupes $p$-adiques II” , 2018
\bibitem{L} R. P. Langlands - “Stable conjugacy : definitions and lemmas” , Canad. J. Math. 31 (1979) no. 4, p. 700-725
\bibitem{LS} R. P. Langlands \& D. Shelstad - “On the definition of transfer factors” , Math. Ann. 278 (1987) no. 1-4, p. 219-271
\bibitem{Lu} G. Lusztig - “Classification of unipotent representations of simple $p$-adic groups” , Internat. Math. Res. Notices (1995) no. 11, p. 517-589
\bibitem{MW} C. Moeglin \& J.-L. Waldspurger - Stabilisation de la formule des traces tordue , Progress in Math. , vol. 316 \& 317 , Birkhäuser/Springer, Cham, 2016
\bibitem{MP1} A. Moy \& G. Prasad - “Unrefined minimal $K$-types for $p$-adic groups” , Invent. Math. 116 (1994) no. 1-3, p. 393-408
\bibitem{MP} A. Moy \& G. Prasad - “Jacquet functors and unrefined minimal $K$-types” , Comment. Math. Helv. 71 (1996) no. 1, p. 98-121
\bibitem{Oi} M. Oi - “Depth preserving property of the local Langlands correspondence for quasi-split classical groups in a large residual characteristic” , 2018
\bibitem{P-R08} G. Pappas \& M. Rapoport - “Twisted loop groups and their affine flag varieties” , Adv. Math. 219 (2008) no. 1, p. 118-198
\bibitem{P} G. Prasad - “Finite group actions on reductive groups and buildings and tamely-ramified descent in Bruhat-Tits theory” , Amer. J. Math. 142 (2020) no. 4, p. 1239-1267
\bibitem{PY} G. Prasad \& J.-K. Yu - “On finite group actions on reductive groups and buildings” , Invent. Math. 147 (2002) no. 3, p. 545-560
\bibitem{S} M. Solleveld - “On unipotent representations of ramified $p$-adic groups” , 2019
\bibitem{T} J. Tits - “Reductive groups over local fields” , in Automorphic forms, representations and $L$-functions (Oregon State Univ., Corvallis, Ore., 1977), Part 1 , Proc. Sympos. Pure Math. , vol. XXXIII , American Mathematical Society, Providence, RI, 1979, p. 29-69
\bibitem{W4} J.-L. Waldspurger - L’endoscopie tordue n’est pas si tordue , Mem. Amer. Math. Soc. , vol. 194 no. 908 , American Mathematical Society, Providence, RI, 2008
\bibitem{W1} J.-L. Waldspurger - “Une formule intégrale reliée à la conjecture locale de Gross-Prasad” , Compositio Math. 146 (2010) no. 5, p. 1180-1290
\bibitem{W3} J.-L. Waldspurger - “Une formule intégrale reliée à la conjecture locale de Gross-Prasad, 2 e partie : extension aux représentations tempérées” , in Sur les conjectures de Gross et Prasad. I , Astérisque , vol. 346 , Société Mathématique de France, Paris, 2012, p. 171-312
\bibitem{W2} J.-L. Waldspurger - “Caractères de représentations de niveau 0” , Ann. Fac. Sci. Toulouse Math. (6) 27 (2018) no. 5, p. 925-984
\end{thebibliography}
\end{document}
